\documentclass[11pt,letterpaper,reqno]{amsart}
\usepackage[T1]{fontenc}
\usepackage{lmodern,microtype}
\usepackage[letterpaper,margin=1in,headheight=15pt]{geometry}
\usepackage{amsmath,amsthm,amssymb,mathtools,mathrsfs}
\usepackage{enumitem,perpage}
\usepackage{xcolor}
\definecolor{manuscriptred}{rgb}{0.6,0.1,0.1}
\definecolor{manuscriptblue}{rgb}{0.1,0.3,0.5}
\definecolor{manuscriptmagenta}{HTML}{60445C}
\usepackage[hypertexnames=false,pdfauthor={}]{hyperref}
\hypersetup{colorlinks=true,linkcolor=manuscriptred,citecolor=manuscriptblue,urlcolor=manuscriptmagenta}
\hypersetup{pdftitle={Massless photons and Maxwell infrared limits in three-dimensional QED}}
\setcounter{tocdepth}{1}
\renewcommand{\thefigure}{\Roman{figure}}
\renewcommand{\thetable}{\Roman{table}}
\MakePerPage{footnote}
\renewcommand{\thefootnote}{\ifnum\value{footnote}>9\arabic{footnote}\else\fnsymbol{footnote}\fi}
\allowdisplaybreaks
\emergencystretch=3em
\newcommand{\namedhead}{}
\theoremstyle{definition}
\newtheorem*{namedinner}{\namedhead}
% Preserve the manually assigned statement identifiers.
\makeatletter
\def\qcm@settype#1 #2\@nil{\def\qcm@type{#1}}
\def\qcm@theorem{Theorem}
\def\qcm@corollary{Corollary}
\newenvironment{named}[1]{%
  \renewcommand{\namedhead}{#1}\begin{namedinner}%
  \qcm@settype#1 \@nil\pushQED{\qed}%
  \ifx\qcm@type\qcm@theorem
    \renewcommand{\qedsymbol}{$\blacksquare\square$}%
  \else\ifx\qcm@type\qcm@corollary
    \renewcommand{\qedsymbol}{$\blacksquare\square$}%
  \else
    \renewcommand{\qedsymbol}{$\blacksquare$}%
  \fi\fi
}{\popQED\end{namedinner}}
\makeatother
\renewcommand{\descriptionlabel}[1]{\hspace\labelsep\upshape\bfseries #1}
\begin{document}
\title[Massless photons and Maxwell infrared limits]{Massless photons and Maxwell infrared limits\newline in three-dimensional QED}
\date{}
\begin{abstract}
We prove further properties of the continuum curvature laws constructed
from noncompact lattice QED with chamber fermions. At any fixed charge,
assume that the positive lattice masses are eventually at most one and
bounded below by a fixed positive multiple of the spacing. A meson
correlator estimate then gives a positive massless photon weight in
almost every translation-ergodic component of every volume-first limit
point. The full infrared law of each such component converges to a free
Maxwell law. The previously established Osterwalder--Schrader axioms
already yield a local relativistic field; we give a direct
realization with quantitative energy and analytic-continuation bounds.
We verify the mass hypothesis for the tuned construction at sufficiently
small charge. Combined with its previously proved positive fourth
cumulant, this gives non-Gaussian component theories with nondegenerate
Maxwell infrared limits.
\end{abstract}
\maketitle
\tableofcontents

\section*{Introduction}

Let $P_{N,T}$ be the all-flux curvature laws of the chamber construction of
\cite{QEDII,QEDIII}, with standard weights, Wilson parameter $0<\rho<1$,
mass parameter $\bar m\in(0,1]$, spacing $\eta=L^{-N}$ and torus side
$8hL^T$, where $L$ is its fixed odd scale factor and
$h=L^{\lfloor\log_L\bar m^{-1}\rfloor}$.
A mass assignment is a real sequence $m_{0,N}$ independent of $T$;
its lattice mass is $\boldsymbol\mu_N=\bar m\eta+m_{0,N}$.
Write $\Pi_{\mathrm c}$ for the orthogonal projection of
$L^2(\mathbb R^3;\Lambda^2)$ onto the closed forms.
Existence of the limit points and the properties in clause~(i) below
are recalled from \cite[Theorems~1.3 and~1.4]{QEDII}.

\begin{named}{Theorem A}
Fix $e>0$ and a mass assignment for which, with a constant
$c_{\boldsymbol\mu}>0$,
\[
c_{\boldsymbol\mu}\bar mL^{-N}\le\boldsymbol\mu_N\le1
\qquad\text{for all sufficiently large }N.
\]
Volume-first limit points exist. Let $\mu$ be any such limit point:
$\mu=\lim_i\mu^{(N_i)}$ weakly on the Hermite--Sobolev space $H_{-4}$,
where $N_i\to\infty$ and each $\mu^{(N_i)}$ is a weak limit of
$P_{N_i,T}$ as $T\to\infty$.
Let $\mu=\int\mu_\xi\,d\nu(\xi)$ be its decomposition under translations
of $\mathbb R^3$. There is a measurable set of full $\nu$-measure on
which every law $\lambda=\mu_\xi$ has the following properties.

\begin{itemize}
\item[(i)] \emph{Euclidean field and OS axioms.}
The law $\lambda$ is a probability measure on the real tempered
distributional two-forms. Its field extends to a continuous linear map
$X:L^2(\mathbb R^3;\Lambda^2)\to L^2(\lambda)$ and obeys
\[
\lambda(\exp X(g))\le\exp\bigl(\tfrac12\|\Pi_{\mathrm c}g\|^2\bigr).
\]
It is concentrated on $dX=0$. Put $S_0=1$. For complex tensor Schwartz
tests $\Psi\in\mathcal S(\mathbb R^{3n};(\Lambda^2\mathbb C^3)^{\otimes n})$,
\[
|S_n^\xi(\Psi)|\le2\,66^n\sqrt{n!}\,q_{6n}(\Psi),\qquad
q_j(\Psi)=\max_{|\alpha|\le j}\sup_y(1+|y|^2)^{j/2}|\partial^\alpha\Psi(y)|,
\]
where the maximum includes tensor components.
Its Schwinger functions are tempered,
symmetric and Euclidean covariant under the two-form action
$((a,R)\cdot g)(y)=(\wedge^2R)g(R^{-1}(y-a))$;
the law's characteristic functional extends to an entire function of every
finite collection of test coefficients.
For every plane with reflection $r$ and either open half space $O$ bounded by that plane,
\[
\lambda(\overline{\widehat r A}\,A)\ge0
\qquad\bigl(A\in L^2(\Sigma_O,\lambda)\bigr),
\]
where $\Sigma_O$ is generated by $X(g)$ with
$g\in C_c^\infty(O;\Lambda^2)$ and
$\widehat r X(g)=X(r\cdot g)$.
For $A,B\in L^2(\lambda)$,
\[
\lambda(\bar A\,U(a)B)\longrightarrow
\overline{\lambda(A)}\lambda(B)\qquad(|a|\to\infty).
\]
These are the regularity, covariance, positivity, symmetry and cluster
axioms for the field strength; the tensor reflection is part of positivity.

\item[(ii)] \emph{Photon weight and Fisher information.}
The two-point function has a unique positive K\"all\'en--Lehmann measure
$\mathrm{kl}_\xi$ on $[0,\infty)$, with the integration variable denoting
mass, determined for real $g,g'\in L^2(\mathbb R^3;\Lambda^2)$ by
\[
S_2^\xi(g,g')=\int_0^\infty
\left\langle\Pi_{\mathrm c}g,\frac{-\Delta}{-\Delta+t^2}
\Pi_{\mathrm c}g'\right\rangle\,d\mathrm{kl}_\xi(t).
\]
Its total mass is at most one, and
$z_\xi:=\mathrm{kl}_\xi(\{0\})>0$.
There are measurable constants $C_\xi$ such that
\[
z_\xi^{-1}\le C_\xi<\infty,\qquad
\int C_\xi\,d\nu\le1+\tfrac12c_F,\qquad
\nu\{z_\xi\le s\}\le\frac{c_Fs}{2(1-s)}\quad(0<s<1),
\]
where
\[
c_F=15\pi^4[3(1+\rho)]^4C_\sharp(\rho)
\frac{e^2}{c_{\boldsymbol\mu}\bar m}
\max\left\{1,\frac{e^2}{c_{\boldsymbol\mu}\bar m}\right\}^3.
\]
Here $C_\sharp(\rho)<\infty$ is the constant of Theorem~5.14.
For every real closed $f\in L^2$,
\[
z_\xi\|f\|^2\le S_2^\xi(f,f)\le\|f\|^2,
\qquad I_\xi(f)\le C_\xi\|f\|^2.
\]
In this formula
$I_\xi(f)=\sup_\Phi[2\lambda(\partial_f\Phi)-\lambda(\Phi^2)]$,
where $\Phi=\phi(X(g_1),\dots,X(g_k))$, $g_j\in L^2$ real,
$\phi\in C_b^\infty(\mathbb R^k)$ (all derivatives bounded), and
$\partial_f\Phi=\sum_j\langle f,g_j\rangle\partial_j\phi$.

\item[(iii)] \emph{Relativistic reconstruction.}
Time is the first Euclidean coordinate. The reflection-positive
reconstruction has a Hilbert space $\mathcal H$, a unique vacuum
$\Omega$, commuting translation generators $(H,P_2,P_3)$ with
$\operatorname{spec}(H,\vec P)\subset\{\omega\ge|\vec p|\}$,
and a strongly continuous unitary representation $\mathcal U$ of
$\mathbb R^{1+2}\rtimes SO_0(1,2)$.
On the invariant dense domain $D=\bigcap_{j\ge0}D(H^j)$ there is a
tempered two-form field $F(h):D\to D$ which is symmetric for real
Schwartz tests, satisfies $dF=0$, is covariant under $\mathcal U$, and
commutes at spacelike separated supports for the metric
$(x^0)^2-(x^2)^2-(x^3)^2$.
Products of the field applied to $\Omega$ span a dense subspace.
Its vacuum functions analytically continue the Euclidean Schwinger
functions, with a factor $i$ for every time index.

Let $E$ be the joint spectral measure of $(H,\vec P)$ and put
$\partial V=\{\omega=|\vec p|\}$ and
$\mathcal H^{(1)}=\allowbreak \overline{\{F(h)\Omega:h\in\mathcal S(\mathbb R^{1+2};\Lambda^2\mathbb C^3)\}}$.
The subspace $\mathcal H_\gamma=E(\partial V)\mathcal H^{(1)}$ is
invariant and irreducible under $\mathcal U$. It is unitarily identified with
\[
L^2\left(\mathbb R^2,\frac{d^2k}{(2\pi)^2\,2|k|}\right),
\]
with $(H,\vec P)$ acting as $(|k|,-k)$. If $B_0(h)$ is the magnetic field
at time zero, as defined in Proposition~4.2, this identification sends
$E(\partial V)[B_0(h)]$ to $z_\xi^{1/2}|k|\hat h(k)$.
Thus $\operatorname{spec}H=[0,\infty)$ and the field produces a
massless particle of multiplicity one.

\item[(iv)] \emph{Infrared limit.}
For every real $g\in L^2(\mathbb R^3;\Lambda^2)$, with
$(\delta_lg)(y)=l^{-3/2}g(y/l)$,
\[
\lambda(\exp iX(\delta_lg))\longrightarrow
\exp\bigl(-\tfrac12z_\xi\|\Pi_{\mathrm c}g\|^2\bigr)
\qquad(l\to\infty).
\]
The limit is the free Maxwell law of weight $z_\xi$.
\end{itemize}
\end{named}

The additional conclusions concern the positive photon weight and
component Fisher bounds in~(ii), the identified massless subspace
in~(iii), and the infrared limit of the full law in~(iv). The local Wightman
field associated with~(i) is supplied by the Osterwalder--Schrader
reconstruction theorem \cite[Section~IV.1]{OS1975}, with the two-form
covariance and reflection conventions used here; see also
\cite[Section~6]{OS1973} for nonscalar fields. Section~3 gives a
direct realization on $D=\bigcap_mD(H^m)$, including the energy bound
of Theorem~3.3(a) and the tube bounds of Theorem~3.3(f). These estimates
also enter the proof of the infrared limit.

The cluster assertion concerns the components. The theorem gives an
integral bound on the reciprocals of their photon weights; it does not give a positive
lower bound independent of $\xi$. The massless subspace in (iii) is
the one generated by a single field applied to the vacuum.

The mass hypothesis holds for every assignment
$c_-\eta\le\boldsymbol\mu_N\le c_+\eta$ eventually, with
$0<c_-\le c_+<\infty$. We call such an assignment physical.
Here this term refers only to the two-sided bounds on the lattice mass.
We prove these bounds for the tuned shifts of \cite{QEDIII} in the
following range and combine them with that construction's
fourth-cumulant estimate.

\begin{named}{Corollary B}
Let $(\beta,x_0,r_\beta)$ be data of \cite[Theorem~44]{QEDIII},
let $f=dd^*\beta$ and $b(f)>0$ be its one-loop coefficient, and let
$e_2>0$ be the charge ceiling of that theorem. Use its tuned mass
assignment. Put
$e_7=(\bar m/(4C_\natural))^{1/2}$, where $C_\natural$ is defined in
Section~6. If $0<e\le\min\{e_2,e_7\}$, Theorem~A applies and, for
every volume-first limit point and almost every component,
\[
\kappa_4^\mu(f)\ge\tfrac12e^4b(f)>0,
\qquad \kappa_4^\xi(f)\ge\tfrac12e^4b(f)>0,
\qquad 0<z_\xi<1.
\]
In particular these component theories are not Gaussian, while their
infrared limits are nondegenerate free Maxwell laws. The tuned lattice
masses obey
$\frac12\bar m\eta\le\boldsymbol\mu_N\le2\bar m\eta$
for every $N\ge N_1(e)$, with $N_1(e)<\infty$.
\end{named}

\subsection*{The proof}
The preceding construction supplies the OS properties, but those
properties alone do not ensure a positive massless photon weight.
The estimate that yields this weight is a bound on the second spatial
moment of the meson correlator in one box of size comparable with the
inverse lattice mass. The determinant supplies a Ward identity, and the
identity turns this moment into a bound on the Fisher information of
curvature shifts. An eigenvalue count for the chamber operator removes
the logarithmic loss in an earlier estimate. With a lattice mass bounded
below by a positive multiple of the spacing, the Fisher bound survives
the limit.

A bound for the limit law must still be transferred to its components.
Invariant functions behave as constants under compactly supported
field shifts. This gives a common score after disintegration. Large
dilates of one test then isolate the atom at zero in the spectral measure.
To obtain a constant valid at all tests in almost every component, we
approximate the indicators of invariant partition cells by weights at
the cutoffs. The lattice clustering estimate controls their ergodic
averages uniformly through the volume and spacing limits. The weighted
inequality survives those limits; taking the supremum over finite
partitions controls the mean of the component constants. This gives
the integral estimate in Theorem~A(ii).

The direct reconstruction in Sections~2 and~3 provides the quantitative
realization used below. Products of the curvature field generate the
probabilistic OS Hilbert space. Bounded insertions in a time slab give
holomorphic point-field vectors with a bound near the boundary of the
forward tube. Their boundary values define the Minkowski field on
$D$ with the energy estimate in Theorem~3.3(a). The construction also
identifies the massless subspace produced by one field.

Finally, dilating an ergodic law sends its spectral measure to its atom
at zero. To identify the law, rather than just its two-point function,
we prove that an ergodic law in this class with a purely massless
spectral measure is Gaussian. Reflection positivity makes the higher
cumulant kernels curl free away from the diagonals. Their harmonic
singularities are removed by the $L^2$ bounds, closedness and clustering.
The contact-term argument then identifies the full law. Applied to
subsequential limits of dilates, this proves (iv).

The general photon theorem uses the mass hypothesis in Theorem~A and
holds at every charge. For Corollary~B, Section~6 verifies that
hypothesis for the tuned masses. The positive fourth cumulant is
imported from \cite[Theorem~44]{QEDIII}; the corollary establishes that
the same non-Gaussian components have the photon and infrared
properties of Theorem~A.

\section*{Notation and the preceding construction}
All Euclidean test forms are two-forms on $\mathbb R^3$.
The real $L^2$ norm is denoted by $\|\cdot\|$, and complexification is
understood when taking moments as multilinear forms.
The Hodge star identifies two-forms with vector fields; under this
identification closed forms are divergence free and
$\Pi_{\mathrm c}$ has Fourier multiplier $\Pi(p)=1-pp^T/|p|^2$.
We use $\hat f(p)=\int e^{-ip\cdot y}f(y)\,dy$.
The law of weight $0\le a\le1$ with characteristic functional
$\exp(-a\|\Pi_{\mathrm c}g\|^2/2)$ is denoted by $\gamma_a$;
$\gamma_0$ is the point mass at zero, and $\gamma=\gamma_1$.

Write $(\mathcal S',\Sigma)$ for the space of tempered distributional
two-forms with its evaluation $\sigma$-algebra, and $X(g)(\omega)=\omega(g)$.
For a law $\lambda$, write
$\Phi_\lambda(g)=\lambda(\exp iX(g))$ for its characteristic functional.
For a law satisfying Gaussian domination, denoted by (G), the field
extends continuously to $L^2$. Condition (dF) means concentration on
closed forms. We write $S_n$ for moments and $\kappa_n$ for cumulants;
in particular $\kappa_4(f)=S_4(f,f,f,f)-3S_2(f,f)^2$.
The translation of a test is $(\tau_ag)(y)=g(y-a)$, and
$U(a)X(g)=X(\tau_ag)$. The action of a Euclidean map on tests is the
two-form action in Theorem~A. Its induced operator on random variables
is denoted by a hat. The local $\sigma$-algebra $\Sigma_O$ is generated by
the smooth tests supported in the open set $O$.

The notation (H) and $\mathcal C_9$ refer to the structural hypotheses
of \cite[Definition~2.38]{QEDII}: Gaussian domination, translation and
cubic invariance, and reflection invariance and positivity at the nine
families of mirrors of the chamber construction. We set
\[
\mathcal K=\{\lambda:\lambda\text{ satisfies (H), (dF), and belongs to }\mathcal C_9\}.
\]
Every law in $\mathcal K$ is Euclidean invariant, reflection positive
at every plane, and has the round spectral cone and the
K\"all\'en--Lehmann representation of Theorem~A(ii)
\cite[Theorem~1.4, Part~I]{QEDII}.
Every volume-first limit point and almost every component belong to
$\mathcal K$, at every charge and for every mass assignment
\cite[Theorem~1.3]{QEDII}. An ergodic law in this class is mixing, with
the local clustering estimate of \cite[Theorem~12.39]{QEDII}.

For these laws the reflection-positive reconstruction at $y_1=0$ is
$(\mathcal H,[\,\cdot\,],H,P_2,P_3,\Omega)$, where
$\langle[A],[B]\rangle=\lambda(\overline{\widehat\theta_1A}\,B)$,
$e^{-tH}[A]=[U(te_1)A]$ for $t\ge0$, and spatial translations are
$e^{i\vec a\cdot\vec P}$. In a picture with time axis $e_j$ we use
subscripts $j$. The symbols $(T_{123})$, $(R_j)$ and $(RP_j)$ mean,
respectively, translation invariance, reflection invariance at $y_j=0$,
and positivity there. The one-plane hypotheses for $(j,i)$ include
(G), (dF), these invariances and positivity, the quarter turn in the
$(j,i)$ plane, and the spectral bound $H_j\ge|P_i|$
\cite[Definition~14.2]{QEDII}.

A probe is a real $\beta\in C_c^\infty(\mathbb R^3;\Lambda^2)$, with
$a=d^*\beta$ and $f=dd^*\beta$; a quotient by $\|f\|$ always requires
$f\ne0$. In the tuned construction,
\[
b(f)=-6\operatorname{Tr}\bigl[(S_{\bar m}\,i\gamma\cdot a)^4\bigr],
\qquad S_{\bar m}=(\bar m+\gamma\cdot\partial)^{-1},
\]
where $\gamma_j=\operatorname{diag}(\sigma_j,-\sigma_j)$ and $\sigma_j$
are the Pauli matrices. The data of Corollary~B are precisely the probe,
separation point and support radius of \cite[Proposition~43.2; Theorem~44]{QEDIII}.
The set $\Xi_\mu$ of tuned components and the exponent $\sigma>0$ are
those of that theorem. Set $g_0=f$ and $g_\pm=f\pm\tau_{x_0}f$. We also write
$\mathrm{Ext}_{x_0}(f)=\allowbreak \kappa_4(f)-3\kappa_4(f,\allowbreak f,\allowbreak \tau_{x_0}f,\allowbreak \tau_{x_0}f)$.
The tuned shifts and all their constants are imported from that
construction; no choice of a new mass prescription is made here.

The sets $\mathcal V_N$ consist of the fixed-spacing volume limits.
All weak limits in the definition of Theorem~A use $H_{-4}$, the
Hermite--Sobolev space of \cite[Sections~2.4 and~7.1.5]{QEDI}.
The separate notation for the finite lattice below is that of
\cite[Sections~2 and~3]{QEDII}.


\section{Contact terms and separated supports}

\emph{Conventions.}
By a \emph{law} is meant a probability measure on $(\mathcal S',\Sigma)$ \cite[Notation~(n1)]{QEDIII}.
If a law has (G), the field $X$ is extended complex linearly to $L^2(\mathbb{R}^3;\Lambda^2)\otimes\mathbb{C}$; for such a law the moments $S_N$ and the joint cumulants $\kappa_N$ are those of \cite[Notation~(n5),~(n8)]{QEDIII}.
Since (G) entails $S_1=0$ \cite[Proposition~45.2\,(1),~(8)]{QEDIII}, it follows that $\kappa_2=S_2$ and that
\begin{multline*}
\kappa_4(f_1,\dots,f_4)=S_4(f_1,\dots,f_4)-S_2(f_1,f_2)\,S_2(f_3,f_4)\\
-S_2(f_1,f_3)\,S_2(f_2,f_4)-S_2(f_1,f_4)\,S_2(f_2,f_3).
\end{multline*}

As in the notation fixed above, a two-form $g$ is identified with the vector field $\ast g$.
For $b_1,\dots,b_N\in L^2(\mathbb{R}^3;\mathbb{C}^3)$ one accordingly writes $S_N(b_1,\dots,b_N)$ and $\kappa_N(b_1,\dots,b_N)$ in place of $S_N(\ast^{-1}b_1,\dots,\ast^{-1}b_N)$ and $\kappa_N(\ast^{-1}b_1,\dots,\ast^{-1}b_N)$.
The first of them is the form $S_N$ of \cite[Lemma~14.3\,(b)]{QEDII}, since $\mathbb{B}(b)=X(\ast^{-1}b)$ \cite[Definition~14.1]{QEDII}.

Condition (H) is that of \cite[Notation~(n3)]{QEDIII}.
The group $O_{\mathrm{cub}}$ which enters this condition is $\mathrm{O}(3)\cap\mathrm{GL}(3,\mathbb{Z})$ \cite[Section~I.1.2]{QEDIII}.
For $j\neq i$ in $\{1,2,3\}$, the rotations $R^{ji}_t$ are those of \cite[Definition~14.1]{QEDII}, and the reconstruction of a law in the direction $e_j$ is called the picture $e_j$ \cite[Section~1.2]{QEDII}.
The cone $(C_{ji})$ is the condition of \cite[Definition~14.2\,(c)]{QEDII} that, in the picture $e_j$, the joint spectral measure of $(\mathsf H_j,\mathsf P_i)$ be carried by the points at which the energy is not smaller than the modulus of the momentum along $e_i$; it is the inequality $\mathsf H_j\ge|\mathsf P_i|$ with the constant one.

The two-dimensional linear subspaces $E$ of $\mathbb{R}^3$ are called \emph{planes}.
Let $p_E$ be the orthogonal projection onto $E$, and let $E_0:=\mathbb{R}e_1+\mathbb{R}e_2$.
Compactly supported forms or vector fields $f_1,\dots,f_N$ are said to be \emph{$E$-separated} when no two of the sets $p_E(\operatorname{supp}f_i)$ intersect; in that case no two of their supports intersect either.
One says that an $N$-linear form \emph{vanishes at non-coinciding arguments} when its value is zero at all real $f_1,\dots,f_N\in C_c^\infty$ with pairwise disjoint supports, and that it \emph{vanishes at $E$-separated arguments} when its value is zero at all real $E$-separated $f_1,\dots,f_N\in C_c^\infty$.
Of these two properties the second is the weaker one.
\begin{named}{Lemma 1.1}
Let $N\ge1$, and let $D$ be a bounded $N$-linear form on $L^2(\mathbb{R}^3;\mathbb{C}^3)$ with the property that $D(b_1,\dots,b_N)=0$ as soon as some entry is equal to $\nabla u$ for a function $u\in C_c^\infty(\mathbb{R}^3;\mathbb{C})$.
Let $E$ be a plane, and let $r>0$.
If $D\neq0$, one can find real $E$-separated $b_1,\dots,b_N\in C_c^\infty(\mathbb{R}^3;\mathbb{R}^3)$, each of them with support in a ball of radius $r$, such that $D(b_1,\dots,b_N)\neq0$.
\end{named}

\begin{proof}
\begin{enumerate}
\item[1.] \emph{Reduction to $E_0$.}
Choose $g\in SO(3)$ with $gE_0=E$, and put $(g\cdot b)(x):=g\,b(g^{-1}x)$ \cite[Definition~14.1]{QEDII}.
The map $b\mapsto g\cdot b$ is an isometry of $L^2(\mathbb{R}^3;\mathbb{C}^3)$; it preserves real fields and the space $C_c^\infty$, it maps every ball onto a ball of the same radius, it maps $\nabla u$ to $\nabla(u\circ g^{-1})$, and it satisfies
\begin{equation*}
p_E(\operatorname{supp}g\cdot b)=g\,p_{E_0}(\operatorname{supp}b).
\end{equation*}
Consequently the hypotheses of the lemma are fulfilled by the form $(b_i)\mapsto D(g\cdot b_1,\dots,g\cdot b_N)$, and the conclusion for this form and the plane $E_0$ coincides with the conclusion for $D$ and the plane $E$.
One may therefore assume $E=E_0$.
\item[2.]
One applies \cite[Proposition~14.13]{QEDII} with $U=\mathbb{R}^2$, so that $L^2_U=L^2(\mathbb{R}^3;\mathbb{C}^3)$.
The hypothesis (b) of that proposition is satisfied and its conclusion $D=0$ is false, so that its hypothesis (a) fails: there exist pairwise disjoint compact sets $K_1,\dots,K_N\subset\mathbb{R}^2$ and fields $b_i\in L^2(\mathbb{R}^3;\mathbb{C}^3)$ such that $b_i=0$ almost everywhere outside $K_i\times\mathbb{R}$ and $D(b_1,\dots,b_N)\neq0$.
\item[3.] \emph{Real fields.}
Write $b_i=\operatorname{Re}b_i+i\operatorname{Im}b_i$ and expand $D(b_1,\dots,b_N)$ by $N$-linearity.
At least one of the resulting $2^N$ terms differs from zero, and each entry of that term vanishes, as the corresponding $b_i$ does, outside $K_i\times\mathbb{R}$.
Real fields $b_i$ may therefore be used.
\item[4.] \emph{Smooth fields.}
Choose $\delta>0$ so that the closed $\delta$-neighbourhoods $V_i$ of the $K_i$ in $\mathbb R^2$ are pairwise disjoint, and take a mollifier $\psi_\varepsilon$ supported in the ball of radius $\varepsilon<\delta$.
The fields
\begin{equation*}
\psi_\varepsilon*\bigl(1_{\{|x_3|\le1/\varepsilon\}}\,b_i\bigr)
\end{equation*}
are real and smooth and have compact supports contained in $V_i\times\mathbb{R}$; hence they are $E_0$-separated.
They tend to $b_i$ in $L^2$ as $\varepsilon\to0$.
The continuity of $D$ implies that its value at these fields is not zero for all sufficiently small $\varepsilon$.
\item[5.] \emph{Small supports.}
Choose a smooth partition of unity $(\varphi_k)$ of $\mathbb{R}^3$ which is locally finite and in which every $\operatorname{supp}\varphi_k$ is contained in a ball of radius $r$.
The sum $b_i=\sum_k\varphi_kb_i$ has only finitely many non-zero terms, and expanding by $N$-linearity shows that at least one of the numbers $D(\varphi_{k_1}b_1,\dots,\varphi_{k_N}b_N)$ differs from zero.
\end{enumerate}
\end{proof}

\begin{named}{Proposition 1.2}
Consider a law $\mathrm{lw}$ with (G) and (dF) and a plane $E$.
\begin{itemize}
\item[(a)] Let $N\ge2$, and let $\kappa_N(g_1,\dots,g_N)\neq0$ for some $g_1,\dots,g_N\in L^2(\mathbb{R}^3;\Lambda^2)\otimes\mathbb{C}$.
Then for each $r>0$ one can find real $E$-separated $f_1,\dots,f_N\in C_c^\infty(\mathbb{R}^3;\Lambda^2)$, each of them supported in some ball of radius $r$, such that $\kappa_N(f_1,\dots,f_N)\neq0$.
The same holds with $S_N$ in place of $\kappa_N$.
\item[(b)] The three conditions below are mutually equivalent:
\begin{itemize}
\item[(i)] the law $\mathrm{lw}$ is Gaussian;
\item[(ii)] $\kappa_N$ vanishes at non-coinciding arguments whenever $N\ge3$;
\item[(iii)] $\kappa_N$ vanishes at $E$-separated arguments whenever $N\ge3$.
\end{itemize}
\item[(c)] Suppose that $\mathrm{lw}'$, a second law with (G) and (dF), is such that, whatever $N\ge2$, the form $S_N^{\mathrm{lw}'}-S_N^{\mathrm{lw}}$ vanishes at $E$-separated arguments.
Then $\mathrm{lw}'=\mathrm{lw}$.
Thus the Schwinger functions at non-coinciding arguments of a law having (G) and (dF) determine that law.
\end{itemize}
\end{named}

\begin{proof}
\begin{enumerate}
\item[1.] \emph{The form.}
For $b_1,\dots,b_N\in L^2(\mathbb{R}^3;\mathbb{C}^3)$ put $D(b_1,\dots,b_N):=\kappa_N(\ast^{-1}b_1,\dots,\ast^{-1}b_N)$.
That this form is $N$-linear and symmetric follows from the display of \cite[Notation~(n5)]{QEDIII}, which is valid also for complex entries \cite[Notation~(n8); Chapter~45-1, Part~F, Lemma~F.1]{QEDIII}.
It is bounded: complex entries included, one has
\begin{equation*}
|D(b_1,\dots,b_N)|\le(\log2)\,\Bigl(\frac{5\exp(1)}{2}\Bigr)^{N}N!\,\prod_{i=1}^{N}\|b_i\|
\end{equation*}
as follows from \cite[Proposition~7.10\,(7)]{QEDI}, applicable to all laws with (G) by \cite[Proposition~45.1\,(1),~(11)]{QEDIII}, combined with the pointwise isometry property of $\ast$ \cite[Definition~14.1]{QEDII}.
\item[2.] \emph{Gradients.}
Let $u\in C_c^\infty(\mathbb{R}^3;\mathbb{C})$; condition (dF) gives $X(\ast^{-1}\nabla u)=0$ almost surely \cite[Definition~14.2\,(a)]{QEDII}, \cite[Proposition~45.1\,(10)]{QEDIII}.
Suppose that $b_\ell=\nabla u$; then each product in the display of \cite[Notation~(n5)]{QEDIII} has among its factors the moment of the block containing $\ell$, and the integrand of that moment is zero almost surely.
Hence $D=0$ at such entries.
\item[3.]
Part~(a) is Lemma~1.1 applied to $D$, with $f_i:=\ast^{-1}b_i$.
In the case of $S_N$, the form is bounded by \cite[Lemma~14.3\,(b)]{QEDII}, and its vanishing on gradients is seen as in step~2.
\item[4.] \emph{Part (b), (i) implies (ii).}
Let $\mathrm{lw}$ be Gaussian \cite[Notation~(n6)]{QEDIII}.
Then every finite family $X(f_1),\dots,X(f_N)$ with real $f_i\in\mathcal S$ is jointly normal \cite[Section~2.5]{QEDI} and centred.
Therefore
\begin{equation*}
\log\mathrm{lw}\Bigl(\exp\sum_az_aX(f_a)\Bigr)=\frac12\sum_{a,c}z_az_c\,S_2(f_a,f_c),
\end{equation*}
and the derivative $\partial_{z_1}\cdots\partial_{z_N}$ of this function at $0$, which is $\kappa_N(f_1,\dots,f_N)$ by \cite[Notation~(n5)]{QEDIII}, vanishes for $N\ge3$.
That (ii) implies (iii) follows from the fact that $E$-separated arguments have pairwise disjoint supports.
\item[5.] \emph{Part (b), (iii) implies (i).}
Lemma~1.1, together with steps 1 and~2, shows that $\kappa_N=0$ identically for each $N\ge3$.
For real $f\in\mathcal S$ the function $w(z):=\log\mathrm{lw}(\exp zX(f))$ is analytic for $|z|\cdot\|f\|\le\tfrac25$, with $w^{(n)}(0)=\kappa_n(f,\dots,f)$; this is \cite[Proposition~7.10\,(7)]{QEDI} with $f_1=\dots=f_n=f$, the function $W$ of that proposition being $W(z)=w(z_1+\dots+z_n)$.
It follows that $w(z)=\tfrac12S_2(f,f)z^2$ for these $z$.
Both sides of the identity
\begin{equation*}
\mathrm{lw}\bigl(\exp zX(f)\bigr)=\exp\bigl(\tfrac12S_2(f,f)z^2\bigr)
\end{equation*}
are entire functions of $z$ \cite[Proposition~7.10\,(1)]{QEDI}, so that the identity holds at $z=it$, and $X(f)$ is normal.
\item[6.] \emph{Part (c).}
Lemma~1.1 applies to the form $S_N^{\mathrm{lw}'}-S_N^{\mathrm{lw}}$, which is bounded \cite[Lemma~14.3\,(b)]{QEDII} and zero on gradients, and shows that this form vanishes identically.
Thus all moments of $\mathrm{lw}'$ and of $\mathrm{lw}$ at real tests in $C_c^\infty$ coincide, the products being $\mathrm{lw}'$-integrable by (G), and \cite[Proposition~45.1\,(8),~(11)]{QEDIII} gives $\mathrm{lw}'=\mathrm{lw}$.
\end{enumerate}
\end{proof}


\section{Cyclicity of the curvature field}

In this section a \emph{law} is a probability measure on $(\mathcal S',\Sigma)$, where $\mathcal S'=\mathcal S'(\mathbb R^3;\Lambda^2)$.
Let $j\ne i$ be two elements of $\{1,2,3\}$.
Following \cite[Definition~14.2\,(f)]{QEDII}, the \emph{one-plane hypotheses for the ordered pair $(j,i)$} on a law $\mu$ are the seven properties listed below:
\begin{itemize}
\item (G) and (dF);
\item $(T_{123})$: the law is invariant under every translation of $\mathbb R^3$;
\item $(R_j)$ and $(RP_j)$: the law is invariant under the reflection $\theta_j$ in the plane $\{x_j=0\}$, and it is reflection positive with respect to this plane;
\item $(Q_{ji})$: the law is invariant under the quarter turn $e_j\mapsto e_i\mapsto-e_j$;
\item $(C_{ji})$: $\mathfrak E_j(\{\omega<|p_i|\})=0$; equivalently, the inequality $\mathsf H_j\ge|\mathsf P_i|$ holds with the constant equal to one.
\end{itemize}
The quintuple $(\mathcal H_j,[\,\cdot\,]_j,\mathsf H_j,V_j,\Omega_j)$ denotes the reconstruction of $\mu$ in the direction $e_j$ \cite[Definition~2.38]{QEDII}.
The space $\mathcal H_j$ is the completion of $L^2(\Sigma_j^+,\mu)$ modulo null vectors, where $\Sigma_j^+=\Sigma_{\{x_j>0\}}$, for the form
\begin{equation*}
\langle[\Phi]_j,[\Phi']_j\rangle=\mu\big((\Theta_j\Phi)\cdot\Phi'\big),\qquad\Theta_j\Phi=\overline{\hat\theta_j\Phi};
\end{equation*}
the operators $\mathsf H_j$ and $V_j(y)$ and the vector $\Omega_j$ are given by
\begin{equation*}
e^{-t\mathsf H_j}[\Phi]_j=[U(te_j)\Phi]_j\quad(t\ge0),\qquad V_j(y)[\Phi]_j=[U(y)\Phi]_j\quad(y\perp e_j),\qquad\Omega_j=[1]_j.
\end{equation*}
For $i\ne j$, let $\mathsf P_i$ denote the generator of $u\mapsto V_j(ue_i)$; one writes $\mathfrak E_j$ for the joint spectral measure of $(\mathsf H_j,(\mathsf P_i)_{i\ne j})$, which is a measure on the space $\mathbb R\times e_j^\perp$ whose points are written $(\omega,p_\perp)$ \cite[Lemma~12.18\,(b)]{QEDII}.
For the hypothesis (H) and the class $\mathcal C_9$ one takes the definitions of \cite[Definition~2.38]{QEDII}, which agree with those of \cite{QEDIII}.

The statements below use the following notation of \cite{QEDII}.
By a \emph{label} one means a pair $\lambda=(\sigma,\chi)$ with $\sigma\in\{+,-,3\}$ and $\chi\in L^2(\mathbb R;\mathbb R)$.
One sets $w(+)=1$, $w(-)=-1$, $w(3)=0$, and $w(\lambda)=\sum_aw(\sigma_a)$ for a tuple $\lambda$ of labels $\lambda_a=(\sigma_a,\chi_a)$ \cite[Definition~14.5]{QEDII}.
The forward tube is
\begin{equation*}
\mathcal W=\{\xi\in\mathbb C^2:\operatorname{Re}\xi^0>|\operatorname{Im}\xi^1|\},
\end{equation*}
with closure $\overline{\mathcal W}$, and one puts
\begin{equation*}
\mathcal W^{(n)}=\{x\in(\mathbb C^2)^n:x_1\in\mathcal W,\ x_{a+1}-x_a\in\mathcal W\}.
\end{equation*}
For real $\beta$, let $\Lambda_\beta$ be the linear map of $\mathbb C^2$ with
\begin{equation*}
(\Lambda_\beta\xi)^0=\xi^0\cosh\beta-i\xi^1\sinh\beta,\qquad(\Lambda_\beta\xi)^1=\xi^1\cosh\beta+i\xi^0\sinh\beta;
\end{equation*}
on configurations it acts entry by entry \cite[Lemma~14.9\,(c)]{QEDII}.
Under $(C_{12})$, the operators $\mathsf T(\kappa)$ of \cite[Lemma~14.17]{QEDII} are given for $\kappa\in\overline{\mathcal W}$ by
\begin{equation*}
\mathsf T(\kappa)=\int d\mathfrak E_1(\omega,p_\perp)\,e^{-\kappa^0\omega+i\kappa^1p_2}.
\end{equation*}
Throughout the section, $\beta$ and $\sigma$ have these meanings; probes, and the exponent $\sigma$ of \cite{QEDIII}, do not occur.

\begin{named}{Theorem 2.1}
Let the law $\mu$ fulfil the one-plane hypotheses for the pair $(1,2)$.
\begin{itemize}
\item[(a)] The linear span of $\Omega_1$ and of the vectors
\begin{equation*}
[X(g_1)\cdots X(g_N)]_1,
\end{equation*}
where $N\ge1$ and $g_1,\dots,g_N\in C_c^\infty(\{x_1>0\};\Lambda^2)$ are real and satisfy
\begin{equation*}
\sup\{x_1:x\in\operatorname{supp}g_a\}<\inf\{x_1:x\in\operatorname{supp}g_{a+1}\}\qquad(1\le a<N),
\end{equation*}
is dense in $\mathcal H_1$.
\item[(b)] One has $\mathcal H^{\Psi}_{12}=\mathcal H_1$, where $\mathcal H^{\Psi}_{12}$ denotes the closed span of the plate vectors $\Psi(x;\lambda)$ with $n\ge0$, $x\in\mathcal W^{(n)}$ and $\lambda$ a tuple of $n$ labels \cite[Definition~14.21; Proposition~14.47\,(c)]{QEDII}.
\item[(c)] The vector $\Omega_1$ is cyclic for the field strength in two forms, one for each of the two families described next.
The first family consists of the point plates $\mathcal B_\lambda(\xi,\xi')$, $\xi,\xi'\in\mathcal W$, each of which stands for $\mathsf T(\xi)\,\mathbb B_\lambda(\delta_0)\,\mathsf T(\xi')$ \cite[Proposition~14.18]{QEDII}.
The second family consists of the operators $\mathcal B_l(f)$, where $l>0$ and where $f\in L^2(\mathbb R^3;\mathbb C^3)$ vanishes outside $\{0\le x_1\le l\}$ \cite[Proposition~14.16]{QEDII}.
For each of the two families, $\Omega_1$ and the vectors obtained by applying finite products of members of the family to $\Omega_1$ form a total set in $\mathcal H_1$.
\end{itemize}
\end{named}

\begin{named}{Corollary 2.2}
\begin{itemize}
\item[(a)] Let the law $\mu$ satisfy the hypotheses of Theorem~2.1.
There is a unique strongly continuous unitary group $(\mathfrak L(\beta))_{\beta\in\mathbb R}$, defined on all of $\mathcal H_1$, such that $\mathfrak L(\beta)\Psi(x;\lambda)=\allowbreak e^{w(\lambda)\beta}\,\allowbreak \Psi(\Lambda_\beta x;\lambda)$ for every plate vector.
For $\beta,u\in\mathbb R$, $t\ge0$, $\kappa\in\overline{\mathcal W}$, $\xi,\xi'\in\mathcal W$ and every label $\lambda_0=(\sigma,\chi)$, the following relations hold on $\mathcal H_1$:
\begin{gather*}
\mathfrak L(\beta)\Omega_1=\Omega_1,\qquad\mathfrak L(\beta)\,\mathsf T(\kappa)\,\mathfrak L(\beta)^{-1}=\mathsf T(\Lambda_\beta\kappa),\\
\mathfrak L(\beta)\,\mathcal B_{\lambda_0}(\xi,\xi')\,\mathfrak L(\beta)^{-1}=e^{w(\sigma)\beta}\,\mathcal B_{\lambda_0}(\Lambda_\beta\xi,\Lambda_\beta\xi'),
\end{gather*}
and
\begin{align}
\mathfrak L(\beta)\,e^{-t\mathsf H_1}\,\mathfrak L(\beta)^{-1}&=e^{-t(\mathsf H_1\cosh\beta+\mathsf P_2\sinh\beta)},\tag{2.1}\\
\mathfrak L(\beta)\,e^{iu\mathsf P_2}\,\mathfrak L(\beta)^{-1}&=e^{iu(\mathsf P_2\cosh\beta+\mathsf H_1\sinh\beta)},\tag{2.2}
\end{align}
the right-hand sides of (2.1) and (2.2) being integrals against $\mathfrak E_1$.
\item[(b)] Under the assumptions of (a), the operators $\mathfrak L(\beta)$ and $V_1(ue_3)$ commute for every $u\in\mathbb R$, and one has
\begin{equation*}
\mathfrak L(\beta)\,\mathfrak E_1(A)\,\mathfrak L(\beta)^{-1}=\mathfrak E_1(\ell_{-\beta}A)
\end{equation*}
for every Borel set $A\subset\mathbb R\times e_1^\perp$, where
\begin{equation*}
\ell_\beta(\omega,p_2,p_3):=(\omega\cosh\beta+p_2\sinh\beta,\ p_2\cosh\beta+\omega\sinh\beta,\ p_3).
\end{equation*}
\item[(c)] Assume that the law $\mu$ belongs to $\mathcal C_9$ and has (H) and (dF).
Then, for every ordered pair $j\ne i$, the law $\mu$ satisfies the one-plane hypotheses for $(j,i)$, one has $\mathcal H^{\Psi}_{ji}=\mathcal H_j$, and the group $\mathfrak L_{ji}$ of \cite[Theorem~1.4, Part~I\,(c)]{QEDII} acts on all of $\mathcal H_j$ and satisfies the relations of (a) and (b) for $(\mathsf H_j,\mathsf P_i)$.
Furthermore, let $\mathbb R\times e_j^\perp$ carry the quadratic form $\omega^2-|p_\perp|^2$; then to each $\Lambda$ in its restricted Lorentz group there corresponds a unitary operator $W$ on $\mathcal H_j$ satisfying $W\Omega_j=\Omega_j$, as well as $W\,\mathfrak E_j(A)\,W^{-1}=\mathfrak E_j(\Lambda A)$ for each Borel set $A$.
In particular, $\mathfrak E_j$ has a Lorentz invariant support, which is a closed subset of the cone $\{\omega\ge|p_\perp|\}$.
\item[(d)] The hypotheses of (c) are satisfied in the following cases:
\begin{itemize}
\item for every charge $e>0$ and every mass assignment: by every volume-first limit point $\mu$ of \cite[Theorem~1.3]{QEDII}, and by $\nu$-almost every component $\mu_\xi$;
\item for every charge $e\in(0,e_2]$, $e_2$ being the ceiling of \cite[Theorem~44]{QEDIII}, and the laws $P_{N,T}$ of that theorem: by every volume-first limit point $\mu$, and by $\mu_\xi$ at every $\xi\in\Xi_\mu$.
\end{itemize}
\end{itemize}
\end{named}

\begin{named}{Corollary 2.3}
Let $\mu$ be as in Theorem~2.1, and suppose that $\ker\mathsf H_1=\mathbb C\,\Omega_1$.
Then the only bounded operators on $\mathcal H_1$ that commute with all point plates $\mathcal B_\lambda(\xi,\xi')$ are the multiples of the identity.
Every $\mu_\xi$ with $\xi\in\Xi_\mu$ falls under the corollary \cite[Theorem~44\,(D1)--(D3)]{QEDIII}, and so does $\nu$-almost every component in \cite[Theorem~1.4, Part~I\,(e)]{QEDII}.
\end{named}

\begin{proof}[Proof of Theorem~2.1]
The notation is that of \cite[Definitions~14.1 and~14.5]{QEDII}.
For $b=\ast g$ one writes $\mathbb B(b)=X(g)$, and $\mathbb B$ is extended complex linearly to $b\in L^2(\mathbb R^3;\mathbb C^3)$.
Points of $\mathbb R^3$ are written $x=(y,x_3)$ with $y=(y^0,y^1)=(x_1,x_2)$, and $\mathbb H=\{y\in\mathbb R^2:y^0>0\}$.
One writes $L^2_{\mathbb H}$ for the subspace formed by the $b$ that are equal to zero almost everywhere outside $\mathbb H\times\mathbb R$.
For a label $\lambda=(\sigma,\chi)$ and $\alpha\in L^2(\mathbb R^2;\mathbb C)$ one puts $(\alpha\otimes\lambda)(x)=\alpha(y)\chi(x_3)e^\sigma$ and $\mathbb B_\lambda(\alpha)=\mathbb B(\alpha\otimes\lambda)$.
The closed linear span of the vectors listed in (a) is denoted by $\mathcal H^{<}$.

\emph{Step 1 (a bound).}
Let $b_1,\dots,b_N\in L^2_{\mathbb H}$.
The product $\prod_a\mathbb B(b_a)$ belongs to $L^2(\Sigma_1^+,\mu)$, and
\begin{equation*}
\Big\|\Big[\prod_a\mathbb B(b_a)\Big]_1\Big\|\le\Big\|\prod_a\mathbb B(b_a)\Big\|_{L^2(\mu)}\le c_{2N}^N\prod_a\|b_a\|.
\end{equation*}
This follows from \cite[Lemma~14.3\,(a)]{QEDII} and H\"older's inequality, together with the fact that the class map $\Phi\mapsto[\Phi]_1$ is a contraction \cite[Proposition~8.19\,(i)]{QEDI}.
Consequently, the class of a product depends $N$-linearly and continuously on its entries.

\emph{Step 2 (plates in time order).}
\emph{2a.} Consider $\alpha_1,\dots,\alpha_N\in L^2(\mathbb R^2;\mathbb C)$ whose supports are time-ordered in $\mathbb H$ \cite[Definition~14.21]{QEDII}.
Cutting the plates $\alpha_a\otimes\lambda_a$ off in $x_3$, mollifying at a scale smaller than the gaps between the time supports and than their distance from $\{y^0=0\}$, and splitting into real and imaginary parts, one exhibits each plate as an $L^2$-limit of fields $\ast g+i\ast g'$ in which $g,g'$ are real, smooth and compactly supported, the time order being preserved.
Step 1 then yields $[\prod_a\mathbb B_{\lambda_a}(\alpha_a)]_1\in\mathcal H^{<}$.

\emph{2b.} Conversely, $\mathcal H^{<}\subset\mathcal H^{\Psi}_{12}$.
This inclusion is used for (b), in Step 6, and the proof of (a) does not rely on it.
Take $g_a$ as in (a) and pick compact rectangles $A_a\subset\mathbb H$ in time order so that $\ast g_a$ vanishes outside $A_a\times\mathbb R$; by \cite[Lemma~14.8\,(a)]{QEDII}, each $\ast g_a$ is then an $L^2$-limit of finite sums of plates $\alpha\otimes\lambda$ in which $\alpha=0$ outside $A_a$.
By \cite[Proposition~14.47\,(c)]{QEDII}, the closed space $\mathcal H^{\Psi}_{12}$ contains the class of every product of such plates, and Step 1 gives the inclusion.

\emph{Step 3 (entries over disjoint compact subsets of the plane).}
\emph{Claim.} Assume that the sets $K_1,\dots,K_N\subset\mathbb H$ are compact and pairwise disjoint, that each $\alpha_a\in L^2(\mathbb R^2;\mathbb C)$ is equal to zero outside $K_a$, and that $\lambda=(\lambda_1,\dots,\lambda_N)$ is a tuple of labels.
Then $[\prod_a\mathbb B_{\lambda_a}(\alpha_a)]_1\in\mathcal H^{<}$.

\emph{3a (decomposition).} Given $d>0$ and $k\in\mathbb Z$, set $\alpha_a^{k}:=\alpha_a\cdot1_{\{kd\le y^0<(k+1)d\}}$; for $k=(k_1,\dots,k_N)\in\mathbb Z^N$, set $\Phi_k:=\prod_a\mathbb B_{\lambda_a}(\alpha_a^{k_a})$.
Then $\prod_a\mathbb B_{\lambda_a}(\alpha_a)=\sum_k\Phi_k$, and the sum is finite.

\emph{3b (distinct intervals).} For pairwise distinct integers $k_a$, one has $[\Phi_k]_1\in\mathcal H^{<}$.
Indeed, let the strips be replaced by $\{kd+\varepsilon\le y^0\le(k+1)d-\varepsilon\}$; since the factors commute, they may be taken in the order of increasing $k_a$, and the resulting functions then have time-ordered supports in $\mathbb H$, so that Step 2 applies.
Letting $\varepsilon\to0$ and using Step 1, one obtains the assertion.

\emph{3c (the rest).} Denote by $\Phi^{\mathrm r}_d$ the sum of the other $\Phi_k$, and by $E_d$ the set of configurations $Y=(y_1,\dots,y_N)\in(\mathbb R^2)^N$ such that two of the times $y_a^0$ fall into a common interval $[kd,(k+1)d)$.
Besides being antilinear, the map $\Theta_1$ is multiplicative; furthermore, according to \cite[Lemma~14.6\,(b)]{QEDII}, one has almost surely
\begin{equation*}
\Theta_1\mathbb B_{\lambda_a}(\alpha)=s(\sigma_a)\,\mathbb B_{\lambda_a}(\overline\alpha\circ\theta_1).
\end{equation*}
It follows that $\|[\Phi^{\mathrm r}_d]_1\|^2=\allowbreak \mu\big((\Theta_1\Phi^{\mathrm r}_d)\cdot\Phi^{\mathrm r}_d\big)$ is $s(\lambda)=\pm1$ times a finite sum of moments $S_{2N}$ whose entries are $(\overline{\alpha_a^{k_a}}\circ\theta_1)\otimes\lambda_a$ and $\alpha_a^{k'_a}\otimes\lambda_a$, $1\le a\le N$.

\emph{3d (no singular part).} The supports of these $2N$ functions are compact subsets of the $2N$ pairwise disjoint sets $\theta_1K_1,\dots,\theta_1K_N,K_1,\dots,K_N$.
Hence \cite[Proposition~14.44]{QEDII}, applied with $2N$ in place of $N$, writes each moment as the integral, against the plate function $\tilde G_{2N}(\,\cdot\,;\lambda,\lambda)$, of the product of these functions.
The plate function is real analytic on the set of configurations of $2N$ distinct points \cite[Lemma~14.35\,(e)]{QEDII}; hence there is a constant $C$, independent of $d$, which bounds it on the compact set $\prod_a\theta_1K_a\times\prod_aK_a$.
Since
\begin{equation*}
\sum_k\prod_a\alpha_a^{k_a}(y_a)=1_{E_d}(Y)\prod_a\alpha_a(y_a),
\end{equation*}
in which $k$ runs through the tuples having a repetition, one obtains after summation
\begin{equation*}
\|[\Phi^{\mathrm r}_d]_1\|^2\le C\cdot\Big(\int_{E_d}dY\,\prod_a|\alpha_a(y_a)|\Big)^2.
\end{equation*}

\emph{3e (limit).} Here $\prod_a|\alpha_a(y_a)|$ is an integrable function of $Y$, and
\begin{equation*}
E_d\subset\{Y:|y_a^0-y_b^0|<d\ \text{for some}\ a\ne b\},
\end{equation*}
where the sets on the right-hand side decrease, as $d\to0$, to a set of measure zero.
The dominated convergence theorem therefore gives $[\Phi^{\mathrm r}_d]_1\to0$.
Together with 3a and 3b this proves the claim, because $\mathcal H^{<}$ is closed.

\emph{Step 4 (no contact terms).}
\emph{4a.} Let $\psi\perp\mathcal H^{<}$ and $N\ge1$, and put
\begin{equation*}
D(b_1,\dots,b_N):=\Big\langle\psi,\Big[\prod_a\mathbb B(b_a)\Big]_1\Big\rangle,\qquad b_a\in L^2_{\mathbb H}.
\end{equation*}
Step 1 shows that $D$ is $N$-linear and bounded.
If some $b_\ell$ equals $\nabla u$ with $u\in C_c^\infty(\mathbb H\times\mathbb R;\mathbb C)$, then $\mathbb B(b_\ell)=0$ almost surely by (dF) \cite[Definition~14.2\,(a)]{QEDII}, and $D$ vanishes.
The form $D$ also vanishes when each $b_a$ is zero outside $K_a\times\mathbb R$, the sets $K_a\subset\mathbb H$ being compact and pairwise disjoint: for plates this is Step 3, and the general case follows from \cite[Lemma~14.8\,(a)]{QEDII} with $A=K_a$, multilinearity and Step 1.

\emph{4b.} These two properties are the hypotheses of \cite[Proposition~14.13]{QEDII} for the open set $U=\mathbb H$.
Consequently $D=0$, for every $N\ge1$.

\emph{Step 5 (polynomials suffice).}
Let $g\in C_c^\infty(\{x_1>0\};\Lambda^2)$ be real.
Each partial sum of $\sum_n(iX(g))^n/n!$ is bounded in modulus by $\exp|X(g)|$, and (G), applied to $\pm2g$, shows that $\exp|X(g)|$ belongs to $L^2(\mu)$.
The partial sums therefore converge to $\exp(iX(g))$ in $L^2(\Sigma_1^+,\mu)$; combined with $\psi\perp\Omega_1$, Step 4 then yields $\langle\psi,[\exp(iX(g))]_1\rangle=0$.
According to \cite[Lemma~8.12\,(i)]{QEDI}, the linear hull of such exponentials is dense in $L^2(\Sigma_1^+,\mu)$, and according to \cite[Proposition~8.19\,(i)]{QEDI}, the class map has a dense image.
Thus $\psi=0$, and (a) is established.

\emph{Step 6 (assertion (b), and assertion (c) for the point plates).}
Combining (a) with 2b gives $\mathcal H_1=\allowbreak \mathcal H^{<}\subset\mathcal H^{\Psi}_{12}\subset\mathcal H_1$.
By \cite[Definition~14.21]{QEDII}, the plate vectors consist of $\Omega_1$ together with the vectors obtained from $\Omega_1$ by applying products of point plates.

\emph{Step 7 (assertion (c) for the operators $\mathcal B_l(f)$).}
With $g_a$ as in (a), write $b_a=\ast g_a$, set $l_0=0$, and pick $l_1<\dots<l_N$ such that each $b_a$ vanishes outside $\{l_{a-1}\le x_1\le l_a\}$.
Iterating the identity $\mathcal B_l(f)[\Phi]_1=[\mathbb B(f)\cdot U(le_1)\Phi]_1$ of \cite[Proposition~14.16]{QEDII}, one finds
\begin{equation*}
[X(g_1)\cdots X(g_N)]_1=\mathcal B_{l_1}(b_1)\,\mathcal B_{l_2-l_1}(\tau_{-l_1e_1}\cdot b_2)\cdots\mathcal B_{l_N-l_{N-1}}(\tau_{-l_{N-1}e_1}\cdot b_N)\,\Omega_1.
\end{equation*}
The assertion then follows from (a).
\end{proof}

\begin{proof}[Proof of Corollary~2.2]
\emph{(a)} The assertion follows from \cite[Proposition~14.47\,(c)]{QEDII}, since $\mathcal H^{\Psi}_{12}=\mathcal H_1$ by Theorem~2.1\,(b).

\emph{(b1) Smeared translations along $e_3$.}
For a tuple $\lambda$ of labels let $\lambda^u$ be the tuple with the profiles $\chi_a(\,\cdot-u)$; then $\tau_{ue_3}\cdot(\alpha\otimes\lambda_a)=\allowbreak \alpha\otimes\lambda_a^u$ and $w(\lambda^u)=w(\lambda)$.
For $\alpha_a$ whose supports are time-ordered in $\mathbb H$, it follows from the identities $V_1(y)[\Phi]_1=[U(y)\Phi]_1$ \cite[Definition~2.38]{QEDII} and $U(a)\mathbb B(b)=\mathbb B(\tau_a\cdot b)$ \cite[Lemma~14.3\,(a)]{QEDII} that
\begin{equation*}
V_1(ue_3)\Big[\prod_a\mathbb B_{\lambda_a}(\alpha_a)\Big]_1=\Big[\prod_a\mathbb B_{\lambda_a^u}(\alpha_a)\Big]_1.
\end{equation*}
Expressed by means of the second identity of \cite[Proposition~14.22\,(b)]{QEDII}, this equality becomes
\begin{equation*}
\int dY\,\prod_a\alpha_a(y_a)\,\big(V_1(ue_3)\Psi(Y;\lambda)-\Psi(Y;\lambda^u)\big)=0.
\end{equation*}

\emph{(b2) The pointwise identity, and the commutation.}
On the convex open set $\mathcal W^{(n)}$, the real points of which are the time-ordered configurations in $\mathbb H$, the vectors $V_1(ue_3)\Psi(x;\lambda)$ and $\Psi(x;\lambda^u)$ depend holomorphically on $x$ \cite[Proposition~14.22\,(c)]{QEDII}.
Applying \cite[Lemma~14.10\,(g), (a)]{QEDII} to the scalar products of these vectors with an arbitrary fixed vector, one therefore obtains $V_1(ue_3)\Psi(x;\lambda)=\Psi(x;\lambda^u)$ on $\mathcal W^{(n)}$.
One deduces
\begin{equation*}
\mathfrak L(\beta)V_1(ue_3)\Psi(x;\lambda)=e^{w(\lambda)\beta}\Psi(\Lambda_\beta x;\lambda^u)=V_1(ue_3)\mathfrak L(\beta)\Psi(x;\lambda),
\end{equation*}
and the set of plate vectors is total.

\emph{(b3) The spectral measure.}
Let $t\ge0$ and $y=u_2e_2+u_3e_3$, and put $F_{t,y}(\omega,p_\perp)=e^{-t\omega+ip_\perp\cdot y}$.
From the relations (2.1) and (2.2), from part (b2), and from the multiplicativity of integrals against $\mathfrak E_1$, one obtains
\begin{equation*}
\mathfrak L(\beta)\,e^{-t\mathsf H_1}V_1(y)\,\mathfrak L(\beta)^{-1}=\int d\mathfrak E_1\,F_{t,y}\circ\ell_\beta.
\end{equation*}
The measure $\mathfrak E_1$ is carried by the set $\{\omega\ge|p_2|\}$, and this set is invariant under the map $\ell_\beta$.
For each $\psi$, the finite positive measures
\begin{equation*}
\langle\psi,\mathfrak L(\beta)\,\mathfrak E_1(\cdot)\,\mathfrak L(\beta)^{-1}\psi\rangle\qquad\text{and}\qquad\langle\psi,\mathfrak E_1(\ell_\beta^{-1}\,\cdot\,)\psi\rangle
\end{equation*}
on $[0,\infty)\times\mathbb R^2$ therefore give rise to one and the same Laplace--Fourier transform.
By \cite[Lemma~12.25]{QEDII} the two measures are equal; moreover, $\ell_\beta^{-1}=\ell_{-\beta}$.

\emph{(c1) All planes.}
The conditions for membership in $\mathcal C_9$ include $(T_{123})$, $(R_j)$, $(RP_j)$, and invariance and reflection positivity at $\{x_j+x_i=0\}$; hence \cite[Theorem~12.21\,(c)]{QEDII} gives $(C_{ji})$.
By \cite[Definition~2.38]{QEDII}, the hypothesis (H) includes (G) and the invariance under $O_{\mathrm{cub}}$, a group which contains the quarter turn of $(Q_{ji})$ \cite[Definition~14.1]{QEDII}; thus $\mu$ has (G) and $(Q_{ji})$, and it has (dF) by assumption.
In view of \cite[Lemma~14.12\,(d)]{QEDII}, the one-plane hypotheses for $(1,2)$ hold for the law $\mu'=(g_{ji}^{-1})_*\mu$, and the map $J$ of that lemma is a unitary map of $\mathcal H_1(\mu')$ onto $\mathcal H_j(\mu)$ which carries $\mathsf H_1'$, $V_1'(y)$, $\mathfrak E_1'$ to $\mathsf H_j$, $V_j(g_{ji}y)$, $\mathfrak E_j$; explicitly,
\begin{equation*}
\mathfrak E_j(A)=J\,\mathfrak E_1'(g_{ji}^{-1}A)\,J^{-1},
\end{equation*}
where $g_{ji}$ also denotes the map $(\omega,p_\perp)\mapsto(\omega,g_{ji}p_\perp)$ of $\mathbb R\times e_1^\perp$ onto $\mathbb R\times e_j^\perp$.
In \cite[Lemma~14.12\,(d)]{QEDII} this relation is stated for Borel subsets $A$ of $[0,\infty)\times e_j^\perp$.
It extends to all Borel subsets of $\mathbb R\times e_j^\perp$: both $\mathfrak E_j$ and $\mathfrak E_1'$ vanish on $\{\omega<0\}$ \cite[Lemma~12.18\,(b)]{QEDII}, while $g_{ji}$ leaves $\omega$ unchanged, so that replacing $A$ by $A\cap\{\omega\ge0\}$ changes neither side.
Since $\mathcal H^{\Psi}_{ji}=J\mathcal H^{\Psi}_{12}(\mu')$ and $\mathfrak L_{ji}=J\mathfrak LJ^{-1}$ by definition, Theorem~2.1\,(b) and parts (a) and (b) for $\mu'$ give the first three assertions of (c).

\emph{(c2) Rotations.}
Here $\mu'=\mu$, because $g_{ji}\in O_{\mathrm{cub}}$, and by (c1) one has $\mathfrak E_j(A)=J\,\mathfrak E_1(g_{ji}^{-1}A)\,J^{-1}$ for all Borel sets $A\subset\mathbb R\times e_j^\perp$; it therefore suffices to take $j=1$.
Indeed, for $\Lambda$ acting on $\mathbb R\times e_j^\perp$, the conjugate $\Lambda':=g_{ji}^{-1}\Lambda g_{ji}$ belongs to the restricted Lorentz group of $\mathbb R\times e_1^\perp$.
Suppose that $W'$ is a unitary operator on $\mathcal H_1$ for which $W'\Omega_1=\Omega_1$ and $W'\mathfrak E_1(A')W'^{-1}=\mathfrak E_1(\Lambda'A')$ for each Borel set $A'$; then $W:=JW'J^{-1}$ fixes the vector $\Omega_j=J\Omega_1$, and one has
\begin{equation*}
W\mathfrak E_j(A)W^{-1}=J\,\mathfrak E_1(\Lambda'g_{ji}^{-1}A)\,J^{-1}=\mathfrak E_j(\Lambda A).
\end{equation*}
The group $\mathrm O(3)$ leaves $\mu$ invariant \cite[Theorem~1.4, Part~I\,(a), (f)]{QEDII}.
Accordingly, \cite[Proposition~14.47\,(a)]{QEDII} provides, for $g\in\mathrm O(3)$ with $ge_1=e_1$, unitary operators $\mathcal U(g)$ such that $\mathcal U(g)\Omega_1=\Omega_1$ and
\begin{equation*}
\mathcal U(g)\,\mathfrak E_1(A)\,\mathcal U(g)^{-1}=\mathfrak E_1(gA),\qquad gA:=\{(\omega,gp_\perp):(\omega,p_\perp)\in A\}.
\end{equation*}
This relation is stated there for Borel subsets of $[0,\infty)\times e_1^\perp$; the argument of (c1) extends it to all Borel subsets of $\mathbb R\times e_1^\perp$, because $(\omega,p_\perp)\mapsto(\omega,gp_\perp)$ leaves $\omega$ unchanged.

\emph{(c3) Lorentz transformations.}
Let $\Lambda$ be given.
Since the point $\Lambda(1,0,0)$ lies on $\{\omega=(1+|p_\perp|^2)^{1/2}\}$, it can be written as $g\ell_\beta(1,0,0)$ with a rotation $g$ of $e_1^\perp$ and a real $\beta$.
As a restricted Lorentz transformation leaving $(1,0,0)$ fixed, the map $g':=(g\ell_\beta)^{-1}\Lambda$ is a rotation.
Parts (b) and (c2) show that the operator $W:=\mathcal U(g)\,\mathfrak L(-\beta)\,\mathcal U(g')$ fixes $\Omega_1$ and that
\begin{equation*}
W\,\mathfrak E_1(A)\,W^{-1}=\mathfrak E_1(g\ell_\beta g'A)=\mathfrak E_1(\Lambda A).
\end{equation*}
Finally, the round cone carries the measure $\mathfrak E_1$ \cite[Theorem~1.4, Part~I\,(c), (f)]{QEDII}.

\emph{(d)} By \cite[Theorem~1.3\,(a), (b)(i)]{QEDII}, the law $\mu$ has (G), (dF), the invariance under $\mathbb R^3\rtimes O_{\mathrm{cub}}$ and reflection positivity at the coordinate planes, and lies in $\mathcal C_9$; by \cite[Theorem~1.3\,(b)(v)]{QEDII} the same holds for almost every component.
In the setting of \cite[Theorem~44]{QEDIII}, the hypothesis (H) for $\mu$ is given by (P1), (P5), (P6), (dF) by (P4), and membership in $\mathcal C_9$ by (C1); for $\xi\in\Xi_\mu$ all three are given by (D1).
\end{proof}

\begin{proof}[Proof of Corollary~2.3]
\emph{Step 1 (two stabilities).}
Let $C$ be a bounded operator that commutes with all point plates, and denote by $\mathcal P$ the set of all finite products of point plates.
If $A\in\mathcal P$, then $A^*$ or $-A^*$ belongs to $\mathcal P$ \cite[Proposition~14.18\,(b)]{QEDII}; hence $C^*$ also commutes with every element of $\mathcal P$.
Applied with $\kappa=-it\,n_0$ (where $n_0=(1,0)$ is the direction of time \cite[Definition~14.5]{QEDII}) and $\kappa'=it\,n_0$, both of which lie in $\overline{\mathcal W}$, the same clause shows that $A\in\mathcal P$ implies $A_t:=e^{it\mathsf H_1}Ae^{-it\mathsf H_1}\in\mathcal P$; here one uses $\mathsf T(\mp it\,n_0)=e^{\pm it\mathsf H_1}$ and $\mathsf T(\kappa)\Omega_1=\Omega_1$ \cite[Lemma~14.17\,(a)]{QEDII}.

\emph{Step 2 (positivity of the energy).}
Let $A\in\mathcal P$, and let $t$ be real; then
\begin{equation*}
\langle C^*\Omega_1,e^{it\mathsf H_1}A\Omega_1\rangle=\langle\Omega_1,CA_t\Omega_1\rangle=\langle\Omega_1,A_tC\Omega_1\rangle=\langle A^*\Omega_1,e^{-it\mathsf H_1}C\Omega_1\rangle.
\end{equation*}
As functions of $t$, the first and the last member of this chain are Fourier transforms of finite complex measures, carried by $[0,\infty)$ and by $(-\infty,0]$ respectively.
Because a finite complex measure is determined by its Fourier transform, both measures are carried by $\{0\}$.
Let $P_0$ denote the projection onto $\ker\mathsf H_1=\mathbb C\,\Omega_1$; evaluation at $t=0$ then yields $\langle C^*\Omega_1,\allowbreak A\Omega_1\rangle=\allowbreak \langle C^*\Omega_1,\allowbreak P_0A\Omega_1\rangle$.

\emph{Step 3 (cyclicity).}
Put $c:=\langle\Omega_1,C^*\Omega_1\rangle$.
By Step 2, the vector $C^*\Omega_1-c\,\Omega_1$ is orthogonal to $\Omega_1$ and to every $A\Omega_1$ with $A\in\mathcal P$, and it vanishes by Theorem~2.1\,(c).
On a total set of vectors one thus has $C^*A\Omega_1=AC^*\Omega_1=c\,A\Omega_1$, whence $C^*=c$.
That $\mu_\xi$ fulfils the hypotheses for every $\xi\in\Xi_\mu$ follows from Corollary~2.2\,(c), (d) and \cite[Theorem~44\,(D3)]{QEDIII}.
\end{proof}


\section{The local relativistic field}

The OS axioms and their growth condition give a local Wightman theory
by \cite[Section~IV.1]{OS1975}. We retain a direct reconstruction to
identify its field on the probabilistic OS Hilbert space and prove
the energy and tube estimates on its common domain in Theorem~3.3. These
estimates are used in Section~4 to identify the full infrared law.

\emph{Conventions.}
We use the conventions above and \cite[(n1)--(n8)]{QEDIII}.
Time is the first coordinate, so that points are written $y=(y_1,\vec y)$, and one sets
\begin{equation*}
\theta:=\theta_1,\qquad \Sigma_+:=\Sigma_{\{y_1>0\}},\qquad L^2:=L^2(\mathbb R^3;\Lambda^2)\otimes\mathbb C.
\end{equation*}
On Minkowski space $\mathbb R^3$, with points $x=(x^0,\vec x)$ and basis $e_0,e_2,e_3$, the quadratic form is $x\cdot x=(x^0)^2-|\vec x|^2$ and the forward cone is $V_+=\{x^0>|\vec x|\}$.
Put
\begin{equation*}
d_{\mathcal T}(\zeta):=\operatorname{Im}\zeta^0-|\operatorname{Im}\vec\zeta|;
\end{equation*}
this function may be compared with the function $d_{\mathcal W}(\xi)=\allowbreak \operatorname{Re}\xi^0-|\operatorname{Im}\xi^1|$ on $\mathbb C^2$ defined in \cite[Section~14.2]{QEDII}, whereas the letters $d$ and $d'$ without index stand for half widths of slabs, as in \cite[Proposition~14.16]{QEDII}.
The sets
\begin{gather*}
\mathcal T:=\mathbb R^3+iV_+=\{d_{\mathcal T}>0\},\\
\mathcal T_n:=\{z\in\mathbb C^{3n}:z_{a+1}-z_a\in\mathcal T\},\qquad \mathcal T_n^+:=\{z\in\mathcal T_n:z_1\in\mathcal T\}
\end{gather*}
are open and convex.
The matrix $J:=\operatorname{diag}(i,1,1)$ sends $y$ to the point $Jy=(iy_1,\vec y)$, called a \emph{Euclidean point}, and sends $e_1$ to $ie_0$.
Every linear map $M$ of $\mathbb C^3$ acts on $b\in\Lambda^2\mathbb C^3$ through $\wedge^2M$, abbreviated $Mb$; it acts componentwise on tuples, and, for $M\in\operatorname{GL}(3,\mathbb R)$, on fields through $(M\cdot g)(y)=M\,g(M^{-1}y)$.
In particular $J(e_1\wedge e_k)=i\,e_0\wedge e_k$ and $|Jb|=|b|$; moreover
\begin{equation}
\overline{Jy}=J\theta y,\qquad \overline{Jb}=J\theta\bar b.
\tag{K}
\end{equation}
The space $\mathcal S(\mathbb R^{3n};(\Lambda^2)^{\otimes n})$ of complex valued Schwartz functions is denoted $\mathrm{Sw}_n$ \cite[(n5)]{QEDIII}.
By a Minkowski test one means an $h\in\mathrm{Sw}_1$; its values lie in $\Lambda^2\mathbb C^3$, and its components are the $h^\alpha$ with $\alpha\in\{02,03,23\}$.
For such an index $\alpha$ let $\mathrm{ti}(\alpha)$ be the number of its entries equal to $0$, and put $\mathrm{ti}(\boldsymbol\alpha):=\sum_a\mathrm{ti}(\alpha_a)$.
The group $\mathcal P_+^\uparrow=\mathbb R^3\rtimes SO_0(1,2)$, and in Proposition~3.5 the group $\mathbb R^3\rtimes O(1,2)$, act on tests by $((a,\Lambda)\cdot h)(x)=\Lambda\,h(\Lambda^{-1}(x-a))$; the matrix groups $SO_0(1,2)$ and $O(1,2)$ are introduced in Fact~6.
Let $j\in\{2,3\}$, and let $R^{(j)}_t:=R^{1j}_t$ with $R^{1j}_t$ as in \cite[Definition~14.1]{QEDII}, so that
\begin{equation*}
R^{(j)}_te_1=\cos t\,e_1+\sin t\,e_j,\qquad R^{(j)}_te_j=-\sin t\,e_1+\cos t\,e_j.
\end{equation*}
This matrix is an entire function of $t$, and $\Lambda^{(j)}_s:=JR^{(j)}_{is}J^{-1}$ is the boost
\begin{equation*}
e_0\mapsto\cosh s\,e_0+\sinh s\,e_j,\qquad e_j\mapsto\sinh s\,e_0+\cosh s\,e_j.
\end{equation*}
Compact sets $K_1,\dots,K_n$ satisfying $\max_{K_a}y_1<\min_{K_{a+1}}y_1$ are said to be \emph{time-ordered}.
Euclidean tests are denoted $g$ and $g_a$; for $y=(y_1,\dots,y_n)$ one puts $g(y):=(g_1(y_1),\dots,g_n(y_n))$, and $h(x)$ has the analogous meaning.

\begin{named}{Proposition 3.1}
\emph{The hypotheses.}
Let the data of \cite[Theorem~44]{QEDIII} be given, together with a charge $e\in(0,e_2]$, $e_2$ denoting the ceiling in that theorem.
Let $\mu$ be a volume-first limit point of the laws $P_{N,T}$ at the cutoffs $(N,T)$, and let $\xi\in\Xi_\mu$.
Then the law $\mathrm{lw}:=\mu_\xi$ enjoys the properties below.
\begin{itemize}
\item[(G), (dF)] Both hold by \cite[Theorem~44\,(D1)]{QEDIII}.
\item[(Inv)] Every Euclidean map, including the reflections, leaves $\mathrm{lw}$ invariant, by \cite[Theorem~44\,(D2) with (C2)]{QEDIII}.
\item[(RP)] For every plane $\{\hat n\cdot y=\mathrm{po}\}$ and all $A\in L^2(\Sigma_{\{\hat n\cdot y>\mathrm{po}\}},\mathrm{lw})$ one has $\langle\hat r_{\hat n,\mathrm{po}}A,A\rangle_{\mathrm{lw}}\ge0$, by \cite[Theorem~44\,(D2) with (C3)]{QEDIII}.
\item[(Cl)] For $A,B\in L^2(\mathrm{lw})$ one has $\mathrm{lw}(\bar A\,\allowbreak U(a)B)\to\overline{\mathrm{lw}(A)}\,\allowbreak \mathrm{lw}(B)$ as $|a|\to\infty$ in any manner, by \cite[Theorem~44\,(D4)]{QEDIII}.
\item[(Rec)] The reconstruction at the plane $\{y_1=0\}$ is denoted
\begin{equation*}
(\mathcal H,[\,\cdot\,],H,P_2,P_3,\Omega):=(\mathcal H_1,[\,\cdot\,]_1,\mathsf H_1,\mathsf P_2,\mathsf P_3,\Omega_1).
\end{equation*}
For $A,B\in L^2(\Sigma_+,\mathrm{lw})$, for $t\ge0$ and for real $a_2,a_3$ it satisfies
\begin{gather*}
\langle[A],[B]\rangle=\langle\hat\theta A,B\rangle_{\mathrm{lw}},\\
e^{-tH}[A]=[U(te_1)A],\qquad H=H^*\ge0,\\
e^{i(a_2P_2+a_3P_3)}[A]=[U(a_2e_2+a_3e_3)A],
\end{gather*}
by \cite[(n4), (n8), Proposition~45.3\,(1), (2), (5)]{QEDIII}.
These statements apply to $\mathrm{lw}$: by \cite[Theorem~44\,(D1)]{QEDIII} it satisfies (dF) and the hypothesis (H) of \cite[(n3)]{QEDIII}, and therefore, by \cite[Proposition~45.1\,(10)]{QEDIII}, the conditions (L1)--(L4) under which \cite[Proposition~45.3]{QEDIII} is stated.
The joint spectral measure of $(H,P_2,P_3)$ is denoted $E$ \cite[Lemma~12.18\,(b)]{QEDII}.
\item[(Sp)] The measure $E$ is carried by $\{(\omega,\vec p):\omega\ge|\vec p|\}$, by \cite[Theorem~44\,(D2) with (C4)]{QEDIII}.
\item[(Vac)] The kernel of $H$ is $\mathbb C\,\Omega$, by \cite[Theorem~44\,(D3)]{QEDIII}.
This also follows from (Cl), as
\begin{equation*}
\langle[A],e^{-tH}[B]\rangle=\mathrm{lw}(\overline{\hat\theta A}\,U(te_1)B)\to\langle[A],\Omega\rangle\langle\Omega,[B]\rangle\qquad(t\to\infty).
\end{equation*}
\end{itemize}
\end{named}

Free use is made of the following facts, taken from \cite[Lemma~14.3\,(a), (b)]{QEDII} and \cite[Section~8.1.2]{QEDI}.
The field $X$ extends to $L^2$, with $\|X(g)\|_{L^N(\mathrm{lw})}\le c_N\|g\|$; and $X(g)$ belongs to $L^2(\Sigma_+,\mathrm{lw})$ whenever $g$ vanishes on $\{y_1<0\}$.
If $\mathrm{Gi}$ is a Euclidean map, then $\widehat{\mathrm{Gi}}$ is a unitary multiplicative operator with $\widehat{\mathrm{Gi}}X(g)=X(\mathrm{Gi}\cdot g)$, and $\widehat{\mathrm{Gi}}L^2(\Sigma_O,\allowbreak \mathrm{lw})=\allowbreak L^2(\Sigma_{\mathrm{Gi}\,O},\allowbreak \mathrm{lw})$ for each open set $O$; these are the third and the fifth items of the list in \cite{QEDI}, combined with (Inv).
Each $S_n$ is a bounded symmetric form on $L^2$; moreover $S_1=0$ and $\overline{S_n(g_1,\dots)}=S_n(\bar g_1,\dots)$.
In addition, (dF) implies that $X(g)=0$ almost surely whenever $\ast g=\nabla u$ for some $u\in\mathcal S(\mathbb R^3;\mathbb C)$ \cite[Definition~14.2\,(a)]{QEDII}.
Among the properties of Proposition~3.1 are the one-plane hypotheses of \cite[Definition~14.2\,(f)]{QEDII} for the pair $(1,2)$.


The hypotheses listed above may instead be imposed on one law.
They hold for almost every component of every volume-first limit
point, at every charge and for every mass assignment
\cite[Theorems~1.3 and~1.4]{QEDII}. In the proofs of Theorems~3.2
and~3.3, clustering is used through the unique-vacuum property.
\emph{Localisation.} If $g\in L^2$ vanishes off an open set $O$,
then $X(g)\in L^p(\Sigma_O)$ for $p<\infty$: approximate $g$ by smooth
tests supported in $O$, use the Gaussian moment bound, and use that
$L^2(\Sigma_O)$ is closed. Products of finitely many such fields
belong to $L^2(\Sigma_O)$ by H\"older's inequality.

Theorems~3.2 and~3.3 apply to every law with the properties listed in
Proposition~3.1. Every ergodic law in $\mathcal K$ has those properties.
For the analytic-kernel statement (f2), Steps~1--10 use only (G),
(Inv), (RP), (Rec) and (Sp); the unique-vacuum property is not needed.

\begin{named}{Theorem 3.2}
\emph{The Poincar\'e group.}
On $\mathcal H$ there is a unique strongly continuous unitary representation $\mathcal U$ of $\mathcal P_+^\uparrow$ satisfying (a), (b) and (c); this representation satisfies (d) and (e) as well.
\begin{itemize}
\item[(a)] One has $\mathcal U(a,1)=e^{i(a^0H+a^2P_2+a^3P_3)}$, and the joint spectrum of $(H,P_2,P_3)$ is contained in $\{\omega\ge|\vec p|\}$.
\item[(b)] For $A\in L^2(\Sigma_+,\mathrm{lw})$ and for every rotation $\varrho$ about $e_1$, that is, on Minkowski space, about $e_0$, one has $\mathcal U(0,\varrho)[A]=[\hat\varrho A]$.
\item[(c)] For $j=2,3$ one has $\mathcal U(0,\Lambda^{(j)}_s)=e^{isL_j}$, where $L_j$ denotes the unique self-adjoint operator with the following property: \emph{whenever $O$ is open, $0<t_0<\frac\pi2$, $R^{(j)}_tO\subset\{y_1>0\}$ for $|t|\le t_0$, and $A\in L^2(\Sigma_O,\mathrm{lw})$, the vector $[A]$ lies in $D(e^{tL_j})$ and $e^{tL_j}[A]=[\hat R^{(j)}_tA]$ for $|t|<t_0$.}
Consequently the map $t\mapsto[\hat R^{(j)}_tA]$ has the holomorphic extension $\zeta\mapsto e^{\zeta L_j}[A]$ to the strip $\{|\operatorname{Re}\zeta|<t_0\}$, whose value at $\zeta=is$ is $e^{isL_j}[A]$; in other words, a rotation through the imaginary angle $is$ is the boost of rapidity $s$.
For every fixed $t_0$ the vectors $[A]$ just described form a total set.
\item[(d)] Each $\mathcal U(a,\Lambda)$ fixes $\Omega$, that is, $\mathcal U(a,\Lambda)\Omega=\Omega$; and every vector that is invariant under all $e^{itH}$ is a scalar multiple of $\Omega$.
\item[(e)] The unitary group $e^{isL_2}$ coincides with the group $\mathfrak L(s)$ of \cite[Proposition~14.47\,(c)]{QEDII}; by Theorem~2.1\,(b), $\mathfrak L(s)$ acts on the whole space $\mathcal H$.
\end{itemize}
\end{named}

Because $\mathcal U$ is a representation, every relation valid in $\mathcal P_+^\uparrow$ becomes an identity between bounded operators; an example is
\begin{equation*}
e^{isL_j}\,\mathcal U(a,1)\,e^{-isL_j}=\mathcal U(\Lambda^{(j)}_sa,1).
\end{equation*}

\begin{named}{Theorem 3.3}
\emph{The field.}
Put $D:=\bigcap_mD(H^m)$ and $\|\psi\|_m:=\|(1+H)^m\psi\|$.
One can find operators $F(h)\colon D\to D$, depending linearly on $h\in\mathrm{Sw}_1$, such that the following hold.
\begin{itemize}
\item[(a)] \emph{Field.}
For $\psi\in D$ and $m\ge0$ one has
\begin{equation*}
\|F(h)\psi\|_m\le|h|_m\|\psi\|_{m+4},\qquad |h|_m:=C\sum_{k\le m+2}\binom{m+2}k\|\partial_0^kh\|_{L^1},
\end{equation*}
where $C$ is an absolute constant; in particular $h\mapsto\langle\psi,F(h)\psi'\rangle$ is tempered.
The relation $\langle\psi,F(h)\psi'\rangle=\langle F(\bar h)\psi,\psi'\rangle$ holds, and therefore $F(h)$ is symmetric when $h$ is real.
Further, $\Omega$ belongs to $D$, and $\mathcal U(a,\Lambda)D=D$.
\item[(b)] \emph{Covariance.}
For $(a,\Lambda)\in\mathcal P_+^\uparrow$ one has $\mathcal U(a,\allowbreak \Lambda)F(h)\,\allowbreak \mathcal U(a,\allowbreak \Lambda)^{-1}=\allowbreak F((a,\allowbreak \Lambda)\cdot h)$; thus $F$ is an antisymmetric tensor field of rank two.
\item[(c)] \emph{Locality.}
The commutator $[F(h),F(h')]$ vanishes on $D$ if $(x-x')\cdot(x-x')<0$ for all $x\in\operatorname{supp}h$ and $x'\in\operatorname{supp}h'$.
\item[(d)] \emph{$dF=0$.}
For $u\in\mathcal S(\mathbb R^3;\mathbb C)$ one has $F(d^*u)=0$ on $D$, where
\begin{equation*}
(d^*u)^{23}=-\partial_0u,\qquad (d^*u)^{03}=\partial_2u,\qquad (d^*u)^{02}=-\partial_3u.
\end{equation*}
\item[(e)] \emph{Cyclicity.}
The vector $\Omega$ and the vectors $F(h_1)\cdots F(h_n)\Omega$ span a dense subspace of $\mathcal H$.
\item[(f)] \emph{Relation to $\mathrm{lw}$.}
Put $\check R:=JRJ^{-1}$ and
\begin{equation*}
\mathcal T_n^{R,\pi}:=\{z:(\check R^{-1}z_{\pi(a)})_a\in\mathcal T_n\},\qquad \mathcal D_n:=\bigcup_{R,\pi}\mathcal T_n^{R,\pi},
\end{equation*}
the union being taken over $R\in SO(3)$ and over all permutations $\pi$.
There exist functions $\mathcal W_n(z;b)$, holomorphic in $z\in\mathcal D_n$ and $n$-linear in $b\in(\Lambda^2\mathbb C^3)^n$, with the properties (f1), (f2), (f3) below.
Here an $n$-linear form in $b$ is identified with a linear form on $(\Lambda^2\mathbb C^3)^{\otimes n}$, so that $\mathcal W_n(z;c)$ makes sense for every $c$ in the tensor product, for instance when $c$ is the value at $x$ of an element of $\mathrm{Sw}_n$.
\begin{itemize}
\item[(f1)] On $\mathcal T_n$ the function $\mathcal W_n$ is a function of the differences $z_{a+1}-z_a$ alone.
For $n\ge2$ it obeys
\begin{equation*}
|\mathcal W_n(z;b)|\le\prod_a(c_0|b_a|)\,\Bigl(\tfrac12\min_ad_{\mathcal T}(z_{a+1}-z_a)\Bigr)^{-3n/2},
\end{equation*}
$c_0$ being an absolute constant, while $\mathcal W_1=0$.
For $\Lambda\in SO_0(1,2)$ one has $\mathcal W_n(\Lambda z;\Lambda b)=\mathcal W_n(z;b)$.
The form $W_n(h_1,\allowbreak \dots,\allowbreak h_n):=\allowbreak \langle\Omega,\allowbreak F(h_1)\cdots F(h_n)\Omega\rangle$ has a continuous extension to $\mathrm{Sw}_n$, and for $\mathbf h\in\mathrm{Sw}_n$
\begin{equation*}
\int dx\,\mathcal W_n(x+i\eta;\mathbf h(x))\to W_n(\mathbf h)\qquad\text{as }\eta\to0\text{ with }\eta_{a+1}-\eta_a\in V_+.
\end{equation*}
\item[(f2)] Whenever the $y_a$ are pairwise distinct, the point $Jy=(Jy_1,\dots,Jy_n)$ lies in $\mathcal D_n$.
Permuting the pairs $(z_a,b_a)$ leaves $\mathcal W_n$ unchanged, and so does the map $(z,b)\mapsto(\check Rz,\check Rb)$.
If $g_1,\dots,g_n\in L^2$ have compact supports which are pairwise disjoint, then
\begin{equation*}
S_n(g_1,\dots,g_n)=\int dy\,\mathcal W_n(Jy;Jg(y)).
\end{equation*}
Written out in components, with the index $1$ identified with $0$, this says that $S_n^{\boldsymbol\alpha}(y)=\allowbreak i^{\mathrm{ti}(\boldsymbol\alpha)}\,\allowbreak \mathcal W_n^{\boldsymbol\alpha}(Jy)$, so that every time index contributes a factor $i$.
It follows that the $S_n$ are real analytic wherever no two arguments coincide, configurations with equal times included.
\item[(f3)] From the $W_n$ one recovers the functions $\mathcal W_n(z;b)$ on $\mathcal D_n$, and from these the $S_n$ at pairwise disjoint compact supports; by Proposition~1.2\,(c), no law with (G) and (dF) other than $\mathrm{lw}$ has these $S_n$ at such supports.
In particular, two laws having the properties of Proposition~3.1 and the same $W_n$ for all $n\ge1$ are equal.
\end{itemize}
\end{itemize}
\end{named}


\begin{proof}[Proof of Theorems~3.2 and~3.3]
We first construct bounded slab insertions and holomorphic point-field
vectors (Steps~1--7). Continuation, totality and the rotation identities
then give the Poincar\'e representation (Steps~8--15). Boundary forms
and Hamiltonian estimates turn the kernels into operators on $D$
(Steps~16--19). Covariance, locality, closedness, cyclicity and the
boundary values are established in Steps~20--25.

\emph{Auxiliary classical facts.}

\textbf{Fact 1} (mean value).
For $u$ holomorphic in a neighbourhood of the closed polydisc $\{a\in\mathbb C^M:|a_i|\le\delta\}$,
\begin{equation*}
u(0)=(\pi\delta^2)^{-M}\int_{\{|a_i|<\delta\}}u\,d\lambda^{2M}.
\end{equation*}

\textbf{Fact 2} (identity theorem on a real slice).
Suppose that $I\subset\mathbb C^M$ is a connected open set, that $L$ is an invertible complex linear map, and that $u$ is a holomorphic function on $I$ with values in a Hilbert space.
If $u=0$ on some non-empty subset of $I\cap L\mathbb R^M$ which is open relative to $I\cap L\mathbb R^M$, then $u$ vanishes identically.

\textbf{Fact 3} (Schwarz reflection).
Suppose that $u$ is continuous on the closed half-plane $\{\operatorname{Im}\lambda\ge0\}$ and holomorphic on the open half-plane $\{\operatorname{Im}\lambda>0\}$, and that $u$ vanishes on some non-empty open interval $I\subset\mathbb R$; then $u=0$.
For the proof, choose an open disc $\Delta$ with centre in $I$ such that $\Delta\cap\mathbb R\subset I$, and let $\tilde u$ be the function on $\Delta$ that coincides with $u$ on $\Delta\cap\{\operatorname{Im}\lambda\ge0\}$ and is zero on $\Delta\cap\{\operatorname{Im}\lambda<0\}$.
As $u=0$ on $\Delta\cap\mathbb R$, the function $\tilde u$ is continuous on $\Delta$; it is holomorphic away from the real axis, and by Morera's theorem it is holomorphic on $\Delta$.
Being zero on the open lower half of $\Delta$, it is zero on all of $\Delta$.
Hence $u$ is zero on the non-empty open set $\Delta\cap\{\operatorname{Im}\lambda>0\}$; the identity theorem, applied on the connected set $\{\operatorname{Im}\lambda>0\}$, and then continuity yield $u=0$.

Restricting to the complex lines $x+\lambda\eta$, $\eta\in V_+$, and treating one variable at a time, one deduces a statement in several variables.
Let $B_1,\dots,B_k$ be non-empty open subsets of $\mathbb R^3$, and let $G$ be a function on $\overline{\mathcal T}^{\,k}$ such that, for each $a$ and for arbitrary values of the other variables in $\overline{\mathcal T}$, $G$ is continuous on $\overline{\mathcal T}$ and holomorphic on $\mathcal T$ as a function of its $a$-th variable.
If $G=0$ on $B_1\times\dots\times B_k$, then $G=0$ on $\overline{\mathcal T}^{\,k}$.
The proof is an induction on $a=1,\dots,k$; at the stage $a$ it is known that $G$ vanishes on $\overline{\mathcal T}^{\,a-1}\times B_a\times\dots\times B_k$.
\begin{itemize}
\item[(1)] Give the variables of index smaller than $a$ arbitrary values in $\overline{\mathcal T}$ and those of index larger than $a$ arbitrary values in $B_{a+1},\dots,B_k$, and let $x\in B_a$ and $\eta\in V_+$.
Inserting $x+\lambda\eta$ as the $a$-th variable, one obtains a function of $\lambda$ that is continuous for $\operatorname{Im}\lambda\ge0$ and holomorphic for $\operatorname{Im}\lambda>0$.
It is zero at each real $\lambda$ with $x+\lambda\eta\in B_a$; these $\lambda$ form an open subset of $\mathbb R$ which contains $0$, and so contains a non-empty open interval.
By the case of one variable, this function of $\lambda$ vanishes identically.
\item[(2)] Taking $\lambda=i$, one sees that $G$, as a function of its $a$-th variable, is zero on $B_a+iV_+$, which is a non-empty open subset of the convex set $\mathcal T$.
By the identity theorem it is zero on $\mathcal T$, and by continuity on $\overline{\mathcal T}$.
Consequently $G$ vanishes on $\overline{\mathcal T}^{\,a}\times B_{a+1}\times\dots\times B_k$, which completes the stage $a$.
\end{itemize}
Steps~24 and~25 and Part~3 of the proof of Corollary~3.4 use the case $B_1=\dots=B_k=\mathbb R^3$; the proof of Proposition~3.6 uses balls.

\textbf{Fact 4} (spectral theorem).
Let $K$ be a self-adjoint operator.
\begin{itemize}
\item[(i)] Suppose $v\in D(e^{tK})$ for all $|t|<t_0$; then $\gamma\mapsto e^{i\gamma K}v$ is holomorphic in the strip $\{|\operatorname{Im}\gamma|<t_0\}$.
\item[(ii)] Conversely, if $s\mapsto e^{isK}v$ has a holomorphic extension $V$ to that strip, then $v\in D(e^{tK})$ and $e^{tK}v=V(-it)$; to see this, one pairs with the vectors of compact spectral support, which constitute a core.
\item[(iii)] Every dense subspace $D'\subset D(e^{tK})$ satisfying $e^{isK}D'\subset D'$ for all real $s$ is a core for $e^{tK}$.
Indeed, if $\phi$ is orthogonal to $D'$ with respect to the graph inner product, then each of the finite measures $(1+e^{2t\lambda})\,d\langle\phi,E_K(\lambda)\psi\rangle$, $\psi\in D'$, has vanishing Fourier transform.
\end{itemize}

\textbf{Fact 5} (nuclear theorem with a bound).
Let $M$ be an $n$-linear form on $\mathrm{Sw}_1$ such that $|M(h_1,\dots,h_n)|\le\prod_ap(h_a)$ for some continuous seminorm $p$.
Then $M(h_1,\dots,h_n)=W(h_1\otimes\dots\otimes h_n)$ for a unique continuous linear form $W$ on $\mathrm{Sw}_n$, and $|W|\le P$, where the continuous seminorm $P$ is determined by $p$ and $n$ alone.
The proof proceeds by expansion in Hermite functions, as in \cite[Proposition~45.2\,(2)]{QEDIII}.

\textbf{Fact 6} (Lorentz group).
In the basis $e_0,e_2,e_3$ the form $x\cdot x$ has the matrix $-J^2=\operatorname{diag}(1,-1,-1)$.
One defines
\begin{gather*}
O(1,2):=\{\Lambda\text{ real}:\Lambda^TJ^2\Lambda=J^2\},\qquad SO_0(1,2):=\{\Lambda\in O(1,2):\det\Lambda=1,\ (\Lambda e_0)^0>0\},\\
so(1,2):=\{Y\text{ real}:Y^TJ^2+J^2Y=0\}.
\end{gather*}
An element $\Lambda\in O(1,2)$ permutes the two convex components $\pm V_+$ of $\{x\cdot x>0\}$; hence $(\Lambda e_0)^0>0$ holds if and only if $\Lambda V_+=V_+$, and $SO_0(1,2)$ is a group.
\begin{itemize}
\item[(i)] For $\Lambda\in O(1,2)$ one has $(\Lambda e_0)\cdot(\Lambda e_0)=1$, whence $|(\Lambda e_0)^0|\ge1$.
Let $Y\in so(1,2)$.
Then $e^{sY}\in O(1,2)$, and both $\det e^{sY}\in\{\pm1\}$ and the sign of $(e^{sY}e_0)^0$ depend continuously on $s$; therefore $e^{sY}\in SO_0(1,2)$.
\item[(ii)] Each $\Lambda\in SO_0(1,2)$ can be written $\Lambda=\varrho\Lambda^{(2)}_s\varrho'$ with $s\ge0$ and rotations $\varrho,\varrho'$ about $e_0$.
To see this, write $\Lambda e_0=\cosh s\,e_0+\sinh s\,u$, with $u$ a unit vector in $\operatorname{span}(e_2,e_3)$, which is possible by (i).
For a rotation $\varrho$ with $\varrho e_2=u$ one has $\varrho\Lambda^{(2)}_se_0=\Lambda e_0$.
The matrix $\varrho':=(\varrho\Lambda^{(2)}_s)^{-1}\Lambda$ preserves the form and fixes $e_0$; it therefore maps $\operatorname{span}(e_2,e_3)$ isometrically onto itself, and since its determinant is $1$, it is a rotation.
\item[(iii)] Each of the three factors in (ii) has the form $e^{sY}$ with $Y\in so(1,2)$.
Hence these one-parameter groups generate $SO_0(1,2)$, and $SO_0(1,2)$ is connected.
\item[(iv)] Let $\theta_2=\operatorname{diag}(1,-1,1)$ and $\theta=\operatorname{diag}(-1,1,1)$.
The group $O(1,2)$ is the disjoint union of the four cosets $SO_0(1,2)$, $SO_0(1,2)\,\theta_2$, $SO_0(1,2)\,\theta$ and $SO_0(1,2)\,\theta_2\theta$.
Indeed, if $(\Lambda e_0)^0<0$ one multiplies $\Lambda$ on the right by $\theta$, and afterwards by $\theta_2$ if the determinant equals $-1$; multiplication by $\theta_2$ leaves $(\Lambda e_0)^0$ unchanged.
\end{itemize}

\textbf{Step 1} (semigroup).
For $\zeta\in\overline{\mathcal T}$ define
\begin{equation*}
T(\zeta):=\int e^{i(\zeta^0\omega+\vec\zeta\cdot\vec p)}\,dE.
\end{equation*}
According to (Sp), the integrand is bounded in modulus by $e^{-d_{\mathcal T}(\zeta)\omega}$.
Consequently $T$ is a contraction semigroup with $T(\zeta)^*=T(-\bar\zeta)$ and $T(\zeta)\Omega=\Omega$, and $T(\mathcal T)\mathcal H\subset D$.
On each $D(H^m)$ equipped with $\|\cdot\|_m$, it is a contraction, strongly continuous on $\overline{\mathcal T}$ and holomorphic on $\mathcal T$.
At real points $a$ one has $T(a)=e^{i(a^0H+\vec a\cdot\vec P)}$, and for $y_1\ge0$ one has $T(Jy)[A]=[U(y)A]$.

\textbf{Step 2} (insertions; the single bound that is uniform in the number of fields).
The hypotheses (G), $(T_{123})$, $(R_1)$, $(RP_1)$ of \cite[Proposition~14.16, (b), (c)]{QEDII} are among those of Proposition~3.1.
That proposition therefore yields, for each $g\in L^2$ vanishing outside $\{|y_1|\le d\}$, a unique bounded operator $\mathcal B_d(g)$, depending linearly on $g$, such that
\begin{equation*}
\mathcal B_d(g)[A]=[X(\tau_{de_1}g)\,U(2de_1)A]\qquad\text{for }A\in\textstyle\bigcap_pL^p(\Sigma_+,\mathrm{lw});
\end{equation*}
in \cite{QEDII} this operator is denoted $\mathcal B^{\mathrm c}_d(\ast g)$.
It satisfies the following.
\begin{itemize}
\item[(i)] $\|\mathcal B_d(g)\|\le k_1\|g\|$ with $k_1=(2\exp1)^{1/2}$.
\item[(ii)] $\mathcal B_d(g)^*=\mathcal B_d(\overline{\theta\cdot g})$.
\item[(iii)] For real $a$ and $s\ge|a_1|$,
\begin{equation*}
T(ise_0+Ja)\,\mathcal B_d(g)\,T(ise_0-Ja)=\mathcal B_{d+s}(\tau_ag).
\end{equation*}
If $a_3=0$, this is part (c) of \cite[Proposition~14.16]{QEDII}.
If $a=a_3e_3$ and $s=0$, it results from the identity that defines $\mathcal B_d$, from (Rec) applied to $e^{ia_3P_3}=T(a_3e_3)$, and from the facts that $U(a_3e_3)$ is multiplicative and that $U(a_3e_3)X(g)=X(\tau_{a_3e_3}g)$.
For an arbitrary real $a$, the semigroup law of Step~1 gives $T(ise_0\pm Ja)=T(\pm a_3e_3)\,T(ise_0\pm J(a_1,a_2,0))$, and one has $\tau_a=\tau_{a_3e_3}\tau_{(a_1,a_2,0)}$; the identity follows on applying the case $a_3=0$ first, and then the case $a=a_3e_3$, $s=0$ with $d+s$ and $\tau_{(a_1,a_2,0)}g$ in place of $d$ and $g$, the latter function vanishing outside $\{|y_1|\le d+s\}$ because $s\ge|a_1|$.
\end{itemize}
The constant $k_1$ is independent of the vectors and so of the number of fields they contain; it is obtained from (G) through repeated reflections and Cauchy's estimate \cite[Lemma~14.15]{QEDII}.
Accordingly, no estimate in what follows involves a sum over $n$.

\textbf{Step 3} (point fields).
Fix $\delta>0$.
For $v\in Q:=(-\delta,\delta)^3$ set $\chi_v(u):=\prod_i1[u_i^2+v_i^2<\delta^2]$, and for $d_{\mathcal T}(\zeta),d_{\mathcal T}(\zeta')\ge3\delta$ set
\begin{equation*}
\Phi_b(\zeta,\zeta'):=(\pi\delta^2)^{-3}\int_Qdv\,T(\zeta-i\delta e_0+iJv)\,\mathcal B_\delta(\chi_vJ^{-1}b)\,T(\zeta'-i\delta e_0-iJv).
\end{equation*}
The arguments of $T$ belong to $\mathcal T$, because $iJv=(-v_1,i\vec v)$.
\begin{itemize}
\item[(a)] Since $\|\chi_v\|_2\le(2\delta)^{3/2}$, one has $\|\Phi_b\|\le2^{9/2}\pi^{-3}k_1|b|\,\delta^{-3/2}$.
\item[(b)] For $\kappa,\kappa'\in\overline{\mathcal T}$ one has $T(\kappa)\Phi_b(\zeta,\allowbreak \zeta')T(\kappa')=\allowbreak \Phi_b(\kappa+\zeta,\allowbreak \zeta'+\kappa')$.
\item[(c)] $\Phi_b(\zeta,\zeta')^*=\Phi_{\bar b}(-\bar\zeta',-\bar\zeta)$; this follows from Step~2\,(ii), from (K), and from the change of variables $v\to-\theta v$.
\end{itemize}

\textbf{Step 4} (the point fields are the kernel of the insertions).
\begin{itemize}
\item[(d)] \emph{If $g\in L^2$ has compact support in $\{|y_1|\le d'\}$ and $d-d'\ge3\delta$, then}
\begin{equation*}
\mathcal B_d(g)=\int dy\,\Phi_{Jg(y)}(ide_0+Jy,\,ide_0-Jy).
\end{equation*}
\end{itemize}
To prove (d), consider
\begin{equation*}
Z(a):=T(i(d-d')e_0+Ja)\,\mathcal B_{d'}(g)\,T(i(d-d')e_0-Ja).
\end{equation*}
This operator depends weakly holomorphically on $a$ in the region $|\operatorname{Re}a_1|+|\operatorname{Im}\vec a|<d-d'$, which contains a neighbourhood of $\{|a_i|\le\delta\}$; moreover $Z(0)=\mathcal B_d(g)$.
One applies Fact~1 with $a=u+iv$.
The imaginary part $iv$ of the shift yields the factors $T(\,\cdot\pm iJv)$ appearing in Step~3, while Step~2\,(iii) converts the real part $u$ into the translation $\tau_ug$.
In $L^2$ one has $\int du\,\chi_v(u)\,\tau_ug=\int dy\,(\tau_y\chi_v)\,g(y)$, and a second use of Step~2\,(iii) converts $\tau_y$ into factors $T$.

\textbf{Step 5} (independence of $\delta$).
Take two values of $\delta$.
The corresponding weakly continuous functions $y\mapsto\Phi_b(ide_0+Jy,ide_0-Jy)$ yield, by (d), the same integral against every $g$, and are therefore equal.
It follows from (b) that the two functions $\Phi_b$ coincide where $\min(d_{\mathcal T}(\zeta),d_{\mathcal T}(\zeta'))>3\delta$, with $\delta$ the larger of the two values, and by continuity also where equality holds.
Hence $\Phi_b$ is a norm holomorphic function on $\mathcal T\times\mathcal T$ which has the properties (b) and (c) of Step~3 and, for compact support in $\{|y_1|<d\}$, the property (d) of Step~4.
Moreover, the choice $3\delta=\min(d_{\mathcal T}(\zeta),d_{\mathcal T}(\zeta'))$ in Step~3\,(a) gives
\begin{equation*}
\|\Phi_b(\zeta,\zeta')\|\le c_0|b|\min(d_{\mathcal T}(\zeta),d_{\mathcal T}(\zeta'))^{-3/2},\qquad c_0=2^{9/2}3^{3/2}\pi^{-3}k_1.
\end{equation*}
From (b) one also sees that $\Phi_b(\cdot,\zeta')\Phi_{b'}(\zeta'',\cdot)$ is a function of $\zeta'+\zeta''$ alone, and that neither $\Phi_b(\zeta,\kappa)\Omega$ nor $\langle\Omega,\Phi_b(\kappa,\zeta)\,\cdot\,\rangle$ depends on $\kappa$.

\textbf{Step 6} (vectors and kernels).
On $\mathcal T_n^+$ put
\begin{equation*}
\Psi_n(z;b):=\Phi_{b_1}\Bigl(z_1,\frac{z_2-z_1}2\Bigr)\Phi_{b_2}\Bigl(\frac{z_2-z_1}2,\frac{z_3-z_2}2\Bigr)\cdots\Phi_{b_n}\Bigl(\frac{z_n-z_{n-1}}2,\kappa\Bigr)\Omega,
\end{equation*}
and put $\Psi_0:=\Omega$ and $D_1:=\operatorname{span}\{\Psi_n(z;b)\}$; on $\mathcal T_n$ put $\mathcal W_n(z;b):=\langle\Omega,\Psi_n(z+\kappa';b)\rangle$.
Here $\kappa\in\mathcal T$ is arbitrary, and $\kappa'$ is any vector with $z_1+\kappa'\in\mathcal T$.
Step~5 implies the following.
\begin{itemize}
\item[(a)] The bound of (f1) holds for $\mathcal W_n$.
For $n=1$, letting $d_{\mathcal T}(\kappa)$ and $d_{\mathcal T}(\kappa')$ tend to $\infty$ in the bound of Step~5 gives $\mathcal W_1=0$; this is also a consequence of Step~7.
\item[(b)] $T(\kappa)\Psi_n(z;b)=\Psi_n(z+\kappa;b)$.
Hence $D_1\subset D$ and $e^{-sH}D_1\subset D_1$, and $\mathcal W_n$ is a function of the differences only.
\item[(c)] For the inner products one finds
\begin{equation*}
\langle\Psi_m(z';b'),\Psi_n(z;b)\rangle=\mathcal W_{m+n}(\bar z'_m,\dots,\bar z'_1,z_1,\dots,z_n;\bar b'_m,\dots,\bar b'_1,b_1,\dots,b_n).
\end{equation*}
\end{itemize}

\textbf{Step 7} (ordered times).
\emph{For $g_a\in L^2$ with time-ordered compact supports in $\{y_1>0\}$,}
\begin{equation*}
\Bigl[\prod_aX(g_a)\Bigr]=\int dy\,\Psi_n(Jy;Jg(y)).
\end{equation*}
Choose slabs $\{|y_1-c_a|<d_a\}$ containing $\operatorname{supp}g_a$ such that $0\le c_1-d_1$ and $c_a+d_a\le c_{a+1}-d_{a+1}$.
By induction over the definition of $\mathcal B$ one obtains
\begin{equation*}
\Bigl[\prod_aX(g_a)\Bigr]=e^{-(c_1-d_1)H}\mathcal B_{d_1}(\tau_{-c_1e_1}g_1)\,e^{-(c_2-d_2-c_1-d_1)H}\mathcal B_{d_2}(\tau_{-c_2e_1}g_2)\cdots\Omega.
\end{equation*}
Substituting Step~4\,(d) for each factor $\mathcal B$ and combining the arguments by means of Step~3\,(b), one sees that the first argument turns into $Jy_1$, and the argument between two consecutive factors into $J(y_{a+1}-y_a)$.
Consequently $S_n(g_1,\dots,g_n)=\int dy\,\mathcal W_n(Jy;Jg(y))$ for time-ordered compact supports wherever they lie, and $\mathcal W_1=0$.

\textbf{Step 8} (rotated and permuted tubes).
Define on $\mathcal T_n^{R,\pi}$
\begin{equation*}
\mathcal W_n^{R,\pi}(z;b):=\mathcal W_n\bigl((\check R^{-1}z_{\pi(a)})_a;(\check R^{-1}b_{\pi(a)})_a\bigr).
\end{equation*}
Since $\mathcal T=\{Jw:\operatorname{Re}w_1>|\operatorname{Im}\vec w|\}$ and since $R$ and $\pi$ are real, the set $\mathcal T_n^{R,\pi}$ is open and convex, and the map $Jw\mapsto J\operatorname{Re}w$ sends it into itself.
The Euclidean points in it are exactly the $Jy$ for which $u\cdot y_{\pi(1)}<\dots<u\cdot y_{\pi(n)}$, with $u=Re_1$.
For $g_a\in L^2$ whose compact supports lie over such points one has
\begin{equation*}
S_n(g_1,\dots,g_n)=S_n(R^{-1}\cdot g_{\pi(1)},\dots,R^{-1}\cdot g_{\pi(n)})=\int dy\,\mathcal W_n^{R,\pi}(Jy;Jg(y)).
\end{equation*}
The first equality holds because the $X(g)$ commute and because of (Inv), tensor indices included; the second follows from Step~7 and the substitution $y\to Ry$.

\textbf{Step 9} (one function; (f2)).
Let two of the tubes have non-empty intersection.
This intersection is convex, and the Euclidean points in it include a relatively open set; there the two functions are both continuous densities of $S_n$, hence equal.
Fact~2, applied with $L=J$, extends the equality to the whole intersection.
The $\mathcal W_n^{R,\pi}$ are therefore restrictions of a single holomorphic function $\mathcal W_n$ on $\mathcal D_n$, whose construction gives it the invariance properties stated in (f2).
If the $y_a$ are pairwise distinct, there is a unit vector $u$ for which the numbers $u\cdot y_a$ are distinct as well; thus $Jy\in\mathcal D_n$.
For disjoint compact supports, one cuts the $g_a$ into pieces by multiplying with indicator functions of small cubes, which is legitimate in $L^2$; by the Lebesgue-number lemma, each product of pieces then lies over a single tube.

\textbf{Step 10} (Lorentz invariance).
Let $z\in\mathcal T_n$.
Owing to the holomorphy in $z$ and the multilinearity in $b$, the derivative $\frac d{ds}\mathcal W_n(e^{sY}z;e^{sY}b)|_{s=0}$ depends complex linearly on the matrix $Y$; by Step~9 it vanishes for $Y\in J\,so(3)J^{-1}$.
The complex span of $J\,so(3)J^{-1}$ contains $so(1,2)$, since $J^2=-\operatorname{diag}(1,-1,-1)$ implies that $Y\in so(1,2)$ if and only if $J^{-1}YJ$ is antisymmetric.
For such $Y$ the matrices $e^{sY}$ preserve $\mathcal T_n$ by Fact~6\,(i), and $\mathcal W_n(e^{sY}z;e^{sY}b)$ is thus independent of $s$; by Fact~6\,(iii), these one-parameter groups generate $SO_0(1,2)$.

\textbf{Step 11} (totality; here (dF) is indispensable).
\emph{The subspace $D_1$ is dense.}
Proposition~3.1 makes Theorem~2.1\,(a) applicable; according to it, $\Omega$ together with the vectors $[\prod_aX(g_a)]$ with time-ordered supports in $\{y_1>0\}$ is a total set, and Step~7 places these vectors in $\overline{D_1}$.
In the proof of Theorem~2.1\,(a), the hypothesis (dF) enters through \cite[Proposition~14.13]{QEDII}, applied to the forms $\langle\psi,[\prod_aX(g_a)]\rangle$, which vanish on gradients by (dF).
Step~7 with shrinking supports, combined with Fact~2 on $\mathcal T_n^+$, shows that a smaller set is total as well: $\Omega$ and the $[\prod_aX(g_a)]$ whose $g_a$ are supported near one time-ordered configuration in $\{y_1>0\}$, one configuration being fixed for each $n$.

\textbf{Step 12} (the representation; Theorem~3.2\,(a), (b), (d)).
Step~10, Step~6\,(c) and the identity $\overline{\Lambda z}=\Lambda\bar z$ show that the Gram matrices of the vectors $\Psi_n(z;b)$ do not change under $(z,b)\mapsto(\Lambda z,\Lambda b)$.
The prescription $\mathcal U(0,\Lambda)\Psi_n(z;b):=\Psi_n(\Lambda z;\Lambda b)$ therefore defines an isometry of $D_1$ onto $D_1$.
On $D_1$ one reads off the group law, the strong continuity and the relation $\mathcal U(0,\allowbreak \Lambda)T(a)\,\allowbreak \mathcal U(0,\allowbreak \Lambda)^{-1}=\allowbreak T(\Lambda a)$.
Consequently $\mathcal U(a,\Lambda):=T(a)\,\mathcal U(0,\Lambda)$ is a representation of the group, and it leaves $\Omega=\Psi_0$ fixed.
What (a) and (d) assert beyond this is contained in (Sp) and (Vac).
Since $J$ commutes with every rotation $\varrho$ about $e_1$, the relation $\mathcal U(0,\varrho)[B]=[\hat\varrho B]$ holds on the total set formed by the time-ordered products of Step~7; moreover $[A]\mapsto[\hat\varrho A]$ is unitary by \cite[Proposition~14.47\,(a)]{QEDII}.

\textbf{Step 13} (continuation in the angle: products).
Let $j$ and $0<t_0<\frac\pi2$ be fixed; write $R_t:=R^{(j)}_t$ and $\Lambda_s:=\Lambda^{(j)}_s$, and let $L$ be the self-adjoint operator, provided by Stone's theorem, for which $\mathcal U(0,\Lambda_s)=e^{isL}$.
A product $B=\prod_aX(g_a)$ is called \emph{admissible} when the $g_a\in L^2$ have compact supports and, for every $|t|\le t_0$, the sets $R_t\operatorname{supp}g_a$ are time-ordered in $\{y_1>0\}$.
Denote by $K$ the union of the supports.
On the compact set $[-t_0,t_0]\times K$ the function $(t,y)\mapsto(R_ty)_1$ is continuous and positive; hence it is bounded below there by some $2\epsilon>0$.
Let $O$ be the open $\epsilon$-neighbourhood of $K$.
Each $R_t$ being an isometry, $(R_ty)_1>\epsilon$ for $y\in O$, so that $R_tO\subset\{y_1>0\}$ for $|t|\le t_0$; and $B\in L^2(\Sigma_O,\mathrm{lw})$ by the statement \emph{Localisation} following Proposition~3.1.
The vectors $[B]$, as well as $\Omega=[1]$, are thus among the vectors $[A]$ of Theorem~3.2\,(c).
Step~11, used at the configuration $(e_1,2e_1,\dots,ne_1)$, shows that these $[B]$ together with $\Omega$ form a total set; this is the totality claimed in Theorem~3.2\,(c).
By Steps~7 and~12, $e^{isL}[B]=V(s)$ with
\begin{equation*}
V(\gamma):=\int dy\,\Psi_n(JR_{i\gamma}y;JR_{i\gamma}g(y)).
\end{equation*}
If $\gamma=s-it$, then $JR_{i\gamma}y=\Lambda_sJR_ty$ lies in $\mathcal T_n^+$.
Therefore $V$ is holomorphic in $\{|\operatorname{Im}\gamma|<t_0\}$ with $V(-it)=[\hat R_tB]$, and Fact~4\,(ii) yields $e^{tL}[B]=[\hat R_tB]$.

\textbf{Step 14} (all $A$; Theorem~3.2\,(c)).
Take $O$, $t_0$ and $A$ as in Theorem~3.2\,(c), an admissible $B$, and $s,t$ with $|s|+|t|<t_0$.
From $\theta R_t=R_{-t}\theta$ and (Inv) one gets
\begin{equation*}
\langle[\hat R_t\hat R_sB],[A]\rangle=\langle\hat\theta\hat R_sB,\hat R_tA\rangle_{\mathrm{lw}},\quad\text{i.e.}\quad\langle e^{(s+t)L}[B],[A]\rangle=\langle e^{sL}[B],[\hat R_tA]\rangle.
\end{equation*}
By Fact~4\,(i) both members are antiholomorphic in $s$; evaluating at $s=is'$ with real $s'$ gives $\langle e^{tL}v,[A]\rangle=\langle v,[\hat R_tA]\rangle$ for $v=e^{is'L}[B]$.
The linear span of these $v$ and of $\Omega$ is a dense subspace contained in $D(e^{tL})$ and stable under every $e^{isL}$; by Fact~4\,(iii) it is a core for the self-adjoint operator $e^{tL}$.
It follows that $[A]\in D(e^{tL})$ and $e^{tL}[A]=[\hat R_tA]$.

\textbf{Step 15} (uniqueness; Theorem~3.2\,(e)).
Suppose that a self-adjoint operator $L'$ has the property stated in Theorem~3.2\,(c).
Applying this property to an admissible $B$ and to the constant function $1$, with the set $O$ of Step~13, one obtains $e^{tL'}[B]=[\hat R_tB]$ and $e^{tL'}\Omega=\Omega$ for $|t|<t_0$.
By Steps~13 and~12, $L$ satisfies the same two identities.
Let $v$ be $[B]$ or $\Omega$; the latter belongs to the total set as well.
By Fact~4\,(i), $e^{i\gamma L'}v$ and $e^{i\gamma L}v$ are holomorphic in the strip, and they coincide at $\gamma=-it$; hence they coincide for real $\gamma$.
Thus $e^{isL'}=e^{isL}$, and Stone's theorem gives $L'=L$.
The uniqueness of $\mathcal U$ follows, because $SO_0(1,2)=\{\varrho\Lambda^{(2)}_s\varrho'\}$ by Fact~6\,(ii).
For (e), consider admissible products $B$ of plates; by Step~11 and \cite[Lemma~14.8]{QEDII}, these $[B]$ and $\Omega$ form a total set.
In \cite[Section~14.2]{QEDII} the symbol $\Lambda_s$ denotes the matrix that is written $R_{is}$ in this section.
For such $B$, the map $s\mapsto\mathfrak L(s)[B]$ continues to the strip and takes the value $[\hat R_tB]$ at $s=-it$, by \cite[Propositions~14.22\,(b), (c), 14.47\,(c), Lemmas~14.6\,(a), 14.9\,(b), (c), Definition~14.21]{QEDII}.
Here, for $\gamma=s-it$ with $|t|<t_0$, one has $R_{i\gamma}=R_{is}R_t$; the matrix $R_t$ carries each point of the product of the supports to a real configuration that is time-ordered in $\{y_1>0\}$ and hence lies in $\mathcal W^{(n)}$, and $R_{is}$ maps $\mathcal W^{(n)}$ into itself.
The conclusion then follows from Fact~4\,(ii), as in the case of $L'$.

\textbf{Step 16} (the field as a form).
For $t,t'>0$ consider the operator
\begin{equation*}
F_{t,t'}(h):=\int dx\,\Phi_{h(x)}(x+ite_0,-x+it'e_0).
\end{equation*}
Its norm does not exceed $c_0\|h\|_{L^1}\min(t,t')^{-3/2}$, and $e^{-sH}F_{t,t'}(h)e^{-s'H}=F_{t+s,t'+s'}(h)$.
Consequently the form
\begin{equation*}
q_h(\psi,\psi'):=\langle e^{tH}\psi,F_{t,t'}(h)e^{t'H}\psi'\rangle\qquad\text{on }D_0:=\bigcup_ME(\{\omega\le M\})\mathcal H
\end{equation*}
is independent of $t$ and $t'$.
Insert $(1+H)^{-2}=\int_0^\infty te^{-t}e^{-tH}dt$ on either side and use $\iint tt'e^{-t-t'}\min(t,t')^{-3/2}\,dt\,dt'\le2\sqrt\pi$; this gives
\begin{equation*}
|q_h(\psi,\psi')|\le C\|h\|_{L^1}\|\psi\|_2\|\psi'\|_2,\qquad C=2\sqrt\pi c_0.
\end{equation*}
The exponent $-3/2$ coincides with that of the free field.
If instead $s$ powers are taken on each side, by means of
\begin{equation*}
(1+H)^{-s}=\int_0^\infty t^{s-1}e^{-t}e^{-tH}dt\Big/\int_0^\infty t^{s-1}e^{-t}dt,
\end{equation*}
the relevant double integral, that of $(tt')^{s-1}e^{-t-t'}\min(t,t')^{-3/2}$, is finite exactly when $s>\frac32$; the estimate above uses $s=2$.
Since the estimate involves $h$ only through $\|h\|_{L^1}$, and since $x\mapsto\Phi_b(x+ite_0,-x+it'e_0)$ is bounded and norm continuous, $q_h$ extends to finite measures $h$.
Put $e_{02}:=e_0\wedge e_2$, $e_{03}:=e_0\wedge e_3$ and $e_{23}:=e_2\wedge e_3$, so that $h=\sum_\alpha h^\alpha e_\alpha$.
For $x\in\mathbb R^3$ and each index $\alpha$, the point field $F_\alpha(x):=q_{\delta_xe_\alpha}$ is therefore a form on $D(H^2)\times D(H^2)$, and
\begin{equation*}
q_h(\psi,\psi')=\sum_\alpha\int dx\,h^\alpha(x)\,F_\alpha(x)(\psi,\psi')\qquad(\psi,\psi'\in D(H^2));
\end{equation*}
as $q_h$ is the form of $F(h)$ (Step~17), this says that $F(h)=\sum_\alpha\int dx\,h^\alpha(x)\,F_\alpha(x)$ in the sense of forms.
For $f\in L^1(\mathbb R^2;\Lambda^2\mathbb C^3)$ the field $F(\delta_c\otimes f)$ at the sharp time $x^0=c$ is likewise a form on $D(H^2)\times D(H^2)$.
Derivatives of $h$ in $x^0$ are needed only in Step~17, where the form is turned into an operator.
No assertion is made on whether $F(\delta_c\otimes f)$ is an operator on $D$.
\textbf{Step 17} (the field as an operator; Theorem~3.3\,(a)).
The relation $e^{isH}F_{t,t'}(h)e^{-isH}=F_{t,t'}(\tau_{se_0}h)$ implies
\begin{equation*}
q_h(H\psi,\psi')-q_h(\psi,H\psi')=q_{i\partial_0h}(\psi,\psi'),
\end{equation*}
and consequently
\begin{equation*}
q_h((1+H)^m\psi,\psi')=\sum_k\binom mk\,q_{(i\partial_0)^kh}(\psi,(1+H)^{m-k}\psi').
\end{equation*}
Applying this identity with $m+2$ in place of $m$ and with $(1+H)^{-2}\psi$ in place of $\psi$, and estimating each term by Step~16, one arrives at
\begin{equation*}
|q_h((1+H)^m\psi,\psi')|\le|h|_m\|\psi\|\,\|\psi'\|_{m+4}.
\end{equation*}
Riesz representation provides $F(h)\psi'$ with $q_h(\psi,\psi')=\langle\psi,F(h)\psi'\rangle$.
Since $D_0$ is a core for $(1+H)^m$, the vector $F(h)\psi'$ lies in $D(H^m)$ and satisfies the bound in (a); by continuity, $F(h)$ extends to $D$.
The equality $\mathcal U(a,\Lambda)D=D$ holds because (Sp) implies
\begin{equation*}
\mathcal U(0,\Lambda)H\,\mathcal U(0,\Lambda)^{-1}=(\Lambda e_0)^0H+\overrightarrow{\Lambda e_0}\cdot\vec P\le2(\Lambda e_0)^0H.
\end{equation*}
More precisely, $\|\mathcal U(0,\Lambda)\psi\|_m\le(2(\Lambda e_0)^0)^m\|\psi\|_m$.
To prove this, write $\mathcal U(0,\Lambda)^{-1}H\,\mathcal U(0,\Lambda)=\int f\,dE$, where $f=\allowbreak (\Lambda^{-1}e_0)^0\omega+\overrightarrow{\Lambda^{-1}e_0}\cdot\vec p$.
On the cone one has $0\le f\le c\,\omega$ with $c=2(\Lambda^{-1}e_0)^0=2(\Lambda e_0)^0\ge2$; here $(\Lambda^{-1}e_0)^0$ and $(\Lambda e_0)^0$ coincide because Fact~6, together with $J^4=1$, gives $\Lambda^{-1}=J^2\Lambda^TJ^2$.
Hence $1+f\le c\,(1+\omega)$, and
\begin{equation*}
\|(1+H)^m\,\mathcal U(0,\Lambda)\psi\|=\Bigl\|\int(1+f)^mdE\,\psi\Bigr\|\le c^m\|(1+H)^m\psi\|.
\end{equation*}

\textbf{Step 18} (rules).
\begin{itemize}
\item[(e1)] $T(\zeta)F(h)T(\zeta')=\int dx\,\Phi_{h(x)}(\zeta+x,\zeta'-x)$ on $\mathcal T\times\mathcal T$.
\item[(e2)] $F_{t,t}(h)^*=F_{t,t}(\bar h)$ by Step~3\,(c), that is, by (K); this is the symmetry in Theorem~3.3\,(a).
\item[(e3)] $T(a)F(h)T(-a)=F(\tau_ah)$ for real $a$.
\item[(e4)] For $\kappa,\kappa_a\in\mathcal T$, the operator $T(\kappa)F(h)\subset(F(\bar h)T(\kappa)^*)^*$ is bounded,
\begin{equation*}
T(\kappa)F(h)\Psi_n(z;b)=\int dx\,\Psi_{n+1}(x+\kappa,z+\kappa;h(x),b),
\end{equation*}
and
\begin{equation*}
T(\kappa_1)F(h_1)\cdots T(\kappa_n)F(h_n)\Omega=\int dx\,\Psi_n(x_1+c_1,\dots,x_n+c_n;h(x)),\qquad c_a=\kappa_1+\dots+\kappa_a.
\end{equation*}
\end{itemize}

\textbf{Step 19} ($D_1$ is a core).
Let $s>0$.
The operator $(1+H)^me^{-sH}$ is bounded, self-adjoint and injective, and hence has dense range; it therefore carries the dense set $D_1$ onto a dense set, which is contained in $(1+H)^mD_1$.
It follows that $D_1$ is dense in $D(H^m)\supset D$ for the norm $\|\cdot\|_m$.

\textbf{Step 20} (covariance; Theorem~3.3\,(b)).
Apply $\mathcal U(0,\Lambda)$ to (e4), change variables by $x\to\Lambda^{-1}x$, and let $\kappa\to0$.
This gives $\mathcal U(0,\allowbreak \Lambda)F(h)=\allowbreak F(\Lambda\cdot h)\,\allowbreak \mathcal U(0,\allowbreak \Lambda)$ on $D_1$.
Both sides are bounded from $\|\cdot\|_4$ to $\mathcal H$: the left one by Theorem~3.3\,(a), the right one by Theorem~3.3\,(a) together with the bound $\|\mathcal U(0,\Lambda)\psi\|_4\le(2(\Lambda e_0)^0)^4\|\psi\|_4$ of Step~17.
By Step~19 the identity therefore extends to $D$.
For translations the statement is (e3).

\textbf{Step 21} (locality: a formula).
Suppose that $h$ and $h'$ have compact supports which are space-like separated.
Put $\Psi'=\Psi_m(z';b')$, $\Psi=\Psi_n(z;b)$, $N=m+n+2$, and
\begin{equation*}
p_\epsilon:=(\bar z'_m,\dots,\bar z'_1,x,x'+i\epsilon e_0,z_1+i\epsilon e_0,\dots,z_n+i\epsilon e_0),
\end{equation*}
which lies in $\mathcal T_N$ when $\epsilon>0$.
Two applications of (e4), together with Step~6\,(b), (c), give
\begin{equation*}
\langle\Psi',F(h)T(i\epsilon e_0)F(h')\Psi\rangle=\iint dx\,dx'\,\mathcal W_N(p_\epsilon;\bar b'_m,\dots,\bar b'_1,h(x),h'(x'),b_1,\dots,b_n).
\end{equation*}
As $\epsilon\to0$, the left-hand side converges to $\langle\Psi',F(h)F(h')\Psi\rangle$.

\textbf{Step 22} ($p_0\in\mathcal D_N$; Theorem~3.3\,(c)).
With one exception, the consecutive differences of the entries of $p_0$ belong to $\mathcal T$; the exception is the real vector $w=x'-x$, which satisfies $|\vec w|>|w^0|$.
Choose a rotation $\varrho$ taking $\vec w$ to $|\vec w|e_2$.
Then
\begin{equation*}
\operatorname{Im}(JR^{(2)}_tJ^{-1}\varrho\,w)=-\sin t\,(|\vec w|,w^0,0)\in V_+
\end{equation*}
for small $t<0$, whereas, $\mathcal T$ being open, the images of the other differences belong to $\mathcal T$ for small $|t|$; therefore $p_0\in\mathcal D_N$.
When $x$ ranges over $\operatorname{supp}h$, $x'$ over $\operatorname{supp}h'$ and $\epsilon$ over $[0,1]$, the points $p_\epsilon$ form a compact subset of $\mathcal D_N$, and $\mathcal W_N$ is continuous there.
It follows that $\langle\Psi',F(h)F(h')\Psi\rangle=\iint\mathcal W_N(p_0;\dots)$, an expression which, by (f2), does not change when $(x,h(x))$ and $(x',h'(x'))$ are interchanged.
The relation $[F(h),F(h')]=0$ extends to $D$ by Step~19 and Theorem~3.3\,(a), and to supports that are not compact by means of cut-offs.

\textbf{Step 23} ($dF=0$; Theorem~3.3\,(d)).
Put $G_b(a):=\Phi_b(ide_0+a,ide_0-a)$.
This function is holomorphic in the convex region $\{|\operatorname{Im}a^0|+|\operatorname{Im}\vec a|<d\}$, which includes $\mathbb R^3$ as well as every point $Jy$ with $|y_1|<d$.
Let $u\in C_c^\infty(\{|y_1|<d\})$ and
\begin{equation*}
g=\partial_1u\,e_2\wedge e_3+\partial_2u\,e_3\wedge e_1+\partial_3u\,e_1\wedge e_2,
\end{equation*}
so that $\ast g=\nabla u$; then (dF) yields $\mathcal B_d(g)=0$.
Combining Step~4\,(d), an integration by parts and the relation $\partial_{y_1}[G(Jy)]=i(\partial_0G)(Jy)$, one deduces
\begin{equation*}
i\,(\partial_0G_{e_2\wedge e_3}+\partial_2G_{e_3\wedge e_0}+\partial_3G_{e_0\wedge e_2})(Jy)=0,
\end{equation*}
where a single factor $i$ is common to the three terms.
Fact~2 shows that the bracket is zero on $\mathbb R^3$.
Hence $F_{d,d}(d^*u)=0$, and so $q_{d^*u}=0$, for every $u\in\mathcal S$.

\textbf{Step 24} (cyclicity; Theorem~3.3\,(e)).
Suppose that $\psi$ is orthogonal to $F(h_1)\cdots F(h_n)\Omega$ for all choices of the $h_a$, and consider
\begin{equation*}
G(c):=\langle\psi,T(c_1)F(h_1)\cdots T(c_n)F(h_n)\Omega\rangle.
\end{equation*}
For real $c$ this vanishes, by (e3).
Step~1 and Theorem~3.3\,(a) show that $G$ fulfils the assumptions of Fact~3 on $\overline{\mathcal T}^{\,n}$, whence $G\equiv0$.
Letting $h_a\to\delta_0b_a$ in (e4), one finds $\psi\perp D_1$, and therefore $\psi=0$.

\textbf{Step 25} (Wightman distributions; (f1), (f3)).
Theorem~3.3\,(a) gives, for $\zeta_a\in\overline{\mathcal T}$,
\begin{equation*}
|\langle\Omega,F(h_1)T(\zeta_1)F(h_2)\cdots T(\zeta_{n-1})F(h_n)\Omega\rangle|\le\prod_a|h_a|_{4n}.
\end{equation*}
Fact~5 then shows that $W_n$ is continuous on $\mathrm{Sw}_n$.
Fix $\eta$ with $\eta_{a+1}-\eta_a\in V_+$, and fix a norm $|c|$ on $(\Lambda^2\mathbb C^3)^{\otimes n}$.
By (f1), $|\mathcal W_n(x+i\eta;c)|\le C_\eta|c|$ for real $x$ and $c\in(\Lambda^2\mathbb C^3)^{\otimes n}$, with a constant $C_\eta$ depending on $\eta$ and $n$; hence the functional $\mathbf h\mapsto\int dx\,\mathcal W_n(x+i\eta;\mathbf h(x))$ on $\mathrm{Sw}_n$ is bounded by a multiple of $\|\mathbf h\|_{L^1}$, and it is continuous.
By (e4), its value at $h_1\otimes\dots\otimes h_n$ is the above form at $\zeta_a=i(\eta_{a+1}-\eta_a)$, an $n$-linear form to which Fact~5 applies with $p=|\cdot|_{4n}$.
By the uniqueness clause of Fact~5, the functional is the continuous linear form $W$ furnished there, and it is therefore bounded by the seminorm $P$, which does not depend on $\eta$.
Since $T(i\eta)\to1$ strongly for every norm $\|\cdot\|_m$, these functionals tend to $W_n$ on products of tests, and consequently on all of $\mathrm{Sw}_n$.
For (f3), observe that by (e3) the forms at real $\zeta_a$ equal $W_n(h_1,\tau_{\zeta_1}h_2,\tau_{\zeta_1+\zeta_2}h_3,\dots)$.
Arguing with Fact~3 as in Step~24, one concludes that the $W_n$ determine these forms on $\overline{\mathcal T}^{\,n-1}$; when two laws are compared, Fact~3 is applied to the difference of their forms.
Therefore the $W_n$ determine $\mathcal W_n$ on $\mathcal T_n$, by (e4) with $h_a\to\delta_0b_a$, and on $\mathcal D_n$, by the definition of the $\mathcal W_n^{R,\pi}$; through (f2) they determine the $S_n$ at pairwise disjoint compact supports.
Proposition~1.2\,(c) completes the argument.
\end{proof}

\begin{named}{Corollary 3.4}
\emph{$F$ is not a generalised free field.}
Take the tuned data of \cite[Theorem~44]{QEDIII}, a charge $0<e\le e_2$, and $\mathrm{lw}=\mu_\xi$ with $\xi\in\Xi_\mu$.
Define $W_n^{\mathrm{tr}}$ in terms of the $W_k$ through the formula of \cite[(n5)]{QEDIII} for $\kappa_n$, keeping the order of the arguments within each block.
Then $W_4^{\mathrm{tr}}\ne0$; in particular, it is not the case that $W_n^{\mathrm{tr}}=0$ for every $n\ne2$.
Furthermore, given non-empty open sets $O_1,\dots,O_4\subset\mathbb R^3$ which are pairwise disjoint, one can choose real $f_a\in C_c^\infty(O_a;\Lambda^2)$ such that $\kappa_4^\xi(f_1,\dots,f_4)\ne0$.
\end{named}

\begin{proof}
\textbf{Part 1.}
Since $\mathcal W_1=0$, the truncated function is $\mathcal W_4^{\mathrm{tr}}:=\allowbreak \mathcal W_4-\sum\mathcal W_2\mathcal W_2$, the sum being over the three pairings.
As $\mathcal T$ is a convex cone, $\mathcal W_4^{\mathrm{tr}}$ is holomorphic on $\mathcal D_4$; and (f2) identifies it, at pairwise disjoint compact supports, with the density of $\kappa_4^\xi$.

\textbf{Part 2.}
The positive fourth cumulant of \cite[Theorem~44(D6)]{QEDIII}, together with Proposition~1.2(a), gives disjoint compact supports at which $\kappa_4^\xi\ne0$.
The function $y\mapsto\mathcal W_4^{\mathrm{tr}}(Jy;\cdot)$ is real analytic on the set of configurations of four distinct points, which is connected; consequently it does not vanish on any non-empty open subset.
This proves the last assertion, and shows that $\mathcal W_4^{\mathrm{tr}}$ is not identically zero on $\mathcal T_4$.

\textbf{Part 3.}
Let $c\in\mathcal T^3$.
By (e4), the integral
\begin{equation*}
\int dx\,\mathcal W_4^{\mathrm{tr}}(x_1,x_2+c_1,x_3+c_1+c_2,x_4+c_1+c_2+c_3;h(x))
\end{equation*}
can be written as a polynomial in the matrix elements $\langle\Omega,F(h_a)T(\cdot)F(h_{a'})\cdots\Omega\rangle$.
As a function of $c$ it extends to $\overline{\mathcal T}^{\,3}$ in the manner required by Fact~3, and at real $c$ its value is $W_4^{\mathrm{tr}}(h_1,\allowbreak \tau_{c_1}h_2,\allowbreak \tau_{c_1+c_2}h_3,\allowbreak \tau_{c_1+c_2+c_3}h_4)$.
If $W_4^{\mathrm{tr}}$ vanished, it would follow, on letting $h_a\to\delta_0b_a$, that $\mathcal W_4^{\mathrm{tr}}=0$ on $\mathcal T_4$.
\end{proof}

The Euclidean fourth-cumulant bound used here is part of the tuned
construction. Corollary~3.4 gives its Minkowski continuation and its
consequence for arbitrary disjoint open sets.

\begin{named}{Proposition 3.5}
\emph{Discrete symmetries and the parity-even field-strength two-point function.}
Let $\mathrm{lw}=\mu_\xi$ be the law of Proposition~3.1.
\begin{itemize}
\item[(a)] \emph{Component reflection invariance, recalled from (Inv).}
Every reflection $r$ of $\mathbb R^3$ satisfies $r_*\mu_\xi=\mu_\xi$.
In other words, although the reflections, under which $\mu$ is invariant, could permute the components of $\mu$, each component with $\xi\in\Xi_\mu$ is fixed by them.
Denoting by $\mathrm{Rl}$ the linear part of $r$, the action of $r$ on $b=\ast f$ is $(r\cdot b)(y)=-\mathrm{Rl}\,b(r^{-1}y)$; thus the dual field is a pseudovector.
\item[(b)] \emph{Implementation.}
The representation $\mathcal U$ extends to a representation of $\mathbb R^3\rtimes O(1,2)$ on $\mathcal H$ fixing $\Omega$.
If $(\Lambda e_0)^0>0$, that is, if $\Lambda$ preserves the direction of time, then $\mathcal U(a,\Lambda)$ is unitary and
\begin{equation*}
\mathcal U(a,\Lambda)F(h)\,\mathcal U(a,\Lambda)^{-1}=F((a,\Lambda)\cdot h);
\end{equation*}
otherwise $\mathcal U(a,\Lambda)$ is antiunitary and the right-hand side is replaced by $F(\overline{(a,\Lambda)\cdot h})$.
When $R\in O(3)$ satisfies $Re_1=e_1$, which on Minkowski space reads $Re_0=e_0$, the operator $\mathcal U(0,R)$ is $[A]\mapsto[\hat RA]$.
For a reflection $R$ of space one has $\mathcal U(0,R)F_{23}(x)\,\mathcal U(0,R)^{-1}=-F_{23}(Rx)$; that is, under parity the magnetic field changes sign.
\item[(c)] \emph{$\mathsf{CT}$, not $\mathsf T$.}
One has $\mathcal U(0,\theta)=\Gamma$, where $\Gamma[A]:=[\bar A]$; moreover $\Gamma H\Gamma=H$, $\Gamma P_i\Gamma=-P_i$ and $\Gamma F_\alpha(x)\Gamma=\allowbreak (-1)^{\mathrm{ti}(\alpha)}F_\alpha(\theta x)$.
Thus $F_{23}$ is even and the electric field $(F_{02},F_{03})$ is odd.
This transformation is $\mathsf{CT}$ and not the time reversal $\mathsf T$, since $\mathsf T$ reverses the magnetic field.
\item[(d)] \emph{Two-point function.}
Let $f',f''$ be real, and put $b'=\ast f'$ and $b''=\ast f''$.
Then
\begin{equation*}
S_2^\xi(f',f'')=(2\pi)^{-3}\int dp\,\overline{\hat b'(p)}\cdot M(p)\,\hat b''(p),\qquad M(p)=\Pi(p)\int d\mathrm{kl}_\xi(t)\,\frac{|p|^2}{|p|^2+t^2}.
\end{equation*}
Symmetry, boundedness on $L^2$, (dF) and invariance under $\mathbb R^3\rtimes SO(3)$, and these alone, yield a bounded measurable $M$, by \cite[proof of Theorem~14.26, Step~7]{QEDII}; they also permit, for almost every $p$,
\begin{equation*}
M(p)=M_+(|p|)\,\Pi(p)+M_-(|p|)\,\epsilon(p),\qquad \epsilon(p)_{ab}:=\sum_c\epsilon_{abc}p_c/|p|,
\end{equation*}
where $M_-$ is purely imaginary in the case of a law.
The term containing $\epsilon(p)$ changes sign under reflections.
For $\mu_\xi$ this term does not occur.
\item[(e)] \emph{Any order of the limits; Wilson fermions.}
Consider the laws of \cite[Theorem~1.3\,(a)]{QEDII}, which include those of \cite[Theorem~44\,(B)]{QEDIII} since $P_{N,T}$ is the law $P^{\mathrm{ch}}_{N,T}$ of \cite[Definition~2.36]{QEDII}, by \cite[Section~I.7.5]{QEDIII}, and the laws of \cite[Theorem~3.1, Part~I\,(b), (c)]{QEDI}, for an arbitrary charge $e>0$ and an arbitrary mass assignment.
Each limit point of either family, with the limits in the cutoff $(N,T)$ taken in an arbitrary manner, is $O_{\mathrm{cub}}$-invariant, by \cite[Theorem~3.1, Part~I\,(b), (P5)]{QEDI} and \cite[Theorem~1.3\,(a)]{QEDII}; and $\nu$-almost all of its components are invariant under the group $\{\pm1\}^3\ni\theta_j,-1$, which is the group $N_3$ of \cite[Section~2.5]{QEDI}, by \cite[Theorem~3.1, Part~I\,(c)]{QEDI} and \cite[Theorem~1.3\,(a)]{QEDII}.
Iterated limits, in either order, are among these limit points, by \cite[Theorem~3.1, Part~I\,(d)]{QEDI} and the last clause of \cite[Theorem~1.3\,(a)]{QEDII}.
Hence for these laws the matrices $M(p)$, which exist by \cite[proof of Theorem~14.26, Step~7]{QEDII}, are real and symmetric for almost every $p$.
\item[(f)] \emph{Charge conjugation.}
Let $\mathrm{Cc}\colon\omega\mapsto-\omega$.
The image law $\mathrm{Cc}_*\mu_\xi$ is a component $\mu_{\xi'}$ as well, with $\xi'\in\Xi_\mu$, and its Schwinger functions are $(-1)^nS_n^\xi$; either it equals $\mu_\xi$, or the two laws are mutually singular.
The following three statements are equivalent.
\begin{itemize}
\item[(i)] $\mathrm{Cc}_*\mu_\xi=\mu_\xi$.
\item[(ii)] $W_n=0$ for every odd $n$.
\item[(iii)] Some unitary operator $\mathsf C$ on $\mathcal H$ satisfies $\mathsf C\Omega=\Omega$, $\mathsf CD=D$ and $\mathsf CF(h)\mathsf C^{-1}=-F(h)$.
\end{itemize}
One says that $\mathsf T$ is \emph{implemented} if some antiunitary operator $\mathsf T_0$ on $\mathcal H$ satisfies $\mathsf T_0\Omega=\Omega$, $\mathsf T_0D=D$ and $\mathsf T_0F(h)\mathsf T_0^{-1}=-F(\overline{\theta\cdot h})$, that is, maps $F_\alpha(x)$ to $-(-1)^{\mathrm{ti}(\alpha)}F_\alpha(\theta x)$; these signs are those of $\mathsf T$.
The operator $\Gamma$ is antiunitary, with $\Gamma^2=1$, $\Gamma\Omega=\Omega$, $\Gamma D=D$ and $\Gamma F(h)\Gamma=F(\overline{\theta\cdot h})$; hence multiplication by $\Gamma$ turns each such $\mathsf T_0$ into an operator $\mathsf C$ as in (iii), and conversely.
Consequently $\mathsf T$ is implemented if and only if (i), (ii), (iii) hold.
\end{itemize}
\end{named}

\begin{proof}
\textbf{Part 1.}
Assertion (a) is a restatement of (Inv); the action on $\ast f$ is described in \cite[Definition~14.1]{QEDII}.

\textbf{Part 2.}
Let $R\in O(3)$ satisfy $Re_1=e_1$.
Then $R$ commutes with $J$, maps $\mathcal T_n$ onto itself and respects the order of the times.
Hence (Inv), Step~7 and Fact~2 yield $\mathcal W_n(Rz;Rb)=\mathcal W_n(z;b)$ on $\mathcal T_n$.
At the Euclidean points, where $\theta\bar z=z$, the relation $\overline{S_n(g_1,\dots)}=S_n(\bar g_1,\dots)$, Step~7 and (K) give $\overline{\mathcal W_n(z;b)}=\allowbreak \mathcal W_n(\theta\bar z;\theta\bar b)$.
Both members are antiholomorphic on $\mathcal T_n$, because $-\theta V_+=V_+$; so the identity holds on all of $\mathcal T_n$.

\textbf{Part 3.}
By Fact~6\,(iv), the group $O(1,2)$ is generated by $SO_0(1,2)$, by one reflection of space and by $\theta$.
Let $\Lambda\in O(1,2)$.
If $\Lambda$ preserves the direction of time, Step~10, Part~2 and Step~6\,(c) show that $\Psi_n(z;b)\mapsto\Psi_n(\Lambda z;\Lambda b)$ extends to a unitary operator; if it reverses that direction, they show that $\Psi_n(z;b)\mapsto\Psi_n(\Lambda\bar z;\Lambda\bar b)$ extends to an antiunitary operator.
On $D_1$ one sees the group law directly, together with the relation $\mathcal U(0,\allowbreak \Lambda)T(a)\,\allowbreak \mathcal U(0,\allowbreak \Lambda)^{-1}=\allowbreak T(\Lambda a)$ for real $a$.
Applied with $a\in\mathbb Re_0$, this relation shows that $\mathcal U(0,\varrho)$ commutes with $H$ whenever $\varrho e_0=\pm e_0$, the operator being antiunitary in the case of the sign $-$.
By Fact~6\,(iv) one may write $\Lambda=\Lambda_0\varrho$ with $\Lambda_0\in SO_0(1,2)$ and $\varrho\in\{1,\theta_2,\theta,\theta_2\theta\}$; then $\Lambda e_0=\pm\Lambda_0e_0$, so that $(\Lambda_0e_0)^0=|(\Lambda e_0)^0|$.
As $\mathcal U(0,\allowbreak \Lambda)=\allowbreak \mathcal U(0,\allowbreak \Lambda_0)\,\allowbreak \mathcal U(0,\allowbreak \varrho)$ and $\mathcal U(0,\varrho)$ preserves each norm $\|\cdot\|_m$, the bound of Step~17, applied to $\Lambda_0$, gives
\begin{equation*}
\|\mathcal U(0,\Lambda)\psi\|_m\le\bigl(2|(\Lambda e_0)^0|\bigr)^m\|\psi\|_m .
\end{equation*}
Hence $\mathcal U(a,\Lambda)D=D$, and (e4) yields the action on $F(h)$ exactly as in Step~20.

\textbf{Part 4.}
For the $R$ of Part~2, Step~7 yields $\mathcal U(0,R)[B]=[\hat RB]$, and also $\mathcal U(0,\theta)[B]=[\bar B]$, for all $B$ in the total set of time-ordered products.
The map $[A]\mapsto[\hat RA]$ is unitary by \cite[Proposition~14.47\,(a)]{QEDII}, and $[A]\mapsto[\bar A]$ is antiunitary because $\hat\theta$ commutes with complex conjugation.
From $\Gamma T(a)\Gamma=T(\theta a)$ one infers $\Gamma H\Gamma=H$ and $\Gamma P_i\Gamma=-P_i$, and the signs are those of $\wedge^2R$ and of $\wedge^2\theta$.
The identities for the components $F_\alpha(x)$ in (b), (c) and (f) are identities between forms; they follow from the corresponding identities for $F(h)$ as $h$ tends to $\delta_xe_\alpha$, by the bound of Step~16.
In (b) and (c) they hold on $D(H^2)\times D(H^2)$, which $\mathcal U(0,R)$, for $Re_0=e_0$, and $\Gamma$ map onto itself, since these operators commute with $H$ by Part~3; in (f) they hold on $D\times D$, since $\mathsf T_0D=D$.

\textbf{Part 5.}
For (d), the formula is \cite[Theorem~44\,(D2) with (C7)]{QEDIII}.
Consider any law with (G), (dF) and invariance under translations.
For such a law, \cite[proof of Theorem~14.26, Step~7]{QEDII} furnishes a measurable Hermitian $M$ with $0\le M\le1$ and $M(p)p=0$ almost everywhere, which represents $S_2$ as in (d).
That step rests on \cite[Chapter~I, Theorem~3.18]{SteinWeiss1971}, according to which a bounded operator on $L^2(\mathbb R^n)$ commuting with all translations is a Fourier multiplier whose symbol is bounded and measurable.
For almost every $p$ one has
\begin{equation*}
M(-p)=M(p)^T,\qquad M(p)p=0,\qquad M(-p)=\overline{M(p)},
\end{equation*}
by symmetry, by (dF) and by reality respectively.
Likewise, for each $R\in O(3)$ leaving the law invariant, $M(Rp)=RM(p)R^T$ for almost every $p$, by invariance, with $\det R$ entering squared; here the exceptional null set depends on $R$.
Recall Tonelli's theorem: if $u\ge0$ is measurable on the product of two $\sigma$-finite measure spaces, then both iterated integrals of $u$ exist in $[0,\infty]$ and coincide with the integral of $u$ over the product.
Assume that the law is invariant under $SO(3)$.
Choose a bounded Borel representative $M$, equip $SO(3)\times\mathbb R^3$ with the product of Haar measure and Lebesgue measure, and let $u$ be the indicator function of $\{(R,p)\in SO(3)\times\mathbb R^3:M(Rp)\ne RM(p)R^T\}$.
For each fixed $R$ the corresponding section of this set is a null set.
By Tonelli's theorem, for almost every $p$ the identity $M(Rp)=RM(p)R^T$ holds for almost every $R$; consequently $\tilde M(p):=\int dR\,R^TM(Rp)R$ coincides with $M(p)$ for almost every $p$.
The right invariance of Haar measure gives $\tilde M(R'p)=R'\tilde M(p)R'^T$ for all $R'\in SO(3)$ and all $p$.
The relations $\tilde M(p)p=0$ and $\tilde M(-p)=\tilde M(p)^T=\overline{\tilde M(p)}$ are valid almost everywhere, on a set that is invariant under rotations; hence they are valid at each $p$ whose length $|p|$ avoids a certain null set of $(0,\infty)$.
At such $p$ the matrix $\tilde M(p)$ annihilates $p$ from both sides and commutes with the rotations about the axis $p$.
On $p^\perp$, the matrices commuting with these rotations are precisely the linear combinations of $\Pi(p)$ and $\epsilon(p)$.
Therefore $\tilde M(p)=M_+(p)\,\Pi(p)+M_-(p)\,\epsilon(p)$.
The exact covariance of $\tilde M$, together with $\epsilon(Rp)=\det R\cdot R\,\epsilon(p)R^T$, shows that $M_+$ and $M_-$ depend on $|p|$ only; and $M_-$ is purely imaginary, since $\tilde M(p)$ is Hermitian whereas $\epsilon(p)$ is real and antisymmetric.
With $\tilde M$ in place of $M$, a single improper $R$ forces $M_-=0$.

\textbf{Part 6.}
For (e), take $R=-1$ in Part~5: it gives $M(-p)=M(p)$ for almost every $p$, so that $M$ is symmetric as well as Hermitian, and therefore real.

\textbf{Part 7.}
For (f): the law $\mu$ is $\mathrm{Cc}$-invariant by \cite[Theorem~44\,(P3)]{QEDIII}, and $\mathrm{Cc}$ commutes with all Euclidean maps.
Therefore $\mathrm{Cc}_*$ maps ergodic laws to ergodic laws, and by the uniqueness statement of \cite[(n7)]{QEDIII} it leaves $\nu$ invariant.
The set $\Xi_\mu$ is defined by the $18$ statements (c1), (c2), (c3) of \cite[Theorem~44\,(B5), (D)]{QEDIII}, and each of them is invariant under $\mathrm{Cc}_*$: (G) and reflection positivity are requirements on all real tests, the condition $\{dF=0\}$ is linear, and the functions $\mathrm{Wk}_{m,\mathrm{ix}}=\allowbreak \mathrm{Pw}_m(X(g_{\mathrm{ix}});\|g_{\mathrm{ix}}\|^2,\allowbreak c^{\mathrm{inf}}_{m,\mathrm{ix}})$ of \cite[Theorem~44\,(B4)]{QEDIII} are even functions of $X(g_{\mathrm{ix}})$, because $\mathrm{Pw}_2$ and $\mathrm{Pw}_4$ are even in their first variable \cite[Section~I.7.6]{QEDIII}.
Two distinct ergodic laws $\mathrm{lw}'$ and $\mathrm{lw}''$ are mutually singular.
Indeed, let $\lambda$ be their mean; one-point sets of laws are measurable by \cite[Proposition~45.6\,(d0)]{QEDIII}, so that $\{\lambda_\omega=\mathrm{lw}'\}$ is an event, and by \cite[Proposition~45.6\,(a), (a$''$)]{QEDIII} this event separates $\mathrm{lw}'$ from $\mathrm{lw}''$.

\textbf{Part 8.}
Assume (i).
On $L^2(\mathrm{lw})$ the operator $A\mapsto A\circ\mathrm{Cc}$ is then unitary; it commutes with $\hat\theta$ and maps $L^2(\Sigma_+,\mathrm{lw})$ onto itself.
The operator $\mathsf C$ it induces on $\mathcal H$ satisfies $\mathsf C\Psi_n=(-1)^n\Psi_n$, by Step~7 and Fact~2, and (e4) yields the other properties required in (iii).
If (iii) holds, then $W_n=(-1)^nW_n$, which is (ii).
Finally assume (ii).
Step~25 gives $\mathcal W_n=0$ for odd $n$, so that $\mathrm{Cc}_*\mathrm{lw}$, a law with (G) and (dF), has the same moments as $\mathrm{lw}$ at disjoint supports; Proposition~1.2\,(c) then yields (i).
\end{proof}

The next proposition is the Reeh--Schlieder consequence of positive
energy and locality \cite{ReehSchlieder1961}. We include its proof for
the common domain $D$ used in Theorem~3.3.

\begin{named}{Proposition 3.6}
\emph{The property of Reeh and Schlieder.}
Let $\mathrm{lw}$ be a law having the properties listed in Proposition~3.1, and let $\mathcal H$, $\Omega$, $D$ and $F$ be attached to it as in Theorems~3.2 and~3.3.
\begin{itemize}
\item[(a)] \emph{Cyclicity.}
Let $O$ be a non-empty open subset of Minkowski space $\mathbb R^3$.
The vector $\Omega$, together with the vectors $F(h_1)\cdots F(h_n)\Omega$ for which $n\ge1$ and each $h_a$ lies in $C_c^\infty(O;\Lambda^2\mathbb C^3)$, spans a dense subspace of $\mathcal H$.
\item[(b)] \emph{Separation.}
Let $O_1\subset\mathbb R^3$ be an open set for which there exists a non-empty open set $O_2$ such that $(x-x')\cdot(x-x')<0$ whenever $x\in O_1$ and $x'\in O_2$; every bounded open set $O_1$ is of this kind.
Consider a finite sum $P=c_0+\sum_kF(h_{k,1})\cdots F(h_{k,n_k})$, where $c_0\in\mathbb C$, each $h_{k,a}$ belongs to $\mathrm{Sw}_1$, and $\operatorname{supp}h_{k,a}\subset O_1$.
If $P\Omega=0$, then $P\chi=0$ for all $\chi\in D$.
\end{itemize}
\end{named}

\begin{proof}
Write $|x|_{\mathrm E}:=((x^0)^2+|\vec x|^2)^{1/2}$ for $x\in\mathbb R^3$.

\textbf{Part 1} (the function $G$).
Suppose that $\psi\in\mathcal H$ is orthogonal to $\Omega$ and to all the vectors described in (a).
Fix $x_0\in\mathbb R^3$ and $r>0$ with $\{|x-x_0|_{\mathrm E}<2r\}\subset O$, let $n\ge1$, take $h_1,\dots,h_n$ in $C_c^\infty(\{|x-x_0|_{\mathrm E}<r\};\Lambda^2\mathbb C^3)$, and define, as in Step~24,
\begin{equation*}
G(c):=\langle\psi,T(c_1)F(h_1)\cdots T(c_n)F(h_n)\Omega\rangle\qquad(c\in\overline{\mathcal T}^{\,n}).
\end{equation*}
According to Step~1, for each norm $\|\cdot\|_m$ the operators $T(\zeta)$ are contractions that depend strongly continuously on $\zeta\in\overline{\mathcal T}$ and holomorphically on $\zeta\in\mathcal T$; according to Theorem~3.3\,(a), $F(h)$ is a bounded map from the norm $\|\cdot\|_{m+4}$ to the norm $\|\cdot\|_m$.
It follows that, in each variable and for arbitrary values of the other variables in $\overline{\mathcal T}$, the function $G$ is continuous on $\overline{\mathcal T}$ and holomorphic on $\mathcal T$.

\textbf{Part 2} (real arguments).
For real $c$, the rule (e3) of Step~18 and the equality $T(a)\Omega=\Omega$ of Step~1 give
\begin{equation*}
G(c)=\langle\psi,F(\tau_{c_1}h_1)F(\tau_{c_1+c_2}h_2)\cdots F(\tau_{c_1+\dots+c_n}h_n)\Omega\rangle.
\end{equation*}
If every $c_a$ satisfies $|c_a|_{\mathrm E}<r/n$, the supports of the tests occurring in this expression lie in $\{|x-x_0|_{\mathrm E}<2r\}\subset O$, and therefore $G(c)=0$.

\textbf{Part 3} (continuation).
Fact~3, applied with the balls $B_a=\{|c_a|_{\mathrm E}<r/n\}$, shows that $G=0$ on $\overline{\mathcal T}^{\,n}$.

\textbf{Part 4} (point fields).
Let $\kappa\in\mathcal T^n$ and $z_a:=\kappa_1+\dots+\kappa_a$.
The rule (e4) of Step~18 expresses $G(\kappa)$ as
\begin{equation*}
G(\kappa)=\int dx\,\langle\psi,\Psi_n(x_1+z_1,\dots,x_n+z_n;h(x))\rangle.
\end{equation*}
Choose $h_a:=\phi_\epsilon(\cdot-x_0)\,b_a$ with $b_a\in\Lambda^2\mathbb C^3$ and $0<\epsilon<r$, where $\phi_\epsilon\ge0$ belongs to $C_c^\infty(\{|x|_{\mathrm E}<\epsilon\})$ and $\int\phi_\epsilon=1$.
Since $\Psi_n$ is norm continuous on $\mathcal T_n^+$ and $n$-linear in $b$ (Steps~5 and~6), the limit $\epsilon\to0$ yields
\begin{equation*}
\langle\psi,\Psi_n(x_0+z_1,\dots,x_0+z_n;b)\rangle=0,
\end{equation*}
and $z=(z_1,\dots,z_n)$ ranges over $\mathcal T_n^+$ as $\kappa$ ranges over $\mathcal T^n$.

\textbf{Part 5} (density).
Step~6\,(b) gives $\Psi_n(x_0+z_1,\dots,x_0+z_n;b)=T(x_0)\Psi_n(z;b)$, and $T(x_0)\Omega=\Omega$.
Thus $T(x_0)^*\psi$ is orthogonal to $D_1$.
The operator $T(x_0)$ being unitary and the subspace $D_1$ being dense by Step~11, one concludes that $\psi=0$, and (a) is proved.

\textbf{Part 6} (a space-like separated set).
If $O_1\subset\{|x|_{\mathrm E}<r'\}$ for some $r'>0$, then, with the basis vector $e_2$ of Minkowski space, the ball $O_2:=\{|x-4r'e_2|_{\mathrm E}<r'\}$ has the property required in (b), since $|x^0-x'^0|<2r'<|\vec x-\vec x'|$ for $x\in O_1$ and $x'\in O_2$.
In the general case, let $O_2$ be as in (b), and let $Q$ be either $1$ or a product $F(h'_1)\cdots F(h'_k)$ with $h'_a\in C_c^\infty(O_2;\Lambda^2\mathbb C^3)$.
By Theorem~3.3\,(c) one has $PQ=QP$ on $D$, a subspace which each of these operators maps into itself; applied to $\Omega$, this yields $PQ\Omega=QP\Omega=0$.

\textbf{Part 7} (the adjoint).
Let $P^+:=\bar c_0+\sum_kF(\bar h_{k,n_k})\cdots F(\bar h_{k,1})$.
Theorem~3.3\,(a) shows that $\langle\phi,P\chi\rangle=\langle P^+\phi,\chi\rangle$ for $\phi,\chi\in D$.
Hence $\langle P^+\phi,Q\Omega\rangle=\langle\phi,PQ\Omega\rangle=0$ for every $Q$ of Part~6, and (a), applied to the set $O_2$, gives $P^+\phi=0$.
Consequently $\langle\phi,P\chi\rangle=0$ for all $\phi,\chi\in D$, and the density of $D$ yields $P\chi=0$.
\end{proof}


\section{The spectral measure and the Maxwell infrared limit}


\subsection{The spectral multiplier}
For $\lambda\in\mathcal K$, let $\mathrm{kl}$ be the measure in
Theorem~A(ii) and put
\[
\mathsf D(s)=\int\frac{s}{s+t^2}\,d\mathrm{kl}(t),\qquad
Y_t(y)=\frac{e^{-t|y|}}{4\pi|y|}.
\]
The measure variable $t$ is mass. The two-point kernel off the origin is
$(-\delta_{ab}\Delta+\partial_a\partial_b)\int Y_t\,\allowbreak d\mathrm{kl}(t)$
after Hodge duality \cite[Theorem~14.26]{QEDII}.
\begin{named}{Proposition 4.1}
Let $0\ne g\in L^2$ be real and closed, and put
$g_l(y)=l^{-3/2}g(y/l)$. Then
\[
\frac{S_2(g_l,g_l)}{\|g\|^2}\downarrow\mathrm{kl}(\{0\})
\quad(l\to\infty),\qquad
\frac{S_2(g_l,g_l)}{\|g\|^2}\uparrow\mathrm{kl}([0,\infty))
\quad(l\downarrow0).
\]
In particular these two masses are the infimum and supremum,
respectively, of $S_2(g,g)/\|g\|^2$ over nonzero real closed tests.
\end{named}
\begin{proof}
The Fourier transform of $g_l$ is $l^{3/2}\hat g(lp)$.
The representation in Theorem~A(ii) therefore gives
\[
\frac{S_2(g_l,g_l)}{\|g\|^2}
=\int\mathsf D(|p|^2/l^2)\,
\frac{|\hat g(p)|^2\,dp}{(2\pi)^3\|g\|^2}.
\]
The last factor is a probability measure. The function $\mathsf D$ is
nondecreasing, bounded between the atom at zero and the total mass,
and tends to those two numbers at zero and infinity. Monotonicity
and dominated convergence give the assertion.

\end{proof}
\subsection{The magnetic field at a fixed time}
\begin{named}{Proposition 4.2}
\emph{Notation.}
Throughout this proposition the letter $t$ stands for a time, while $s$ serves as the variable of $\mathrm{kl}$.
Momenta are written $p=(p_1,k)$ with $k\in\mathbb R^2$, and $\omega_s(k):=(|k|^2+s^2)^{1/2}$.
Section~2 supplies the conditions $(T_{123})$, $(R_1)$, $(RP_1)$, the reflection $\theta_1$ in the plane $\{x_1=0\}$, its operator $\hat\theta_1$ on functions on $\mathcal S'$, the antilinear map $\Theta_1$ given by $\Theta_1\Phi=\overline{\hat\theta_1\Phi}$, the reconstruction $(\mathcal H_1,[\,\cdot\,]_1,\mathsf H_1,V_1,\Omega_1)$ along $e_1$, and the joint spectral measure $\mathfrak E_1$ of this reconstruction, written in the variables $(\omega,p_\perp)$; see \cite[Definition~2.38; Lemma~12.18\,(b)]{QEDII}.
For a real $h\ne0$ in $C_c^\infty(\mathbb R^2)$ put
\begin{equation*}
\|h\|_{1/2}^2:=\int\frac{d^2k}{(2\pi)^2}\,|k|\,|\hat h(k)|^2.
\end{equation*}
Given a real number $t$, let the real closed form $g_{t,h}$ in $L^2(\mathbb R^3;\Lambda^2)$ be determined by
\begin{equation*}
\hat g_{t,h}(p)=\exp(-ip_1t)\,\hat h(k)\,\Pi(p)\,e_1,
\end{equation*}
so that $\|g_{t,h}\|^2=\frac12\|h\|_{1/2}^2$, and put $B_t(h):=X(g_{t,h})$.
Fix an even function $\varrho\ge0$ in $C_c^\infty(\mathbb R)$ with $\int\varrho=1$, set $\varrho_\varepsilon(x_1):=\allowbreak \varepsilon^{-1}\varrho(x_1/\varepsilon)$, and define the two-form $u_{\varepsilon,t}$ by $\ast u_{\varepsilon,t}(x)=\allowbreak \varrho_\varepsilon(x_1-t)\,\allowbreak h(x_2,\allowbreak x_3)\,\allowbreak e_1$; among its components only the component $23$ differs from zero.
The letter $\lambda$ stands, in (a), (b) and (c), for a probability measure on $(\mathcal S',\Sigma)$.
\begin{itemize}
\item[(a)] Suppose that $\lambda$ satisfies (G) and (dF).
As $\varepsilon\to0$, one has $X(u_{\varepsilon,t})\to B_t(h)$ in $L^p(\lambda)$ for every $p<\infty$.
In this sense, $B_t(h)$ is the smearing with $h$, at the time $t$, of the component $F_{23}$, that is, of the magnetic field of the picture $e_1$.
Moreover $\lambda(\exp B_t(h))\le\exp\bigl(\frac14\|h\|_{1/2}^2\bigr)$.
\item[(b)] If, besides (G) and (dF), the measure $\lambda$ satisfies $(T_{123})$, $(R_1)$ and $(RP_1)$, then $B_0(h)$ belongs to both $L^2(\Sigma_{\{x_1>0\}},\lambda)$ and $L^2(\Sigma_{\{x_1<0\}},\lambda)$, and $\Theta_1B_0(h)=B_0(h)$.
The symbol $\mathcal E_0$ denotes the $L^2(\lambda)$-closure of the linear space generated by all bounded Borel functions of a finite number of the variables $B_0(h)$.
Restricted to $\mathcal E_0$, the class map is an isometry: $\|[\Phi]_1\|=\|\Phi\|_{L^2(\lambda)}$.
\item[(c)] Suppose moreover that $S_2$ has the form (C7) of \cite[Theorem~44]{QEDIII}, with some positive measure $\mathrm{kl}$ whose mass does not exceed $1$, and set $\psi_h:=[B_0(h)]_1$.
For $t\ge0$ and $y\perp e_1$,
\begin{equation}
\begin{split}
&\bigl\langle\psi_h,\exp(-t\,\mathsf H_1)\,V_1(y)\,\psi_h\bigr\rangle\\
&\qquad=\int d\mathrm{kl}(s)\int\frac{d^2k}{(2\pi)^2}\,|\hat h(k)|^2\,|k|^2\,\frac{\exp(-t\,\omega_s(k)+ik\cdot y)}{2\,\omega_s(k)}.
\end{split}
\tag{4.1}\label{eq:magnetic-semigroup}
\end{equation}
Consequently $\langle\psi_h,\mathfrak E_1(\cdot)\,\psi_h\rangle$ coincides with the image of the measure
\begin{equation*}
\mathrm{kl}(ds)\,\frac{|\hat h(k)|^2\,|k|^2}{2\,\omega_s(k)}\,\frac{d^2k}{(2\pi)^2}
\end{equation*}
under the map $(s,k)\mapsto(\omega_s(k),k)$.
Furthermore, $\psi_h$ is in the domain of the quadratic form of $\mathsf H_1$, and
\begin{equation}
2\,\langle\psi_h,\mathsf H_1\,\psi_h\rangle=\mathrm{kl}([0,\infty))\,\|\nabla h\|_{L^2(\mathbb R^2)}^2.
\tag{4.2}\label{eq:magnetic-energy}
\end{equation}
\item[(d)] For a law satisfying (a)--(c), suppose that the cone $\{\omega\ge|p_\perp|\}$ carries its measure $\mathfrak E_1$ \cite[Theorem~44\,(P1), (P4)--(P6), (C4), (C7); (D1), (D2)]{QEDIII}.
The mass operator is defined as $\mathsf M:=\int(\omega^2-|p_\perp|^2)^{1/2}\,d\mathfrak E_1$.
In this situation $\psi_h\perp\Omega_1$, the operator $\mathsf M$ is reduced by the subspace $\Omega_1^\perp$, and $\mathsf M$ has in $\psi_h$ the spectral measure $c_h(s)\,\mathrm{kl}(ds)$, with
\begin{equation*}
c_h(s):=\int\frac{d^2k}{(2\pi)^2}\,\frac{|\hat h(k)|^2\,|k|^2}{2\,\omega_s(k)}\in(0,\infty).
\end{equation*}
Consequently the spectrum of $\mathsf M$ on $\Omega_1^\perp$ contains the support of $\mathrm{kl}$;
one has $\|\mathfrak E_1(\{\omega=|p_\perp|\})\,\allowbreak \psi_h\|^2=\allowbreak \frac12\|h\|_{1/2}^2\,\allowbreak \mathrm{kl}(\{0\})$;
\end{itemize}
\end{named}

\begin{proof}
\begin{enumerate}
\item[1.] \emph{Assertion (a).}
One has
\begin{gather*}
\hat u_{\varepsilon,t}(p)=\exp(-ip_1t)\,\hat\varrho(\varepsilon p_1)\,\hat h(k)\,e_1,\qquad
e_1\cdot\Pi(p)\,e_1=\frac{|k|^2}{|p|^2},\\
\int\frac{dp_1}{2\pi}\,\frac{1}{p_1^2+|k|^2}=\frac{1}{2|k|}.
\end{gather*}
These formulas yield the norm of $g_{t,h}$.
The reality of $h$ implies that of $g_{t,h}$, and $\Pi^2=\Pi$ shows that $g_{t,h}$ is closed.
As $|\hat\varrho|\le1$ and $\hat\varrho(\varepsilon p_1)\to1$, the convergence $\Pi_{\mathrm c}u_{\varepsilon,t}\to g_{t,h}$ in $L^2$ follows by dominated convergence.
If (G) and (dF) hold, then every real $f'\in L^2$ satisfies $X(f')=X(\Pi_{\mathrm c}f')$ almost surely and $\lambda(\exp X(f'))\le\exp\bigl(\frac12\langle f',\allowbreak \Pi_{\mathrm c}f'\rangle\bigr)$ \cite[Proposition~45.1\,(4), (11)]{QEDIII}, \cite[Remark~14.4]{QEDII}.
Moreover $\|X(f')\|_{L^N(\lambda)}\le c_N\,\|f'\|$ for $N=1,2,\dots$ \cite[Lemma~14.3\,(a)]{QEDII}.
\item[2.] \emph{Assertion (b).}
When $t>0$ and $\varepsilon$ is small, $u_{\varepsilon,t}$ is zero on $\{x_1<0\}$; hence $X(u_{\varepsilon,t})\in L^2(\Sigma_{\{x_1>0\}},\lambda)$ \cite[Lemma~14.3\,(a)]{QEDII}.
This subspace is closed.
By (a) it therefore contains $B_t(h)$; it contains $B_0(h)$ as well, since dominated convergence gives $\|g_{t,h}-g_{0,h}\|\to0$ as $t\to0$.
The same holds for the half space $\{x_1<0\}$.

With respect to $\theta_1$, the component of $\ast f'$ in the direction normal to $\{x_1=0\}$ is even \cite[Section~14.1, after Definition~14.1]{QEDII}, and the function $\varrho$ has the same parity.
Hence $\theta_1\cdot u_{\varepsilon,0}=u_{\varepsilon,0}$; as $u_{\varepsilon,0}$ is real, $\Theta_1X(u_{\varepsilon,0})=\allowbreak X(\theta_1\cdot u_{\varepsilon,0})=\allowbreak X(u_{\varepsilon,0})$ \cite[Section~14.1, after Definition~14.1]{QEDII}.
By $(R_1)$ the map $\Theta_1$ is isometric on $L^2(\lambda)$, and the limit $\varepsilon\to0$ gives $\Theta_1B_0(h)=B_0(h)$.

Further, $\hat\theta_1$ preserves products, and it commutes with the operation of composing with a function \cite[Section~8.1.2]{QEDI}; under (G) and $(R_1)$ it maps functions that agree $\lambda$-almost surely to functions that agree $\lambda$-almost surely \cite[Section~14.1, after Lemma~14.3]{QEDII}.
It follows that $\hat\theta_1\Phi=\Phi$ for every bounded Borel function $\Phi$ of finitely many $B_0(h)$, and thus on $\mathcal E_0$; moreover $\mathcal E_0\subset L^2(\Sigma_{\{x_1>0\}},\lambda)$.
Therefore $\|[\Phi]_1\|^2=\allowbreak \lambda((\Theta_1\Phi)\cdot\Phi)=\allowbreak \lambda(|\Phi|^2)$.
\item[3.] \emph{Assertion (c).}
Combining (b) with the definition of the reconstruction and with the relation $U(a)X(f')=X(\tau_a\cdot f')$ \cite[Lemma~14.3\,(a)]{QEDII}, one finds that the left side of \eqref{eq:magnetic-semigroup} equals
\begin{equation*}
\lambda\bigl(B_0(h)\cdot X(\tau_{te_1+y}\cdot g_{0,h})\bigr)=S_2(g_{0,h},\tau_{te_1+y}\cdot g_{0,h}).
\end{equation*}
Using (C7), one can write this as
\begin{equation*}
(2\pi)^{-3}\int dp\,|\hat h(k)|^2\,\frac{|k|^2}{|p|^2}\,\exp(-ip_1t-ik\cdot y)\int d\mathrm{kl}(s)\,\frac{|p|^2}{|p|^2+s^2}.
\end{equation*}
A majorant of the integrand is $|\hat h(k)|^2\,|k|^2/|p|^2$, and step 1 shows that it is integrable.
Fubini's theorem, the formula
\begin{equation*}
\int\frac{dp_1}{2\pi}\,\frac{\exp(-ip_1t)}{p_1^2+\omega^2}=\frac{\exp(-t\omega)}{2\omega}
\end{equation*}
and the equality $|\hat h(-k)|=|\hat h(k)|$ give \eqref{eq:magnetic-semigroup}.

The two sides of \eqref{eq:magnetic-semigroup} are the Laplace--Fourier transforms of $\langle\psi_h,\mathfrak E_1(\cdot)\,\psi_h\rangle$ \cite[Lemma~12.18\,(b)]{QEDII} and of the image measure, respectively.
Both being finite measures on $[0,\infty)\times\mathbb R^2$, they are equal by \cite[Lemma~12.25]{QEDII}.
The integral of $\omega$ against the image measure equals $\frac12\,\allowbreak \mathrm{kl}([0,\allowbreak \infty))\int\frac{d^2k}{(2\pi)^2}\,\allowbreak |k|^2\,\allowbreak |\hat h(k)|^2<\infty$; this is \eqref{eq:magnetic-energy}.
\item[4.] \emph{Assertion (d).}
By the definition of the reconstruction, $\exp(-t\,\mathsf H_1)\,V_1(y)\,\Omega_1=\Omega_1$.
Hence, by \cite[Lemma~12.25]{QEDII}, the measure $\langle\Omega_1,\mathfrak E_1(\cdot)\,\Omega_1\rangle$ is the unit mass at the origin, and $\mathfrak E_1(A)\,\Omega_1=1_A(0,0)\,\Omega_1$; consequently $\Omega_1^\perp$ reduces $\mathsf M$.
Further, $\langle\Omega_1,\psi_h\rangle=\lambda(B_0(h))=0$ \cite[Lemma~14.3\,(a)]{QEDII}.
Under the map $(\omega,p_\perp)\mapsto(\omega^2-|p_\perp|^2)^{1/2}$ the measure of (c) is carried onto $c_h(s)\,\mathrm{kl}(ds)$.
Since $\hat h$ is not identically zero, $c_h$ is a continuous positive function; consequently this measure and $\mathrm{kl}$ have the same support.
The point $0$ has the preimage $\{\omega=|p_\perp|\}$ under this map, and $c_h(0)=\frac12\|h\|_{1/2}^2$.
\end{enumerate}
\end{proof}

The magnetic variable is defined on the Euclidean probability space.
The proposition identifies its vector in the reconstructed space; it
does not define an operator at a sharp real time.
\subsection{A purely massless spectral measure and the infrared limit}
The following harmonic and probability facts will be used.
\begin{itemize}
\item[(K20)] (Harmonic Laurent expansion; \cite[Theorems~10.1 and~1.21; Chapter~5]{AxlerBourdonRamey}.) Let $u$ satisfy $\operatorname{div}u=\operatorname{curl}u=0$ on $\{0<|z|<R\}\subset\mathbb R^3$.
Then, with locally uniform convergence,
\begin{equation*}
u=\sum_{l\ge1}\nabla P_l+\sum_{l\ge0}\nabla H_l,
\end{equation*}
where each $P_l$ is a homogeneous harmonic polynomial of degree $l$, and $H_l=|z|^{-2l-1}\tilde P_l$ for a polynomial $\tilde P_l$ of the same kind;
every derivative of order $m$ of an $H_l$ has the form of an $H_{l+m}$.
Let $P_l,P_k$ be homogeneous harmonic polynomials of the degrees $l,k$, and let $H_l,H_k$ be functions of the form just described, taken from the same expansion or from different ones.
On every sphere $\{|z|=r\}$ the fields $\nabla P_l$ and $\nabla H_k$ are orthogonal in $L^2$ for all $l$ and $k$, and so are $\nabla P_l$ and $\nabla P_k$, as well as $\nabla H_l$ and $\nabla H_k$, whenever $l\ne k$.
A harmonic function on $\mathbb R^3$ whose square is integrable outside some ball is zero.
(Poisson.) Let $\varphi$ be harmonic in a neighbourhood of the closed ball $\{|z|\le s\}\subset\mathbb R^3$;
then
\begin{equation*}
\varphi(z)=\int_{|z'|=s}\frac{s^2-|z|^2}{4\pi s\,|z-z'|^3}\,\varphi(z')\,dS(z')\qquad(|z|<s).
\end{equation*}
\item[(K21)] (Minlos; \cite[Theorem~1]{FageotAminiUnser2014}.) On $\mathcal S$, every continuous positive definite function $\Phi$ with $\Phi(0)=1$ is the characteristic functional of a unique law on $(\mathcal S',\Sigma)$.
(L\'evy.) Let $Q$ and $Q_k$, $k\ge1$, be probability measures on $\mathbb R$;
if the characteristic function of $Q_k$ converges at every point to that of $Q$, then $Q_k\to Q$ weakly.
(Fatou.) If $Q_k\to Q$ weakly and $\varphi\ge0$ is continuous, then $Q(\varphi)\le\liminf_kQ_k(\varphi)$.
Let finitely many random variables have moments of every order, and let them fall into two non-empty groups such that the two groups are independent of each other;
then the joint cumulant of all of them vanishes.
\end{itemize}
\begin{named}{Theorem 4.3}
\emph{The massless part of the field.}
\begin{itemize}
\item[(a)] An ergodic law $\mathrm{lw}$ of $\mathcal K$ with $\mathrm{kl}=c_0\delta_0$ equals $\gamma_{c_0}$.
\item[(b)] Let $\mathrm{lw}$ be an ergodic law of $\mathcal K$; at every charge $e>0$ and for every mass assignment, $\nu$-almost all components of each volume-first limit point $\mu$ are laws of this kind.
Put $c_0:=\mathrm{kl}(\{0\})$ and $(\delta_lg)(y):=l^{-3/2}g(y/l)$.
Then, for real $g\in L^2$,
\begin{equation*}
\mathrm{lw}\bigl(\exp iX(\delta_lg)\bigr)\longrightarrow\exp\bigl(-\tfrac{c_0}2\|\Pi_{\mathrm c}g\|^2\bigr)\qquad(l\to\infty).
\end{equation*}
\end{itemize}
\end{named}

\begin{proof}
We first prove (a), assuming $\mathrm{kl}=c_0\delta_0$; Steps~7 and~8 prove (b).
Throughout, two-forms are identified with vector fields by means of $\ast$; a closed form then corresponds to a divergence free field, and $\Pi_{\mathrm c}\operatorname{curl}=\operatorname{curl}$.

\emph{Step 1} (kernels).
By Corollary~2.2\,(c), every ergodic law of $\mathcal K$ satisfies the one-plane hypotheses for $(1,2)$; it is, moreover, invariant under $SO(3)$ \cite[Theorem~1.4, Part~I\,(a), (f)]{QEDII}.
For pairwise disjoint compact supports, Theorem~3.3\,(f2) therefore represents $S_n$, and with it $\kappa_n$ \cite[(n5)]{QEDIII}, by integration against a kernel which is real analytic wherever the $n$ points are distinct.
Write $\mathsf k_n$ for the kernel of $\kappa_n$.
Step~1 of the proof of Proposition~1.2 gives $|\kappa_n(b_1,\dots,b_n)|\le C_n\prod_i\|b_i\|$ on $L^2$.
Condition (dF) yields $\operatorname{div}_{y_a}\mathsf k_n=0$.

\emph{Step 2} ($\operatorname{curl}_{y_a}\mathsf k_n=0$).
Let $n\ge2$, and let $a,b_2,\dots,b_n\in C_c^\infty$ be supported in balls of radius $\delta$ any two of whose centres are farther apart than $4\delta$.
Set $a':=a(\cdot-y)$ and $f(y):=S_n(\operatorname{curl}a',b_2,\dots,b_n)$.
The function $f$ is real analytic on the set of those $y$ for which the closed balls do not meet; this set is the complement of finitely many disjoint closed balls and is therefore connected.
Suppose that some plane, with reflection $\mathsf r$, has the first ball on one side and the other balls on the other side; the points $y$ with this property fill a neighbourhood of infinity.
From reflection positivity at this plane, invariance under $\mathsf r$ \cite[Theorem~1.4, Part~I\,(b), (f)]{QEDII} and two applications of the Schwarz inequality one obtains
\begin{equation*}
|f(y)|^2\le S_2(\mathsf r\cdot\operatorname{curl}a',\operatorname{curl}a')\,\Bigl\|\prod_jX(b_j)\Bigr\|_2^2,
\end{equation*}
where $\mathsf r\,\cdot$ denotes the action of $\mathsf r$ on two-forms \cite[Section~1.2]{QEDII}.
Because $\mathrm{kl}=c_0\delta_0$ and the two supports are disjoint,
\begin{equation*}
S_2(\mathsf r\cdot\operatorname{curl}a',\operatorname{curl}a')=c_0\,\langle\mathsf r\cdot\operatorname{curl}a',\operatorname{curl}a'\rangle=0
\end{equation*}
\cite[Theorem~1.4, Part~I\,(d), (f)]{QEDII}.
Thus $f$ is zero on a non-empty open subset of the connected open set on which it is real analytic, and the identity theorem for real analytic functions \cite[Lemma~14.10\,(b)]{QEDII} shows that $f$ vanishes identically.
In the formula for the cumulant, the block containing $\operatorname{curl}a$ is either a singleton, whose factor is $S_1=0$, or has two elements or more, in which case its factor is the value at $y=0$ of the function $f$ of that order, which vanishes as well; hence $\kappa_n(\operatorname{curl}a,b_2,\dots,b_n)=0$.
Since $a$, the $b_j$ and $\delta$ were arbitrary, $\operatorname{curl}_{y_1}\mathsf k_n=0$, and the symmetry of $\mathsf k_n$ gives the same for the other arguments.

\emph{Step 3} (the singularity at $y_1=y_2$).
Fix $n\ge3$ and $b_3,\dots,b_n$ supported in disjoint closed balls, and let $O$ be the complement of these balls, which is connected.
Define
\begin{equation*}
k_{ij}(y_1,y_2):=\int\mathsf k_n\,b_3\cdots b_n,\qquad B(b_1,b_2):=\kappa_n(b_1,\dots,b_n).
\end{equation*}
To $z\mapsto k_{\cdot j}(y_2+z,y_2)$ apply (K20), and call $H_{l,j}(z;y_2)$ the terms of the resulting expansion.
The orthogonality in (K20) expresses each coefficient of these terms as the integral, over a small sphere $|z|=s$, of $k_{\cdot j}(y_2+z,y_2)$ paired with a fixed field, so that the coefficients depend real analytically on $y_2$.
Fix a point $y_*\in O$, let $\varepsilon_0(y_*)>0$ be such that the closed ball of radius $3\varepsilon_0(y_*)$ about $y_*$ lies in $O$, let $0<\varepsilon\le\varepsilon_0(y_*)$, and set
\begin{equation*}
b_i^\varepsilon:=\varepsilon^{-3/2}\beta_i\bigl((\cdot-y_*)/\varepsilon\bigr),
\end{equation*}
where $\beta_2=\chi e_j$ for some $0\le\chi\in C_c^\infty$ with $\int\chi=1$ and support in $|\eta|<\frac12$, and where $\beta_1=\psi(|\eta|)\nabla H_{l,j}(\cdot;y_*)$ for some $0\le\psi\in C_c^\infty((1,2))$ with $\psi\ne0$.
Then $b_2^\varepsilon$ is supported in $\{|y-y_*|<\varepsilon/2\}$ and $b_1^\varepsilon$ in $\{\varepsilon<|y-y_*|<2\varepsilon\}$;
on the product of these supports $|y_1-y_2|<5\varepsilon/2$, which is less than the distance from $y_2$ to the complement of $O$, so that the expansion about $y_2$ applies there.
Uniformly for $|\eta_2|\le\frac12$ and $1\le|\eta_1|\le2$, expand $H_{l'}(\eta_1-\eta_2)$ in homogeneous terms of degree $m$ in $\eta_2$.
In $\eta_1$, the term of degree $m$ is an $H_{l'+m}$, so that about the centre $0$ there occur only orders $\ge l'$; and the regular part, once shifted by $\eta_2$, becomes a series of fields $\nabla P_k(\eta_1)$.
From the orthogonality on spheres one thus gets
\begin{equation*}
B(b_1^\varepsilon,b_2^\varepsilon)=\sum_{l'\le l}\varepsilon^{1-l'}\int\chi\,F_{l'},
\end{equation*}
where each $F_{l'}=F_{l'}(\eta_2;\varepsilon)$ extends continuously to $\varepsilon=0$, and $F_l\to\int\psi\,|\nabla H_{l,j}(\cdot;y_*)|^2$ as $\varepsilon\to0$.
As $\|b_i^\varepsilon\|=\|\beta_i\|$, the bound on $B$ of Step~1 then forces $H_{l,j}(\cdot;y_*)=0$ for every $l\ge2$;
since the point $y_*\in O$ and the index $j$ were arbitrary, $H_{l,j}(\cdot;y_2)$ vanishes for $l\ge2$ at each $y_2\in O$ and for each $j$.

\emph{Step 4} ($H_{0,j}=0$).
Write $H_{0,j}=q_j(y_2)/(4\pi|z|)$.
Take $b_2\in C_c^\infty(O)$, and $u\in C_c^\infty(O)$ with $u=1$ near $\operatorname{supp}b_2$.
Condition (dF) and the divergence theorem yield
\begin{equation*}
0=B(\nabla u,b_2)=\sum_j\int b_{2,j}(y_2)\,q_j(y_2)\,dy_2,
\end{equation*}
since the flux of $k_{\cdot j}(\cdot,y_2)$ across small spheres centred at $y_2$ equals $-q_j(y_2)$, and $\nabla H_{0,j}$ is the only term of the expansion with a non-zero flux.
As this is true for every $b_2$, one concludes that $q_j\equiv0$.
Consequently, with $G(z):=\frac1{4\pi|z|}$,
\begin{equation*}
k_{ij}=\sum_mT_{mj}(y_2)\,\partial_i\partial_mG(y_1-y_2)+R_{ij},
\end{equation*}
where $z\mapsto R_{\cdot j}(y_2+z,y_2)$ is the regular part.
On each sphere $|z|=s$ the three fields $\nabla\partial_mG$ are mutually orthogonal, and $\int_{|z|=s}|\nabla\partial_mG|^2\,dS=1/(2\pi s^4)$; therefore
\begin{equation*}
T_{mj}(y_2)=2\pi s^4\int_{|z|=s}\sum_ik_{ij}(y_2+z,y_2)\,\partial_i\partial_mG(z)\,dS,
\end{equation*}
and $T_{mj}$ is real analytic.
For $|z|<s$, Poisson's formula for the ball $\{|z|<s\}$ of $\mathbb R^3$, applied to each component of the harmonic field $z\mapsto R_{\cdot j}(y_2+z,y_2)$, gives
\begin{equation*}
R_{ij}(y_2+z,y_2)=\int_{|z'|=s}\frac{s^2-|z|^2}{4\pi s\,|z-z'|^3}\Bigl[k_{ij}(y_2+z',y_2)-\sum_mT_{mj}(y_2)\,\partial_i\partial_mG(z')\Bigr]\,dS(z'),
\end{equation*}
an expression which is smooth in $(z,y_2)$; thus $R$ is smooth on $O\times O$ as a function of both variables jointly.
By symmetry, $k$ is free of divergence and of curl also in the pair $(y_2,j)$.
Write $y=y_2$ and $z=y_1-y$. For $n\in S^2$ put
\[
Q(n)=3nn^T-1,\qquad
\partial_i\partial_mG(rn)=\frac{Q(n)_{im}}{4\pi r^3}.
\]
The degree $-4$ part of
$\partial_{y_a}k_{ib}-\partial_{y_b}k_{ia}=0$ is
\[
\sum_m\bigl(T_{mb}\partial_i\partial_m\partial_aG
             -T_{ma}\partial_i\partial_m\partial_bG\bigr)=0.
\]
Contracting in $i$ with $n_i$ gives
$Q(n)T=T^TQ(n)$ for every $n\in S^2$.
For $a\ne b$, taking $n=e_a$ and $n=e_b$ yields
$2T_{ab}=-T_{ba}$ and $-T_{ab}=2T_{ba}$, so both off-diagonal entries vanish.
Taking $n=(e_a+e_b)/\sqrt2$ then gives $T_{aa}=T_{bb}$.
Thus $T=\tau1$.
Since $\Delta G=0$ away from zero and the derivatives of $R$ are bounded near $z=0$, the degree $-3$ part of
$\sum_j\partial_{y_j}k_{ij}=0$ now gives
\[
Q(n)\nabla\tau=0\qquad(n\in S^2).
\]
The eigenvalues of $Q(n)$ are $2,-1,-1$, so $\nabla\tau=0$.
As $O$ is connected, $\tau$ is constant there.

\emph{Step 5} ($\tau=0$).
The functions $\partial\partial G$ are homogeneous of degree $-3$, whereas $R$ is a bounded function near $(y_*,y_*)$.
It follows that, whenever the supports of $\beta_1$ and $\beta_2$ are disjoint,
\begin{equation*}
B(b_1^\varepsilon,b_2^\varepsilon)=\tau\,c_\beta+O(\varepsilon^3),\qquad c_\beta:=\sum_{i,m}\iint\beta_{1,i}(\eta_1)\,\partial_i\partial_mG(\eta_1-\eta_2)\,\beta_{2,m}(\eta_2)\,d\eta_1\,d\eta_2.
\end{equation*}
Choose the $\beta_i$ as in Step~3, but with $\partial_jG$ instead of $H_{l,j}(\cdot;y_*)$; the computation made in Step~3 yields
\begin{equation*}
c_\beta=\int\psi\,|\nabla\partial_jG|^2=\frac1{2\pi}\int\psi(x)\,x^{-4}\,dx>0.
\end{equation*}
Assume that the balls carrying the $b_j$, $j\ge3$, are contained in the cube $(-r,r)^3$, that $|y_*|_\infty\ge4r$, and that $\varepsilon<r/2$; the functions $b_i^\varepsilon(\cdot+y_*)$ are then supported in that cube as well.
Every mixed moment occurring in the cumulant formula is to be replaced by the product of its far part and its near part.
Since $\mathrm{lw}$ is ergodic, \cite[Theorem~12.39; Theorem~1.4, Part~II\,(E4)]{QEDII} bounds the error so made by $Cr/|y_*|_\infty$ with a constant $C$ independent of $\varepsilon$ and of $y_*$:
that bound is expressed through the $L^2(\mathrm{lw})$ norms of the two parts, H\"older's inequality and \cite[Lemma~14.3\,(a)]{QEDII} give, for real $g_1,\dots,g_k\in L^2$,
\begin{equation*}
\Bigl\|\prod_{a=1}^kX(g_a)\Bigr\|_{L^2(\mathrm{lw})}\le c_{2k}^k\prod_{a=1}^k\|g_a\|,
\end{equation*}
and $\|b_i^\varepsilon\|=\|\beta_i\|$.
After the replacement one is left with the joint cumulant of $(X(b_1^\varepsilon),X(b_2^\varepsilon))$ and of an independent copy of $(X(b_3),\dots,X(b_n))$, neither family being empty since $n\ge3$, and the moments of every order of these variables being finite \cite[Lemma~14.3\,(a)]{QEDII};
by the independent-groups cumulant identity in (K21), this cumulant is $0$.
Hence $|\tau|\,c_\beta\le Cr/|y_*|_\infty+O(\varepsilon^3)$.
Letting first $\varepsilon\to0$ and afterwards $|y_*|_\infty\to\infty$, which is legitimate because $\tau$ is constant on $O$, one arrives at $\tau=0$.

\emph{Step 6.}
With $\tau=0$ one has $k=R$, which is smooth on $O\times O$ and harmonic in $y_1$; so $k$ has a harmonic extension across $y_1=y_2$.
Let $b_2,\dots,b_n\in C_c^\infty$ have supports in pairwise disjoint closed balls $\bar B_j$, denote by $V$ the complement of the union of these balls, put $O_j:=V\cup\bar B_j$, and define on $V$
\begin{equation*}
w(y_1):=\int\mathsf k_n(y_1,\cdot)\,b_2\cdots b_n.
\end{equation*}
Steps~3--5, applied with $O=O_2$, show that $\int R(\cdot,y_2)\,b_2(y_2)\,dy_2$ is harmonic on $O_2$ and agrees with $w$ on $V$.
Since $\kappa_n$ is symmetric, Steps~3--5 apply to every pair $(1,j)$, the other entries being held fixed, and they furnish a harmonic extension of $w$ to $O_j$.
The sets $O_j$ cover $\mathbb R^3$, and any two of them intersect in $V$; so $w$ extends to a harmonic field on all of $\mathbb R^3$.
For $b_1\in C_c^\infty(V)$, Step~1 gives $|\int w\cdot b_1|\le C_n\|b_1\|\prod_{j\ge2}\|b_j\|$; hence $w$ is square integrable outside the balls, and (K20) gives $w=0$.
A partition of unity shows that $\kappa_n$ vanishes at non-coinciding arguments for each $n\ge3$, and Proposition~1.2\,(b) shows that $\mathrm{lw}$ is Gaussian.
Its mean is zero and its covariance is $c_0\langle\cdot,\Pi_{\mathrm c}\cdot\rangle$, as in Step~2.
The definition of $\gamma_{c_0}$ then gives $\Phi_{\mathrm{lw}}=\Phi_{\gamma_{c_0}}$ on $\mathcal S$, and (K21) gives $\mathrm{lw}=\gamma_{c_0}$.
This establishes (a).

\emph{Step 7} (part (b)).
For $\omega\in\mathcal S'$ put $\delta_l^*\omega:=\omega\circ\delta_l$, and let $\mathrm{lw}_l:=(\delta_l^*)_*\mathrm{lw}$, that is, the law of $g\mapsto X(\delta_lg)$; then $\Phi_{\mathrm{lw}_l}=\Phi_{\mathrm{lw}}\circ\delta_l$.
The map $\delta_l$ is an isometry, it sends gradients to gradients, it commutes with $\Pi_{\mathrm c}$, and it normalises the Euclidean group:
\begin{equation*}
\delta_l(\mathrm{Gi}\cdot f)=\mathrm{Gi}^{(l)}\cdot\delta_lf,\qquad\text{where }\mathrm{Gi}^{(l)}\colon y\mapsto\mathrm{Rl}\,y+la\ \text{ for }\ \mathrm{Gi}\colon y\mapsto\mathrm{Rl}\,y+a.
\end{equation*}
In particular, $\delta_l$ transforms the reflection in $\{\hat n\cdot y=\mathrm{po}\}$, and the forms supported in $\{\hat n\cdot y>\mathrm{po}\}$, into the corresponding objects for the plane $\{\hat n\cdot y=l\,\mathrm{po}\}$; within each of the nine families, $\mathrm{po}$ takes every real value.
The hypothesis (H), condition (dF) and membership in $\mathcal C_9$ reduce to a single inequality and to relations among values of the characteristic functional \cite[Section~1.2, Definition~2.38, Remark~14.4]{QEDII}:
\begin{itemize}
\item[(i)] (G) and (dF), taken together, are equivalent to the inequality $\mathrm{lw}(\exp X(f))\le\exp(\frac12\|\Pi_{\mathrm c}f\|^2)$ for $f\in\mathcal S$;
\item[(ii)] $\Phi_{\mathrm{Gi}_*\mathrm{lw}}(f)=\allowbreak \Phi_{\mathrm{lw}}(\mathrm{Gi}^{-1}\cdot f)$, and therefore $\mathrm{Gi}_*\mathrm{lw}=\mathrm{lw}$ holds exactly when $\Phi_{\mathrm{lw}}(\mathrm{Gi}\cdot f)=\Phi_{\mathrm{lw}}(f)$ for all $f\in\mathcal S$ (K21); this applies to the translations, to $O_{\mathrm{cub}}$ and to the reflections in the planes of the nine families;
\item[(iii)] reflection positivity at a plane, including $(RP_j)$ of Section~2, is by its definition a family of inequalities among values of $\Phi_{\mathrm{lw}}$.
\end{itemize}
Hence $\mathrm{lw}_l\in\mathcal K$.
This law is ergodic, because $\Psi\mapsto\Psi\circ\delta_l^*$ is a unitary map of $L^2(\mathrm{lw}_l)$ onto $L^2(\mathrm{lw})$ under which $U(a)$ corresponds to $U(la)$.

According to \cite[Lemma~14.3\,(a)]{QEDII}, $|\Phi_{\mathrm{lw}_l}(g)-\Phi_{\mathrm{lw}_l}(g')|\le\|g-g'\|$, so that every sequence $l\to\infty$ admits a subsequence $l_k$ such that $\Phi_{\mathrm{lw}_{l_k}}$ has a limit at every point of real $L^2$.
This limit, $\Phi_\infty$, is a positive definite function with $\Phi_\infty(0)=1$, Lipschitz with respect to $\|\cdot\|$, and hence continuous on $\mathcal S$.
By (K21) (Minlos), there is one law $\mathrm{lw}_\infty$ with $\Phi_\infty=\Phi_{\mathrm{lw}_\infty}$ on $\mathcal S$.
For $f\in\mathcal S$, the law of $X(f)$ under $\mathrm{lw}_{l_k}$ converges weakly to its law under $\mathrm{lw}_\infty$, by L\'evy's continuity theorem, since $\Phi_{\mathrm{lw}_{l_k}}(tf)\to\Phi_{\mathrm{lw}_\infty}(tf)$ for every real $t$.
Applied to $\mathrm{lw}_{l_k}$ at $2f$, (i) gives $\sup_k\mathrm{lw}_{l_k}(\exp2X(f))\le\exp(2\|\Pi_{\mathrm c}f\|^2)$, so that uniform integrability of the exponentials yields $\mathrm{lw}_\infty(\exp X(f))=\allowbreak \lim_k\mathrm{lw}_{l_k}(\exp X(f))\le\exp(\frac12\|\Pi_{\mathrm c}f\|^2)$;
in this way (i) carries over to $\mathrm{lw}_\infty$, and (ii) and (iii) are preserved under pointwise limits.
Thus $\mathrm{lw}_\infty\in\mathcal K$.
As a law of $\mathcal K$, $\mathrm{lw}_\infty$ has (G), so that $\Phi_{\mathrm{lw}_\infty}$ is defined on $L^2$; there it coincides with $\Phi_\infty$, both functions being Lipschitz and $\mathcal S$ being dense.

\emph{Step 8.}
Let $S_2^{(l)}$ denote the second moment of $\mathrm{lw}_l$, put $b:=\ast g$, let $P^\perp(p):=1-pp^T/|p|^2$, and let
\begin{equation*}
\mathsf D(s):=\int d\mathrm{kl}(t)\,\frac{s}{s+t^2},
\end{equation*}
a function which decreases to $\mathrm{kl}(\{0\})$ as $s\downarrow0$.
By \cite[Theorem~1.4, Part~I\,(d), (f)]{QEDII},
\begin{equation*}
S_2^{(l)}(g,g)=(2\pi)^{-3}\int|P^\perp\hat b|^2\,\mathsf D(|p|^2/l^2)\,dp\longrightarrow c_0\,\|\Pi_{\mathrm c}g\|^2.
\end{equation*}
Consider a law with (G), a real $g\in L^2$ and $t>0$.
The elementary inequality $|u^2-2(1-\cos u)|\le u^4/12$ for real $u$, together with \cite[Lemma~14.3\,(a)]{QEDII} and $c_4^4=64$, yields
\begin{equation*}
\bigl|S_2(g,g)-2t^{-2}\bigl(1-\operatorname{Re}\Phi(tg)\bigr)\bigr|\le\frac{t^2}{12}\,\mathrm{lw}\bigl(X(g)^4\bigr)\le\frac{c_4^4}{12}\,t^2\|g\|^4.
\end{equation*}
Every $\mathrm{lw}_l$, and $\mathrm{lw}_\infty$ too, satisfies (G) with the constant one, and $\Phi_{\mathrm{lw}_{l_k}}(tg)\to\Phi_{\mathrm{lw}_\infty}(tg)$.
It follows that $\limsup_k|S_2^{(l_k)}(g,\allowbreak g)-S_2^{\mathrm{lw}_\infty}(g,\allowbreak g)|\le\frac{c_4^4}6t^2\|g\|^4$ for each $t>0$, which means $S_2^{\mathrm{lw}_\infty}(g,g)=c_0\|\Pi_{\mathrm c}g\|^2$.
Polarisation and the uniqueness assertion of \cite[Theorem~1.4, Part~I\,(d), (f)]{QEDII} lead to $\mathrm{kl}_\infty=c_0\delta_0$.
The constants in the bound of \cite[Theorem~12.39]{QEDII} are absolute:
\begin{equation*}
\bigl|\Phi_{\mathrm{lw}_l}\bigl(g'(\cdot-a)-g\bigr)-\overline{\Phi_{\mathrm{lw}_l}(g)}\,\Phi_{\mathrm{lw}_l}(g')\bigr|\le C_0\,r/|a|_\infty
\end{equation*}
for every $l$, every $r>0$, all $g,g'\in C_c^\infty((-r,r)^3)$ and all $a$ with $|a|_\infty\ge4r$.
This inequality carries over to $\mathrm{lw}_\infty$.
Accordingly $\mathrm{lw}_\infty$ clusters on the linear span of the functions $\exp iX(g)$, $g\in C_c^\infty$, a span which is dense in $L^2(\mathrm{lw}_\infty)$; so $\mathrm{lw}_\infty$ is mixing and ergodic \cite[Section~8.1.4, Lemma~8.14\,(a)]{QEDI}.
Part (a) then yields $\mathrm{lw}_\infty=\gamma_{c_0}$, whichever subsequence was chosen, and this proves (b).
\end{proof}
For laws $\mathrm{lw}_1,\mathrm{lw}_2$ and $0<\epsilon<1$, let
$\mathrm{lw}_1\ast_\epsilon\mathrm{lw}_2$ be the law of
$(1-\epsilon)^{1/2}F+\epsilon^{1/2}G$, where $F$ and $G$ are independent
with these respective laws \cite[Section~21.11.1]{QEDI}.
The constants $C_U$, $C_{\bar\kappa}$, $C'_{\bar\kappa}$,
$c_{\mathrm{cl}}$ and $\mathfrak e_2$ in the next proposition are those of
\cite[Theorem~44]{QEDIII}.

\subsection{The properties (D1)--(D6) do not determine a law}

\begin{named}{Proposition 4.4}
With the data of \cite[Theorem~44]{QEDIII} held fixed, let $e\in(0,e_2]$ be a charge, $\mu$ a volume-first limit point, and $\xi\in\Xi_\mu$.
Define
\begin{equation*}
\epsilon_0:=\frac{\bigl(C'_{\bar\kappa}-\frac12C_{\bar\kappa}\bigr)\cdot e^{\sigma}}{2\max_{\mathrm{ix}}|b(g_{\mathrm{ix}})|};
\end{equation*}
then $0<\epsilon_0\le\frac1{20}$.
Set $\mu_\xi^{0}:=\mu_\xi$, and $\mu_\xi^{\epsilon}:=\mu_\xi\ast_\epsilon\gamma$ when $0<\epsilon\le\epsilon_0$.
For each $\epsilon\in(0,\epsilon_0]$ the following assertions are true.
\begin{itemize}
\item[(a)] The single law $\mu_\xi^\epsilon$ enjoys all of the properties (D1)--(D6) listed in \cite[Theorem~44]{QEDIII}, with $g_{\mathrm{ix}}$ as tests, $x_0$ as point and $c_{\mathrm{cl}}$, $C'_{\bar\kappa}$, $\sigma$, $\mathfrak e_2$ as constants, exactly as in that theorem.
\item[(b)] For $n\ge3$ the cumulants of order $n$ of $\mu_\xi^\epsilon$ are $(1-\epsilon)^{n/2}\kappa_n^\xi$.
In particular, $\kappa_4(f)=(1-\epsilon)^2\kappa_4^\xi(f)$ for the fourth cumulant of $\mu_\xi^\epsilon$; hence the laws $\mu_\xi^\epsilon$, $0\le\epsilon\le\epsilon_0$, differ from one another.
\item[(c)] The K\"all\'en--Lehmann measure of $\mu_\xi^\epsilon$ is $(1-\epsilon)\,\mathrm{kl}_\xi+\epsilon\,\delta_0$.
\item[(d)] Consider a positive measure $\varrho$, defined on $[0,\infty)$, of total mass at most $1$; the symbol $\varrho$ is unrelated to the Wilson parameter $\rho$.
Let $\gamma^\varrho$ denote the centred Gaussian law with covariance
\begin{equation*}
C_\varrho(g,g'):=(2\pi)^{-3}\int\mathrm dp\,\overline{\widehat{\ast g}(p)}\cdot\Pi(p)\,\widehat{\ast g'}(p)\int\mathrm d\varrho(t)\,\frac{|p|^2}{|p|^2+t^2};
\end{equation*}
thus $\gamma^{a\delta_0}=\gamma_a$.
For $\gamma^\varrho$ the properties (H) and (dF) hold; this law belongs to $\mathcal C_9$, is ergodic, and its K\"all\'en--Lehmann measure is $\varrho$.
Assume in addition that
\begin{equation*}
\epsilon\cdot\max_{\mathrm{ix}}\bigl(\|g_{\mathrm{ix}}\|^2-C_\varrho(g_{\mathrm{ix}},g_{\mathrm{ix}})\bigr)\le\bigl(C'_{\bar\kappa}-2C_U\bigr)\cdot e^{2-\mathfrak e_2}.
\end{equation*}
The law $\mu_\xi\ast_\epsilon\gamma^\varrho$ then fulfils (a) together with the first assertion of (b), and its K\"all\'en--Lehmann measure equals $(1-\epsilon)\,\mathrm{kl}_\xi+\epsilon\,\varrho$.
In the case $\varrho=a\delta_0$, $0<a\le1$, the assumption takes the form
\begin{equation*}
\epsilon\,(1-a)\max_{\mathrm{ix}}\|g_{\mathrm{ix}}\|^2\le(C'_{\bar\kappa}-2C_U)\,e^{2-\mathfrak e_2},
\end{equation*}
and it is fulfilled for all $\varrho$ as soon as
\begin{equation*}
\epsilon\max_{\mathrm{ix}}\|g_{\mathrm{ix}}\|^2\le(C'_{\bar\kappa}-2C_U)\,e^{2-\mathfrak e_2}.
\end{equation*}
\end{itemize}
\end{named}

The analogue for the class $\mathbf K(C)$ of \cite{QEDI} is \cite[Proposition~21.40\,(a)]{QEDI}, where $C$ is replaced by $C+\boldsymbol\delta$; Proposition~4.4, on the other hand, enlarges no constant.

\begin{proof}
\emph{Step 1: the product formula.}
For real $g\in\mathcal S$, the definition of $\ast_\epsilon$ \cite[Section~21.11.1]{QEDI} and the independence of the two summands give
\begin{equation*}
\Phi_{\mu_\xi^\epsilon}(g)=\Phi_{\mu_\xi}\bigl((1-\epsilon)^{1/2}g\bigr)\cdot\Phi_\gamma\bigl(\epsilon^{1/2}g\bigr);
\end{equation*}
the same factorisation is valid for $\mathrm{lw}(\exp X(g))$.

\emph{Step 2: (H), (dF) and $\mathcal C_9$.}
The two factors share the following properties: (G) and (dF); invariance under $\mathbb R^3$ and under $O_{\mathrm{cub}}$; and, for each plane in one of the nine families, invariance under the reflection in that plane together with reflection positivity at it.
For $\mu_\xi$ these are contained in (D1), and for $\gamma$ in \cite[Lemma~12.29]{QEDII}.
Step~2 of the proof of \cite[Proposition~21.40\,(a)]{QEDI} transfers them to $\mu_\xi^\epsilon$, and the argument there rests on three facts.
First, the map $(F,G)\mapsto(1-\epsilon)^{1/2}F+\epsilon^{1/2}G$ commutes with every Euclidean map and sends closed forms to closed forms.
Second, $\frac{1-\epsilon}2\|g\|^2+\frac\epsilon2\|g\|^2=\frac12\|g\|^2$.
Third, Step~1 exhibits the matrix $\bigl(\Phi(f_q-r\cdot f_p)\bigr)_{p,q}$ of \cite[(n3)]{QEDIII} as the entrywise product of two positive semidefinite matrices.
Although the transfer argument is written for the planes $\{y_j=0\}$, it makes no use of anything special to these planes.

\emph{Step 3: ergodicity.}
Take real $g,g'\in C_c^\infty$.
By Step~1, $\mu_\xi^\epsilon\bigl(\exp(iX(g))\cdot U(a)\exp(iX(g'))\bigr)$ splits into two factors: the corresponding expression for $\mu_\xi$, evaluated at $(1-\epsilon)^{1/2}g$ and $(1-\epsilon)^{1/2}g'$, and the corresponding expression for $\gamma_\epsilon$, evaluated at $g$ and $g'$.
Each factor factorises when $|a|\to\infty$ in any manner: for the first this is the second sentence of (D4), for the second it is \cite[Lemma~12.29]{QEDII}.
From \cite[Lemma~8.14\,(a)]{QEDI} and the first part of its proof one concludes that $\mu_\xi^\epsilon$ is ergodic, and that it clusters on the entire space $L^2(\mu_\xi^\epsilon)$ when $|a|\to\infty$ in any manner.

\emph{Step 4: (D1)--(D5).}
Steps~2 and~3 together establish (H) and (dF) for $\mu_\xi^\epsilon$, its membership in $\mathcal C_9$, and its ergodicity.
These properties constitute (D1), and they coincide with the display (7.1) of \cite[Chapter~44]{QEDIII}; in paragraph~7.3 of that chapter, (D2)--(D5) are deduced from this display for a single law.
More precisely, (D2) is \cite[Theorem~1.4, Part~I\,(a)--(d), (f)]{QEDII}, (D3) is Part~II\,(E8) of the same theorem, the rate in (D4) is \cite[Theorem~12.39]{QEDII}, and (D5) is Part~II\,(E0$'$)--(E4) of \cite[Theorem~1.4]{QEDII}.

\emph{Step 5: moments and cumulants.}
For a sum of independent random variables, the logarithm of the generating function equals the sum of the logarithms of the generating functions of the summands; furthermore, all cumulants of $\gamma$ of order $3$ or more vanish.
Hence (b) holds, and the second moment $S_2^\epsilon$ of $\mu_\xi^\epsilon$ is given by
\begin{equation*}
S_2^{\epsilon}(g,g')=(1-\epsilon)\,S_2^\xi(g,g')+\epsilon\,\langle g,\Pi_{\mathrm c}g'\rangle
\end{equation*}
(Step~4 of the proof of \cite[Proposition~21.40\,(a)]{QEDI}).
If $g$ is closed, then $\|\Pi_{\mathrm c}g\|=\|g\|$ \cite[Lemma~9.12\,(a)]{QEDI}.

\emph{Step 6: (D6).}
As $\xi\in\Xi_\mu$, the statements (c1) and (c3) hold at $\xi$ \cite[Theorem~44\,(D)]{QEDIII}.
Fix $\mathrm{ix}$, put $z':=\|g_{\mathrm{ix}}\|^2$, and let $\mathrm{dv}_m:=\mu_\xi(\mathrm{Wk}_{m,\mathrm{ix}})$ for $m=2,4$, so that $|\mathrm{dv}_m|\le\mathsf u^{\mathrm{inf}}_{m,\mathrm{ix}}$ by (c3).
Here $\mathrm{Wk}_{m,\mathrm{ix}}=\allowbreak \mathrm{Pw}_m\bigl(X(g_{\mathrm{ix}});z',\allowbreak c^{\mathrm{inf}}_{m,\mathrm{ix}}\bigr)$ is the function of (B4), formed with the polynomials of \cite[Section~I.7.6]{QEDIII},
\begin{equation*}
\mathrm{Pw}_2(z;z',\mathrm{cf})=e^{-2}(z^2-z')-\mathrm{cf},\qquad\mathrm{Pw}_4(z;z',\mathrm{cf})=e^{-4}(z^4-6z'z^2+3z'^2)-\mathrm{cf}.
\end{equation*}
By (c1) the law $\mu_\xi$ has (G), hence $S_1^\xi=0$ \cite[Proposition~45.2\,(1), (8)]{QEDIII}, and the formula for $\kappa_4$ in \cite[(n5)]{QEDIII} reduces to $\kappa_4^\xi(g)=S_4^\xi(g,g,g,g)-3S_2^\xi(g,g)^2$.
Integrating $\mathrm{Wk}_{2,\mathrm{ix}}$ and $\mathrm{Wk}_{4,\mathrm{ix}}$ against $\mu_\xi$, one finds
\begin{equation*}
e^{-2}\bigl(S_2^\xi(g_{\mathrm{ix}},g_{\mathrm{ix}})-z'\bigr)=c^{\mathrm{inf}}_{2,\mathrm{ix}}+\mathrm{dv}_2,\qquad e^{-4}\kappa_4^\xi(g_{\mathrm{ix}})=c^{\mathrm{inf}}_{4,\mathrm{ix}}+\mathrm{dv}_4-3\bigl(c^{\mathrm{inf}}_{2,\mathrm{ix}}+\mathrm{dv}_2\bigr)^2.
\end{equation*}
Since $0<e\le e_2\le1$, $\sigma>0$ and $\mathfrak e_2\in[0,1)$ \cite[Proposition~43.2]{QEDIII}, one has $e^{2\sigma}\le e^\sigma\le1\le e^{-\mathfrak e_2}$.
The four inequalities in the first display of \cite[Theorem~44\,(B4)]{QEDIII} give $\mathsf u^{\mathrm{inf}}_{2,\mathrm{ix}}\le C_Ue^\sigma$ and $|c^{\mathrm{inf}}_{2,\mathrm{ix}}|\,\allowbreak \mathsf u^{\mathrm{inf}}_{2,\mathrm{ix}}\le C_Ue^\sigma$, and with them
\begin{equation*}
\bigl|S_2^\xi(g_{\mathrm{ix}},g_{\mathrm{ix}})-\|g_{\mathrm{ix}}\|^2\bigr|\le e^2\bigl(|c^{\mathrm{inf}}_{2,\mathrm{ix}}|+\mathsf u^{\mathrm{inf}}_{2,\mathrm{ix}}\bigr)\le e^2\bigl(C_Ue^{-\mathfrak e_2}+C_Ue^\sigma\bigr)\le2C_U\,e^{2-\mathfrak e_2},
\end{equation*}
\begin{align*}
\bigl|e^{-4}\kappa_4^\xi(g_{\mathrm{ix}})-b(g_{\mathrm{ix}})\bigr|&\le\bigl|c^{\mathrm{inf}}_{4,\mathrm{ix}}-3\bigl(c^{\mathrm{inf}}_{2,\mathrm{ix}}\bigr)^2-b(g_{\mathrm{ix}})\bigr|+|\mathrm{dv}_4|+6\,|c^{\mathrm{inf}}_{2,\mathrm{ix}}|\,|\mathrm{dv}_2|+3\,|\mathrm{dv}_2|^2\\
&\le(1+1+6)\,C_Ue^\sigma+3C_U^2e^{2\sigma}\le\bigl(8C_U+3C_U^2\bigr)e^\sigma=\tfrac12C_{\bar\kappa}\,e^\sigma,
\end{align*}
the last equality being the definition $C_{\bar\kappa}=2(8C_U+3C_U^2)$ \cite[Proposition~43.2, Constants]{QEDIII}.
Thus $|\kappa_4^\xi(g_{\mathrm{ix}})-e^4b(g_{\mathrm{ix}})|\le\frac12C_{\bar\kappa}e^{4+\sigma}$, and the fourth cumulant $(1-\epsilon)^2\kappa_4^\xi(g_{\mathrm{ix}})$ of $\mu_\xi^\epsilon$ at $g_{\mathrm{ix}}$ (Step~5) satisfies
\begin{equation*}
\bigl|(1-\epsilon)^2\kappa_4^\xi(g_{\mathrm{ix}})-e^4b(g_{\mathrm{ix}})\bigr|\le\tfrac12C_{\bar\kappa}e^{4+\sigma}+2\epsilon\,e^4|b(g_{\mathrm{ix}})|\le C'_{\bar\kappa}e^{4+\sigma}.
\end{equation*}
As $g_{\mathrm{ix}}=dd^*\beta_{\mathrm{ix}}$ is closed, Step~5 leads to
\begin{equation*}
S_2^\epsilon(g_{\mathrm{ix}},g_{\mathrm{ix}})-\|g_{\mathrm{ix}}\|^2=(1-\epsilon)\bigl(S_2^\xi(g_{\mathrm{ix}},g_{\mathrm{ix}})-\|g_{\mathrm{ix}}\|^2\bigr).
\end{equation*}

\emph{Step 7: (D6), continued.}
Every translation invariant law with (G) whose fourth cumulant obeys the bound established for $\mu_\xi^\epsilon$ in Step~6, namely
\begin{equation*}
\bigl|\kappa_4(g_{\mathrm{ix}})-e^4b(g_{\mathrm{ix}})\bigr|\le C'_{\bar\kappa}e^{4+\sigma}\qquad(\mathrm{ix}=0,+,-),
\end{equation*}
falls under \cite[Chapter~44, Lemma~3.2]{QEDIII}, which yields $\mathrm{Ext}_{x_0}(f)\ge\frac12e^4b(f)$ and $\kappa_4(f)\ge\frac12e^4b(f)>0$.
As for $\kappa_4(\mathrm{Gi}\cdot f)=\kappa_4(f)$, it follows from (C2), which is valid by Step~4.

\emph{Step 8: the size of $\epsilon_0$.}
One has $C'_{\bar\kappa}=\allowbreak 8C_{\bar\kappa}+3C_{\bar\kappa}^2>\frac12C_{\bar\kappa}>0$ \cite[Proposition~43.2, Constants]{QEDIII}, $\max_{\mathrm{ix}}|b(g_{\mathrm{ix}})|\ge b(f)>0$, and $C'_{\bar\kappa}e^\sigma\le\frac1{10}\,b(f)$ \cite[Chapter~44, (0.3)]{QEDIII}.

\emph{Step 9: (c).}
Assertion (c) follows on inserting (C7) for $\mu_\xi$ and $\mathrm{kl}_\gamma=\delta_0$ \cite[Theorem~14.26\,(c)]{QEDII} into the formula of Step~5, and on using the uniqueness in (C7) for $\mu_\xi^\epsilon$.

\emph{Step 10: (d), the law $\gamma^\varrho$.}
Restricted to real tests, $C_\varrho$ is bilinear, symmetric and positive semidefinite, and it satisfies $C_\varrho(g,g)\le\|\Pi_{\mathrm c}g\|^2\le\|g\|^2$.
Its invariance under all Euclidean maps follows from $\ast(\mathrm{Gi}\cdot g)=\allowbreak \det(\mathrm{Rl})\cdot\mathrm{Rl}\,\allowbreak (\ast g)(\mathrm{Gi}^{-1}\,\allowbreak \cdot)$ \cite[Lemma~9.12]{QEDI} together with $\Pi(\mathrm{Rl}\,p)=\mathrm{Rl}\,\Pi(p)\,\mathrm{Rl}^T$.
By \cite[Lemma~9.11\,(a)--(c)]{QEDI}, there is one and only one law $\gamma^\varrho$ with $\Phi_{\gamma^\varrho}=\exp(-\frac12C_\varrho)$.
It is a centred Gaussian law whose covariance is $C_\varrho$, and $\gamma^\varrho(\exp X(g))=\exp(\frac12C_\varrho(g,g))$.
Moreover, it is invariant under all Euclidean maps; and since it is mixing, it is ergodic.
By \cite[Remark~14.4]{QEDII}, (G) and (dF) hold for it.

\emph{Step 11: (d), reflection positivity.}
Let $g\in C_c^\infty(\{y_1>0\};\Lambda^2)$ be real, and set $j:=\operatorname{curl}\ast g$.
For $t\ge0$ and $k=(k_2,k_3)\in\mathbb R^2$ define $E_t:=(|k|^2+t^2)^{1/2}$ and
\begin{equation*}
v(k):=\int\mathrm dy\,\exp\bigl(-i(k_2y_2+k_3y_3)-E_ty_1\bigr)\,j(y)\in\mathbb C^3.
\end{equation*}
The identity $\overline{(p\times u)}\cdot(p\times u')=\allowbreak \bar u\cdot(|p|^2-pp^T)\,\allowbreak u'$ shows that the form whose multiplier is $\Pi(p)\,\frac{|p|^2}{|p|^2+t^2}$ can be written as
\begin{equation*}
(g,g')\mapsto\bigl\langle\operatorname{curl}\ast g,(-\Delta+t^2)^{-1}\operatorname{curl}\ast g'\bigr\rangle.
\end{equation*}
Further, $\operatorname{curl}\ast(\theta_1\cdot g)=\allowbreak \theta_1\,\allowbreak j(\theta_1\,\allowbreak \cdot)$, and one has
\begin{equation*}
\int\frac{\mathrm dp_1}{2\pi}\,\frac{\exp(ip_1s)}{p_1^2+E_t^2}=\frac{\exp(-E_t|s|)}{2E_t}.
\end{equation*}
Combining these facts, one arrives at
\begin{equation*}
C_\varrho(\theta_1\cdot g,g)=\int\mathrm d\varrho(t)\int\frac{\mathrm d^2k}{(2\pi)^2}\,\frac{|v_2|^2+|v_3|^2-|v_1|^2}{2E_t}.
\end{equation*}
Since $\operatorname{div}j=0$, integrating by parts yields $E_tv_1+i(k_2v_2+k_3v_3)=0$; hence $|v_1|^2\le\frac{|k|^2}{E_t^2}\bigl(|v_2|^2+|v_3|^2\bigr)$, and the integrand cannot be negative.
Reflection positivity of $\gamma^\varrho$ at $\{y_1=0\}$ is then a consequence of \cite[Lemma~9.11\,(d)]{QEDI}.
In the proof of \cite[Lemma~12.29]{QEDII}, the paragraph headed ``The law $\gamma$'' uses nothing about $\gamma$ except this property and the invariance under all Euclidean maps, and its conclusion is reflection positivity at every plane.
Accordingly $\gamma^\varrho$ has (H) and belongs to $\mathcal C_9$; by \cite[Theorem~1.4, Part~I\,(d), (f)]{QEDII} it satisfies (C7), and the uniqueness contained in that theorem gives $\mathrm{kl}_{\gamma^\varrho}=\varrho$.

\emph{Step 12: (d), the law $\mu_\xi\ast_\epsilon\gamma^\varrho$.}
Concerning $\gamma$, Steps~1--5, 7 and~9 require nothing beyond the properties proved for $\gamma^\varrho$ in Steps~10 and~11 and the vanishing of the cumulants of order at least $3$.
In Step~3 the second factor is the corresponding expression for $\gamma^{\epsilon\varrho}$, which factorises by \cite[Lemma~9.11\,(c)]{QEDI}; in Steps~5 and~9 one puts $C_\varrho(g,g')$ and $\varrho$ in the places of $\langle g,\Pi_{\mathrm c}g'\rangle$ and $\delta_0$.
The bound of Step~6 for $(1-\epsilon)^2\kappa_4^\xi(g_{\mathrm{ix}})$ holds for this law as well.
For the second moment, Step~6 gives $|S_2^\xi(g_{\mathrm{ix}},\allowbreak g_{\mathrm{ix}})-\|g_{\mathrm{ix}}\|^2|\le2C_U\,\allowbreak e^{2-\mathfrak e_2}$, and therefore the second moment $S_2$ of $\mu_\xi\ast_\epsilon\gamma^\varrho$ satisfies
\begin{equation*}
\bigl|S_2(g_{\mathrm{ix}},g_{\mathrm{ix}})-\|g_{\mathrm{ix}}\|^2\bigr|\le(1-\epsilon)\,2C_U\,e^{2-\mathfrak e_2}+\epsilon\bigl(\|g_{\mathrm{ix}}\|^2-C_\varrho(g_{\mathrm{ix}},g_{\mathrm{ix}})\bigr)\le C'_{\bar\kappa}\,e^{2-\mathfrak e_2}.
\end{equation*}
Here $C'_{\bar\kappa}\ge8C_{\bar\kappa}\ge128\,C_U>2C_U$ by the definitions of $C_U$, $C_{\bar\kappa}$ and $C'_{\bar\kappa}$ in \cite[Proposition~43.2, Constants]{QEDIII} (the inequality $2C_U\le C'_{\bar\kappa}$ is \cite[Chapter~44, (0.2)]{QEDIII}), and $0\le\|g\|^2-C_\varrho(g,g)\le\|g\|^2$.
\end{proof}


\section{The lattice Fisher estimate}

In the fixed-spacing lattice statements below, the lattice has unit spacing, and the notation is that of \cite[Section~2]{QEDII}.
Thus the cochains $C^k$, $C^k_{\mathrm c}$, the operators $d$, $d^{*}$, the pairing $\langle\cdot,\cdot\rangle$ and the norm $\|\cdot\|$ are those of \cite[Definition~2.1]{QEDII};
the space $\Omega$ and the condition (G) are those of \cite[Definition~2.5]{QEDII};
$\mathrm{Conf}(\Gamma)$, $\mathrm d(A,z)$ and $F=dA+\frac{2\pi}{q}z$ are those of \cite[Definition~2.6]{QEDII};
the sites $\mathbb X$, their neighbours $x^{(i)}$ and the map $\pi$ are those of \cite[Definition~2.4]{QEDII};
the operator $D_A$ is that of \cite[Definition~2.10]{QEDII};
the tori $\mathbb T=(\Gamma,\chi)$ and $C(n)$ are those of \cite[Definition~2.11]{QEDII};
and $Z_{\mathbb T}$, $\mathbb E_{\mathbb T}$, $\mu_{\mathbb T}$ are those of \cite[Definition~2.13]{QEDII}.
In the fixed-spacing statements, $\beta$ denotes a lattice two-cochain and $\nu$ a law on $\Omega$.
The lattice symmetry group is $W=\mathbb Z^3\rtimes O_{\mathrm{cub}}$.
The lattice charge is $q>0$, and $\mathrm{sw}:=\sum_iw_i$.
For $\rho\ge0$, the norm of a type~3 hopping matrix is
$\mathsf h_3=w_3(1+\rho)/2$ \cite[Lemma~3.8(a)]{QEDII};
for standard weights it equals $3(1+\rho)$.
The lattice mass, denoted $m$ in \cite[Definition~2.10]{QEDII} and $\boldsymbol\mu$ in \cite[(I.2.5)]{QEDIII}, is written
\begin{equation*}
\mathbf m:=M-\frac{\rho}{2}\,\mathrm{sw}.
\end{equation*}
The operator $D_A$ here includes $\mathbf m$; the massless convention
will be stated explicitly where it is used. For the continuum cutoffs,
$q=e\eta^{1/2}$ and $\mathbf m=\boldsymbol\mu_N$.
The symbol $d_{\mathbb T}$ denotes the graph distance of the fermion graph on $\mathbb X/\Gamma$, or on $\mathbb X$ when there is no torus.
Of all bonds, only those of type~$3$ change $\pi$, and they change it by a unit vector \cite[Proposition~3.6\,(c)]{QEDII}.
Hence
\begin{equation*}
d_{\mathbb T}(x,y)\ge|\pi(x)-\pi(y)|_1,
\end{equation*}
the right side being the $\ell^1$ distance of $\mathbb Z^3/\Gamma$.


\subsection{The positive-mass inverse}
\begin{named}{Lemma 5.1}
Let $\rho\ge0$, $w_0,\dots,w_3>0$, $q>0$, and let $M\in\mathbb R$ be such that $\mathbf m>0$.
Consider a real one-cochain $A$ and a space $\mathcal H$ from among those of \cite[Proposition~3.18]{QEDII}; thus $\mathcal H$ is $\ell^2(\mathbb X;\mathbb C^4)$ or the space of spinors of a pair $(\Gamma,\chi)$, and in the latter case $A$ is assumed to be $\Gamma$-periodic.
Put
\begin{equation*}
\vartheta:=\log\Bigl(1+\frac{\mathbf m}{(1+\rho)\,\mathrm{sw}}\Bigr).
\end{equation*}
Then $D_A$ is invertible on $\mathcal H$, $\|D_A^{-1}\|\le\mathbf m^{-1}$, and for all sites $x,y$ the $4\times4$ blocks of the inverse satisfy
\begin{equation*}
\|D_A^{-1}(x,y)\|\le\frac{2}{\mathbf m}\,e^{-\vartheta\,d_{\mathbb T}(x,y)}.
\end{equation*}
The same statements hold for $\check D_A$ and its $2\times2$ blocks.
The constants are independent of $A$, $q$, $\Gamma$ and $\chi$.
\end{named}

\begin{proof}
\emph{Step 1.}
By \cite[Proposition~3.18\,(b)]{QEDII} one has $\operatorname{Re}D_A\ge\mathbf m$.
Therefore
\begin{equation*}
\mathbf m\|\psi\|^2\le\operatorname{Re}\langle\psi,D_A\psi\rangle\le\|\psi\|\,\|D_A\psi\|,
\end{equation*}
and the same inequalities hold with $D_A^*$ in place of $D_A$.
It follows that $D_A$ is invertible and that $\|D_A^{-1}\|\le\mathbf m^{-1}$.
Only the lower bound on the real part enters this argument.

\emph{Step 2.}
Fix a site $y$ and $N<\infty$, and let $\phi:=\min(d_{\mathbb T}(\cdot,y),N)$.
This function is bounded and $\Gamma$-periodic, and $|\phi(x)-\phi(x^{(i)})|\le1$.
One has $e^{\vartheta\phi}D_Ae^{-\vartheta\phi}=D_A-R$, where
\begin{equation*}
(R\psi)(x)=\sum_i\bigl(e^{\vartheta(\phi(x)-\phi(x^{(i)}))}-1\bigr)\,w(x,x^{(i)})\,U_A(x,x^{(i)})\,\psi(x^{(i)}).
\end{equation*}
Repeating the computation in the proof of \cite[Proposition~3.18\,(c)]{QEDII} with the additional factors, each of modulus at most $e^{\vartheta}-1$, one obtains
\begin{equation*}
\|R\|\le(e^{\vartheta}-1)\cdot\frac{1+\rho}{2}\,\mathrm{sw}=\frac{\mathbf m}{2}.
\end{equation*}

\emph{Step 3.}
Consequently $\operatorname{Re}(D_A-R)\ge\frac{\mathbf m}{2}$, and Step~1 gives $\|(D_A-R)^{-1}\|\le\frac{2}{\mathbf m}$.
Since $\phi(y)=0$ and
\begin{equation*}
D_A^{-1}(x,y)=e^{-\vartheta(\phi(x)-\phi(y))}\,(D_A-R)^{-1}(x,y),
\end{equation*}
the bound on the blocks follows on letting $N\to\infty$.
The statements for $\check D_A$ follow, since $\check D_A$ is a diagonal block of $D_A$ \cite[Proposition~3.18\,(d)]{QEDII}.
\end{proof}

Conjugation with exponential weights is the method of \cite[Lemma~15.23]{QEDI} and \cite[Theorem~12.15]{QEDI}.
Those two results deal with real symmetric positive matrices and with blocked Wilson operators at zero field, and neither of them applies to $D_A$ in an arbitrary real field.

\subsection{The finite-mass polarisation bound}
\begin{named}{Lemma 5.2}
Let $\rho\ge0$, $w_0,\dots,w_3>0$ and $\mathbf m>0$.
For $q>0$, an integer $n\ge4$ and $\Gamma=n\mathbb Z^3$, let $\ell_n(A):=\log\det D^{C(n)}_A$, and define the polarisation form
\begin{equation*}
K_n(u,u'):=-\,\mathbb E_{C(n)}\bigl[\partial_u\partial_{u'}\ell_n(A)\bigr]\qquad(u,u'\in C^1(\Gamma)),
\end{equation*}
where $\partial_u$ is the derivative along $A\mapsto A+tu$.
There exists $c_*=c_*(\mathbf m,\rho,w)<\infty$, independent of $q$ and of $n$, such that
\begin{equation*}
K_n(u,u)\le c_*\,q^2\Bigl(\|du\|^2+n^{-5}\,\Bigl|\sum_xu(x)\Bigr|^2\Bigr)
\end{equation*}
for all $u\in C^1(\Gamma)$, where $\sum_xu(x)\in\mathbb R^3$ denotes the vector with the components $\sum_{x\in\mathbb Z^3/\Gamma}u_\mu(x)$.
\end{named}


\begin{proof}
Lemma~5.8 gives
\[
|K_n(u,u)|\le\mathsf c_n\left(\frac{\pi^2}{4}\|du\|^2
+2\pi^2n^{-5}\left|\sum_xu(x)\right|^2\right).
\]
Apply Theorem~5.7(i). One may take
$c_*=1152\pi^4\mathsf h_3^4\mathbf m^{-4}$.
\end{proof}

\subsection{Shift identities at a fixed spacing}
\begin{named}{Proposition 5.3}
Let $0<\rho<1$, $w_0,\dots,w_3>0$, $q>0$, and let $M\in\mathbb R$ be such that $\mathbf m>0$.
Fix integers $n_i\ge4$ with $n_i\to\infty$ and, as in \cite[Theorem~1.2]{QEDII}, a weak limit point $\nu$ of the laws $\mu_{C(n_i)}$.
On the space of square summable two-cochains, let $P_{\mathrm{ex}}$ be the orthogonal projection whose range is the closure of $dC^1_{\mathrm c}$.
Let $c_*$ be the constant of Lemma~5.2.
\begin{itemize}
\item[(a)] For every $\beta\in C^2_{\mathrm c}$,
\begin{equation*}
(1+c_*q^2)^{-1}\,\|P_{\mathrm{ex}}\beta\|^2\le\int\langle\beta,F\rangle^2\,\mathrm d\nu\le\|P_{\mathrm{ex}}\beta\|^2.
\end{equation*}
The right side is the variance of $\langle\beta,F\rangle$ under the Gaussian weight $\exp(-\frac12\|dA\|^2)$.
\item[(b)] For every plaquette $\sigma$,
\begin{equation*}
\sum_{v\in\mathbb Z^3}\Bigl|\int F_\sigma F_{\sigma+v}\,\mathrm d\nu\Bigr|=\infty.
\end{equation*}
Exponential decay of the two-point function is thereby excluded.
\item[(c)] Let $T_\mu$ be the transfer operator of \cite[Definition~10.11]{QEDII}, where $\mu=1,2,3$.
For every $\varepsilon\in(0,1)$ the spectral projection of $T_\mu$ belonging to the interval $(1-\varepsilon,1)$ is different from zero.
In particular $1$ is a limit point of the spectrum of $T_\mu$.
\end{itemize}
\end{named}

In the proof below, Steps~6 and~7 give the upper bound in (a), Steps~1 to~5 the lower bound, and Steps~8 to~15 parts (b) and (c).

\begin{proof}[Proof of Proposition~5.3]
\emph{Step 1} (flow on $C(n)$).
In Steps~1 to~3 the assumption $\rho\ge0$ suffices.
The density of $Z_{C(n)}$ is $e^{-V}$, where $V:=\frac12\|F\|^2-\ell_n(A)$.
Let $u\in C^1(\Gamma)$ and $t\in\mathbb R$.
The shift $(A,z)\mapsto(A+tu,z)$ commutes with each of the maps $(A,z)\mapsto(A+\frac{2\pi}{q}k,\,z-dk)$, $k\in C^1(\Gamma;\mathbb Z)$, by which the equivalence relation defining $\mathrm{Conf}(\Gamma)$ is given, and its inverse is the shift by $-tu$;
consequently it maps each equivalence class onto an equivalence class, and it descends to $\mathrm{Conf}(\Gamma)$.
Since Lebesgue measure is invariant under translations and $z$ is left unchanged, the shift preserves the product of $\mathrm dA$ with the counting measure on the closed $z$.
For each closed $z$, the section at $z$ of the image of the fundamental domain given by the box is the box translated by $tu$;
this image is thus a Borel set of pairs, and each equivalence class meets it in a single pair, because the fundamental domain given by the box has this property and the shift is a bijection of the pairs which maps equivalence classes onto equivalence classes.
The clause of \cite[Definition~2.6]{QEDII} on such sets allows it to take the place of the fundamental domain given by the box in the definition of $\mathrm d(A,z)$, and therefore the shift preserves $\mathrm d(A,z)$.

\emph{Step 2} (integration by parts).
Let $\Phi$ be smooth in $A$, and let $\Phi$ and $\partial_u\Phi$ grow at most polynomially in $F$.
The integral $\int\mathrm d(A,z)\,(\Phi e^{-V})(A+tu,z)$ is independent of $t$.
For $|t|\le1$ the $t$-derivative of the integrand is bounded by a constant multiple of $\prod_\sigma\exp(-\frac14F_\sigma^2)$, which is integrable \cite[Lemma~3.14]{QEDII}.
Differentiating at $t=0$ one obtains $\mathbb E[\partial_u\Phi]=\mathbb E[\Phi\,\partial_uV]$.

\emph{Step 3} (the Schwarz inequality).
For $\Phi=\langle\beta,F\rangle$ and for $\Phi=\partial_uV=\langle F,du\rangle-\partial_u\ell_n$, Step~2 reads
\begin{equation*}
\mathbb E[\langle\beta,F\rangle\,\partial_uV]=\langle\beta,du\rangle,\qquad\mathbb E[(\partial_uV)^2]=\|du\|^2+K_n(u,u).
\end{equation*}
Consequently, for all $\beta\in C^2(\Gamma)$ and $u\in C^1(\Gamma)$,
\begin{equation*}
\langle\beta,du\rangle^2\le\mathbb E_{C(n)}\langle\beta,F\rangle^2\cdot\bigl(\|du\|^2+K_n(u,u)\bigr).
\end{equation*}

\emph{Step 4} (limit of the moments).
Let $\beta\in C^2_{\mathrm c}$ and $u\in C^1_{\mathrm c}$.
Once $i$ is large, the periodisations of $\beta$ and $u$ have the same norms and pairings as $\beta$ and $u$ themselves, and $\langle\beta,\widetilde F\rangle$ coincides with the pairing on the torus \cite[Lemma~2.14\,(b)]{QEDII}.
On every $C(n)$ one has $\mathbb E\langle\beta,F\rangle^4\le\frac32\exp(2)\,\|\beta\|^4$ \cite[Theorem~6.27\,(b)]{QEDII}, and therefore the second moments converge along the sequence.

\emph{Step 5} (lower bound in (a)).
Since $\sum_xu(x)$ is fixed, Lemma~5.2 gives $\limsup_iK_{n_i}(u,u)\le c_*q^2\|du\|^2$.
Hence
\begin{equation*}
\langle\beta,du\rangle^2\le(1+c_*q^2)\,\|du\|^2\int\langle\beta,F\rangle^2\,\mathrm d\nu,
\end{equation*}
and taking the supremum over $u\in C^1_{\mathrm c}$ yields the lower bound.

\emph{Step 6} (upper bound in (a)).
By \cite[Theorem~1.2]{QEDII}, the law $\nu$ satisfies (G) and is symmetric under $F\mapsto-F$.
Expanding (G) for $t\beta$ to the order $t^2$ one finds $\int\langle\beta,F\rangle^2\,\mathrm d\nu\le\|\beta\|^2$.
Since $\langle d^*\gamma,F\rangle=0$ on $\Omega$ for every $\gamma\in C^3_{\mathrm c}$ \cite[Lemma~2.14\,(a)]{QEDII}, the cochain $\beta$ can be replaced by $\beta-d^*\gamma$.

\emph{Step 7} (Hodge decomposition).
The infimum of $\|\beta-d^*\gamma\|^2$ over $\gamma$ equals $\|P_{\mathrm{ex}}\beta\|^2$.
Put $\mathsf s_\mu(p)=2\sin(p_\mu/2)$.
Indeed, the square summable two-cochains form the orthogonal sum of the closures of $dC^1_{\mathrm c}$ and of $d^*C^3_{\mathrm c}$, because a square summable $\beta$ with $d\beta=0$ and $d^*\beta=0$ satisfies $|\mathsf s(p)|^2\hat\beta(p)=0$.

\emph{Step 8} (Fourier form).
Let $\beta$ be identified with the triple $(\beta_{23},-\beta_{13},\beta_{12})$, and let Fourier transforms be taken with respect to the centres of the plaquettes.
Then $P_{\mathrm{ex}}$ acts as multiplication by $1-\mathsf s\,\mathsf s^{T}/|\mathsf s|^2$.
Write $\varpi:=1-\mathsf s_1^2/|\mathsf s|^2$.

\emph{Step 9} (reduction of (b)).
By the invariance under $W$ it suffices to take $\sigma=(0;23)$; let $g(v):=\int F_\sigma F_{\sigma+v}\,\mathrm d\nu$.
Suppose that $g$ is summable.
Then $\hat g$ is continuous, and (a), applied to $\beta=\sum_vc_v1_{\sigma+v}$, becomes
\begin{equation*}
(1+c_*q^2)^{-1}\int|\hat c|^2\varpi\le\int|\hat c|^2\hat g\le\int|\hat c|^2\varpi.
\end{equation*}

\emph{Step 10} (conclusion for (b)).
Fej\'er kernels for $|\hat c|^2$ then give $(1+c_*q^2)^{-1}\varpi\le\hat g\le\varpi$ on $[-\pi,\pi]^3\setminus\{0\}$.
However, $\varpi=0$ on the $p_1$-axis and $\varpi=1$ on the plane $p_1=0$, so that $\hat g$ cannot have a limit at $0$.

\emph{Step 11} (set-up for (c)).
It suffices to treat $\mu=1$, the other directions following by the invariance under $W$.
For integers $r\ge1$ and $s$, let $\beta^s$ be equal to $1$ on the plaquettes $((s,y);23)$ with $y\in\{0,\dots,r-1\}^2$ and to $0$ elsewhere, and put
\begin{equation*}
\Phi_s:=\langle\beta^s,F\rangle,\qquad g(t):=\int\Phi_0\Phi_t\,\mathrm d\nu.
\end{equation*}
The reflection $\theta_1$ of \cite[Definition~10.11]{QEDII} sends $\Phi_s$ to $\Phi_{-1-s}$, and the translation $\tau_{e_1}$ sends $\Phi_s$ to $\Phi_{s+1}$.

\emph{Step 12} (spectral representation).
By \cite[Lemma~10.9\,(a)]{QEDII} the class map extends to the $L^2(\nu)$-closure of $\mathcal O^+_1$, to which $\Phi_0$ belongs.
By \cite[Lemma~10.9\,(b), (d)]{QEDII},
\begin{equation*}
g(t)=\bigl\langle[\Phi_0],T_1^{|t|-1}[\Phi_0]\bigr\rangle=\int\lambda^{|t|-1}\,\mathrm d\mathfrak e(\lambda)\ge0\qquad(|t|\ge1),
\end{equation*}
where $\mathfrak e$ is the spectral measure of $[\Phi_0]$ on $[0,1]$; its mass is $g(1)\le g(0)$.

\emph{Step 13} (the sum of $g$).
For $N\ge1$,
\begin{equation*}
\sum_{|t|<N}\Bigl(1-\frac{|t|}{N}\Bigr)\,g(t)=\frac1N\int\Bigl\langle\sum_{s=0}^{N-1}\beta^s,F\Bigr\rangle^2\,\mathrm d\nu.
\end{equation*}
By (a) this quantity lies between $(1+c_*q^2)^{-1}J_N$ and $J_N$, where
\begin{equation*}
J_N:=\frac1N\Bigl\|P_{\mathrm{ex}}\sum_{s<N}\beta^s\Bigr\|^2\le r^2,
\end{equation*}
and $J_N\to r^2$ by Step~8 and Fej\'er's theorem in the variable $p_1$.
As $g\ge0$, it follows that $(1+c_*q^2)^{-1}r^2\le\sum_tg(t)\le r^2$; in particular $\mathfrak e(\{1\})=0$.

\emph{Step 14} (the value $g(0)$).
Write $k=(p_2,p_3)$, and let $1_r$ be the indicator function of $\{0,\dots,r-1\}^2$.
By (a) and Step~8,
\begin{align*}
g(0)\le\|P_{\mathrm{ex}}\beta^0\|^2
&=\int\frac{\mathrm d^2k}{(2\pi)^2}\,|\hat{1}_r(k)|^2\,|\mathsf s(k)|\,\bigl(4+|\mathsf s(k)|^2\bigr)^{-1/2}\\
&\le\frac12\|1_r\|\,\|\nabla 1_r\|=r^{3/2}.
\end{align*}

\emph{Step 15} (conclusion for (c)).
Suppose that the spectral projection of $(1-\varepsilon,1)$ vanishes.
Then
\begin{equation*}
\sum_tg(t)=g(0)+2\int(1-\lambda)^{-1}\,\mathrm d\mathfrak e(\lambda)\le\Bigl(1+\frac2\varepsilon\Bigr)\,g(0)\le\Bigl(1+\frac2\varepsilon\Bigr)\,r^{3/2},
\end{equation*}
which contradicts Step~13 as soon as $r>(1+c_*q^2)^2(1+\frac2\varepsilon)^2$.
\end{proof}

\subsection{The flavoured meson}
\begin{named}{Proposition 5.4}
Let $\rho,M\in\mathbb R$, $w_0,\dots,w_3>0$ and $q>0$, and let $\mathbb T=(\Gamma,\chi)$ be any torus.
Write $\psi(x)=(\psi^{(1)}(x),\psi^{(2)}(x))$ and $\bar\psi(x)=(\bar\psi^{(1)}(x),\bar\psi^{(2)}(x))$ according to the decomposition $\mathbb C^4=\mathbb C^2\oplus\mathbb C^2$ in which $\gamma_\mu=\operatorname{diag}(\sigma_\mu,-\sigma_\mu)$ \cite[Definition~2.7]{QEDII}, and define
\begin{equation*}
\mathfrak M(x):=\bar\psi^{(1)}(x)\,\psi^{(2)}(x),\qquad\mathfrak M^\star(x):=-\,\bar\psi^{(2)}(x)\,\psi^{(1)}(x).
\end{equation*}
The elements $\mathfrak M(x)$ and $\mathfrak M^\star(x)$ of $\operatorname{Gr}_{\mathbb T}\subset\mathcal B(\mathbb T)$ are even, $\Gamma$-periodic, and invariant under the gauge transformations of \cite[Definition~2.10]{QEDII}.
\begin{itemize}
\item[(a)] For every mirror $\Pi$ and every site $x$ one has $\Theta_\Pi\,\mathfrak M(x)=\mathfrak M^\star(\theta_\Pi x)$ in $\operatorname{Gr}(\mathbb X)$, with $\Theta_\Pi$ as in \cite[Definition~2.9]{QEDII}.
On the torus, $\Theta_\Pi$ maps $\operatorname{Gr}_{\mathbb T}$ onto $\operatorname{Gr}_{\theta_\Pi\mathbb T}$ \cite[Lemma~3.12]{QEDII}, and the identity holds for the classes, its right side being formed in $\operatorname{Gr}_{\theta_\Pi\mathbb T}$.
If $\theta_\Pi\mathbb T=\mathbb T$, as for a torus adapted to the family of $\Pi$ \cite[Section~1.1]{QEDII} and for the cubic torus $C(n)$ of side $n$, $n\ge4$ an integer, both sides lie in $\operatorname{Gr}_{\mathbb T}$.
\item[(b)] Let $\operatorname{adj}$ denote the adjugate matrix and $(y,x)$ its $2\times2$ block.
For all sites $x,y$,
\begin{equation*}
Z_{\mathbb T}\bigl[\mathfrak M^\star(y)\,\mathfrak M(x)\bigr]=\int_{\mathrm{Conf}(\Gamma)}\mathrm d(A,z)\,e^{-\frac12\|F\|^2}\,\bigl\|\operatorname{adj}\bigl(\check D^{\mathbb T}_A\bigr)(y,x)\bigr\|_{\mathrm{HS}}^2\ge0.
\end{equation*}
\item[(c)] If $\rho\ge0$ and $\mathbf m>0$, then, with $\vartheta$ as in Lemma~5.1,
\begin{gather*}
\mathbb E_{\mathbb T}\bigl[\mathfrak M^\star(y)\,\mathfrak M(x)\bigr]=\mathbb E_{\mathbb T}\bigl\|(\check D^{\mathbb T}_A)^{-1}(y,x)\bigr\|_{\mathrm{HS}}^2\le\frac{8}{\mathbf m^2}\,e^{-2\vartheta\,d_{\mathbb T}(x,y)},\\
\frac{2}{\bigl(|M|+\frac{1+\rho}{2}\,\mathrm{sw}\bigr)^2}\le\sum_{y\in\mathbb X/\Gamma}\mathbb E_{\mathbb T}\bigl[\mathfrak M^\star(y)\,\mathfrak M(x)\bigr]\le\frac{2}{\mathbf m^2}.
\end{gather*}
\end{itemize}
The constants do not depend on $\mathbb T$ or on $q$.
\end{named}

\begin{proof}
\emph{Step 1} (proof of (a)).
One has $\mathfrak M(x)=\bar\psi(x)\,Y\psi(x)$, where the only non-zero block of $Y$ is $Y_{12}=1$.
By \cite[Lemma~3.10\,(c)]{QEDII},
\begin{equation*}
\Theta_\Pi\mathfrak M(x)=\bar\psi(\theta_\Pi x)\,\gamma_nY^*\gamma_n\,\psi(\theta_\Pi x),
\end{equation*}
and $\gamma_nY^*\gamma_n=-Y^*$, since $\gamma_n=\operatorname{diag}(\sigma_n,-\sigma_n)$ and $\sigma_n^2=1$.
On the torus, \cite[Lemma~3.12]{QEDII} gives $\Theta_\Pi[\mathfrak M(x)]_{\mathbb T}=\allowbreak [\Theta_\Pi\mathfrak M(x)]_{\theta_\Pi\mathbb T}$, and the identity passes to the classes.

\emph{Step 2} (proof of (b)).
Write $\check D=\check D^{\mathbb T}_A$.
By \cite[Proposition~3.18\,(d)]{QEDII},
\begin{equation*}
S_{\mathbb T}(A)=\bar\psi^{(1)}\check D\,\psi^{(1)}+\bar\psi^{(2)}\check D^{*}\psi^{(2)}.
\end{equation*}
Rearranging even brackets, one obtains
\begin{equation*}
\mathfrak M^\star(y)\,\mathfrak M(x)=\sum_{a,b}\bigl[\psi^{(1)}_b(y)\,\bar\psi^{(1)}_a(x)\bigr]\bigl[\psi^{(2)}_a(x)\,\bar\psi^{(2)}_b(y)\bigr].
\end{equation*}

\emph{Step 3.}
Since the Berezin integral picks out the coefficient of a product of even pairs \cite[Definition~2.8]{QEDII}, it factorises on the product of an even element of the first flavour with an even element of the second.
For every square matrix $B$, differentiating $\int\exp(-\bar\psi B\psi)=\det B$ with respect to $B_{ji}$ gives
\begin{equation*}
\int\exp(-\bar\psi B\psi)\,\psi_i\bar\psi_j=\frac{\partial\det B}{\partial B_{ji}}=\operatorname{adj}(B)_{ij}.
\end{equation*}

\emph{Step 4.}
Applying this with $B=\check D$ and with $B=\check D^*$, and using $\operatorname{adj}(\check D^*)=\operatorname{adj}(\check D)^*$, one finds
\begin{equation*}
\int_{\mathbb T}e^{-S_{\mathbb T}(A)}\,\mathfrak M^\star(y)\,\mathfrak M(x)=\sum_{a,b}\bigl|\operatorname{adj}(\check D)_{(y,b),(x,a)}\bigr|^2.
\end{equation*}
Convergence of the integral over $\mathrm{Conf}(\Gamma)$ follows from \cite[Lemma~3.13\,(b)]{QEDII}.

\emph{Step 5} (proof of (c)).
Lemma~5.1 shows that $\check D$ is invertible.
Since $\operatorname{adj}(\check D)=\det\check D\cdot\check D^{-1}$ and $|\det\check D|^2=\det D^{\mathbb T}_A$ \cite[Proposition~3.18\,(d)]{QEDII}, the identity follows.
The upper bound is a consequence of Lemma~5.1 and of the inequality $\|B\|_{\mathrm{HS}}^2\le2\|B\|^2$ for $2\times2$ blocks $B$.

\emph{Step 6.}
Finally,
\begin{equation*}
\sum_y\|\check D^{-1}(y,x)\|_{\mathrm{HS}}^2=\operatorname{tr}_{\mathbb C^2}\bigl[(\check D\check D^*)^{-1}(x,x)\bigr],
\end{equation*}
and $\|\check D\|^{-2}\le(\check D\check D^*)^{-1}\le\mathbf m^{-2}$ by \cite[Proposition~3.18\,(c)]{QEDII} and Lemma~5.1.
\end{proof}

For the continuum the notation fixed above is used: $\beta$ denotes a probe, $a=d^*\beta$ and $f=dd^*\beta$, the charge is $e>0$, and $\|\cdot\|$, $\langle\cdot,\cdot\rangle$ stand, on forms, for the norm and the inner product of $L^2(\mathbb R^3;\Lambda^2)$.
The exponential function is written $\exp$.

\subsection{The variance identity}

\begin{named}{Lemma 5.5}
In this lemma only, $D_A$ denotes the massless chamber operator of
\cite[(I.2.3)]{QEDIII}, with standard weights, and $E_n=\mathbb E_{C(n)}$ is formed at
lattice mass $\boldsymbol\mu$. Thus $D_A+\boldsymbol\mu$ is the
massive operator denoted by $D_A$ elsewhere in this section.

Suppose that $-1\le\rho\le1$, $q>0$, $\boldsymbol\mu\in\mathbb R$ and $n\ge4$, and that $u$ is a real one-cochain on $\mathbb Z^3/n\mathbb Z^3$ with $du\ne0$.
Write $\partial_u$ for the derivative along $A\mapsto A+tu$, and define
\begin{equation*}
\ell(A):=\log\det(D_A+\boldsymbol\mu),\qquad x:=\frac{E_n\langle du,F\rangle^2}{\|du\|^2},\qquad j:=\frac{E_n(\partial_u\ell)^2}{\|du\|^2}.
\end{equation*}
Then $j<\infty$, $x\le1$, and
\begin{equation*}
E_n\bigl(\langle du,F\rangle-\partial_u\ell\bigr)^2=(2-x+j)\,\|du\|^2,\qquad(1-x)^2\le j\,x,\qquad x\ge1-\sqrt j .
\end{equation*}
\end{named}

\begin{proof}
\begin{itemize}
\item[1.] Let $\check D_A$ be the upper block of the operator, mass included; then $\det(D_A+\boldsymbol\mu)=|\det\check D_A|^2$ \cite[Proposition~3.18\,(d)]{QEDII}.
This is a trigonometric polynomial in $qA$, which does not vanish identically, since $Z_n[1]>0$ \cite[Theorem~4.26]{QEDII}.
Hence $\ell$ is smooth off a null set, and
\begin{equation*}
\frac{\bigl(\partial_u\det(D_A+\boldsymbol\mu)\bigr)^2}{\det(D_A+\boldsymbol\mu)}\le4\,|\partial_u\det\check D_A|^2
\end{equation*}
is bounded; therefore $j<\infty$.
\item[2.] The density $\varrho:=\exp(-\frac12\|F\|^2)\det(D_A+\boldsymbol\mu)$ is smooth, and almost everywhere one has $\partial_u\varrho=-\varrho\,s$, where $s:=\langle du,F\rangle-\partial_u\ell$.
The argument of Steps~1 and~2 of the proof of Proposition~5.3, with $\varrho$ substituted for $\exp(-V)$, yields $E_n[\partial_u\Phi]=E_n[\Phi\,s]$ for every $\Phi$ smooth in $A$ such that $\Phi$ and $\partial_u\Phi$ grow at most polynomially in $F$; it makes no use of the sign of $\boldsymbol\mu$.
\item[3.] For $\Phi=\langle du,F\rangle$ this reads $E_n[\langle du,F\rangle\,s]=\|du\|^2$, that is, $E_n[\langle du,F\rangle\,\partial_u\ell]=(x-1)\|du\|^2$.
Expanding $s^2$, one obtains the first identity.
By the Schwarz inequality, $\|du\|^4\le x\|du\|^2\cdot E_ns^2$, that is, $1\le x(2-x+j)$, which is the inequality $(1-x)^2\le jx$.
\item[4.] The bound $x\le1$ is \cite[Theorem~6.27\,(b)]{QEDII}; hence $1-x\le\sqrt{jx}\le\sqrt j$.
\qedhere
\end{itemize}
\end{proof}

\subsection{Scores and the continuum limit}
\begin{named}{Proposition 5.6}
Let $-1\le\rho\le1$, $w_0,\dots,w_3>0$, $q>0$ and $M\in\mathbb R$; no assumption is made on the sign of $\mathbf m$.
Let $n\ge4$ be an integer and $\Gamma=n\mathbb Z^3$.
Denote by
\begin{equation*}
\varrho:=\exp(-\tfrac12\|F\|^2)\,\det D^{C(n)}_A\ge0
\end{equation*}
the density of $Z_{C(n)}$ with respect to $\mathrm d(A,z)$.
Given $u\in C^1(\Gamma)$, define $s_u:=\partial_u\varrho/\varrho$ on the set where $\varrho>0$, and $s_u:=0$ on the set where $\varrho=0$.
The Fisher information of the shift along $u$ is
\begin{equation*}
I_n(u):=\mathbb E_{C(n)}[s_u^2].
\end{equation*}
\begin{itemize}
\item[(a)] \emph{One torus.}
One has $I_n(u)<\infty$, $\mathbb E_{C(n)}[s_u]=0$,
\begin{equation*}
\|du\|^4\le\mathbb E_{C(n)}\langle du,F\rangle^2\cdot I_n(u)\qquad\text{and}\qquad I_n(u)\ge\|du\|^2.
\end{equation*}
On the set where $\varrho>0$ put $\ell_n:=\log\det D^{C(n)}_A$; the set where $\varrho=0$ is a null set for $\mathrm d(A,z)$.
Whatever $u,u'\in C^1(\Gamma)$, the function $\partial_u\partial_{u'}\ell_n$ is integrable with respect to $Z_{C(n)}[1]^{-1}\varrho\,\mathrm d(A,z)$, a probability measure, and the formula
\begin{equation*}
K_n(u,u'):=-\,\mathbb E_{C(n)}\bigl[\partial_u\partial_{u'}\ell_n\bigr]
\end{equation*}
defines on $C^1(\Gamma)$ a bilinear form that is symmetric and positive semidefinite; for $\rho\ge0$ and $\mathbf m>0$ it coincides with the form of Lemma~5.2.
For every real $M$, that is, whatever the real lattice mass $\mathbf m$, and for all $u,u'\in C^1(\Gamma)$, those with $du=0$ included,
\begin{gather*}
\mathbb E_{C(n)}[s_u\,s_{u'}]=\langle du,du'\rangle+K_n(u,u'),\qquad I_n(u)=\|du\|^2+K_n(u,u),\\
K_n(u,u)=\|du\|^2-\mathbb E_{C(n)}\langle du,F\rangle^2+\mathbb E_{C(n)}(\partial_u\ell_n)^2\ge\mathbb E_{C(n)}(\partial_u\ell_n)^2.
\end{gather*}
Suppose in addition that the weights are standard and that $du\ne0$, and let $x$ and $j$ be the numbers that Lemma~5.5 attaches to $u$ at $\boldsymbol\mu=\mathbf m$.
Then
\begin{equation*}
I_n(u)=(2-x+j)\,\|du\|^2,\qquad K_n(u,u)=(1-x+j)\,\|du\|^2,
\end{equation*}
and consequently $j\le K_n(u,u)/\|du\|^2\le j+\min\{1,\sqrt j\}$.
\end{itemize}
For (b), (c) and (d), assume $0<\rho<1$ and standard weights, fix $\bar m\in(0,1]$, $e>0$ and an arbitrary mass assignment $N\mapsto m_{0,N}$, and let $\mu$ be a limit point of $(P_{N,T})$ along a defining sequence $(N_i,T_i)$.
Set $\eta_i=L^{-N_i}$, $n_i=8hL^{T_i+N_i}$ and $R_i=8hL^{T_i}$, and form $E_i:=\mathbb E_{C(n_i)}$ and $I_{n_i}$ with $q=e\,\eta_i^{1/2}$ and $\mathbf m=\boldsymbol\mu_{N_i}=\bar m\eta_i+m_{0,N_i}$.
Let $\beta$ be a probe.
For every index $i$ with $R_i>D_\beta$, the lattice probe $\hat{\mathfrak a}_i:=\eta_i^{1/2}a_{\eta_i}[\beta]$ of \cite[Proposition~47.1, grade~1 at $v=0$]{QEDIII} is defined on the torus having $n_i$ sites per side.
Put
\begin{equation*}
\bar I(\beta):=\liminf_{i\to\infty}I_{n_i}(\hat{\mathfrak a}_i)\in\bigl[\,\|f\|^2,\infty\,\bigr].
\end{equation*}
\begin{itemize}
\item[(b)] \emph{Limit points.}
$S_2(f,f)\ge\|f\|^4/\bar I(\beta)$.
\item[(c)] \emph{Integration by parts.}
Suppose that $\bar I(\beta)<\infty$.
Then some $\mathsf s_\beta\in L^2(\mu)$ with $\|\mathsf s_\beta\|_2^2\le\bar I(\beta)$ satisfies $\mu(\partial_f\Phi)=\mu(\Phi\,\mathsf s_\beta)$ for every $\Phi=\phi(X(g_1),\dots,X(g_k))$ in which $\phi$ is smooth and bounded together with all its derivatives and the $g_j\in L^2(\mathbb R^3;\Lambda^2)$ are real; here
\begin{equation*}
\partial_f\Phi:=\sum_j\langle g_j,f\rangle\,(\partial_j\phi)(X(g_1),\dots,X(g_k)).
\end{equation*}
If $J(a)$ is defined by $J(a):=e^{-1}(X(f)-\mathsf s_\beta)$, this identity takes the form $\mu(\partial_f\Phi)=\mu\bigl(\Phi\cdot(X(f)-e\,J(a))\bigr)$.
Moreover, the image $\mu_t$ of $\mu$ under $\omega\mapsto\omega+t\,f$ satisfies $|\mu_t(\Phi)-\mu(\Phi)|\le|t|\,\bar I(\beta)^{1/2}$ for every real $t$ and every $\Sigma$-measurable $\Phi$ with $|\Phi|\le1$.
\item[(d)] \emph{The atom at zero.}
Assume that $\mu$ is volume-first and that $\bar I(\beta)\le(1+c)\,\|f\|^2$ for all probes $\beta$, with a constant $c<\infty$.
Then $\mathrm{kl}_\mu(\{0\})\ge(1+c)^{-1}$.
\end{itemize}
\end{named}

\begin{proof}
\emph{Step 1} (finiteness in (a)).
Let $\check D:=\check D^{C(n)}_A$.
Then $\det D^{C(n)}_A=|\det\check D|^2$ \cite[Proposition~3.18\,(d)]{QEDII}, and $\det\check D$ is a trigonometric polynomial in the variables $qA(b)$.
In particular $\varrho$ is smooth, and both $\det\check D$ and $\partial_u\det\check D$ are bounded.
At a point where $\varrho=0$ the non-negative function $\varrho$ attains a minimum, so that $\partial_u\varrho=\varrho\,s_u$ everywhere.
Where $\varrho>0$ one has $s_u=\allowbreak -\langle du,\allowbreak F\rangle+2\operatorname{Re}(\partial_u\det\check D/\det\check D)$, and consequently
\begin{equation*}
\varrho\,s_u^2\le2\exp(-\tfrac12\|F\|^2)\bigl(\langle du,F\rangle^2\,|\det\check D|^2+4\,|\partial_u\det\check D|^2\bigr).
\end{equation*}
The right side is integrable \cite[Lemma~3.14]{QEDII}, and $Z_{C(n)}[1]>0$ \cite[Theorem~4.26]{QEDII}.

\emph{Step 2} (the inequalities in (a)).
The argument of Steps~1 and~2 in the proof of Proposition~5.3, carried out with $\varrho$ in place of $\exp(-V)$, gives
\begin{equation*}
\mathbb E[\partial_u\Phi]=-\mathbb E[\Phi\,s_u]
\end{equation*}
for every $\Phi$ smooth in $A$ such that $\Phi$ and $\partial_u\Phi$ grow at most polynomially in $F$; the sign of $\mathbf m$ plays no part here.
Taking $\Phi=1$ one finds $\mathbb E[s_u]=0$.
With $\Phi=\langle du,F\rangle$ one obtains $\|du\|^2=-\mathbb E[\langle du,F\rangle\,s_u]$, and the Schwarz inequality yields the first inequality.
The second follows from it together with $\mathbb E\langle du,F\rangle^2\le\|du\|^2$ \cite[Theorem~6.27\,(b)]{QEDII}.
The identities of (a) in which $K_n$ appears are established in Step~2a below.
For the identity $I_n(u)=(2-x+j)\,\|du\|^2$, let the weights be standard and $du\ne0$.
By \cite[(I.2.5), Proposition~46\,(b)]{QEDIII}, the operator $D^{C(n)}_A$ is then the operator $D_A+\boldsymbol\mu$ of Lemma~5.5 at $\boldsymbol\mu=\mathbf m$, and $\mathbb E_{C(n)}$ is the expectation $E_n$ of that Lemma.
Where $\varrho>0$ one has $s_u=-(\langle du,F\rangle-\partial_u\ell)$, with the function $\ell$ of Lemma~5.5, so that $I_n(u)=\mathbb E_{C(n)}(\langle du,F\rangle-\partial_u\ell)^2$; the identity $I_n(u)=(2-x+j)\,\|du\|^2$ is therefore the first identity in the conclusion of Lemma~5.5.

\emph{Step 2a} (the identities of (a) that involve $K_n$; every real $M$).
Write $D:=D^{C(n)}_A$, $X_u:=\langle du,F\rangle$, $Z:=Z_{C(n)}[1]$ and $\mathbb E:=\mathbb E_{C(n)}$.
\begin{itemize}
\item[(i)] (The null set.)
The $A$ for which $\det D=0$ form a closed set of Lebesgue measure zero \cite[Theorem~4.25\,(c)]{QEDII}, and there are countably many closed $z$ \cite[Definition~2.6]{QEDII}.
Therefore $\{\varrho=0\}$ is a null set for $\mathrm d(A,z)$, and for each closed $z$ the $A$ with $\varrho(A,z)>0$ form an open set.
On that open set $\ell_n$ depends smoothly on $A$, and
\begin{equation*}
s_u=-X_u+\partial_u\ell_n,\qquad\partial_us_{u'}=-\langle du,du'\rangle+\partial_u\partial_{u'}\ell_n.
\end{equation*}
\item[(ii)] (Integrability.)
At the points where $\det D>0$,
\begin{equation*}
\det D\cdot\partial_u\partial_{u'}\ell_n=\partial_u\partial_{u'}\det D-\frac{(\partial_u\det D)(\partial_{u'}\det D)}{\det D},
\end{equation*}
and, because $\det D=|\det\check D|^2$ \cite[Proposition~3.18\,(d)]{QEDII},
\begin{equation*}
|\partial_u\det D|=2\,\bigl|\operatorname{Re}\bigl(\overline{\det\check D}\;\partial_u\det\check D\bigr)\bigr|\le2\,(\det D)^{1/2}\,|\partial_u\det\check D|.
\end{equation*}
It follows that
\begin{gather*}
|\det D\cdot\partial_u\partial_{u'}\ell_n|\le\sup|\partial_u\partial_{u'}\det D|+4\sup|\partial_u\det\check D|\,\sup|\partial_{u'}\det\check D|,\\
\det D\cdot(\partial_u\ell_n)^2\le4\sup|\partial_u\det\check D|^2,
\end{gather*}
and these suprema are finite, since $\det D$ and $\det\check D$ are trigonometric polynomials in the variables $qA(b)$ (Step~1).
The function $\exp(-\frac12\|F\|^2)$ being integrable \cite[Lemma~3.14]{QEDII}, both $\partial_u\partial_{u'}\ell_n$ and $(\partial_u\ell_n)^2$ are integrable for the probability measure of density $\varrho/Z$ \cite[Definition~2.13]{QEDII}.
Thus $K_n$ is well defined, bilinear and symmetric.
\item[(iii)] (The integral of $\partial_u\partial_{u'}\varrho$ is zero.)
Each of the functions $\partial_{u'}\varrho$ and $\partial_u\partial_{u'}\varrho$ is $\exp(-\frac12\|F\|^2)$ times a polynomial in $F$ of degree at most two, whose coefficients are trigonometric polynomials in the $qA(b)$.
By Step~1 of the proof of Proposition~5.3 and \cite[Definition~2.6]{QEDII}, the integral $\int\mathrm d(A,z)\,(\partial_{u'}\varrho)(A+tu,z)$ does not depend on $t$.
When $|t|\le1$, the derivative in $t$ of the integrand, namely $(\partial_u\partial_{u'}\varrho)(A+tu,z)$, is bounded in modulus by $C\prod_\sigma\exp(-\frac14F_\sigma^2)$ for some constant $C$, and this function is integrable, as in Step~2 of that proof \cite[Lemma~3.14]{QEDII}.
Differentiating at $t=0$, one obtains
\begin{equation*}
\int\mathrm d(A,z)\,\partial_u\partial_{u'}\varrho=0.
\end{equation*}
\item[(iv)] (The identity for $K_n(u,u')$.)
By Step~1, $\partial_{u'}\varrho=\varrho\,s_{u'}$.
On $\{\varrho>0\}$, part (i) then gives
\begin{equation*}
\partial_u\partial_{u'}\varrho=\varrho\,s_us_{u'}-\varrho\,\langle du,du'\rangle+\varrho\,\partial_u\partial_{u'}\ell_n.
\end{equation*}
These four terms are integrable, by Step~1 combined with the Schwarz inequality, by (ii) and by (iii).
One integrates over $\mathrm{Conf}(\Gamma)$, applies (iii) and divides by $Z>0$ \cite[Theorem~4.26]{QEDII}; the result is $\mathbb E[s_us_{u'}]=\langle du,du'\rangle+K_n(u,u')$.
\item[(v)] (The expression for $K_n(u,u)$.)
Step~2 gives $\mathbb E[X_us_u]=-\|du\|^2$, which is the same as $\mathbb E[X_u\,\partial_u\ell_n]=\mathbb EX_u^2-\|du\|^2$.
On expanding $s_u^2=(X_u-\partial_u\ell_n)^2$ one finds, for arbitrary weights,
\begin{equation*}
I_n(u)=2\|du\|^2-\mathbb EX_u^2+\mathbb E(\partial_u\ell_n)^2.
\end{equation*}
Combined with (iv), this is the expression for $K_n(u,u)$ stated in (a); the inequality $\mathbb EX_u^2\le\|du\|^2$ \cite[Theorem~6.27\,(b)]{QEDII} then yields $K_n(u,u)\ge\mathbb E(\partial_u\ell_n)^2$ and, in particular, $K_n(u,u)\ge0$.
\item[(vi)] (Two special cases; the range of $\rho$.)
If $\rho\ge0$ and $\mathbf m>0$, Lemma~5.1 gives $\varrho>0$ everywhere, and $s_u=-\partial_uV$; hence $K_n$ is the form of Lemma~5.2, and (iv) with $u'=u$ is the second identity of Step~3 of the proof of Proposition~5.3.
For standard weights and $du\ne0$, (v) takes the form $K_n(u,u)=(1-x+j)\|du\|^2$, and $0\le1-x\le\min\{1,\sqrt j\}$, because Lemma~5.5 gives $0<x\le1$ and $x\ge1-\sqrt j$.
Parts (i)--(iv) and the identities of (v) are valid for every real $\rho$, the zero set of $\det D$ in (i) being described for $\rho^2\ne1$ by \cite[Theorem~4.13\,(e)]{QEDII}; the assumption $-1\le\rho\le1$ enters the two inequalities of (v).
\end{itemize}

\emph{Step 3} (lattice probes).
Let $X_i:=\langle d\hat{\mathfrak a}_i,F\rangle$ and $\boldsymbol\sigma_i:=\|d\hat{\mathfrak a}_i\|$.
By \cite[Proposition~47.1\,(b2$'$), (b3), (b5)]{QEDIII}, $X_i=X_{\eta_i,R_i}(f^{(\eta_i)})$ with $\|f^{(\eta)}-f\|\le\eta^2\|\beta\|_{H^4}$, and $\boldsymbol\sigma_i\to\|f\|$.
The map $\beta\mapsto\hat{\mathfrak a}_i$ is linear, so that \cite[Proposition~47.1\,(b5)]{QEDIII} applied to $\beta\pm\beta'$ gives $\langle d\hat{\mathfrak a}_i,\allowbreak d\hat{\mathfrak a}'_i\rangle\to\langle f,\allowbreak f'\rangle$ for any two probes $\beta,\beta'$, the primes referring to $\beta'$.
For all $e>0$ and all mass assignments, \cite[Chapter~44, Lemma~2.1\,(a)--(c)]{QEDIII} gives
\begin{equation*}
E_i\bigl(X_i-X_{\eta_i,R_i}(f)\bigr)^2\le\|f^{(\eta_i)}-f\|^2\to0
\end{equation*}
as soon as $(-\frac{R_i}2,\frac{R_i}2)^3$ contains $\operatorname{supp}\beta+[-1,1]^3$.

\emph{Step 4} (convergence).
Under $E_i$, the family $X_{\eta_i,R_i}(g)$, $g\in\mathcal S$, has the same joint laws as the family $\omega(g)$ under $P_{N_i,T_i}$.
These laws converge weakly, their generating functions converge uniformly on compact sets, and all their moments converge \cite[Theorem~1.3\,(a)]{QEDII}, \cite[Chapter~44, Lemma~2.2]{QEDIII}.
Combined with Step~3, this yields $E_iX_i^2\to S_2(f,f)$.
Moreover, let $\beta_1,\dots,\beta_k$ be probes and $g_j=dd^*\beta_j$; then the variables $X_i^{(j)}:=\langle d\hat{\mathfrak a}_i^{(j)},F\rangle$ converge jointly in law to the variables $X(g_j)$ under $\mu$.

\emph{Step 5} (proof of (b)).
For $u=\hat{\mathfrak a}_i$, part (a) reads
\begin{equation*}
\boldsymbol\sigma_i^4\le E_iX_i^2\cdot I_{n_i}(\hat{\mathfrak a}_i),\qquad I_{n_i}(\hat{\mathfrak a}_i)\ge\boldsymbol\sigma_i^2.
\end{equation*}
Letting $i\to\infty$ along a subsequence that realises $\bar I(\beta)$ gives (b).

\emph{Step 6} (proof of (c) for test forms of probes).
Let $g_j=dd^*\beta_j$ and $\Phi_i:=\phi(X_i^{(1)},\dots,X_i^{(k)})$.
Then
\begin{equation*}
\partial_{\hat{\mathfrak a}_i}\Phi_i=\sum_j\langle d\hat{\mathfrak a}_i^{(j)},d\hat{\mathfrak a}_i\rangle\,(\partial_j\phi)(X_i^{(1)},\dots,X_i^{(k)}),
\end{equation*}
and Step~2 together with the Schwarz inequality yields
\begin{equation*}
|E_i[\partial_{\hat{\mathfrak a}_i}\Phi_i]|\le(E_i\Phi_i^2)^{1/2}\,I_{n_i}(\hat{\mathfrak a}_i)^{1/2}.
\end{equation*}
Along the subsequence chosen in Step~5, Steps~3 and~4 give in the limit
\begin{equation*}
|\mu(\partial_f\Phi)|\le\bar I(\beta)^{1/2}\,\|\Phi\|_2.
\end{equation*}
Hence the linear functional $\Phi\mapsto\mu(\partial_f\Phi)$ is well defined on classes in $L^2(\mu)$ and bounded there, and the Riesz representation theorem furnishes $\mathsf s_\beta$.

\emph{Step 7} (proof of (c) for all test forms).
Every closed form in $L^2$ is a limit in $L^2$ of forms $dd^*\beta$.
Indeed, the dual fields of the forms $dd^*\beta$ have the Fourier transforms $|p|^2P^\perp(p)\,\widehat{\ast\beta}(p)$, and a field $w=P^\perp w\in L^2$ orthogonal to all of them satisfies $|p|^2w=0$.
Write $\Pi_{\mathrm c}$ for the projection of $L^2$ onto the subspace of closed forms.
Every real $g\in L^2$ satisfies $X(g)=X(\Pi_{\mathrm c}g)$ almost surely and $\|X(g)\|_2\le\|g\|$ \cite[Remark~14.4, Lemma~14.3\,(a)]{QEDII}, as well as $\langle g,f\rangle=\langle\Pi_{\mathrm c}g,f\rangle$.
The function $\phi$ and its partial derivatives $\partial_j\phi$ are bounded and Lipschitz continuous; hence each side of the identity in (c) is a continuous function of $(g_1,\dots,g_k)\in(L^2)^k$, and the identity extends from the case of Step~6 to the general case.

\emph{Step 8} (the shift).
For $g_j\in\mathcal S$ put
\begin{equation*}
\Phi_t(\omega):=\Phi(\omega+tf)=\phi\bigl(\omega(g_1)+t\langle g_1,f\rangle,\dots\bigr).
\end{equation*}
Then $\frac{\mathrm d}{\mathrm dt}\mu(\Phi_t)=\allowbreak \mu(\partial_f\Phi_t)=\allowbreak \mu(\Phi_t\,\allowbreak \mathsf s_\beta)$, whose modulus is at most $\sup|\phi|\cdot\|\mathsf s_\beta\|_2$.
Integrating in $t$ one obtains the bound for these $\Phi$; this is enough, because the maps $\omega\mapsto\omega(g)$, $g\in\mathcal S$, generate $\Sigma$.

\emph{Step 9} (proof of (d)).
A volume-first limit point is, in particular, a limit point \cite[Theorem~1.3\,(a)]{QEDII}.
Part (b) therefore gives $S_2(f,f)\ge(1+c)^{-1}\|f\|^2$ for every probe, and Step~7, combined with the continuity of $S_2$ on $L^2$, extends this to $S_2(g,g)\ge(1+c)^{-1}\|g\|^2$ for every real closed $g\in L^2$.
Fix such a $g\ne0$, and set $b=\ast g$ and $g_l(x):=g(x/l)$.
The measure $\mathrm{kl}_\mu$ is denoted $\mathfrak K_\mu$ in \cite[Theorem~1.4, Part~I\,(d)]{QEDII}, and that theorem gives
\begin{equation*}
\frac{S_2(g_l,g_l)}{\|g_l\|^2}=\frac{\int|\hat b(k)|^2\,\mathsf D(|k|^2/l^2)\,\mathrm dk}{\int|\hat b(k)|^2\,\mathrm dk},\qquad\mathsf D(r):=\int\mathrm d\,\mathrm{kl}_\mu(t)\,\frac{r}{r+t^2},
\end{equation*}
and $\mathsf D(r)$ decreases to $\mathrm{kl}_\mu(\{0\})$ as $r\downarrow0$.
As $l\to\infty$, dominated convergence shows that the left side converges to $\mathrm{kl}_\mu(\{0\})$.
\end{proof}
Given a charge $e>0$ and a mass assignment, $P_{N,T}$ stands in this section for the law $P^{\mathrm{ch}}_{N,T}$ of \cite[Section~1.2]{QEDII} formed with them, and, for a volume-first limit point $\mu$ of these laws, $\mathrm{kl}$ stands for the measure $\mathrm{kl}_\mu$ of Proposition~5.6\,(d), which is the measure $\mathfrak K_\mu$ of \cite[Theorem~1.4, Part~I\,(d)]{QEDII}.
If $0<e\le e_2$ and the mass assignment is $N\mapsto m_{0,N}(e)$, these laws and this measure coincide with those of the notation fixed above, by \cite[Section~I.7.5]{QEDIII} and by the uniqueness of the measure in \cite[Theorem~44\,(C7)]{QEDIII}.
For the remaining fixed-spacing statements of this section, $n\ge4$ is an integer, $\Gamma=n\mathbb Z^3$ and $\mathbb X_n:=\mathbb X/\Gamma$; moreover $D:=D^{C(n)}_A$ and $E_n:=\mathbb E_{C(n)}$, and $P_n$ is the probability that belongs to $E_n$.
Put
\begin{equation*}
\mathsf M_n(x,y):=E_n\|\check D_A^{-1}(x,y)\|^2_{\mathrm{HS}},\qquad\mathsf N(x,y):=E_n\|D^{-1}(x,y)\|^2_{\mathrm{HS}};
\end{equation*}
the quantity $\mathsf M_n$ is the expectation of Proposition~5.4\,(c).
Let $\mathsf a:=2\mathbf m/\mathsf h_3$ and $\varepsilon_*:=\frac12\log(1+\mathsf a)$.
For real $a$ write $|a|_n:=\operatorname{dist}(a,n\mathbb Z)$; for a class $s$ of $(\frac12\mathbb Z)^3$ modulo $\Gamma$ write $|s|^2:=\sum_j|s_j|_n^2$ and $|s|_\infty:=\max_j|s_j|_n$.
For a link $b$, let $V(b)$ be the set of the $16$ end points of the eight bonds over $b$.
The \emph{cell} of $k\in\mathbb Z^3/\Gamma$ is the set $\pi^{-1}(k)$, which has $48$ sites.
Further,
\begin{equation*}
L_t:=\{y\in\mathbb X_n:\pi(y)_1\equiv t\},\qquad v(x,y):=\pi(y)-\pi(x),\qquad k(b,b'):=K_n(\delta_b,\delta_{b'});
\end{equation*}
$s(b,b')$ denotes the midpoint of $b'$ minus the midpoint of $b$, and $\mathrm{sv}_j(\cdot)$ are the singular values of an operator.

Define
\begin{gather*}
\mathsf c_n:=\max_b\sum_{b'}|s(b,b')|^2\,|k(b,b')|,\qquad G_n(t):=\sum_{\pi(x)=0}\ \sum_{y\in L_t}\mathsf M_n(x,y),\qquad\mathsf S_n:=\sum_{t\in\mathbb Z/n}|t|_n^2\,G_n(t),\\
N_A(\lambda):=\#\{j:\mathrm{sv}_j(D_A)^2\le\lambda\},\qquad\mathsf T_n:=n^{-3}E_n\operatorname{Tr}\bigl[(D^*D)^{-2}\bigr].
\end{gather*}
Accordingly, $G_n(t)$ is the meson correlator summed over the sites of the cell $0$ and over the time slice $L_t$.
With $\mathsf k_2:=\max_{\mu,\nu}\sum_s|s|^2|k_{\mu\nu}(s)|$ one has $\mathsf k_2\le\mathsf c_n\le3\mathsf k_2$.
Condition \textbf{(FI)} asserts that there is a constant $C$, independent of $N$, $T$ and $u$, such that
\begin{equation*}
I_n(u)\le\bigl(1+C\,e^2/\bar m\bigr)\|du\|^2\qquad\text{for all $u$ with $\sum_xu(x)=0$.}
\end{equation*}

The following classical facts enter the proofs.
\begin{itemize}
\item[(K1)] (Min-max principle.) Let $h$ be a Hermitian form on an inner product space of finite dimension, and let $h(\psi,\psi)\le\mu\|\psi\|^2$ for all $\psi$ in a subspace of dimension $d$.
Then $h$ has at least $d$ eigenvalues $\le\mu$.
\item[(K2)] (H\"older's inequality for Schatten norms.) For complex square matrices,
\begin{equation*}
|\operatorname{tr}(XBYB')|\le\|X\|_{\mathrm{HS}}\|B\|\,\|Y\|_{\mathrm{HS}}\|B'\|,\qquad\|XYZ\|^2_{\mathrm{HS}}\le\|Y\|^2\bigl(\operatorname{Tr}[(X^*X)^2]\,\operatorname{Tr}[(Z^*Z)^2]\bigr)^{1/2}.
\end{equation*}
\item[(K3)] Every bounded invertible operator $T$ on a Hilbert space admits a polar decomposition $T=\Omega|T|$ with $\Omega$ unitary; moreover $|T^*|=\Omega|T|\Omega^*$.
\item[(K4)] On $\mathbb Z^3/n\mathbb Z^3$ the operator $\Delta:=dd^*+d^*d$ is, componentwise, the scalar Laplacian; it commutes with $d$ and with $d^*$, and its kernel consists of the constant cochains.
On a box $\prod_\mu[a_\mu,b_\mu]\cap\mathbb Z^3$ every closed one-cochain or two-cochain is exact, and an integer valued one has an integer valued primitive.
\item[(K5)] Every sequence of positive measures of bounded mass on a compact interval has a weakly convergent subsequence (Helly), and each such measure is determined by its moments (Weierstrass).
\item[(K6)] On a separable metric space, weak convergence of Borel probability measures is metrisable (L\'evy, Prokhorov).
\end{itemize}
\subsection{The Hessian, the Ward identity and the meson moment}
\begin{named}{Theorem 5.7}
\emph{Bounds on the Fisher information; an atom at zero when the fermion is heavy.}
Let $0<\rho<1$, let the weights be standard, and let $q>0$, $\mathbf m>0$ and $n\ge4$.
For every $u\in C^1(\Gamma)$ with $\sum_xu(x)=0$,
\begin{equation*}
I_n(u)\le\Bigl(1+\frac{\pi^2}{4}\,\mathsf c_n\Bigr)\|du\|^2,
\end{equation*}
and
\begin{equation*}
\mathsf c_n\le24\,\mathsf h_3^2q^2\,\mathsf S_n\le3\pi^2\mathsf h_3^4q^2\,\mathsf T_n.
\end{equation*}
There exist constants $C_A,C_B,C_*<\infty$, depending on $\rho$ alone, with the following properties.
\begin{itemize}
\item[(i)] $\mathsf c_n\le576\,\pi^2\mathsf h_3^4\,q^2\,\mathbf m^{-4}$; this holds also for every $\rho\ge0$ and for arbitrary weights $w_i>0$.
\item[(ii)] $\mathsf c_n\le C_A\,\allowbreak q^2\bigl(\mathbf m^{-5/2}+n^{-3}\mathbf m^{-4}\bigr)$.
\item[(iii)] For every $\mathbf m'\in[\mathbf m,1]$ such that $n\ge20/\mathbf m'$ and $\mathbf m'\ge C_*q^2\log^2\frac2{\mathbf m}$,
\begin{equation*}
\mathsf c_n\le C_B\,\frac{q^2}{\mathbf m}\Bigl(\frac{\mathbf m'}{\mathbf m}\Bigr)^3.
\end{equation*}
\end{itemize}
The bounds (i) and (ii) rest on estimates valid in every real field, and (iii) on Gaussian domination.
For the continuum scaling, fix an arbitrary charge $e>0$ and a mass assignment with $m_{0,N}\ge0$; then $\mathbf m=\boldsymbol\mu_N\ge\bar m\eta$.
\begin{itemize}
\item[(a)] Let $m_{0,N}=0$.
For all $T$ and all $N$ with
\begin{equation*}
\max\Bigl\{\tfrac52L\bar m\eta,\ C_*e^2\eta\log^2\frac{2}{\bar m\eta}\Bigr\}\le1
\end{equation*}
one has
\begin{equation*}
\mathsf c_n\le C_B\,\frac{e^2}{\bar m}\,\max\Bigl\{\tfrac52L,\ C_*\frac{e^2}{\bar m}\log^2\frac{2}{\bar m\eta}\Bigr\}^3.
\end{equation*}
Consequently $\mathsf k_2\le C_\kappa\,\allowbreak \frac{q^2}{\bar m\eta}\,\allowbreak (\bar m\eta)^{-\kappa}$ for every $\kappa>0$; the bound (i) corresponds to $\kappa=3$ and the bound (ii) to $\kappa=\frac32$.
The bound on $\mathsf k_2$ holds for all $T$ and the values of $N$ just specified, with
\begin{equation*}
C_\kappa:=C_B\max\Bigl\{\bigl(\tfrac52L\bigr)^3,\ C_*^3\Bigl(\frac{e^2}{\bar m}\Bigr)^3\,2^\kappa\Bigl(\frac{6}{\kappa\exp 1}\Bigr)^6\Bigr\},
\end{equation*}
a constant that depends only on $(\kappa,\rho,L,e^2/\bar m)$.
Let $\mathsf x_0=\mathsf x_0(\rho,L,e^2/\bar m)\in(0,1]$ be the largest number such that
\begin{equation*}
\max\Bigl\{\tfrac52L\mathsf x,\ C_*\frac{e^2}{\bar m}\,\mathsf x\log^2\frac{2}{\mathsf x}\Bigr\}\le1\qquad\text{for all }\mathsf x\in(0,\mathsf x_0].
\end{equation*}
The other values of $N$ are finitely many and satisfy $\bar m\eta>\mathsf x_0$; for them (i) gives
\begin{equation*}
\mathsf c_n\le576\,\pi^2\mathsf h_3^4\,\mathsf x_0^{-3}\,\frac{q^2}{\bar m\eta}.
\end{equation*}
Hence the bound on $\mathsf k_2$ holds for all $N$ and $T$ once $576\,\pi^2\mathsf h_3^4\,\mathsf x_0^{-3}$ is added to $C_\kappa$.
\item[(b)] Let $\varkappa>0$, and let $m_{0,N}\ge\varkappa\,e^2\eta\log^{3/2}\frac2\eta$ for all large $N$.
Then
\begin{equation*}
\limsup_{N\to\infty}\,\sup_T\,\mathsf c_n\le c_\varkappa:=C_BC_*^3\varkappa^{-4},
\end{equation*}
and $\mathrm{kl}(\{0\})\ge(1+\frac{\pi^2}{4}c_\varkappa)^{-1}>0$ for every volume-first limit point.
\item[(c)] Suppose that $m_{0,N}/(\eta\log^{3/2}\frac2\eta)\to\infty$.
Then $P_{N,T}$ converges to the free Maxwell law $\gamma$ as $N,T\to\infty$ in an arbitrary manner, and $\mathrm{kl}_\gamma=\delta_0$.
\end{itemize}
\end{named}

In (b) and (c) the bare mass of the fermion in physical units, $\boldsymbol\mu_N/\eta$, tends to infinity, so that the fermion is not retained in the limit.
The proof of Theorem~5.7 is given after Proposition~5.11.

\begin{named}{Lemma 5.8}
\emph{From the kernel to the form.}
Let $\rho\ge0$, $w_0,\dots,w_3>0$, $q>0$, $\mathbf m>0$ and $n\ge4$.
The matrix $k$ is real and symmetric.
The entry $k((x,\mu),(y,\nu))=k_{\mu\nu}(s)$ depends only on $\mu$, $\nu$ and the class $s=y-x+\frac12(e_\nu-e_\mu)$, and
\begin{equation*}
k_{\mu\nu}(-s)=k_{\mu\nu}(s)=k_{\nu\mu}(s).
\end{equation*}
For every link $b$ and every $\lambda\in C^0(\Gamma)$ one has $\sum_{b'}k(b,b')\,(d\lambda)(b')=0$.
Finally, for all $u\in C^1(\Gamma)$,
\begin{equation*}
|K_n(u,u)|\le\mathsf c_n\Bigl(\frac{\pi^2}{4}\|du\|^2+2\pi^2n^{-5}\Bigl|\sum_xu(x)\Bigr|^2\Bigr).
\end{equation*}
Combined with Theorem~5.7\,(i), this yields Lemma~5.2 with
\begin{equation*}
c_*=1152\,\pi^4\mathsf h_3^4\,\mathbf m^{-4};
\end{equation*}
the assertions of Proposition~5.3 are therefore valid with this value of $c_*$.
\end{named}

\begin{proof}
\emph{Step 1} (the kernel).
Lemma~5.1 gives $\|D^{-1}\|\le\mathbf m^{-1}$, and $\det D=|\det\check D|^2>0$ by \cite[Proposition~3.18\,(d)]{QEDII}; hence $\ell_n$ is smooth.
A phase is carried only by the bonds of type~$3$; each of them lies over one link, and every site is an end point of exactly one of them \cite[Proposition~3.6\,(c), (d)]{QEDII}.
It follows that $\partial_b\partial_{b'}D=\delta_{bb'}D_{bb}$, with $D_b:=\partial_bD$ and $D_{bb}:=\partial_b^2D$ with derivatives taken along the indicated links, and that the sets $V(b)$ are pairwise disjoint and cover $\mathbb X_n$.
Moreover the non-zero blocks of $D_b$ are the sixteen blocks
\begin{equation*}
D_b(x_1,x_2)=\pm iq\,D(x_1,x_2),\qquad x_1\in V(b),\quad x_2=x_1^{(3)},
\end{equation*}
each of norm $q\mathsf h_3$.
Jacobi's formula gives
\begin{equation*}
-\partial_b\partial_{b'}\ell_n=\operatorname{Tr}\bigl(D^{-1}D_bD^{-1}D_{b'}\bigr)-\delta_{bb'}\operatorname{Tr}\bigl(D^{-1}D_{bb}\bigr);
\end{equation*}
this is a bounded function of the phases of the links.

\emph{Step 2} (symmetries).
Let $g\in W$.
Since $\vec g$ is a signed permutation matrix and $\chi(nv)=(-1)^{v_1+v_2+v_3}$, one has $gC(n)=C(n)$.
Further $\ell_n((g^{-1})^*A)=\ell_n(A)$ \cite[Proposition~3.16\,(b)]{QEDII}, and $(A,z)\mapsto(g^*A,g^*z)$ leaves $\mathrm d(A,z)$ and $\|F\|$ unchanged \cite[Definition~2.6]{QEDII}.
The law on $\mathrm{Conf}(\Gamma)$ is therefore invariant under $W$, and $K_n(g^*u,g^*u')=K_n(u,u')$.
The identity $\ell_n(A-d\lambda)=\ell_n(A)$ \cite[Proposition~3.16\,(a)]{QEDII} gives $\partial_u\partial_{d\lambda}\ell_n=0$ in every field.

\emph{Step 3} (evenness).
The translations yield the form $k_{\mu\nu}(s)$ of the kernel.
For $g=-1$ one has $(g^*u)_\mu(x)=-u_\mu(-x-e_\mu)$ \cite[Definition~2.1]{QEDII}; the two signs cancel and the midpoints of the links are carried to their negatives, whence $k_{\mu\nu}(-s)=k_{\mu\nu}(s)$.
In terms of $k_{\mu\nu}$, the symmetry of $k$ reads $k_{\mu\nu}(s)=k_{\nu\mu}(-s)$.

\emph{Step 4} (Parseval's identity with midpoints).
Let $\mathcal B:=\frac{2\pi}{n}\mathbb Z^3\cap(-\pi,\pi]^3$, and for $p\in\mathcal B$ let
\begin{equation*}
\hat u_\mu(p):=\sum_xu_\mu(x)\exp\bigl(-ip\cdot(x+\tfrac12e_\mu)\bigr),\qquad\mathsf s_\mu:=2\sin\frac{p_\mu}{2},\qquad\hat k_{\mu\nu}(p):=\sum_sk_{\mu\nu}(s)\cos(p\cdot s);
\end{equation*}
the cosine is a function of the class $s$, because $\exp(ip\cdot ne_j)=1$.
By Fourier inversion on $(\mathbb Z/n)^3$ and by Step~3,
\begin{equation*}
K_n(u,u)=n^{-3}\sum_{p\in\mathcal B}\hat u^*\hat k\hat u,
\end{equation*}
where each $\hat k(p)$ is a real symmetric matrix.
If Fourier transforms of two-cochains are taken with respect to the centres of the plaquettes, then $\widehat{du}_{\mu\nu}=\allowbreak i(\mathsf s_\mu\hat u_\nu-\mathsf s_\nu\hat u_\mu)$; moreover $\widehat{d\lambda}_\mu=i\mathsf s_\mu\hat\lambda$ \cite[Definition~2.1]{QEDII}.
Lagrange's identity therefore gives $\|du\|^2=n^{-3}\sum_p|\mathsf s|^2|\hat u^\perp|^2$, where $\hat u^\perp$ denotes the component of $\hat u$ orthogonal to $\mathsf s$.

\emph{Step 5} (transversality; the mode $p=0$).
In Fourier variables, Step~2 reads $\hat k(p)\mathsf s(p)=0$.
Put $m_{\mu\nu}:=\sum_s|s|^2|k_{\mu\nu}(s)|$.
Choosing for each class $s$ a shortest representative, one obtains
\begin{equation*}
|\hat k_{\mu\nu}(p)-\hat k_{\mu\nu}(0)|\le\tfrac12|p|^2m_{\mu\nu};
\end{equation*}
by the evenness of the kernel there is no term linear in $p$.
At $p=\frac{2\pi}{n}e_\nu$ transversality gives $\hat k_{\mu\nu}(p)=0$.
Consequently $|\hat k_{\mu\nu}(0)|\le2\pi^2n^{-2}m_{\mu\nu}$, and $|\hat k_{\mu\nu}(p)|\le|p|^2m_{\mu\nu}$ for every $p\in\mathcal B\setminus\{0\}$.

\emph{Step 6} (conclusion).
The norm of a real symmetric matrix is bounded by the largest of the sums of the absolute values of the entries in a row, and $\max_\mu\sum_\nu m_{\mu\nu}=\mathsf c_n$.
For $p\ne0$ only $\hat u^\perp$ contributes, and $|p|\le\frac\pi2|\mathsf s(p)|$ on $\mathcal B$; finally $\hat u(0)=\sum_xu(x)$.
\end{proof}

\begin{named}{Lemma 5.9}
\emph{From the meson to the kernel.}
Let the hypotheses of Lemma~5.8 be satisfied.
\begin{itemize}
\item[(a)] For every $g\in W$,
\begin{equation*}
\mathsf N(x,y)=\mathsf M_n(x,y)+\mathsf M_n(y,x)=\mathsf N(y,x)=\mathsf N(gx,gy),
\end{equation*}
and $\mathsf M_n(gx,gy)=\mathsf M_n(x,y)$ for every even $g\in W$.
Consequently $G_n(-t)=G_n(t)$, the function $G_n$ is the same for each of the three coordinate directions,
\begin{equation*}
\sum_{\pi(x)=0}\ \sum_{y\in L_t}\mathsf M_n(y,x)=G_n(-t),
\end{equation*}
and $\sum_y|v(x,y)|^2\mathsf N(x,y)=\frac18\mathsf S_n$ for every site $x$.
In addition $\mathsf M_n(x,y)=\mathsf M_n(y,x)$ (charge conjugation), so that $\mathsf N=2\mathsf M_n$ and $\mathsf M_n$ is invariant under the whole of $W$.
\item[(b)] For all $b\ne b'$, that is, for all $s\ne0$,
\begin{equation*}
|k(b,b')|\le\mathsf h_3^2q^2\sum_{x'\in V(b)}\ \sum_{y'\in V(b')}\mathsf N(x',y')=2\mathsf h_3^2q^2\sum_{x'\in V(b)}\ \sum_{y'\in V(b')}\mathsf M_n(x',y'),
\end{equation*}
a sum over $256$ pairs.
The term $\operatorname{Tr}(D^{-1}D_{bb})$ occurs only at $b=b'$, where the weight $|s|^2$ is $0$; it does not contribute to $\mathsf c_n$.
For $b\ne b'$ the vertices $V(b)$ and $V(b')$ never overlap.
\item[(c)] (Ward identity.) One has $\mathsf c_n\le24\,\mathsf h_3^2q^2\,\mathsf S_n$.
Links that share an end point enter $\mathsf c_n$ through pairs of sites in a single cell; no $G_n(t)$ with $t\ne0$ sees these pairs, and such links are dominated by the other links.
\end{itemize}
\end{named}

\begin{proof}
\emph{Step 1} (proof of (a)).
By \cite[Proposition~3.18\,(d)]{QEDII}, $D^{-1}=\operatorname{diag}(\check D^{-1},(\check D^{-1})^*)$, so that in every field
\begin{equation*}
\|D^{-1}(x,y)\|^2_{\mathrm{HS}}=\|\check D^{-1}(x,y)\|^2_{\mathrm{HS}}+\|\check D^{-1}(y,x)\|^2_{\mathrm{HS}}.
\end{equation*}
According to \cite[Proposition~3.16\,(b)]{QEDII},
\begin{equation*}
D^{-1}_{(g^{-1})^*A}(gx,gy)=\tilde g\,D_A^{-1}(x,y)\,\tilde g^{-1},
\end{equation*}
where $\tilde g$ is a product of unitary matrices $\gamma_n\gamma_5$, each of which interchanges the two summands $\mathbb C^2$ \cite[Definition~2.7]{QEDII}; for even $g$ the matrix $\tilde g$ is block diagonal.
Combined with Step~2 of the proof of Lemma~5.8 and with $\pi(gx)=g\pi(x)$ \cite[Proposition~3.6\,(c)]{QEDII}, this proves the invariances.
For the consequences, apply them to $\operatorname{diag}(-1,-1,1)$, to the cyclic permutations and to the translations; the group $W$ acts transitively on $\mathbb X$, and summing $\sum_y|v|^2\mathsf N$ over the sites of a cell gives $6\mathsf S_n$.
For the last assertion of (a), the identity of \cite[Proposition~3.16\,(c)]{QEDII} is block diagonal with respect to $\mathbb C^4=\mathbb C^2\oplus\mathbb C^2$, and on every torus its upper block is
\begin{equation*}
\sigma_2(\check D_A)^{\top}\sigma_2^{-1}=\check D_{-A};
\end{equation*}
hence $\|\check D_{-A}^{-1}(x,\allowbreak y)\|_{\mathrm{HS}}=\allowbreak \|\check D_A^{-1}(y,\allowbreak x)\|_{\mathrm{HS}}$.
The map $(A,z)\mapsto(-A,-z)$ is compatible with the equivalence relation and with the product of Lebesgue measure and counting measure, and therefore preserves $\mathrm d(A,z)$ \cite[Definition~2.6]{QEDII}; it preserves $\|F\|$ and $\det D_A$ too.

\emph{Step 2} (proof of (b)).
Step~1 of the proof of Lemma~5.8 gives
\begin{equation*}
\operatorname{Tr}\bigl(D^{-1}D_bD^{-1}D_{b'}\bigr)=\sum\operatorname{tr}\bigl[D^{-1}(y_2,x_1)\,D_b(x_1,x_2)\,D^{-1}(x_2,y_1)\,D_{b'}(y_1,y_2)\bigr],
\end{equation*}
the sum being taken over $x_1\in V(b)$ and $y_1\in V(b')$, with $x_2=x_1^{(3)}$ and $y_2=y_1^{(3)}$.
The assertion follows from (K2), the inequality $2\xi\zeta\le\xi^2+\zeta^2$, part (a), and the fact that $x_1\mapsto x_2$ permutes $V(b)$.

\emph{Step 3} (proof of (c): the bad links).
Fix a link $b$, parallel to $e_\mu$.
A link $b'\ne b$ is called \emph{bad} if it has an end point in common with $b$, and \emph{good} otherwise.
There are $10$ bad links, and $|s|^2\le1$ for each of them; for good links $|s|^2\ge1$, since the value $|s|^2=\frac12$ occurs only at corners.
For $\lambda=\delta_y$ the identity of Lemma~5.8 becomes
\begin{equation*}
\sum_\nu\bigl[k(b,(y-e_\nu,\nu))-k(b,(y,\nu))\bigr]=0.
\end{equation*}
Let $b'$ be bad, and let $p'$ be the end point of $b'$ that is not on $b$.
The vector from the end point of $b$ not shared with $b'$ to $p'$ is $\pm2e_\mu$ or $\pm e_{\mu'}\pm e_\mu$ with $\mu'\ne\mu$; since $n\ge4$, it is neither $0$ nor a unit vector modulo $n$.
Hence the five links at $p'$ other than $b'$ are good, and $|k(b,b')|$ does not exceed the sum of $|k(b,\cdot)|$ over these five links.
The map $b'\mapsto p'$ is injective, and every good link arises at most twice in this way.
Therefore
\begin{equation*}
\sum_{b'\ne b}|s|^2|k|\le3\sum_{\mathrm{good}}|s|^2|k|.
\end{equation*}

\emph{Step 4} (proof of (c): the good links).
For a good link $b'$, parallel to $e_\nu$, and sites $x'\in V(b)$, $y'\in V(b')$ one has $v(x',y')\ne0$ and $s=v(x',y')\pm\frac12e_\mu\pm\frac12e_\nu$, whence $|s|^2\le4|v(x',y')|^2$.
From (b), from the disjointness of the sets $V(b')$ and from (a) it follows that
\begin{equation*}
\sum_{\mathrm{good}}|s|^2|k(b,b')|\le4\mathsf h_3^2q^2\sum_{x'\in V(b)}\ \sum_{y'}|v(x',y')|^2\mathsf N(x',y')=8\mathsf h_3^2q^2\mathsf S_n.
\end{equation*}
\end{proof}

\begin{named}{Lemma 5.10}
Let $\rho\ge0$, $w_i>0$, $q>0$, $\mathbf m>0$ and $n\ge4$, and let $A$ be real and $\Gamma$-periodic.
\begin{itemize}
\item[(a)] One has
\begin{equation*}
\sum_{x,y\in\mathbb X_n}|v(x,y)|^2\,\|D^{-1}(x,y)\|^2_{\mathrm{HS}}\le\frac{3\pi^2\mathsf h_3^2}{4}\,\operatorname{Tr}\bigl[(D^*D)^{-2}\bigr],
\end{equation*}
and consequently $\mathsf S_n\le\frac{\pi^2\mathsf h_3^2}{8}\,\mathsf T_n$.
\item[(b)] If $0<\rho<1$ and the weights are standard, then
\begin{equation*}
\operatorname{Tr}\bigl[(D^*D)^{-2}\bigr]\le\operatorname{Tr}\bigl[(\mathbf m+\tfrac\rho2\nabla_A^*\nabla_A)^{-4}\bigr]\le C_\rho\bigl(\mathbf m^{-4}+n^3\mathbf m^{-5/2}\bigr).
\end{equation*}
\item[(c)] \emph{Energy inequality.}
Given a site $x$ and $\alpha\in\{1,2\}$, set
\begin{equation*}
\psi:=\check D_A^{-1}\delta_{x,\alpha},\qquad f_A(t):=\sum_{|v(x,y)_1|_n=t}|\psi(y)|^2,\qquad F_A(t):=\sum_{t'\ge t}f_A(t').
\end{equation*}
Then, for every integer $t\ge1$,
\begin{equation*}
\mathbf m\,F_A(t)\le\mathsf h_3\sqrt{f_A(t)\,f_A(t-1)}.
\end{equation*}
Let
\begin{equation*}
\mathsf f(t):=\begin{cases}G_n(t)+G_n(-t)&\text{for }0<t<\frac n2,\\G_n(t)&\text{for }t\in\{0,\frac n2\},\\0&\text{for }t>\lfloor n/2\rfloor,\end{cases}\qquad\mathsf F(t):=\sum_{t'\ge t}\mathsf f(t').
\end{equation*}
As a consequence,
\begin{equation}
\tag{E1}
\begin{aligned}
\mathsf a\,\mathsf F(t)&\le\mathsf f(t)+\mathsf f(t-1)=\mathsf F(t-1)-\mathsf F(t+1)\qquad(t\ge1),\\
\mathsf F(t+1)&\le\exp(-2\varepsilon_*)\,\mathsf F(t-1)\qquad(t\ge1),\\
\mathsf F(0)&\le96\,\mathbf m^{-2},
\end{aligned}
\end{equation}
and
\begin{equation}
\tag{E2}
\sum_{t\ge t_1}t^2\,\mathsf f(t)\le\mathsf F(t_1-1)\bigl[t_1^2+2(2t_1+1)\,\mathsf a^{-1}+8\,\mathsf a^{-2}\bigr]\qquad(t_1\ge1).
\end{equation}
The inequalities (E1) and (E2) carry no multiplicative constant and impose no further condition on $n$: the periodicity of the torus is incorporated in their statement.
\end{itemize}
\end{named}

\begin{proof}
\emph{Step 1} (proof of (a)).
Let $\phi_j:=\exp(2\pi i\,\pi(\cdot)_j/n)$.
By \cite[Proposition~3.6\,(c)]{QEDII} the commutator $[D,\phi_j]$ involves hops of type~$3$ only, each of its blocks has norm at most $2\sin\frac\pi n\,\mathsf h_3$, and $x\mapsto x^{(3)}$ is a bijection.
Hence
\begin{equation*}
\|[D,\phi_j]\|\le\frac{2\pi}{n}\,\mathsf h_3.
\end{equation*}
The blocks of $[\phi_j,D^{-1}]=D^{-1}[D,\phi_j]D^{-1}$ are the matrices $(\phi_j(x)-\phi_j(y))\,D^{-1}(x,y)$, and
\begin{equation*}
|\phi_j(x)-\phi_j(y)|\ge\frac4n\,|v(x,y)_j|_n.
\end{equation*}
The first inequality of (a) is obtained from (K2) on summing over $j$.
By Lemma~5.9\,(a), the expectation of its left side equals $6n^3\mathsf S_n$, which gives the bound on $\mathsf S_n$.

\emph{Step 2} (proof of (b)).
Put $R_A:=\frac\rho2\nabla_A^*\nabla_A$.
For $\psi$ in the span of the eigenvectors of $D^*D$ with eigenvalues $\le\lambda$, \cite[Proposition~3.18\,(b)]{QEDII} gives
\begin{equation*}
\langle\psi,(\mathbf m+R_A)\psi\rangle=\operatorname{Re}\langle\psi,D\psi\rangle\le\sqrt\lambda\,\|\psi\|^2.
\end{equation*}
According to (K1), at least $N_A(\lambda)$ eigenvalues of $\mathbf m+R_A$ are therefore $\le\sqrt\lambda$.
This implies
\begin{equation*}
\operatorname{Tr}\bigl[(D^*D)^{-2}\bigr]\le\operatorname{Tr}\bigl[(\mathbf m+R_A)^{-4}\bigr]=\frac16\int_0^\infty t^3\exp(-\mathbf mt)\,\operatorname{Tr}\exp(-tR_A)\,\mathrm dt.
\end{equation*}

\emph{Step 3} (diamagnetic comparison for (b)).
One has $R_A=\frac\rho2(\mathrm{sw}-H_A)$, and every block of $H_A$ is the product of $w_i$, a complex number of modulus one and $1_4$ \cite[Proposition~3.18\,(a), (b)]{QEDII}.
Expanding $\exp(\frac12t\rho H_A)$ as a power series yields
\begin{equation*}
\operatorname{Tr}\exp(-tR_A)\le4\operatorname{Tr}\exp\bigl(-\tfrac12t\rho L_0\bigr),
\end{equation*}
where $L_0$ is the Laplacian of the weighted graph $\mathbb X_n$, acting on periodic functions with scalar values.
This graph is connected and invariant under $\mathbb Z^3$ \cite[Proposition~3.6\,(d)]{QEDII}, so that $x$ and $x+e_\mu$ are joined, covariantly, by at most $\ell_0$ bonds.
Consequently
\begin{equation*}
\sum_{\mu,x}|\psi(x+e_\mu)-\psi(x)|^2\le24\,\ell_0^2\,\langle\psi,L_0\psi\rangle.
\end{equation*}
In view of (K1) and of $4\sin^2\frac{\pi j}{n}\ge16j^2/n^2$, it follows that $\operatorname{Tr}\exp(-sL_0)\le48\,(1+c\,ns^{-1/2})^3$.

\emph{Step 4} (proof of (c)).
Let $Y:=\{y:|v(x,y)_1|_n\ge t\}$ and $P:=1_Y$.
Since $x\notin Y$, one has $P\check D\psi=0$; moreover $\operatorname{Re}\check D_A\ge\mathbf m$ by \cite[Proposition~3.18\,(b), (d)]{QEDII}.
Therefore
\begin{equation*}
\mathbf m\|P\psi\|^2\le\operatorname{Re}\langle P\psi,\check DP\psi\rangle\le|\langle P\psi,\check D(1-P)\psi\rangle|.
\end{equation*}
By \cite[Proposition~3.6\,(c), (d)]{QEDII}, the only bonds leaving $Y$ are bonds of type~$3$; each of them connects distance $t$ with distance $t-1$, at most one of them meets a given site, and their blocks have norm $\mathsf h_3$.
The Schwarz inequality bounds the right side by $\mathsf h_3\sqrt{f_A(t)f_A(t-1)}$, and $\|P\psi\|^2=F_A(t)$.

\emph{Step 5} (proof of (E1) and (E2)).
Bounding the right side of the inequality of (c) by means of $2\sqrt{\xi\zeta}\le\xi+\zeta$, summing over $\alpha$ and over the $48$ sites of the cell $0$, and taking the expectation $E_n$, one obtains the first inequality of (E1) from Lemma~5.9\,(a).
The second inequality of (E1) follows from the first together with $\mathsf F(t)\ge\mathsf F(t+1)$.
The bound $\mathsf F(0)\le48\cdot2\,\mathbf m^{-2}$ is a consequence of Proposition~5.4\,(c).
Summing the first inequality of (E1) over $t\ge t_1$, first without weight and then with the weight $t$, one finds
\begin{equation*}
\mathsf a\sum_{t\ge t_1}\mathsf F(t)\le2\,\mathsf F(t_1-1),\qquad\mathsf a\sum_{t\ge t_1}t\,\mathsf F(t)\le(2t_1+1)\,\mathsf F(t_1-1)+2\sum_{t\ge t_1}\mathsf F(t).
\end{equation*}
Together with Abel's summation
\begin{equation*}
\sum_{t\ge t_1}t^2\,\mathsf f(t)=t_1^2\,\mathsf F(t_1)+\sum_{t>t_1}(2t-1)\,\mathsf F(t),
\end{equation*}
these two inequalities give (E2).
\end{proof}

\subsection{Localization and the first singular-value estimate}
\begin{named}{Proposition 5.11}
Let $0<\rho<1$, and let the weights be standard.
There exist constants $C_*,C_B'<\infty$, depending only on $\rho$, with the following property.
If $q>0$, $0<\mathbf m\le\mathbf m'\le1$, $n\ge20/\mathbf m'$ and $\mathbf m'\ge C_*q^2\log^2\frac2{\mathbf m}$, then
\begin{equation*}
n^{-3}E_nN_A(\lambda)\le C_B'\bigl(\sqrt\lambda+\mathbf m'\bigr)^3\quad(\lambda>0),\qquad\mathsf T_n\le20\,C_B'\,\mathbf m'^3\,\mathbf m^{-4}.
\end{equation*}
\end{named}

\begin{proof}
Every constant denoted by $c$ with a subscript depends on $\rho$ alone.
Let $\ell:=\lceil1/\mathbf m'\rceil$.
Then $\ell^{-1}\le\mathbf m'$, and $10\ell<n$ because $\ell<1/\mathbf m'+1\le2/\mathbf m'$.
Put $J:=\lceil\log_2(4/\mathbf m')\rceil$, and let $r\ge1$.

\emph{Step 1} (localisation).
Divide $\mathbb Z/n$ into arcs $[t_i,t_{i+1}]$ whose lengths lie in $[\ell,2\ell)$, and define
\begin{equation*}
\mathsf j_i(t):=\begin{cases}\sin\dfrac{\pi(t-t_{i-1})}{2(t_i-t_{i-1})}&\text{on }[t_{i-1},t_i],\\[2mm]\cos\dfrac{\pi(t-t_i)}{2(t_{i+1}-t_i)}&\text{on }[t_i,t_{i+1}],\\[1mm]0&\text{elsewhere,}\end{cases}\qquad\mathsf j_a(k):=\prod_\mu\mathsf j_{a_\mu}(k_\mu).
\end{equation*}
These functions satisfy $\sum_a\mathsf j_a^2=1$.
Each support $B_a:=\operatorname{supp}\mathsf j_a$ is a box with at most $4\ell$ sites per side, the number of these boxes does not exceed $(n/\ell)^3$, and for neighbouring $k,k'$
\begin{equation*}
\sum_a|\mathsf j_a(k)-\mathsf j_a(k')|^2\le\frac{\pi^2}{4\ell^2}.
\end{equation*}
The argument of Step~1 of the proof of Lemma~5.10 gives
\begin{equation*}
\sum_a\|D\,\mathsf j_a\psi\|^2\le2\|D\psi\|^2+c_J\,\ell^{-2}\|\psi\|^2,\qquad c_J=\frac{\pi^2}{2}\,\mathsf h_3^2.
\end{equation*}
Since the map $\psi\mapsto(\mathsf j_a\psi)_a$ is isometric, (K1) yields
\begin{equation*}
N_A(\lambda)\le\sum_aN_a(\lambda'),\qquad\lambda':=2\lambda+c_J\,\ell^{-2},
\end{equation*}
where $N_a(\lambda')$ counts the eigenvalues not exceeding $\lambda'$ of the quadratic form $\phi\mapsto\|D\phi\|^2$, restricted to spinors with support over $B_a$.

\emph{Step 2} (a gauge in the box, linear in $F$).
Fix $a$ and write $B=B_a$.
Let $B^+$ be the $1$-neighbourhood of $B$, and $B^{++}$ the $(\ell+1)$-neighbourhood of $B^+$.
These boxes have at most $6\ell+4\le10\ell<n$ sites per side.
Hence they, the cochains on them and the spinors over $B$ lift to $\mathbb Z^3$ and to $\mathbb X$, and the character $\chi$ does not enter.
As $dF=0$, (K4) yields
\begin{equation*}
F=dA^{\mathrm C}+F_{\mathrm h},\qquad A^{\mathrm C}:=d^*\mathcal GF,\qquad\mathcal G:=\int_0^\infty\exp(-t\Delta)\,P^\perp\,\mathrm dt,
\end{equation*}
where $P^\perp$ is the projection onto the orthogonal complement of the constants and $F_{\mathrm h}$ is constant.
On $B^{++}$ define
\begin{equation*}
A^{\mathrm h}_\nu(k):=\frac12\sum_\mu(F_{\mathrm h})_{\mu\nu}\,(k-k_B)_\mu,\qquad A':=A^{\mathrm C}+A^{\mathrm h},
\end{equation*}
with $k_B$ the centre of the box.
Then $d(A-A')=-\frac{2\pi}{q}z$, and (K4) furnishes an integer cochain $\mathfrak k$ such that $A-A'+\frac{2\pi}{q}\mathfrak k=d\lambda$.
Put
\begin{equation*}
\zeta(k):=\max\{0,1-\operatorname{dist}_\infty(k,B^+)/\ell\},\qquad A''_\mu:=\zeta A'_\mu.
\end{equation*}
For $\phi$ over $B$, \cite[Proposition~3.16\,(a)]{QEDII} gives $\|D_A\phi\|=\|D_{A''}(\exp(iq\,\lambda\circ\pi)\,\phi)\|$.

\emph{Step 3} (a test against the polar part of the free operator).
Denote by $D_0'$ the free operator with the mass $\mathbf m'$ on $\ell^2(\mathbb X;\mathbb C^4)$, and write $D_0'=\Omega|D_0'|$ with $\Omega$ unitary, according to (K3).
Let $\Lambda:=(-\Delta_{\mathbb Z^3}+\mathbf m'^2)^{1/2}$ act in the variable $k=\pi(x)$, for fixed $\vec g$.
On functions of the form $\psi(k+\vec gv_0)=\exp(ip\cdot k)\,u(\vec g)$ the operator $D_0'$ acts as the symbol $D^{(p)}$, as a comparison of \cite[Proposition~3.18\,(a)]{QEDII} with \cite[Proposition~37.1\,(a)]{QEDIII} shows.
By \cite[Proposition~37.1\,(f)]{QEDIII},
\begin{equation*}
|(D^{(p)})^{-1}|_2\le c_E^{-1}\,(|p|_2+\mathbf m')^{-1}\quad\text{on }[-\pi,\pi]^3,\qquad c_E:=(\sqrt2+7C_E)^{-1}.
\end{equation*}
The symbol of $\Lambda$ is a scalar not larger than $|p|_2+\mathbf m'$; therefore $|D_0'|\ge c_E\Lambda$ and $|D_0'^*|\ge c_E\Lambda$.
With $V:=D_{A''}-D_0$, suppose that for all $\phi_1,\phi_2$
\begin{equation}
\tag{5.1}\label{eq:polar-form-bound}
|\langle\phi_1,V\phi_2\rangle|\le\frac{c_E}{2}\,\|\Lambda^{1/2}\phi_1\|\,\|\Lambda^{1/2}\phi_2\|.
\end{equation}
One has $D_{A''}=D_0'+V-(\mathbf m'-\mathbf m)$ and $c_E\|\Lambda^{1/2}\Omega\phi\|^2\le\langle\phi,|D_0'|\phi\rangle$.
It follows that
\begin{equation*}
\|D_{A''}\phi\|\,\|\phi\|\ge\operatorname{Re}\langle\Omega\phi,D_{A''}\phi\rangle\ge\frac12\langle\phi,|D_0'|\phi\rangle-\mathbf m'\|\phi\|^2\ge\frac{c_E}{2}\langle\phi,\Lambda\phi\rangle-\mathbf m'\|\phi\|^2.
\end{equation*}

\emph{Step 4} (the count, given \eqref{eq:polar-form-bound}).
Let $\mathcal W$ be a space of spinors $\phi$ over $B$ on which $\|D_{A''}\phi\|^2\le\lambda'\|\phi\|^2$.
By Step~3, every $\phi\in\mathcal W$ satisfies
\begin{equation*}
\langle\phi,\Lambda\phi\rangle\le\mathsf e\,\|\phi\|^2,\qquad\mathsf e:=\frac{2}{c_E}\bigl(\sqrt{\lambda'}+\mathbf m'\bigr),
\end{equation*}
and hence $\langle\phi,\allowbreak 1[\Lambda\le2\mathsf e]\phi\rangle\ge\frac12\|\phi\|^2$.
Consequently
\begin{equation*}
\tfrac12\dim\mathcal W\le\operatorname{Tr}\bigl(1_B\,1[\Lambda\le2\mathsf e]\,1_B\bigr)\le192\,|B|\,\frac{\pi\mathsf e^3}{6},
\end{equation*}
and so $N_a(\lambda')\le64\pi(4\ell)^3\mathsf e^3$.

\emph{Step 5} (\eqref{eq:polar-form-bound} from scalar multipliers).
With a matrix $w$ of norm $\|w\|=\mathsf h_3$ \cite[Definition~2.10]{QEDII} one has
\begin{equation*}
(V\psi)(k,\vec g)=-w\,\bigl(\exp(iqA''(k,k+\vec ge_1))-1\bigr)\,\psi\bigl(k+\vec ge_1,\vec g\operatorname{diag}(-1,1,1)\bigr).
\end{equation*}
The translations commute with $\Lambda$, and the multipliers attached to the directions $-e_\mu$ are translates of the complex conjugates of those attached to $+e_\mu$.
It is therefore enough for \eqref{eq:polar-form-bound} that
\begin{equation*}
\mathsf y:=\max_\mu\bigl\|\Lambda^{-1/2}\bigl(\exp(iq\zeta A'_\mu)-1\bigr)\Lambda^{-1/2}\bigr\|\le\frac{c_E}{2\mathsf h_3}\qquad\text{on }\ell^2(\mathbb Z^3).
\end{equation*}
From $|\exp(i\theta)-1-i\theta|\le\frac12\theta^2$ and $\Lambda\ge\mathbf m'$ one obtains
\begin{equation*}
\mathsf y\le q\,\|\Lambda^{-1/2}\zeta A^{\mathrm C}_\mu\Lambda^{-1/2}\|+\frac{q}{\mathbf m'}\,\|\zeta A^{\mathrm h}\|_\infty+\frac{q^2}{2\mathbf m'}\,\|\zeta A'\|^2_\infty.
\end{equation*}

\emph{Step 6} (scales; the good event).
For $\tau>0$ one has $A^{\mathrm C}=d^*c^\tau+A^{>\tau}$, where
\begin{equation*}
c^\tau:=\int_0^\tau\exp(-t\Delta)\,P^\perp F\,\mathrm dt,\qquad A^{>\tau}:=d^*\int_\tau^\infty\exp(-t\Delta)\,P^\perp F\,\mathrm dt.
\end{equation*}
Every value of these cochains has the form $\langle\beta,F\rangle$ for some $\beta\in C^2(\Gamma)$.
The diagonal of the scalar heat kernel obeys, uniformly in $n$,
\begin{equation*}
p_s(0)-n^{-3}\le c_h\min\{1,s^{-3/2}\},
\end{equation*}
as follows from $4\sin^2\frac{\pi j}{n}\ge16j^2/n^2$.
Together with (K4) and $d^*d\le\Delta$, this bounds $\|\beta\|^2$ by $4c_h\sqrt\tau$ for $c^\tau_{\mu\nu}(k)$, by $\sqrt2\,c_h\tau^{-1/2}$ for $A^{>\tau}_\mu(k)$ and by $3c_h$ for $A^{\mathrm C}_\mu(k)$, while $\|\beta\|^2=n^{-3}$ for $(F_{\mathrm h})_{\mu\nu}$.
Let $\mathcal G_a$ be the event on which $|\langle\beta,F\rangle|\le r\|\beta\|$ holds for all these $\beta$, the cells ranging over $B^{++}$ and $\tau$ over $\{4^i:0\le i\le J\}$.
Applying \cite[Theorem~6.27\,(c)]{QEDII} to one $\beta$ at a time, one finds
\begin{equation*}
P_n(\mathcal G_a^c)\le20\,(J+1)\,(10\ell)^3\exp\bigl(-\tfrac12r^2\bigr).
\end{equation*}
On $\mathcal G_a$ the following bounds hold in $B^{++}$:
\begin{equation*}
|A^{>4^i}|\le c_1r\,2^{-i/2},\qquad|c^{4^i}|\le c_2r\,2^{i/2},\qquad|A^{\mathrm C}|\le c_0r,\qquad|A^{\mathrm h}|\le5\ell rn^{-3/2}\le r\,\mathbf m'^{1/2}.
\end{equation*}

\emph{Step 7} (dyadic sum).
The spectrum of $\Lambda$ is contained in $[\mathbf m',\sqrt{13}]$.
For $0\le j\le J$ let $\Pi_j:=1[\Lambda\in(2^{1-j},2^{2-j}]]$, and for $f\in\ell^2(\mathbb Z^3)$ put $f_j:=\Pi_jf$ and $x_j:=\|\Lambda^{1/2}f_j\|$; then
\begin{equation*}
\|f_j\|\le2^{(j-1)/2}x_j,\qquad\|\nabla f_j\|\le2^{2-j}\|f_j\|.
\end{equation*}
The quantities $g_j$ and $y_j$ are defined in the same way for a second function $g$.
For $i\le j$ the cochain $A^{\mathrm C}$ is split at $\tau=4^i$.
The part $A^{>\tau}$ is estimated by its supremum.
In the part $d^*c^\tau$ a summation by parts shifts the difference onto $\bar f_ig_j\zeta$, and the contribution is bounded by
\begin{equation*}
2\sup|c^\tau|\bigl(\|\nabla f_i\|\,\|g_j\|+\|f_i\|\,\|\nabla g_j\|+\ell^{-1}\|f_i\|\,\|g_j\|\bigr).
\end{equation*}
Since $\ell^{-1}<2^{3-J}$, this gives
\begin{equation*}
|\langle f_i,\zeta A^{\mathrm C}_\mu g_j\rangle|\le(c_1+32c_2)\,r\,2^{-i/2}\,\|f_i\|\,\|g_j\|\le\tfrac12c_3\,r\,2^{j/2}x_iy_j,
\end{equation*}
and the case $i>j$ is symmetric.
The Schwarz inequality yields
\begin{equation*}
\sum_{i\le j}2^{j/2}x_iy_j\le\bigl((J+1)\,2^{J+1}\bigr)^{1/2}\|x\|\,\|y\|,
\end{equation*}
and $2^{J+1}<16/\mathbf m'$.
Consequently, on $\mathcal G_a$,
\begin{equation*}
\|\Lambda^{-1/2}\zeta A^{\mathrm C}_\mu\Lambda^{-1/2}\|\le4c_3\,r\sqrt{(J+1)/\mathbf m'}.
\end{equation*}

\emph{Step 8} (assembly).
Put $\mathsf t:=qr\sqrt{(J+1)/\mathbf m'}$.
Steps~5 to~7 give $\mathsf y\le c_4(\mathsf t+\mathsf t^2)$; hence there is $\mathsf t_0(\rho)>0$ such that \eqref{eq:polar-form-bound} holds on $\mathcal G_a$ whenever $\mathbf m'\ge\mathsf t_0^{-2}q^2r^2(J+1)$.
Choose
\begin{equation*}
r^2:=2\log\bigl(20\,(J+1)\,(10\ell)^3\,\mathbf m^{-3}\bigr).
\end{equation*}
With this choice $P_n(\mathcal G_a^c)\le\mathbf m^3$ and $r^2(J+1)\le c_5\log^2\frac2{\mathbf m}$, and one sets $C_*:=c_5\mathsf t_0^{-2}$.
On $\mathcal G_a^c$ one uses $N_a\le192(4\ell)^3$.
Steps~1 and~4 then yield
\begin{equation*}
n^{-3}E_nN_A(\lambda)\le4096\,\pi\,\mathsf e^3+12288\,\mathbf m^3.
\end{equation*}
By Lemma~5.1, $N_A$ vanishes below $\mathbf m^2$, so that
\begin{equation*}
\operatorname{Tr}\bigl[(D^*D)^{-2}\bigr]=2\int_{\mathbf m^2}^\infty\lambda^{-3}N_A(\lambda)\,\mathrm d\lambda;
\end{equation*}
combined with $(\xi+\zeta)^3\le4(\xi^3+\zeta^3)$, this gives the bound on $\mathsf T_n$.
\end{proof}

\begin{proof}[Proof of Theorem~5.7]
\emph{Step 1.}
By Proposition~5.6\,(a), $I_n(u)=\|du\|^2+K_n(u,u)$; combined with Lemma~5.8, Lemma~5.9\,(c) and Lemma~5.10\,(a), this identity gives the first assertion of the theorem.
Assertion (i) follows from $\|D^{-1}\|\le\mathbf m^{-1}$, since $D$ acts on a space of dimension $192n^3$.
Assertion (ii) is a consequence of Lemma~5.10\,(b), and (iii) of Proposition~5.11.

\emph{Step 2} (proof of (a)).
From \cite[Section~1.2]{QEDII} and the notation fixed above one has $n\bar m\eta=8\bar mhL^T>8/L$.
Hence the number
\begin{equation*}
\mathbf m':=\max\Bigl\{\tfrac52L\mathbf m,\;C_*q^2\log^2\frac2{\mathbf m}\Bigr\}
\end{equation*}
satisfies $n\ge20/\mathbf m'$.
Inserting $q^2=e^2\eta$ and $\mathbf m=\bar m\eta$ into (iii), one arrives at the bound of (a).
The constant $C_\kappa$ results from $\frac{q^2}{\bar m\eta}=\frac{e^2}{\bar m}$, from $\bar m\eta\le1$ and from
\begin{equation*}
\sup_{0<\mathsf x<2}\mathsf x^\kappa\log^6\frac2{\mathsf x}=2^\kappa\Bigl(\frac{6}{\kappa\exp1}\Bigr)^6,
\end{equation*}
the supremum being attained at $\mathsf x=2\exp(-6/\kappa)$.

\emph{Step 3} (the bound in (b)).
Let $\mathbf m=\boldsymbol\mu_N$.
Suppose first that $\mathbf m\le\eta^{1/3}$ and that $N$ is so large that $\frac52L\eta^{1/3}\le1$ and $C_*e^2\eta\log^2\frac2{\bar m\eta}\le1$; then the $\mathbf m'$ of Step~2 satisfies $\mathbf m'\le1$.
Part (iii) with this $\mathbf m'$, combined with $\frac{q^2}{\mathbf m}\le\varkappa^{-1}\log^{-3/2}\frac2\eta$ and with $\log\frac2{\mathbf m}\le\log\frac2{\bar m\eta}$, gives
\begin{equation*}
\mathsf c_n\le C_B\max\Bigl\{\bigl(\tfrac52L\bigr)^3\varkappa^{-1}\log^{-3/2}\frac2\eta,\;C_*^3\varkappa^{-4}\Bigl(\log\frac2{\bar m\eta}\Big/\log\frac2\eta\Bigr)^6\Bigr\}.
\end{equation*}
Suppose next that $\mathbf m>\eta^{1/3}$.
In this case (ii), together with $n\ge8/\eta$, gives $\mathsf c_n\le C_Ae^2(\eta^{1/6}+\eta^{8/3})$.

\emph{Step 4} (the atom in (b)).
By \cite[Section~1.2]{QEDII}, every real sequence qualifies as a mass assignment.
By \cite[Proposition~47.1\,(b1), (b5)]{QEDIII}, each lattice probe is orthogonal to the constant cochains, and $\boldsymbol\sigma_i\to\|f\|$.
Along every defining sequence one therefore has $\bar I(\beta)\le(1+\frac{\pi^2}{4}c_\varkappa)\|f\|^2$ for all probes $\beta$, and Proposition~5.6\,(d) applies.

\emph{Step 5} (proof of (c)).
Part (b), applied for every $\varkappa$, gives $\bar I(\beta)=\|f\|^2$ for every limit point $\mu$.
Let $\mathsf s_\beta$ be as in Proposition~5.6\,(c).
Taking $\Phi=r\tanh(X(f)/r)$ and letting $r\to\infty$, one obtains
\begin{equation*}
\|f\|^2=\mu\bigl(X(f)\,\mathsf s_\beta\bigr)\le\|X(f)\|_2\,\|\mathsf s_\beta\|_2\le\|f\|^2,
\end{equation*}
whence $\mathsf s_\beta=X(f)$.
The choices $\Phi=\cos tX(f)$ and $\Phi=\sin tX(f)$ then lead to
\begin{equation*}
\frac{\mathrm d}{\mathrm dt}\,\mu\bigl(\exp(itX(f))\bigr)=-t\,\|f\|^2\,\mu\bigl(\exp(itX(f))\bigr).
\end{equation*}
By Step~7 of the proof of Proposition~5.6 it follows that $\mu=\gamma$, and the tightness of the laws \cite[Theorem~1.3\,(a)]{QEDII} yields the convergence.
\end{proof}

\subsection{A finite-torus moment representation}
\begin{named}{Proposition 5.12}
Let $0\le\rho\le1$, $w_i>0$, $q>0$, $\mathbf m>0$ and $n\ge4$, and set $k_A:=\lfloor\frac n2\rfloor$.
\begin{itemize}
\item[(a)] As matrices,
\begin{equation*}
\bigl(G_n(i+j+1)\bigr)_{0\le i,j\le k_A-1}\ge0\qquad\text{and}\qquad\bigl(G_n(i+j+2)\bigr)_{0\le i,j\le k_A-2}\ge0,
\end{equation*}
the first inequality for the link plane $x_1=\frac12$ and the second for the site plane $x_1=0$.
\item[(b)] One can find a positive measure $\sigma_n$ on $[0,\infty)$, having at most $k_A-1$ atoms, for which
\begin{equation*}
G_n(t)=\int\lambda^{t-1}\,\mathrm d\sigma_n(\lambda)
\end{equation*}
holds for every integer $t$ with $1\le t\le2k_A-2$; this range contains every integer $t$ with $1\le t\le\frac n2$.
One writes $\lambda=\exp(-E)$, with $E\in(-\infty,\infty]$.
This represents every integer step in the stated range. The value $t=0$ is excluded, since the site plane cuts through the slice $L_0$.
\item[(c)] Let $0<E_0<\varepsilon_*$.
For $1\le t\le k_A$ one has
\begin{equation*}
G_n(t)=\int_{[E_0,\infty]}\exp(-E(t-1))\,\mathrm d\varrho_n(E)+R_n(t),
\end{equation*}
with a measure $\varrho_n\ge0$ and a remainder subject to
\begin{equation*}
0\le R_n(t)\le\frac{96}{\mathbf m^2}\exp\bigl(\varepsilon_*-(\varepsilon_*-E_0)k_A-E_0t\bigr).
\end{equation*}
If the weights are standard and $\mathbf m\le(1+\rho)\,\mathrm{sw}$, then $2\vartheta<\varepsilon_*$, and $E_0=2\vartheta$ is an admissible choice.
If $n\ge6$ is even and $G_n(2)\ne0$, then no positive measure carried by $\{E>0\}$ represents $G_n$ on the range of (b).
\item[(d)] Any sequence $n_i\to\infty$ possesses a subsequence along which $G_{n_i}(t)$ converges, for every $t\ge1$, to a limit $G_\infty(t)$, and there is one positive measure $\varrho_\infty$ such that
\begin{equation*}
G_\infty(t)=\int_{[\varepsilon_*,\infty]}\exp(-E(t-1))\,\mathrm d\varrho_\infty(E)\qquad(t\ge1).
\end{equation*}
\end{itemize}
\end{named}

The proof is given in Appendix~\ref{app:finite-torus-proof}.


The finite-torus representation in Proposition~5.12(c) has a remainder.
The local criterion below uses (E1) and (E2) directly and does not require
an estimate for that remainder.

\subsection{A local meson criterion}
\begin{named}{Proposition 5.13}
Let $\rho\ge0$, $w_i>0$, $q>0$, $\mathbf m>0$ and $n\ge8$.
Let $t_0$ be an integer with $4\le t_0\le\frac n2$, and define
\begin{equation*}
\varsigma:=\mathsf a\,t_0=\frac{2\mathbf m\,t_0}{\mathsf h_3},\qquad c(\varsigma):=\frac{45}{\varsigma}+\frac{157}{\varsigma^2}+\frac{228}{\varsigma^3}.
\end{equation*}
\begin{itemize}
\item[(a)] \emph{The far part is dominated by the near part.}
One has
\begin{equation*}
\mathsf S_n\le\bigl(1+c(\varsigma)\bigr)\sum_{0<|t|_n\le t_0}|t|_n^2\,G_n(t).
\end{equation*}
\item[(b)] \emph{One box.}
If $\varsigma\ge190$, then $\mathsf S_n\le\frac56B_n(t_0)$, where
\begin{equation*}
B_n(t_0):=\sum_{\pi(x)=0}\ \sum_{|\pi(y)|_\infty\le t_0}|\pi(y)|^2\,\mathsf M_n(x,y).
\end{equation*}
\item[(c)] \emph{Mirror pairs.}
Assume moreover $\rho\le1$.
For $1\le d\le n-1$ denote by $\mathsf p_n(d)$ the largest value of
\begin{equation*}
\mathsf M_n(\theta_{x_j,c}x,x)=\mathbb E\bigl[\Theta_{x_j,c}(\mathfrak M(x))\,\mathfrak M(x)\bigr]
\end{equation*}
over the directions $j$, the sites $x\in\mathbb X$ and the numbers $c\in\frac12\mathbb Z$ with $2(\pi(x)_j-c)=d$, and set $\bar{\mathsf p}_n(r):=\max\{\mathsf p_n(r),\mathsf p_n(n-r)\}$.
Then $\mathsf M_n(x,y)\le\bar{\mathsf p}_n(|v(x,y)|_\infty)$ whenever $v(x,y)\ne0$, and
\begin{equation*}
B_n(t_0)\le1.8\cdot10^5\sum_{r=1}^{t_0}r^4\,\bar{\mathsf p}_n(r).
\end{equation*}
In addition $\mathsf p_n(n-r)=\mathsf p_n(r)$, so that $\bar{\mathsf p}_n=\mathsf p_n$.
\item[(d)] Suppose that $G_n(t)\le C_0t^{-2}$ for $1\le t\le t_0$ and that $c_0\le\mathbf mt_0\le c_0'$.
Then $\mathsf c_n\le C_1q^2/\mathbf m$, where $C_1=48\,\mathsf h_3^2C_0c_0'\bigl(1+c(2c_0/\mathsf h_3)\bigr)$.
\item[(e)] \emph{Continuum.}
Let $0<\rho<1$, let the weights be standard, and let $e>0$ be arbitrary.
Let $c_{\boldsymbol\mu}>0$, and assume that $\boldsymbol\mu_N\ge c_{\boldsymbol\mu}\,\bar m\eta$ and $\limsup_{T\to\infty}\boldsymbol\mu_N\mathsf S_n\le C_M$ for all large $N$.
Then
\begin{equation*}
\limsup_T\,\sup_u\frac{I_n(u)}{\|du\|^2}\le1+\frac{Ce^2}{\bar m},
\end{equation*}
the supremum being taken over the $u$ with $\sum_xu(x)=0$ and $du\ne0$, with the constant $C=6\pi^2\mathsf h_3^2C_M/c_{\boldsymbol\mu}$.
Moreover every volume-first limit point satisfies $\mathrm{kl}(\{0\})\ge(1+Ce^2/\bar m)^{-1}$.
The ratio $e^2/\bar m$ need not be small.
\end{itemize}
\end{named}

\begin{proof}
\emph{Step 1} (proof of (a)).
Let $N_0:=\sum_{1\le t\le t_0}t^2\mathsf f(t)$; this is the sum that appears on the right of (a).
The set $\mathcal I:=\{t:\frac{t_0}2<t\le t_0\}$ has at least $\frac{t_0}2$ elements, and each $t\in\mathcal I$ satisfies $t-1\ge\frac38t_0$.
Hence
\begin{equation*}
\sum_{t\in\mathcal I}\bigl(\mathsf f(t)+\mathsf f(t-1)\bigr)\le\frac{128}{9t_0^2}\,N_0,
\end{equation*}
and (E1) provides some $t_1\in\mathcal I$ with $\mathsf a\mathsf F(t_1)\le\frac{256}{9t_0^3}N_0$.
Since $t_1+1\le\frac54t_0$, the inequality (E2) at the point $t_1+1$ gives
\begin{equation*}
\sum_{t>t_0}t^2\mathsf f(t)\le\frac{256\,N_0}{9\,\mathsf a\,t_0^3}\Bigl[\frac{25}{16}\,t_0^2+\frac{11\,t_0}{2\mathsf a}+\frac{8}{\mathsf a^2}\Bigr]\le c(\varsigma)\,N_0.
\end{equation*}

\emph{Step 2} (proof of (b)).
For each direction $j$, split
\begin{equation*}
\sum_{\pi(x)=0}\sum_y|\pi(y)_j|_n^2\,\mathsf M_n(x,y)=\mathrm{near}_j+\mathrm{tail}_j
\end{equation*}
at the value $|\pi(y)_j|_n=t_0$, the terms with $|\pi(y)_j|_n\le t_0$ forming $\mathrm{near}_j$.
The sum over $j$ of the left sides is $3\mathsf S_n$, and $\mathrm{tail}_j\le c(\varsigma)\,\mathrm{near}_j$ by (a) and Lemma~5.9\,(a).
Within $\mathrm{near}_1$, the sites $y$ with $\max(|\pi(y)_2|_n,|\pi(y)_3|_n)>t_0$ contribute no more than $\mathrm{tail}_2+\mathrm{tail}_3$.
Therefore
\begin{equation*}
\sum_j\mathrm{near}_j\le B_n(t_0)+2c(\varsigma)\sum_j\mathrm{near}_j,
\end{equation*}
and, since $c(\varsigma)\le c(190)<\frac14$, one arrives at $3\mathsf S_n\le\bigl(1+c(\varsigma)\bigr)\,2B_n(t_0)$.

\emph{Step 3} (proof of (c)).
Let $x,y$ be sites with $v(x,y)\ne0$, and let $\pi(y)_j-\pi(x)_j\equiv d\in[1,n-1]$.
Choose lifts of $x$ and $y$ for which this congruence is an equality, and put $c:=\frac12(\pi(x)_j+\pi(y)_j)$ and $y':=\theta_cx$.
Then $y$ and $y'$ both satisfy $2(\pi(\cdot)_j-c)=d$, and both lie in the open slab, since $y_j-c\in(\frac d2-\frac12,\frac d2+\frac12)$.
In the direction $j$, the Schwarz inequality for (RP) gives
\begin{equation*}
\mathsf M_n(x,y)^2=\mathsf M_n(\theta_cy',y)^2\le\mathsf M_n(\theta_cy',y')\,\mathsf M_n(\theta_cy,y)\le\mathsf p_n(d)^2.
\end{equation*}
At most $24r^2+2\le26r^2$ cells $k$ satisfy $|k|_\infty=r$; each has $|k|^2\le3r^2$, and to each of them correspond $48^2$ pairs $(x,y)$ with $\pi(x)=0$ and $\pi(y)=k$, since a cell has $48$ sites.
To prove the final claim of (c), assume $2(\pi(x)_j-c)=n-r$ and set $c':=c+\frac n2\in\frac12\mathbb Z$.
Since $\theta_{c'}$ is $\theta_c$ followed by the translation by $ne_j$, the two reflections induce one and the same map of $\mathbb X_n$.
For $x':=\theta_{c'}x$ one has $2(\pi(x')_j-c')=r$ and $\mathsf M_n(\theta_cx,x)=\mathsf M_n(x',\theta_{c'}x')$.
Take $i\ne j$ and $c_i\in\frac12\mathbb Z$.
Then $g:=\theta_{x_j,c'}\theta_{x_i,c_i}$ is an even element of $W$, and $x'':=\theta_{x_i,c_i}x'$ satisfies
\begin{equation*}
gx'=\theta_{c'}x'',\qquad g\theta_{c'}x'=x'',\qquad\pi(x'')_j=\pi(x')_j.
\end{equation*}
Lemma~5.9\,(a) therefore yields $\mathsf M_n(x',\allowbreak \theta_{c'}x')=\allowbreak \mathsf M_n(\theta_{c'}x'',\allowbreak x'')\le\mathsf p_n(r)$.
This proves $\mathsf p_n(n-r)\le\mathsf p_n(r)$, and the reverse inequality is the same statement with $n-r$ put for $r$.

\emph{Step 4} (proof of (d)).
Part (a) gives $\mathsf S_n\le\bigl(1+c(\varsigma)\bigr)\,2C_0t_0$, and the function $c(\cdot)$ is decreasing.
Lemma~5.9\,(c) completes the proof of (d).

\emph{Step 5} (proof of (e)).
Theorem~5.7 gives $I_n(u)\le(1+6\pi^2\mathsf h_3^2q^2\mathsf S_n)\,\|du\|^2$, and $q^2/\boldsymbol\mu_N\le e^2/(c_{\boldsymbol\mu}\bar m)$.
Let $\mu=\lim\mu^{(N_i)}$, where $\mu^{(N_i)}=\lim_jP_{N_i,T_{ij}}$.
By (K6), applied on the separable Hilbert space $H_{-4}$, the limit point $\mu$ admits defining sequences whose $T_i$ can be taken arbitrarily large.
Along such a sequence \cite[Proposition~47.1\,(b1), (b5)]{QEDIII} yields $\bar I(\beta)\le(1+Ce^2/\bar m)\|f\|^2$ for all probes.
The assertion on $\mathrm{kl}(\{0\})$ then follows from Proposition~5.6\,(d).
\end{proof}

\emph{Conventions.}
Put $g:=q^2/\mathbf m$.
A mass assignment is called \emph{physical} if there are numbers $0<c\le C<\infty$ such that $c\,\eta\le\boldsymbol\mu_N\le C\eta$ for all large $N$.
In Theorem~5.14 and its proof, the letters $C$, $C_\bullet$ and $c_\bullet$ stand for numbers that depend on $\rho$ alone;
the exceptions are the datum $c_{\boldsymbol\mu}$ and the numbers $C_M$, $C_F$ and $\mathsf g_*$, which are defined in Theorem~5.14.
The norm of $L^{\mathsf p}(P_n)$ is written $\|\cdot\|_{L^{\mathsf p}}$.
The symbols $J$, $J^*$, $\mathrm{Rsy}$, $\mathrm{Vx}_\nu$ and $\mathrm{Ehp}^\pm_\nu$ denote the objects of the fibre of \cite[Section~G5.6.0]{QEDIII};
in the proof of Theorem~5.16, however, $J$ and $J^*$ are the cluster maps of \cite[(I.3.1)]{QEDIII} on spinor fields, which are given by the same formulas.
In the proofs of Theorems~5.14 and~5.16 the letter $J$ is different from the integer $J$ of the proof of Proposition~5.11.
In the same way, the shells $\Pi_j$ of the proof of Theorem~5.14 are different from those of Step~7 of the proof of Proposition~5.11.
Throughout the proof of Theorem~5.14, the symbol $D^{(p)}$ of the fibre, its inverse $\mathrm{Rsy}(p)$ and the remainder $\mathrm{Er}$ introduced below are formed with the mass $\mathbf m'$ of Setting~(S) in the place of the free number $\ge0$ of \cite[Section~G5.6.0]{QEDIII}, and not with the lattice mass $\mathbf m$ of the law.
For $0<\rho<1$, \cite[Proposition~37.1\,(f)]{QEDIII} provides a number $C_E=C_E(\rho)<\infty$ such that, for every $\mathbf m'>0$ and every $p\in[-\pi,\pi]^3$, the symbol $D^{(p)}$ formed with $\mathbf m'$ is invertible and
\begin{equation*}
\mathrm{Rsy}(p)=J\bigl(\mathbf m'+i\gamma\cdot p\bigr)^{-1}J^*+\mathrm{Er}(p,\mathbf m'),\qquad\bigl|\mathrm{Er}(p,\mathbf m')\bigr|_2\le C_E,
\end{equation*}
where the first relation defines the remainder $\mathrm{Er}(p,\mathbf m')$ and $|\cdot|_2$ is the operator norm on the fibre;
one abbreviates $\mathrm{Er}(p):=\mathrm{Er}(p,\mathbf m')$.
The vector $\mathsf s(p)$ is the one of Step~4 of the proof of Lemma~5.8, and in the proof of Theorem~5.14 the letter $\mathsf s$ stands for its length $|\mathsf s(p)|$.

\emph{Classical facts.}
\begin{itemize}
\item[(K7)] Let $\mathsf g$ be a standard Gaussian vector in $\mathbb R^N$, let $s\in\mathbb R^N$ and $\theta\ge0$, and let $G\ge0$ be a real symmetric matrix such that $2\theta\|G\|<1$.
Then
\begin{equation*}
E\exp(s\cdot\mathsf g)=\exp\tfrac12|s|^2,\qquad E\exp\bigl(\theta\langle\mathsf g,G\mathsf g\rangle\bigr)=\det(1-2\theta G)^{-1/2}.
\end{equation*}
\item[(K8)] (Poisson summation.)
Let $K\in\ell^1(\mathbb Z^3)$ be scalar or matrix valued, and let $N\ge1$ be an integer.
A function $\psi$ on $\mathbb Z^3/N\mathbb Z^3$ is read as an $N\mathbb Z^3$-periodic function on $\mathbb Z^3$, and the momenta $p$ run through $\frac{2\pi}N\mathbb Z^3$.
Then, on $\mathbb Z^3/N\mathbb Z^3$, the Fourier multiplier whose symbol is $\sum_zK(z)\exp(ip\cdot z)$ acts as $\psi\mapsto\sum_zK(z)\,\psi(\cdot+z)$.
(Plancherel's formula on $\mathbb Z^3$.)
Every $K\in\ell^1(\mathbb Z^3)$ satisfies
\begin{equation*}
\sum_z|K(z)|^2=(2\pi)^{-3}\int_{[-\pi,\pi]^3}\Bigl|\sum_zK(z)\exp(ip\cdot z)\Bigr|^2\,\mathrm dp.
\end{equation*}
\item[(K9)] The inequalities of H\"older, Minkowski and Markov, Tonelli's theorem, and Lyapunov's inequality
\begin{equation*}
\|Z\|_{L^{\mathsf p}}\le\|Z\|_{L^2}^{\vartheta}\,\|Z\|_{L^{2\mathsf p}}^{1-\vartheta},\qquad\vartheta=\frac1{\mathsf p-1},\quad\mathsf p\ge2.
\end{equation*}
\item[(K10)] Let all spaces be of finite dimension, let $\beta\ne0$, and let $X$ and $X+\beta R^*R$ be invertible.
Then
\begin{equation*}
\bigl(\beta-\beta^2R(X+\beta R^*R)^{-1}R^*\bigr)\bigl(\beta^{-1}+RX^{-1}R^*\bigr)=1\qquad\text{(Woodbury)}.
\end{equation*}
In finite dimension one has, moreover, $\det(1-XY)=\det(1-YX)$ (Sylvester).
\item[(K11)] If a real analytic function on $\mathbb R^d$ does not vanish identically, its zero set has Lebesgue measure zero.
Two rational functions of one variable which coincide on an interval coincide at every point where both of them are regular.
\end{itemize}

The earlier localized estimate loses a factor $\log^2(2/\mathbf m)$
in the auxiliary mass. The next theorem removes this loss. We first
factor the localized hopping perturbation so that its high--low terms
are absorbed in the coercive free part. The remaining quadratic block
term is decomposed into its eight coordinate-reflection sectors.
In the sector even under all three reflections, the odd singular part
of the free resolvent cancels. Each of the other seven sectors is odd
under a reflection, and Lemma~5.15 controls its spatial average.
These two estimates are summable over frequency scales. The last
three proof steps pass through auxiliary masses on the same box and
then to the global eigenvalue count and the continuum Fisher bound.

\subsection{The uniform singular-value estimate}
\begin{named}{Theorem 5.14}
\emph{Proposition~5.11 with no logarithmic factor; conditions (M$_{\mathrm{loc}}$) and (FI), and an atom at zero for every charge.}
Let $0<\rho<1$, and let the weights be standard.
There is a constant $C_\sharp=C_\sharp(\rho)<\infty$ with the following property.
Whenever $q>0$, $0<\mathbf m\le\mathbf m'\le1$, $\mathbf m'\ge q^2$ and $n\ge20/\mathbf m'$, one has
\begin{gather*}
n^{-3}E_nN_A(\lambda)\le C_\sharp\bigl(\sqrt\lambda+\mathbf m'\bigr)^3\quad(\lambda>0),\qquad\mathsf T_n\le20\,C_\sharp\,\mathbf m'^3\,\mathbf m^{-4},\\
\mathsf c_n\le C_B^\sharp\,g\,\Bigl(\frac{\mathbf m'}{\mathbf m}\Bigr)^3,\qquad C_B^\sharp:=60\,\pi^2\mathsf h_3^4\,C_\sharp.
\end{gather*}
For the continuum scaling, let $e>0$ and $\bar m\in(0,1]$ be arbitrary, let $c_{\boldsymbol\mu}>0$, and let the mass assignment satisfy $c_{\boldsymbol\mu}\bar m\eta\le\boldsymbol\mu_N\le1$ for all large $N$;
every physical mass assignment is of this kind, for instance $m_{0,N}=0$, where no mass shift is subtracted.
Put $\mathsf g_*:=\max\{1,e^2/(c_{\boldsymbol\mu}\bar m)\}$.
Then the following assertions hold for all large $N$ and for all $T$ such that $n\boldsymbol\mu_N\ge20$.
\begin{itemize}
\item[(a)] For every integer $t_0\ge0$,
\begin{equation*}
\boldsymbol\mu_NB_n(t_0)\le3\boldsymbol\mu_N\mathsf S_n\le\frac{3\pi^2\mathsf h_3^2}{8}\,\boldsymbol\mu_N\mathsf T_n\le C_M:=\frac{15}{2}\,\pi^2\mathsf h_3^2\,C_\sharp\,\mathsf g_*^3,
\end{equation*}
whereas in Proposition~5.13 the integer $t_0$ satisfies $4\le t_0\le\frac n2$.
This is \textbf{(M$_{\mathrm{loc}}$)}.
\item[(b)] For these $T$ condition \textbf{(FI)} holds in the form
\begin{equation*}
K_n(u,u)\le C_F\,\frac{e^2}{\bar m}\,\|du\|^2\quad\text{whenever }\sum_xu(x)=0,\qquad C_F:=\frac{15\,\pi^4\mathsf h_3^4\,C_\sharp\,\mathsf g_*^3}{c_{\boldsymbol\mu}}.
\end{equation*}
If $e^2\ge c_{\boldsymbol\mu}\bar m$, then
\begin{equation*}
C_F\,\frac{e^2}{\bar m}=15\,\pi^4\mathsf h_3^4\,C_\sharp\Bigl(\frac{e^2}{c_{\boldsymbol\mu}\bar m}\Bigr)^4,
\end{equation*}
so that the bound is of the fourth order in $e^2/\bar m$;
a bound linear in $e^2/\bar m$ is not proved.
Since $n\boldsymbol\mu_N\ge8c_{\boldsymbol\mu}\bar mhL^T$, the $T$ with $n\boldsymbol\mu_N\ge20$ include every $T$ with $8c_{\boldsymbol\mu}\bar mhL^T\ge20$.

\item[(c)] Every volume-first limit point satisfies $\mathrm{kl}(\{0\})\ge(1+C_Fe^2/\bar m)^{-1}>0$, that is, it has a massless photon.
\end{itemize}
\end{named}

The ratio $e^2/\bar m$ need not be small, and if $e^2\le c_{\boldsymbol\mu}\bar m$, the constants depend on $\rho$ and $c_{\boldsymbol\mu}$ alone.
Here $E_n$ denotes expectation for the cutoff law formed with the mass assignment under consideration.
In Theorem~5.14 the law $P_{N,T}$ is the law $P^{\mathrm{ch}}_{N,T}$ of \cite[Section~1.2]{QEDII}, and $\mathrm{kl}$ is the measure $\mathrm{kl}_\mu$ of Proposition~5.6\,(d), which is the measure $\mathfrak K_\mu$ of \cite[Theorem~1.4, Part~I\,(d)]{QEDII}.
Parts (b) to (d) of Proposition~5.6 hold for an arbitrary mass assignment and every $e>0$.
The hypotheses of Theorem~5.14 exclude the mass assignments with $\boldsymbol\mu_N\le0$ and those with $\boldsymbol\mu_N=o(\eta)$.

\begin{named}{Lemma 5.15}
\emph{Local functionals odd under mirrors: a law of large numbers from reflection positivity.}
Let $-1\le\rho\le1$, $w_i>0$, $M\in\mathbb R$ and $q>0$.
Let $\ell\ge2$, $n>10\ell$ and $1\le\mathsf R\le\ell$ be integers, and put $a_0:=2\mathsf R+4$;
let $\mathsf Z\subset\mathbb Z^3$ be a box with at most $7\ell$ sites per side, and let $\lambda\in\{1,2,3\}$.
For each $y\in\mathsf Z$ let $\mathsf o_y$ be a real polynomial in the variables $F_\sigma$ subject to the following conditions.
\begin{itemize}
\item[\textbf{(L)}] The polynomial $\mathsf o_y$ depends only on the $F_\sigma$ with $\sigma\subset y+[-\mathsf R-\frac12,\mathsf R+\frac12]^3$.
\item[\textbf{(O)}] One has $\mathsf o_{\theta y}(\theta^*F)=-\mathsf o_y(F)$ for every reflection $\theta=\theta_{x_\lambda,c}$ with $c\in\frac12\mathbb Z$ and every $y$ with $y,\theta y\in\mathsf Z$.
\item[\textbf{(S)}] With a number $\mathsf t$, one has $\|\sum_yv(y)\mathsf o_y\|_{L^2}\le\mathsf t\|v\|_1$ for every real $v$ on $\mathsf Z$.
\end{itemize}
Let $w$ be a real function carried by $\mathsf Z$, extended by zero to $\mathbb Z^3$, so that a jump of $w$ at the faces of $\mathsf Z$ enters $\partial w$;
let $x^*$ be an integer in the projection of $\mathsf Z$ on the $\lambda$-th coordinate, and define
\begin{equation*}
\partial w(y):=\begin{cases}w(y)-w(y+e_\lambda)&\text{if }y_\lambda\ge x^*,\\w(y)-w(y-e_\lambda)&\text{if }y_\lambda<x^*.\end{cases}
\end{equation*}
Then
\begin{equation*}
\Bigl\|\sum_yw(y)\,\mathsf o_y\Bigr\|_{L^2}\le3\,\mathsf t\,a_0^{1/2}\sum_y\bigl(|y_\lambda-x^*|+1+a_0\bigr)^{1/2}\,|\partial w(y)|.
\end{equation*}
\end{named}

\begin{proof}[Proof for $\lambda=1$]
Points of $\mathbb Z^3$ are written $y=(y_1,y_\perp)$.

\emph{Step 1} (the mirror).
For $\alpha,t\in\mathbb Z$ put $I_t:=[\alpha+ta_0,\alpha+(t+1)a_0-1]$, and consider only those $t$ for which $I_t$ lies in the projection of $\mathsf Z$.
For a real function $\omega$ of $y_\perp$ put
\begin{equation*}
\Psi_t:=\sum_{y_1\in I_t}\ \sum_{y_\perp}\omega(y_\perp)\,\mathsf o_y.
\end{equation*}
Let $t'\ge t+2$, and set $c:=\alpha+\frac12\bigl((t+t'+1)a_0-1\bigr)$; this number lies in $\frac12\mathbb Z=\mathrm{Adm}(x_1)$ \cite[Definition~2.2]{QEDII}.
The reflection $\theta:=\theta_{x_1,c}$ carries $I_{t'}\times\{y_\perp\}$ onto $I_t\times\{y_\perp\}$.
By (L), the plaquettes on which $\Psi_{t'}$ depends satisfy
\begin{equation*}
2\le x_1-c\le\tfrac72\ell+\mathsf R+\tfrac12<\tfrac n2,
\end{equation*}
which places them in the open slab.

\emph{Step 2} (the sign).
Consequently the polynomial $\Psi_{t'}$ belongs to the class $\mathcal G^+(C(n),x_1,c)$ of \cite[Definition~3.33]{QEDII}.
By the paragraph \emph{Reflections of integrands} of \cite[Section~2.3]{QEDII}, the map $\Theta_{x_1,c}$ of \cite[Definition~2.9]{QEDII} acts on scalar functions of $F$ as complex conjugation composed with $\theta^*$;
on the real polynomial $\Psi_{t'}$ it therefore acts by $\Theta_{x_1,c}\Psi_{t'}=\Psi_{t'}\circ\theta^*$, and (O) gives $\Psi_{t'}\circ\theta^*=-\Psi_t$.
Hence \cite[Theorem~1.1\,(c), (d)]{QEDII} yields
\begin{equation*}
E_n[\Psi_t\Psi_{t'}]=-E_n\bigl[\Theta_{x_1,c}(\Psi_{t'})\,\Psi_{t'}\bigr]\le0.
\end{equation*}

\emph{Step 3} (the intervals).
Divide an interval $I$ of the projection of $\mathsf Z$ into $N_I\le|I|/a_0$ intervals $I_t$ and a remainder of length less than $a_0$.
Step~2 and (S) give
\begin{equation*}
E_n\Bigl(\sum_{t\ \mathrm{even}}\Psi_t\Bigr)^2\le\sum_{t\ \mathrm{even}}E_n\Psi_t^2\le\frac{N_I+1}{2}\,\bigl(\mathsf t\,a_0\|\omega\|_1\bigr)^2,
\end{equation*}
and the odd $t$ obey the same bound;
by (S), the sum over the remainder has norm at most $\mathsf t\,a_0\|\omega\|_1$.
As $a_0(N_I+1)^{1/2}\le(a_0(|I|+a_0))^{1/2}$ and $\sqrt2+1\le3$, one obtains
\begin{equation*}
\Bigl\|\sum_{y_1\in I}\ \sum_{y_\perp}\omega(y_\perp)\,\mathsf o_y\Bigr\|_{L^2}\le3\,\mathsf t\,\|\omega\|_1\bigl(a_0(|I|+a_0)\bigr)^{1/2}.
\end{equation*}

\emph{Step 4} (the weights).
Let $y_1\ge x^*$.
The sum $\sum_{s\ge y_1}\partial w(s,y_\perp)$ telescopes, and its boundary term vanishes because $w$ is extended by zero;
hence it equals $w(y)$, and
\begin{equation*}
\sum_{y_1\ge x^*}w(y)\,\mathsf o_y=\sum_{s\ge x^*}\ \sum_{y_1\in[x^*,s]}\ \sum_{y_\perp}\partial w(s,y_\perp)\,\mathsf o_y.
\end{equation*}
This sum is bounded by Minkowski's inequality over $s$ together with Step~3 for $\omega=\partial w(s,\cdot)$.
The part with $y_1<x^*$ is bounded by the same argument.
\end{proof}

\begin{proof}[Proof of Theorem~5.14]
\textbf{Setting (S) of Steps~1 to~27.}
Let $\ell\ge4$ and $n>10\ell$ be integers, and let $0<\mathbf m\le\mathbf m'\le1$ satisfy $\ell\mathbf m'\ge1$.
No upper bound is imposed on the product $\ell\mathbf m'$;
in particular, the pair $(\ell,2^k\mathbf m')$ fulfils these conditions for every integer $k\ge0$ with $2^k\mathbf m'\le1$.
Let $B$ be a box with at most $4\ell$ sites per side.
The boxes $B^+$ and $B^{++}$, the second of which has at most $7\ell$ sites per side, the cut-off $\zeta$, and the cochains $A^{\mathrm C}$, $A^{\mathrm h}$, $A'$ and $A''=\zeta A'$ are formed as in Step~2 of the proof of Proposition~5.11, with the present value of $\ell$;
that construction is valid for every integer $\ell$ with $6\ell+4<n$.
Put
\begin{equation*}
\mathsf n:=10\ell,\qquad\mathbb T':=\mathbb Z^3/\mathsf n\mathbb Z^3\supset B^{++},\qquad\mathcal V:=\mathbb C^{48}\otimes\mathbb C^4,\qquad\mathcal H':=\ell^2(\mathbb T')\otimes\mathcal V,
\end{equation*}
the site $k+\vec gv_0$ being identified with the pair $(k,\vec g)$ \cite[Definition~2.10]{QEDII}.
By \cite[Proposition~46\,(a), (b)]{QEDIII}, the sites, the hops together with their directions, and the operator of \cite{QEDIII}, expressed in these coordinates, coincide with the corresponding objects of \cite{QEDII}.
On the fibre of \cite{QEDIII} the inner product is $\frac1{48}\sum_{\vec g}$, whereas on $\mathcal V$ it is the plain sum over $\vec g$;
the two inner products give the same operator norms, adjoints and polar parts on $\mathcal V$.
Write $\mathsf q(p):=\sum_\mu4\sin^2\frac{p_\mu}{2}$ and $\mathsf s:=\mathsf q^{1/2}$.
Let $D_0'$ denote the free operator with the mass $\mathbf m'$ on $\mathcal H'$, which is the Fourier multiplier with the symbol $D^{(p)}$, and set
\begin{equation*}
S:=D_0'^{-1},\qquad\Lambda:=\bigl(\mathsf q+\mathbf m'^2\bigr)^{1/2},\qquad\|f\|_\sigma:=\|\Lambda^\sigma f\|.
\end{equation*}
As in Step~3 of the proof of Proposition~5.11, \cite[Proposition~37.1\,(f)]{QEDIII} gives $|D_0'|\ge c_E\Lambda$ and $|D_0'^*|\ge c_E\Lambda$;
moreover $D_0'=\Omega|D_0'|$ with a unitary multiplier $\Omega$, and $\Omega$ commutes with $\Lambda$.

\emph{Shells.}
Put $\kappa_j:=2^j\mathbf m'$ and
\begin{equation*}
\Pi_0:=1[\Lambda\le\kappa_1],\qquad\Pi_j:=1[\kappa_j<\Lambda\le\kappa_{j+1}]\quad(j\ge1),
\end{equation*}
and let $\mathsf J$ be the largest $j$ such that $\kappa_j\le\sqrt{13}$.
Since $\Lambda\le\sqrt{13}<\kappa_{\mathsf J+1}$, one has $\sum_{j=0}^{\mathsf J}\Pi_j=1$.
On $\operatorname{ran}\Pi_j$ the inequalities $\Lambda\ge\kappa_j$ and $\mathsf s\le\kappa_{j+1}$ hold, and so does $\mathsf s>\frac{\sqrt3}{2}\kappa_j$ when $j\ge1$.

\emph{Filters.}
Fix $\varpi\in C^\infty(\mathbb R;[0,1])$ with $\varpi=0$ on $(-\infty,\frac14]$ and $\varpi=1$ on $[\frac9{16},\infty)$, and set $\varpi_k:=\varpi(\mathsf q/\kappa_k^2)$;
thus $\varpi_k=0$ for $k\ge\mathsf J+2$.
Define
\begin{equation*}
\varphi_k:=\varpi_k-\varpi_{k+1},\qquad\Phi_k:=1-\varpi_{k+1},\qquad\tilde\varphi_k:=\varpi_{k-2}-\varpi_{k+3}.
\end{equation*}
Here $\varphi_k$ is carried by $\frac12\kappa_k\le\mathsf s\le\frac32\kappa_k$;
next, $\Phi_k$ equals $1$ where $\mathsf s\le\kappa_k$ and is carried by $\mathsf s\le\frac32\kappa_k$;
finally, $\tilde\varphi_k$ equals $1$ where $\frac3{16}\kappa_k\le\mathsf s\le4\kappa_k$ and is carried by $\frac18\kappa_k\le\mathsf s\le6\kappa_k$.
All these functions are real, even and invariant under $O_{\mathrm{cub}}$.
Their convolution kernels on $\mathbb Z^3$ are denoted by the same letters, and so are the kernels periodised according to (K8) on $\mathbb T'$ and on $\mathbb Z^3/n\mathbb Z^3$.

\emph{Dependence on $\ell$.}
The number $\mathsf e$ of Step~3 satisfies $\mathsf n\mathsf e\ge40\ell\mathbf m'/c_E$, and the numbers $L_j$ of (BL) satisfy $L_j\le\ell$;
the bounds of Steps~8, 11 and~18 contain the factor $(\ell\mathbf m')^3$.
Moreover $\ell^{-1}\le\mathbf m'$, one has $(\kappa_k\ell)^{-1}\le2^{-k}$ for every integer $k\ge0$, the integer $\mathsf R$ of Step~24 satisfies $\mathsf R\le\ell$, and $\frac2\ell\le2\kappa_m$ for every integer $m\ge0$.
The periods obey $\mathsf n\ge1/\kappa_k$ and $n\ge1/\kappa_k$ for every $k\ge0$, which are the conditions on the period in (Ker) and in (Per) at $\kappa=\kappa_k$.
Apart from the factor $(\ell\mathbf m')^3$, each of these relations either follows from $\ell\mathbf m'\ge1$ alone or becomes stronger as $\ell\mathbf m'$ increases.

\emph{Step 1} (factorisation).
Let $V=V^{\mathrm{hl}}+V^{\mathrm{lh}}+V^{\mathrm d}$ be a splitting of an operator $V$ on $\mathcal H'$, and put
\begin{equation*}
\mathsf A:=SV^{\mathrm{hl}},\qquad\mathsf B:=V^{\mathrm{lh}}S,\qquad\mathcal T:=V^{\mathrm{lh}}SV^{\mathrm{hl}}.
\end{equation*}
Expanding the product $(1+\mathsf B)D_0'(1+\mathsf A)$, one obtains the identity \textbf{(F)}
\begin{equation}
\tag{F}
D_0'+V=(1+\mathsf B)D_0'(1+\mathsf A)+V^{\mathrm d}-\mathcal T,
\end{equation}
in which no term of higher order occurs.

\emph{Step 2} (an auxiliary vector).
Let $0\le s<\frac12$ and $D:=D_0'+V-(\mathbf m'-\mathbf m)$;
throughout Steps~1 to~27 the letter $D$ stands for this operator on $\mathcal H'$, and not for $D^{C(n)}_A$.
Suppose that the following two conditions hold.
\begin{itemize}
\item[\textbf{(C1)}] With respect to $\|\cdot\|_0$, and also with respect to $\|\cdot\|_s$, the operators $\mathsf A$ and $\mathsf B^*$ have norm at most $\frac18$, where $\mathsf B^*$ denotes the adjoint of $\mathsf B$ in $\mathcal H'$.
\item[\textbf{(C2)}] For all $f,g\in\mathcal H'$,
\begin{equation*}
|\langle f,(V^{\mathrm d}-\mathcal T)g\rangle|\le\frac{c_E}{4}\,\mathbf m'^{1-2s}\,\|f\|_s\,\|g\|_s.
\end{equation*}
\end{itemize}
To each $\phi\in\mathcal H'$ associate
\begin{equation*}
\phi^\sharp:=(1+\mathsf A)\phi,\qquad\phi_1:=(1+\mathsf B^*)^{-1}\Omega\phi^\sharp.
\end{equation*}
Then (F) yields
\begin{equation*}
\langle\phi_1,D\phi\rangle=\langle\phi^\sharp,|D_0'|\phi^\sharp\rangle+\langle\phi_1,(V^{\mathrm d}-\mathcal T)\phi\rangle-(\mathbf m'-\mathbf m)\langle\phi_1,\phi\rangle.
\end{equation*}
From (C1) and $[\Omega,\Lambda]=0$ one gets $\|\phi\|_\sigma\le\frac87\|\phi^\sharp\|_\sigma$ and $\|\phi_1\|_\sigma\le\frac87\|\phi^\sharp\|_\sigma$ for $\sigma=0$ and for $\sigma=s$.
Taking the real part and using $\mathbf m'^{1-2s}\Lambda^{2s}\le\Lambda$ and $0\le\mathbf m'-\mathbf m\le\mathbf m'$, one arrives at
\begin{equation*}
\frac87\,\|\phi^\sharp\|\,\|D\phi\|\ge\frac{c_E}{2}\,\langle\phi^\sharp,\Lambda\phi^\sharp\rangle-2\mathbf m'\|\phi^\sharp\|^2.
\end{equation*}

\emph{Step 3} (the count).
Assume (C1) and (C2), and put $C_3:=768\,(12/c_E)^3$.
Then the number of eigenvalues not exceeding $\lambda'$ of the quadratic form $\|D\phi\|^2$, restricted to the spinors over $B$, is at most $C_3|B|(\sqrt{\lambda'}+\mathbf m')^3$.
To prove this, let $\mathcal W$ be a space of spinors $\phi$ over $B$ on which $\|D\phi\|^2\le\lambda'\|\phi\|^2$.
For $\phi\in\mathcal W$, Step~2 gives
\begin{equation*}
\langle\phi^\sharp,\Lambda\phi^\sharp\rangle\le\mathsf e\,\|\phi^\sharp\|^2,\qquad\mathsf e:=\frac{4}{c_E}\bigl(\sqrt{\lambda'}+\mathbf m'\bigr).
\end{equation*}
Consequently the projection $\mathsf E:=1[\Lambda\le4\mathsf e]$ satisfies
\begin{equation*}
\|\mathsf E\phi\|\ge\|\mathsf E\phi^\sharp\|-\|\mathsf A\phi\|\ge\Bigl(\frac{\sqrt3}{2}\cdot\frac78-\frac18\Bigr)\|\phi\|\ge\frac12\,\|\phi\|.
\end{equation*}
It follows that
\begin{equation*}
\frac14\dim\mathcal W\le\operatorname{Tr}\bigl(1_B\mathsf E1_B\bigr)=192\,|B|\,\mathsf n^{-3}\,\#\{p:\Lambda(p)\le4\mathsf e\}\le192\,|B|\,(3\mathsf e)^3,
\end{equation*}
where the last inequality rests on $|p|\le\frac\pi2\mathsf s$ and on $\mathsf n\mathsf e\ge40\ell\mathbf m'/c_E\ge40$, with $c_E=(\sqrt2+7C_E)^{-1}\le1/\sqrt2$.
The assertion results on choosing for $\mathcal W$ the span of the eigenvectors whose eigenvalues do not exceed $\lambda'$.

\emph{Step 4} (condition (G) for quadratic forms).
\textbf{(G1)} Let $\beta_i\in C^2(\Gamma)$ be finitely many real cochains, let $w_i\ge0$ and $Q:=\sum_iw_i\langle\beta_i,F\rangle^2$, let $\mathsf t>0$ be such that $\mathsf t\ge\sum_iw_i\|\beta_i\|^2$, and let $\mathsf p\ge1$.
Then $E_n\exp(Q/4\mathsf t)\le\exp\frac12$, whence
\begin{equation*}
P_n(Q\ge4\mathsf t\mathsf r)\le\exp\bigl(\tfrac12-\mathsf r\bigr),\qquad\|Q\|_{L^{\mathsf p}}\le7\mathsf p\mathsf t.
\end{equation*}
Indeed, let $\mathsf g_i$ be the components of the Gaussian vector of (K7).
Applying (K7), Tonelli's theorem and \cite[Theorem~6.27\,(a)]{QEDII} at $\beta=\sqrt{2\theta}\sum_i\sqrt{w_i}\,\mathsf g_i\beta_i$, one finds
\begin{equation*}
E_n\exp\theta Q\le\det(1-2\theta G)^{-1/2},\qquad G_{ii'}=\sqrt{w_iw_{i'}}\,\langle\beta_i,\beta_{i'}\rangle,\qquad\operatorname{tr}G\le\mathsf t,
\end{equation*}
and the inequality $-\log(1-\xi)\le\frac{\xi}{1-\xi}$, used with $\theta=\frac1{4\mathsf t}$, gives the assertion.

\textbf{(G2)} For real numbers $w_i$,
\begin{equation*}
\Bigl\|\sum_iw_i\langle\beta_i,F\rangle\langle\beta_i',F\rangle\Bigr\|_{L^{\mathsf p}}\le7\mathsf p\sum_i|w_i|\,\|\beta_i\|\,\|\beta_i'\|.
\end{equation*}
The terms with $\|\beta_i\|\,\|\beta_i'\|=0$ vanish.
For the other terms, $|\xi\zeta|\le\frac12(\vartheta_i\xi^2+\vartheta_i^{-1}\zeta^2)$ with $\vartheta_i:=\|\beta_i'\|/\|\beta_i\|$ produces two positive forms, each of which admits $\mathsf t=\sum_i|w_i|\,\|\beta_i\|\,\|\beta_i'\|$;
by (G1), the norm is at most $\frac12(7\mathsf p\mathsf t+7\mathsf p\mathsf t)$.
Complex cochains are treated with the bilinear pairing and with $\|\beta\|^2=\allowbreak \|\operatorname{Re}\beta\|^2+\|\operatorname{Im}\beta\|^2$.
The product of two complex linear forms is then a combination of four real products, and
\begin{equation*}
\bigl(\|\operatorname{Re}\beta\|+\|\operatorname{Im}\beta\|\bigr)\bigl(\|\operatorname{Re}\beta'\|+\|\operatorname{Im}\beta'\|\bigr)\le2\,\|\beta\|\,\|\beta'\|;
\end{equation*}
hence (G2) is valid for complex cochains with $14\mathsf p$ in place of $7\mathsf p$.

\emph{Step 5} (functionals of $A^{\mathrm C}$).
To a function $\mathsf k$ on $\mathbb Z^3/n\mathbb Z^3$ associate
\begin{equation*}
A^{\mathrm C}_\nu(\mathsf k):=\sum_y\mathsf k(y)\,A^{\mathrm C}_\nu(y)=\langle\beta_{\mathsf k},F\rangle,\qquad\beta_{\mathsf k}=\mathcal Gd(\mathsf ke_\nu).
\end{equation*}
Since $d^*d\le\Delta$ by (K4), one has the bound \textbf{(B0)}
\begin{equation}
\tag{B0}
\|\beta_{\mathsf k}\|^2\le\langle\mathsf k,\mathcal G\mathsf k\rangle=n^{-3}\sum_{p\ne0}\frac{|\hat{\mathsf k}(p)|^2}{\mathsf q(p)}.
\end{equation}
The bound on the heat kernel in Step~6 of the proof of Proposition~5.11 gives \textbf{(B1)}
\begin{equation}
\tag{B1}
\langle\mathsf k,\mathcal G\mathsf k\rangle\le\int_0^\infty\min\bigl\{\|\mathsf k\|_2^2,\ c_h\|\mathsf k\|_1^2\,t^{-3/2}\bigr\}\,\mathrm dt=3\,c_h^{2/3}\,\|\mathsf k\|_1^{4/3}\,\|\mathsf k\|_2^{2/3}.
\end{equation}
In particular $\|\mathsf k\|_1\le c$ and $\|\mathsf k\|_2^2\le c\kappa^3$ imply $\|\beta_{\mathsf k}\|^2\le c'\kappa$.

\textbf{(Per)} For $\mathsf k\in\ell^1(\mathbb Z^3)$ let
\begin{equation*}
A^{\mathrm C}_\nu(\mathsf k):=\sum_{y\in\mathbb Z^3}\mathsf k(y)\,A^{\mathrm C}_\nu(y)=A^{\mathrm C}_\nu(\operatorname{per}_n\mathsf k),\qquad\operatorname{per}_n\mathsf k:=\sum_{\mathsf m\in\mathbb Z^3}\mathsf k(\cdot+n\mathsf m).
\end{equation*}
Then $\|\operatorname{per}_n\mathsf k\|_1\le\|\mathsf k\|_1$.
Assume further that $|\mathsf k(y')|\le\mathsf k_0(1+\kappa|y'-y|_\infty)^{-8}$, with $n\kappa\ge1$.
On the fundamental domain $y'-y\in(-\frac n2,\frac n2]^3$ the image of index $\mathsf m\ne0$ does not exceed $\mathsf k_0(\frac12\kappa n|\mathsf m|_\infty)^{-8}$, and Minkowski's inequality over the images yields
\begin{equation*}
\|\operatorname{per}_n\mathsf k\|_2\le\|\mathsf k\|_{\ell^2(\mathbb Z^3)}+C\mathsf k_0\,\kappa^{-3/2}(\kappa n)^{-13/2}.
\end{equation*}

\emph{Step 6} (frequencies add).
On $(\mathbb R/2\pi)^3$ one has $\mathsf s(p+p')\le\mathsf s(p)+\mathsf s(p')$, by $|\sin\frac{a+b}{2}|\le|\sin\frac a2|+|\sin\frac b2|$ and Minkowski's inequality.
Therefore, if $\hat g$ is carried by $\mathsf s\le a$, the component of $ug$ with $b_1<\mathsf s\le b_2$ involves only those modes of the multiplier $u$ that have $b_1-a<\mathsf s\le b_2+a$;
this is exact, and no aliasing occurs.
The mode of $u$ which takes $p$ to $P$ is $p'=P-p$, and $\mathsf s$ is even;
hence
\begin{equation*}
|\mathsf s(p)-\mathsf s(P)|\le\mathsf s(p')\le\mathsf s(p)+\mathsf s(P).
\end{equation*}
For an input with $a_1\le\mathsf s\le a_2$ and an output with $\mathsf s\le b<a_1$ it follows that only the modes with $a_1-b\le\mathsf s\le a_2+b$ are selected.
In the left factor of Step~16 the inner momenta form the input and the outer momenta the output, and the selected modes are those with $\frac14\kappa_k\le\mathsf s\le\frac74\kappa_k$, on which $\tilde\varphi_k=1$.

\emph{Step 7.}
\textbf{(Ker)} Let $0<\kappa\le64$, and let $\mathsf f$ be a smooth, periodic, matrix valued function carried by $\mathsf s\le6\kappa$ such that $|\partial^\gamma\mathsf f|\le\mathsf f_0\kappa^{-|\gamma|}$ for $|\gamma|\le8$.
The kernel $K$ of $\mathsf f$ obeys
\begin{equation*}
|K(z)|\le C\mathsf f_0\,\kappa^3\,(1+\kappa|z|_\infty)^{-8},
\end{equation*}
as one sees from the volume of the support and from eight integrations by parts;
the same bound holds after periodisation with a period $\ge1/\kappa$, $|z|_\infty$ being then the distance on the torus.
At $\kappa=\kappa_k$ this applies to the filters $\varphi_k$, $\Phi_k$ and $\tilde\varphi_k$ with $\mathsf f_0=C$, because $|\partial_\mu\mathsf q|\le2\mathsf s$, the higher derivatives of $\mathsf q$ are at most $2$, and $\kappa\le64$;
to $\tilde\varphi_k\,\mathsf q^{-1}(1-\exp(\mp ip_\mu))$ and to $\varphi_k\mathrm{Rsy}$ with $\mathsf f_0=C\kappa_k^{-1}$;
to the product of these two with $\mathsf f_0=C\kappa_k^{-2}$;
and to $\Phi_k(1-\exp(\mp ip_\mu))$ with $\mathsf f_0=C\kappa_k$, since $|1-\exp(\mp ip_\mu)|\le\mathsf s$ and the derivatives of $1-\exp(\mp ip_\mu)$ have modulus at most $1$.
For $\mathrm{Rsy}$ the required bounds are those of \cite[Proposition~37.1\,(g)]{QEDIII}, which applies because $\operatorname{dist}_2(p,\allowbreak 2\pi\mathbb Z^3)\ge\mathsf s\ge\frac12\kappa_k$ on the support of $\varphi_k$.

\emph{Step 8.}
\textbf{(BL)} Let $0\le j\le\mathsf J+1$ and $L_j:=\lceil1/\kappa_j\rceil$, and denote by $\mathcal C_j$ a partition of $\mathbb T'$ of product type, into boxes all of whose sides lie in $[L_j,2L_j)$.
Such a partition exists, since $L_j\le\lceil1/\mathbf m'\rceil\le\ell<\mathsf n$, and $\#\mathcal C_j\le(\mathsf n/L_j)^3\le(\mathsf n\kappa_j)^3$.
For a box $c$ put $\langle W\rangle_c:=|c|^{-1}\sum_{y\in c}W(y)$.
If $W\ge0$ and $\hat g$ is carried by $\mathsf s\le\kappa_j$, then
\begin{equation*}
\langle g,Wg\rangle\le C_{\mathrm{BL}}\,\max_{c\in\mathcal C_j}\langle W\rangle_c\,\|g\|^2.
\end{equation*}
In fact $g=\Phi_j*g$, so that the Schwarz inequality gives $|g(x)|^2\le\|\Phi_j\|_1\sum_y|\Phi_j(x-y)|\,|g(y)|^2$;
and (Ker), applied with the period $\mathsf n\ge10/\kappa_j$, gives $\|\Phi_j\|_1\le C$ and $\sum_c|c|\max_{x\in c}|\Phi_j(y-x)|\le C$.

\emph{Step 9} (the potential).
Let $e$ range over $\{\pm e_\nu\}$, and write $(\tau_ef)(k):=f(k+e)$ and $\mathrm{Ehp}_{\pm e_\nu}:=\mathrm{Ehp}^\pm_\nu$.
The operator $V:=D_{A''}-D_0$ takes the form
\begin{equation*}
V=-\sum_eu_e\,\tau_e\otimes\mathrm{Ehp}_e,\qquad u_e(k):=\exp\bigl(iqA''(k,k+e)\bigr)-1
\end{equation*}
\cite[Definition~2.10]{QEDII};
the hop points in the direction $\vec ge_1$, because $x^{(3)}-x=(1-2v_1)\vec ge_1$ and $0<v_1<\frac12$.
Further $\mathrm{Ehp}_e=1_e\mathrm{Ehp}_e1_{-e}$, where $1_e$ is the projection onto $\{\vec ge_1=e\}$, and $\|\mathrm{Ehp}_e\|=\mathsf h_3$.
Consequently every family of multipliers $v_e$ satisfies
\begin{equation*}
\Bigl\|\sum_ev_e\,\tau_e\otimes\mathrm{Ehp}_e\,g\Bigr\|^2\le\mathsf h_3^2\sum_e\|v_e\,\tau_e1_{-e}g\|^2.
\end{equation*}
Extend $A''$ by zero.
As $D_0'-D_0=\mathbf m'-\mathbf m$, the operator $D$ of Step~2 is the chamber operator on $\mathcal H'$ with the mass $\mathbf m$ and the field $A''$.
For $\phi$ over $B$ the value $(D\phi)(k,\vec g)$ involves $\phi$ only at $(k,\vec g)$, at $(k,\vec gs_i)$ with $i\le2$, and at $(k+\vec ge_1,\vec g\operatorname{diag}(-1,1,1))$;
thus $D\phi$ lives over $B^+$.
Since $7\ell<10\ell$, the quotient map $\mathbb Z^3\to\mathbb T'$ is injective on $B^{++}$.
Hence $\|D\phi\|_{\mathcal H'}=\|D_{A''}\phi\|$, and by Step~2 of the proof of Proposition~5.11 this equals $\|D_A\exp(-iq\,\lambda\circ\pi)\phi\|$.
The phase factor is unitary on the spinors over $B$, so that the forms $\|D\phi\|^2$ and $\|D_A\phi\|^2$ over $B$ have the same eigenvalues.

\emph{Step 10} (the splitting).
Let $\mathsf u_\nu:=\zeta A^{\mathrm C}_\nu$, and let $V_1$ be the operator obtained by putting
\begin{equation*}
u^{(1)}_{e_\nu}:=iq\,\mathsf u_\nu,\qquad u^{(1)}_{-e_\nu}:=-iq\,\tau_{-e_\nu}\mathsf u_\nu
\end{equation*}
in the places of the $u_e$;
the adjoint $V_1^*$ has the same multipliers, with $\mathrm{Ehp}_{-e}^*$ in place of $\mathrm{Ehp}_e$.
Set $V_{\mathrm r}:=V-V_1$ and
\begin{equation*}
V^{\mathrm{hl}}:=\sum_{i\ge j+3}\Pi_iV_1\Pi_j,\qquad V^{\mathrm{lh}}:=\sum_{j\ge i+3}\Pi_iV_1\Pi_j,\qquad V^{\mathrm d}:=\sum_{|i-j|\le2}\Pi_iV_1\Pi_j+V_{\mathrm r}.
\end{equation*}
With this splitting, $\mathcal T$ is exactly quadratic in $F$.

\emph{Step 11} (the linear event).
For $r\ge1$ put $\mathsf r_{kj}:=r+3j+2\log(2+k-j)$, and let $\mathcal G^{\mathrm{lin}}(r)$ be the event on which the following inequalities hold for all $\mu$, $\nu$, all $0\le j\le\mathsf J+1$ and all $c\in\mathcal C_j$:
\begin{itemize}
\item[(a)] $\langle|\tilde\varphi_k*\mathsf u_\nu|^2\rangle_c\le c_1\kappa_k\mathsf r_{kj}$ for $j\le k\le\mathsf J+1$;
\item[(b)] $\langle|\Phi_{j+3}*\mathsf u_\nu|^2\rangle_c\le c_1\kappa_j\mathsf r_{jj}$;
\item[(c)] $\langle|\mathsf u_\nu|^2\rangle_c\le c_1\mathsf r_{jj}$;
\item[(d)] $|(F_{\mathrm h})_{\mu\nu}|\le(2r)^{1/2}n^{-3/2}$.
\end{itemize}
In (a), (b) and (c) each value of $\tilde\varphi_k*\mathsf u_\nu$, of $\Phi_{j+3}*\mathsf u_\nu$ and of $\mathsf u_\nu$ is of the form $A^{\mathrm C}_\nu(\mathsf k)$, with $\mathsf k$ a kernel multiplied by $\zeta$.
By (Ker) and (B1) the three quadratic forms admit $\mathsf t\le c\kappa_k$, $\mathsf t\le c\kappa_j$ and $\mathsf t\le c$ respectively, $c$ being the largest of the three constants;
in (b) this constant contains the ratio $\kappa_{j+3}/\kappa_j=8$, and (Ker) is applicable there because $\kappa_{j+3}\le16\sqrt{13}<64$.
Put $c_1:=4c$;
the thresholds in (a), (b) and (c) are then the numbers $4\mathsf t\mathsf r$ of (G1).
Combining (G1), \cite[Theorem~6.27\,(c)]{QEDII} for a single $\beta$, and
\begin{equation*}
\sum_j(10\ell\kappa_j)^3\exp(-3j)\sum_k(2+k-j)^{-2}\le C(\ell\mathbf m')^3,
\end{equation*}
one obtains
\begin{equation*}
P_n\bigl(\mathcal G^{\mathrm{lin}}(r)^c\bigr)\le C(\ell\mathbf m')^3\exp(-r).
\end{equation*}

\emph{Step 12} (blocks on the event $\mathcal G^{\mathrm{lin}}(r)$).
\begin{itemize}
\item[(i)] Let $\hat g$ be carried by $\mathsf s\le\kappa_j$, and let $j\le k$.
By (BL),
\begin{equation*}
\|(\tilde\varphi_k*\mathsf u_\nu)g\|^2\le C_{\mathrm{BL}}\,c_1\kappa_k\mathsf r_{kj}\,\|g\|^2.
\end{equation*}
This extends to translates of $\mathsf u_\nu$.
Indeed, for $w\in\mathbb Z^3$ and a multiplier $U$ one has $\|(\tau_wU)g\|=\|U\tau_{-w}g\|$, and $\tau_{-w}g$ has the same norm and the same Fourier support as $g$;
consequently the inequality of (i) is valid with $\tilde\varphi_k*\mathsf u_\nu$ replaced by any of its translates.
\item[(ii)] Let $i\ge j+3$.
The input of $\Pi_iV_1\Pi_j$ has $\mathsf s\le\frac14\kappa_i$, and its output has $\frac{\sqrt3}{2}\kappa_i<\mathsf s\le2\kappa_i$.
According to Step~6, only the modes of $\mathsf u_\nu$ with $0.6\,\kappa_i<\mathsf s\le2.25\,\kappa_i$ contribute, and on them $\tilde\varphi_i=1$.
Step~9 therefore gives
\begin{equation*}
\|\Pi_iV_1\Pi_j\|,\ \|\Pi_iV_1^*\Pi_j\|\le c_2\,q\,(\kappa_i\mathsf r_{i,j+1})^{1/2},\qquad c_2:=\mathsf h_3\,(2C_{\mathrm{BL}}c_1)^{1/2}.
\end{equation*}
\item[(iii)] Let $|i-j|\le2$.
Only the modes with $\mathsf s\le\kappa_{j+4}$ contribute, and (b) at the index $j+1$ gives
\begin{equation*}
\|\Pi_iV_1\Pi_j\|\le c_2\,q\,(\kappa_j\mathsf r_{j+1,j+1})^{1/2};
\end{equation*}
the factor $2$ in $c_2$ comes from $\kappa_{j+1}=2\kappa_j$.
\end{itemize}

\emph{Step 13} (the condition (C1)).
For an operator $X$ one has $\|X\|^2\le\sum_{i,j}\|\Pi_iX\Pi_j\|^2$, and $\|S\Pi_i\|\le(c_E\kappa_i)^{-1}$.
Let $0\le\sigma<\frac12$.
From (ii),
\begin{align*}
\|\Lambda^\sigma\mathsf A\Lambda^{-\sigma}\|^2&\le\sum_{i\ge j+3}\kappa_{i+1}^{2\sigma}\,(c_E\kappa_i)^{-2}\,c_2^2q^2\,\kappa_i\mathsf r_{i,j+1}\,\kappa_j^{-2\sigma}\\
&\le C_\sigma q^2\sum_j\kappa_j^{-1}(r+3j+3)\le C_\sigma'\,q^2(r+1)/\mathbf m'.
\end{align*}
Since $\sigma<\frac12$, the summation over $i$ is that of a geometric series;
in $\mathsf r_{i,j+1}=r+3j+3+2\log(1+i-j)$ the term $2\log(1+i-j)$ is absorbed into $C_\sigma$, as $r+3j+3\ge4$.
The operator $\mathsf B^*=S^*(V^{\mathrm{lh}})^*$ satisfies the same bound.

\emph{Step 14} (the operator $V^{\mathrm d}$).
Henceforth $s=\frac38$, $s'=\frac5{16}$ and $\mathsf p=16$.
By Step~12\,(iii), the near-diagonal part of $V^{\mathrm d}$ contributes to $\langle f,V^{\mathrm d}g\rangle$ a term of modulus at most
\begin{equation*}
Cq\,\sup_j\kappa_j^{-1/4}\mathsf r_{j+1,j+1}^{1/2}\,\|f\|_s\|g\|_s\le Cq(r+1)^{1/2}\mathbf m'^{-1/4}\|f\|_s\|g\|_s.
\end{equation*}
Moreover
\begin{equation*}
|u_e-u_e^{(1)}|\le q^2|\zeta A^{\mathrm C}|^2+q^2|\zeta A^{\mathrm h}|^2+q|\zeta A^{\mathrm h}|,
\end{equation*}
where for $e=-e_\nu$ the right side is to be replaced by its translate by $-e_\nu$; the argument of Step~12\,(i) applies to such translates.
From (BL) and Step~11\,(c) one obtains
\begin{equation*}
\bigl\||\mathsf u_\nu|f\bigr\|\le\sum_jC\mathsf r_{j+1,j+1}^{1/2}\kappa_j^{-s}\|\Pi_jf\|_s\le C(r+1)^{1/2}\mathbf m'^{-s}\|f\|_s.
\end{equation*}
Step~11\,(d), together with $n>10\ell$ and $\ell\mathbf m'\ge1$, gives $|\zeta A^{\mathrm h}|\le(r\mathbf m')^{1/2}$, and therefore $q^2|\zeta A^{\mathrm h}|^2\le q^2r$, because $\mathbf m'\le1$.
The summation over the six directions $e$ is controlled by
\begin{equation*}
\sum_e\|1_ef\|_s\|1_{-e}g\|_s\le\|f\|_s\|g\|_s,
\end{equation*}
a consequence of Step~9 and the Schwarz inequality.
It follows that there is a number $c_6>0$ with the following property: on $\mathcal G^{\mathrm{lin}}(r)$, whenever $q^2(r+1)\le c_6\mathbf m'$, condition (C1) is satisfied and
\begin{equation*}
|\langle f,V^{\mathrm d}g\rangle|\le\frac{c_E}{8}\,\mathbf m'^{1-2s}\|f\|_s\|g\|_s.
\end{equation*}

\emph{Step 15} (blocks of $\mathcal T$).
Let $m:=\max(i,j)+3$.
Then
\begin{equation*}
\Pi_i\mathcal T\Pi_j=\Pi_iV_1S\,1[\Lambda>\kappa_m]\,V_1\Pi_j,
\end{equation*}
and this block is zero unless $m\le\mathsf J$.
One has $1[\Lambda>\kappa_m]=\varpi_m-\varpi_m\Pi_{m-1}$, for $\Lambda>\kappa_m$ implies $\mathsf q\ge\frac9{16}\kappa_m^2$, while $\varpi_m\ne0$ implies $\Lambda>\kappa_{m-1}$.
By Step~12\,(ii), (iii), the second term contributes at most $Cq^2\mathsf r_{mm}$ to $\|\Pi_i\mathcal T\Pi_j\|$.
Indeed, the input index of the left factor $\Pi_iV_1\Pi_{m-1}$ is $m-1$.
If $i=m-3$, the bound (iii) applies, Step~11\,(b) being used at the index $m\le\mathsf J+1$, and this yields $\mathsf r_{mm}$; if $i<m-3$, the left factor is the adjoint of $\Pi_{m-1}V_1^*\Pi_i$, a block covered by (ii).
Finally, $\varpi_m=\sum_{k=m}^{\mathsf J+1}\varphi_k$.

\emph{Step 16} (filtered fields).
Let $\tilde S^{(k)}:=\varphi_kS$; thus
\begin{equation*}
(\tilde S^{(k)}\psi)(y)=\sum_z\tilde S^{(k)}(z)\,\psi(y+z),
\end{equation*}
where $z$ runs through $\mathbb Z^3$ and $\tilde S^{(k)}(z)$ is the kernel on $\mathbb Z^3$ of (K8), the same kernel as in the definition of $\mathsf b_{e'}$ in Step~20.
According to (Ker), $|\tilde S^{(k)}(z)|\le C\kappa_k^2(1+\kappa_k|z|_\infty)^{-8}$.
The inner momenta satisfy $\frac12\kappa_k\le\mathsf s\le\frac32\kappa_k$, and the outer ones $\mathsf s\le\frac14\kappa_k$.
By Step~6, the right factor involves only those modes of its multiplier for which $\frac14\kappa_k\le\mathsf s\le\frac74\kappa_k$.
The left factor, whose input is the inner momentum $p$ and whose output is the outer momentum $P$, has the same property: a mode $p'$ carrying $p$ to $P$ equals $P-p$, so that $\mathsf s(p)-\mathsf s(P)\le\mathsf s(p')\le\mathsf s(p)+\mathsf s(P)$.
On these modes $\tilde\varphi_k=1$.
Consequently $\Pi_iV_1\tilde S^{(k)}V_1\Pi_j=\Pi_i\mathcal X_k\Pi_j$, where $\mathcal X_k$ is built from the multipliers $u_{e[k]}$, in which $\tilde\varphi_k*\mathsf u_\nu$ stands in place of $\mathsf u_\nu$:
\begin{equation*}
(\mathcal X_kg)(y)=\sum_{e,e',z}u_{e[k]}(y)\,u_{e'[k]}(y+e+z)\,\mathrm{Ehp}_e\tilde S^{(k)}(z)\mathrm{Ehp}_{e'}\,g(y+e+z+e').
\end{equation*}

\emph{Step 17} ($\mathcal T$ is a multiplication operator modulo convergent errors).
Denote by $\Xi_k(y)\in\operatorname{End}\mathcal V$ the sum of Step~16 with the factor $g(y+e+z+e')$ deleted:
\begin{equation*}
\Xi_k(y):=\sum_{e,e',z}u_{e[k]}(y)\,u_{e'[k]}(y+e+z)\,\mathrm{Ehp}_e\tilde S^{(k)}(z)\mathrm{Ehp}_{e'}.
\end{equation*}
On $\operatorname{ran}\Pi_j$ one has $\|(\tau_v-1)g\|\le|v|_2\kappa_{j+1}\|g\|$.
Together with the Schwarz inequality and Step~12\,(i), this gives
\begin{equation*}
\|\Pi_i(\mathcal X_k-\Xi_k)\Pi_j\|\le Cq^2\kappa_k\mathsf r_{km}\kappa_{j+1}\sum_z|\tilde S^{(k)}(z)|\,(|z|_2+2)\le Cq^2\mathsf r_{km}\,\frac{\kappa_j}{\kappa_k},
\end{equation*}
because $\sum_z|\tilde S^{(k)}(z)|(|z|_2+2)\le C(\kappa_k^{-2}+\kappa_k^{-1})\le C'\kappa_k^{-2}$, by $\kappa_k\le2\sqrt{13}$.
Summed over $k\ge m$, these bounds amount to at most $Cq^2\mathsf r_{mm}$.
Here Step~12\,(i) bounds $\|\overline{u_{e[k]}}f\|$, where $f\in\operatorname{ran}\Pi_i$ has $\mathsf s\le\kappa_{i+1}$, and, for every $z$, $\|(\tau_{e+z}u_{e'[k]})(\tau_v-1)g\|$ with $v=e+z+e'$, whose multiplier is a translate.
As $i+1,j+1\le m\le k$ and $\mathsf r_{kj}$ is an increasing function of $j$, both bounds carry the factor $\mathsf r_{km}$.
By Step~6, the block $\Pi_i\Xi_k\Pi_j$ involves only the modes of $\Xi_k$ with $\mathsf s\le\kappa_{m-1}$.
Put
\begin{equation*}
Y_m:=\Phi_m*\sum_{k=m}^{\mathsf J+1}\Xi_k,\qquad\|Y_m\|_\infty:=\sup_{x\in\mathbb T'}\|Y_m(x)\|,
\end{equation*}
the norm of $Y_m(x)$ being the operator norm of $\operatorname{End}\mathcal V$.
Then
\begin{equation*}
\|\Pi_i\mathcal T\Pi_j\|\le Cq^2\mathsf r_{mm}+\|Y_m\|_\infty.
\end{equation*}

\emph{Step 18} (the quadratic event).
Let $\epsilon_3>0$ be a number depending on $\rho$ alone, whose value is chosen in Step~27.
This positive number is distinct from the component $\varepsilon_3$ of the sign vectors $\varepsilon\in\{\pm1\}^3$ of Step~21.
Define
\begin{equation*}
\mathcal G^{\mathrm{qu}}:=\bigl\{\|Y_m\|_\infty\le\epsilon_3\mathbf m'2^{s'm}\ \text{for}\ 3\le m\le\mathsf J\bigr\}.
\end{equation*}
On $\mathcal G^{\mathrm{lin}}(r)\cap\mathcal G^{\mathrm{qu}}$ one has
\begin{equation*}
|\langle f,\mathcal Tg\rangle|\le C\bigl[q^2(r+1)+\epsilon_3\mathbf m'\bigr]\,\mathbf m'^{-2s}\|f\|_s\|g\|_s,
\end{equation*}
since $\sum_{i,j}4^{-s(i+j)}(1+m)^24^{s'm}<\infty$;
the number $C$ in this bound does not depend on $\epsilon_3$.
If $q^2(r+1)=c_*\mathbf m'$ with $0<c_*\le c_6$, then on $\mathcal G^{\mathrm{lin}}(r)\cap\mathcal G^{\mathrm{qu}}$ Step~14 yields (C1), and Step~14 together with this bound yields
\begin{equation*}
|\langle f,(V^{\mathrm d}-\mathcal T)g\rangle|\le\Bigl[\frac{c_E}{8}+C(c_*+\epsilon_3)\Bigr]\,\mathbf m'^{1-2s}\|f\|_s\|g\|_s,
\end{equation*}
with the same number $C$.
Consequently (C2) holds on this event whenever $C(c_*+\epsilon_3)\le\frac{c_E}{8}$, because $\frac{c_E}{8}+\frac{c_E}{8}=\frac{c_E}{4}$.
As $Y_m=\Phi_{m+1}*Y_m$, H\"older's inequality and (Ker) yield $\|Y_m\|_\infty\le C\kappa_m^{3/\mathsf p}\|Y_m\|_{\ell^{\mathsf p}(\mathbb T')}$, where $\|Y_m\|_{\ell^{\mathsf p}(\mathbb T')}$ is the $\ell^{\mathsf p}$ norm of the function $x\mapsto\|Y_m(x)\|$.
Suppose that condition \textbf{(H)} is satisfied:
\begin{equation*}
\bigl\|\,\|Y_m(x)\|\,\bigr\|_{L^{\mathsf p}}\le\mathsf C_{\mathsf p}q^2\qquad\text{for all }x\in\mathbb T'\text{ and }3\le m\le\mathsf J.
\end{equation*}
Markov's inequality then gives
\begin{equation*}
P_n\bigl((\mathcal G^{\mathrm{qu}})^c\bigr)\le\sum_mC(\ell\mathbf m')^3\,2^{(3-\mathsf ps')m}\Bigl(\frac{\mathsf C_{\mathsf p}q^2}{\epsilon_3\mathbf m'}\Bigr)^{\mathsf p}\le C(\ell\mathbf m')^3\Bigl(\frac{\mathsf C_{\mathsf p}q^2}{\epsilon_3\mathbf m'}\Bigr)^{16},
\end{equation*}
as $\mathsf ps'=5$.
Step~27 establishes (H) with $\mathsf C_{\mathsf p}=C\mathsf p^2$, and $\epsilon_3$ depends on $\rho$ alone;
the factor $(\mathsf C_{\mathsf p}/\epsilon_3)^{16}$ of the last bound is therefore a number depending on $\rho$ alone, and it is absorbed into the constant $C_{11}$ of Step~27.

\emph{Step 19} (covariance).
Let $g\in W$.
Then $U_g=T_g\otimes\mathsf U_{\vec g}$, where $(T_gf)(k):=f(g^{-1}k)$ and where the unitary operator $\mathsf U_{\vec g}$ acts by $(\mathsf U_{\vec g}u)(\vec h):=\tilde g\,u(\vec g^{-1}\vec h)$; furthermore $U_gD_AU_g^{-1}=D_{g_*A}$ with $g_*:=(g^{-1})^*$, and $\tilde g\gamma_v\tilde g^{-1}=\gamma_{\vec gv}$ \cite[Proposition~3.16\,(b)]{QEDII}.
The matrix $\tilde g$ is formed with a word $\vec g=r_{n_1}\cdots r_{n_k}$.
For a different word it changes at most by a sign (an even word gives $\operatorname{diag}(X,X)$ with $X\in SU(2)$), and hence $\operatorname{Ad}\mathsf U_{\vec g}$ depends on $\vec g$ alone.
Since
\begin{equation*}
(\mathrm{Ehp}_eu)(\vec h)=1[\vec he_1=e]\,3(\rho-\gamma\cdot e)\,u(\vec hs_3),
\end{equation*}
with the linear map $s_3=\operatorname{diag}(-1,1,1)$ of \cite{QEDIII}, one finds $\mathsf U_{\vec g}\mathrm{Ehp}_e\mathsf U_{\vec g}^{-1}=\allowbreak \mathrm{Ehp}_{\vec ge}$.
At $A=0$ the relation $\mathsf U_{\vec g}D^{(p)}\mathsf U_{\vec g}^{-1}=D^{(\vec gp)}$ holds, whence $\mathsf U_{\vec g}\tilde S^{(k)}(z)\mathsf U_{\vec g}^{-1}=\allowbreak \tilde S^{(k)}(\vec gz)$.

\emph{Step 20} (the local functional).
Fix an integer $\mathsf R$ with $1\le\mathsf R\le\ell$; its value is chosen in Step~24 and depends on $(k,m)$.
Write $\xi_k$ for $\xi_k^{(\mathsf R)}$.
On $\mathbb Z^3$, with the periodic field, define the one-cochain $a^{[k]}:=\tilde\varphi_k*A^{\mathrm C}$ and
\begin{gather*}
\mathsf b_{e'}(x):=\sum_z\tilde S^{(k)}(z)\,a^{[k]}(x+z,x+z+e'),\\
\xi_k^\infty(y;F):=\sum_{e,e'}a^{[k]}(y,y+e)\,\mathrm{Ehp}_e\,\mathsf b_{e'}(y+e)\,\mathrm{Ehp}_{e'}.
\end{gather*}
Each of the two factors is a sum $\sum_\sigma\mathsf K(\sigma-y)F_\sigma$ over plaquettes, and $\xi_k$ is obtained from $\xi_k^\infty$ by restricting both sums to the plaquettes whose centres satisfy $|\sigma-y|_\infty\le\mathsf R$.
The operators $d^*$ and $\mathcal G$, as well as the $O_{\mathrm{cub}}$-invariant filters, commute with $g_*$ \cite[Definition~2.1]{QEDII}.
Therefore $a^{[k]}[g_*F](gx,gx+\vec ge)=a^{[k]}[F](x,x+e)$.
The centre of $g\sigma$ is the image under $g$ of the centre of $\sigma$, and $|\vec gv|_\infty=|v|_\infty$.
Accordingly, the identity is valid also when the sums are cut off at $|\sigma-y|_\infty\le\mathsf R$.
In each of the two factors this cut-off is centred at $y$, which gives the property (L) of Lemma~5.15 for this $\mathsf R$.
In detail, (K8) gives on $\mathbb Z^3$ the representation $a^{[k]}(b)=\langle\mathsf K_b,F\rangle$ for a bond $b$, with $\mathsf K_b:=d(\mathsf G_k*\delta_b)\in\ell^1$ and $\langle\delta_b,a\rangle:=a(b)$, where $\mathsf G_k$ denotes the kernel of the symbol $\tilde\varphi_k/\mathsf q$.
This symbol is $O_{\mathrm{cub}}$-invariant, and it is smooth because $\tilde\varphi_k$ vanishes near $p=0$.
Moreover $g_*\delta_b=\delta_{gb}$, and $g_*$ commutes with $d$ and with this convolution.
Hence $g_*\mathsf K_b=\mathsf K_{gb}$ and $g_*(\mathsf K_b1_{y,\mathsf R})=\mathsf K_{gb}1_{gy,\mathsf R}$, where $1_{y,\mathsf R}$ is the indicator of the plaquettes with $|\sigma-y|_\infty\le\mathsf R$.
Since $g_*$ is an orthogonal map, $\langle\mathsf K_{gb}1_{gy,\mathsf R},g_*F\rangle$ coincides with $\langle\mathsf K_b1_{y,\mathsf R},F\rangle$ for every bounded $F$.
Change the summation variables in the sums defining $\xi_k^\infty(gy;g_*F)$ and $\xi_k(gy;g_*F)$, which converge absolutely, to $\tilde e$, $\tilde z$, $\tilde e'$, where $e=\vec g\tilde e$, $z=\vec g\tilde z$, $e'=\vec g\tilde e'$, and apply both identities of Step~19.
The outcome is
\begin{equation*}
\xi_k(gy;g_*F)=\mathsf U_{\vec g}\,\xi_k(y;F)\,\mathsf U_{\vec g}^{-1},
\end{equation*}
together with the same identity for $\xi_k^\infty$.

\emph{Step 21} (sectors).
Let $r_\lambda$ be the reflection $x_\lambda\mapsto-x_\lambda$; then $\tilde r_\lambda=\gamma_\lambda\gamma_5$ (formed with the word $r_{e_\lambda}$), and one puts $\operatorname{Ad}_\lambda:=\allowbreak \operatorname{Ad}\mathsf U_{r_\lambda}$.
These maps are commuting isometric involutions of $\operatorname{End}\mathcal V$, because $(\gamma_\lambda\gamma_5)^2=-1$ and $\gamma_\lambda\gamma_5$ anticommutes with $\gamma_\mu\gamma_5$ for $\mu\ne\lambda$ \cite[Definition~2.7]{QEDII}.
For $\varepsilon\in\{\pm1\}^3$ set
\begin{equation*}
\mathsf P^\varepsilon:=\prod_\lambda\tfrac12\bigl(1+\varepsilon_\lambda\operatorname{Ad}_\lambda\bigr),
\end{equation*}
so that $\sum_\varepsilon\mathsf P^\varepsilon=1$.
For $\theta=\theta_{x_\lambda,c}$ with $c\in\frac12\mathbb Z$, Step~20 implies
\begin{equation*}
(\mathsf P^\varepsilon\xi_k)(\theta y;\theta^*F)=\varepsilon_\lambda\,(\mathsf P^\varepsilon\xi_k)(y;F).
\end{equation*}

\emph{Step 22} (sizes).
A matrix valued function $\mathsf z(y)$ is said to have \emph{size} $\mathsf N$ if every entry of $\mathsf z(y)$ can be written as a sum of no more than $C$ products $\langle\beta,F\rangle\langle\beta',F\rangle$, where $\beta,\beta'$ are complex cochains such that $\|\beta\|\|\beta'\|\le\mathsf N$.
Let $\mathsf p'\ge1$ be an arbitrary real number.
This exponent is independent of the number $\mathsf p=16$ fixed in Step~14, and its prime does not mark a conjugate exponent.
Let further $v$ be a real function of $y$ with finite support, as in condition (S) of Lemma~5.15 (a weight, not a vector), and put $\|v\|_1:=\sum_y|v(y)|$.
For a function $\mathsf z$ of size $\mathsf N$, the inequality (G2), with $\mathsf p'$ in the place of $\mathsf p$, gives
\begin{equation*}
\Bigl\|\sum_yv(y)\mathsf z(y)\Bigr\|_{L^{\mathsf p'}}\le C\mathsf p'\mathsf N\|v\|_1,
\end{equation*}
because the product of two complex linear forms is a linear combination of four products of real ones, built from the real and the imaginary parts of $\beta$ and $\beta'$.
In this bound the norm of $L^{\mathsf p'}$ is applied to each real component of the matrix on the left, that is, to the real part and to the imaginary part of each of its entries;
for $C$ one may then take the factor $14$ of (G2) for complex cochains times the number of products admitted in the definition of size.
With another $C$, the bound holds as well when the matrix is measured by its operator norm in $\operatorname{End}\mathcal V$.
Steps~26 and~27 use this bound at $\mathsf p'=2$, $2\mathsf p$ and $\mathsf p$.
The size is preserved by $\mathsf P^\varepsilon$.
Indeed, $\gamma_\lambda\gamma_5$ is a monomial matrix \cite[(I.c6)]{QEDIII}, so that $\mathsf U_{r_\lambda}$ is monomial too, and each entry of $\operatorname{Ad}_\lambda\mathsf z$ equals an entry of $\mathsf z$ times a unimodular number.
If the symbol of a kernel is supported in $\frac18\kappa_k\le\mathsf s\le6\kappa_k$ and has modulus at most $\mathsf f_0$, then (B0) gives $\|\beta\|^2\le C\mathsf f_0^2\kappa_k$.
For the factors of $\xi_k^\infty$ this yields $\|\beta\|\le C\kappa_k^{1/2}$ and, as $\tilde\varphi_k\varphi_k=\varphi_k$ and $|\mathrm{Rsy}|\le2/(c_E\kappa_k)$, also $\|\beta'\|\le C\kappa_k^{-1/2}$; so $\xi_k^\infty$ has size $C$ for every $k$.
Denote by $\mathsf K$ the kernel on $\mathbb Z^3$ of either factor, and by $\hat{\mathsf k}$ the symbol in front of $d^*\mathcal G$; it is supported in $\frac18\kappa_k\le\mathsf s\le6\kappa_k$ and obeys $|\hat{\mathsf k}|\le\mathsf f_0$.
The truncated factor has the cochain $\mathsf K(\cdot-y)1[|\cdot-y|_\infty\le\mathsf R]$.
Because $2\mathsf R+2\le2\ell+2<n$, no two of its plaquettes are congruent modulo $\Gamma$, so that its norm in $C^2(\Gamma)$ does not exceed $\|\mathsf K\|_{\ell^2(\mathbb Z^3)}$.
Plancherel's formula on $\mathbb Z^3$ (K8) and $d^*d\le\Delta$ give
\begin{equation*}
\|\mathsf K\|^2_{\ell^2(\mathbb Z^3)}\le(2\pi)^{-3}\int_{[-\pi,\pi]^3}|\hat{\mathsf k}|^2\,\mathsf q^{-1}\,\mathrm dp\le C\mathsf f_0^2\kappa_k.
\end{equation*}
This is (B0), with the integral replacing the Riemann sum; the factors of $\xi_k$ are therefore subject to the same two bounds.
For a lattice difference of a factor, the bound acquires an additional factor $C\kappa_k$.

\emph{Step 23} (cut-off and period).
One has
\begin{equation*}
(\tilde\varphi_k*\mathsf u_\nu)(y)-\zeta(y)\,a^{[k]}_\nu(y)=A^{\mathrm C}_\nu(\mathsf k_1+\mathsf k_2),
\end{equation*}
where
\begin{equation*}
\mathsf k_1:=\tilde\varphi_k(y-\cdot)\,\bigl(\zeta-\zeta(y)\bigr),\qquad\mathsf k_2:=\sum_{\mathsf m\ne0}\tilde\varphi_k(y+\mathsf n\mathsf m-\cdot)\,\zeta.
\end{equation*}
In this identity $y\in\mathbb T'$ is lifted to a point of $\mathbb Z^3$ with $|y-k_B|_\infty\le5\ell$, where $k_B$ is the centre of $B^{++}$ (Step~2 of the proof of Proposition~5.11).
On $B^{++}$ one has $|y'-k_B|_\infty\le\frac72\ell$, and so the arguments of $\tilde\varphi_k$ in $\mathsf k_2$ have $|\cdot|_\infty\ge10\ell-5\ell-\frac72\ell=\frac32\ell$.
The function $\mathsf k_1$ lives on $\mathbb Z^3$ and is read through (Per), with $\mathsf k_0=C\kappa_k^3$ for the first factor and $\mathsf k_0=C\kappa_k^2$ for the second, whose kernel is $\tilde S^{(k)}$ since $\tilde\varphi_k\varphi_k=\varphi_k$.
From $|\zeta(y')-\zeta(y)|\le\min\{1,|y-y'|_\infty/\ell\}$, $(\kappa_kn)^{-13/2}\le(\kappa_k\ell)^{-1}$ and (Ker) one obtains for the periodisation of $\mathsf k_1$
\begin{align*}
\|\cdot\|_1&\le C(\kappa_k\ell)^{-1},&\|\cdot\|_2^2&\le C\kappa_k^3(\kappa_k\ell)^{-2}&&\text{in the first factor},\\
\|\cdot\|_1&\le C\kappa_k^{-1}(\kappa_k\ell)^{-1},&\|\cdot\|_2^2&\le C\kappa_k(\kappa_k\ell)^{-2}&&\text{in the second factor};
\end{align*}
the function $\mathsf k_2$ has finite support, and the same bounds hold for it.
If $e=-e_\nu$, as well as in the second factor, the centre of the kernel is a point $y_0$ such that $|y_0-y|_\infty\le2$, namely $y-e_\nu$, respectively $y+e$ or $y+e-e_\nu$, while $\zeta$ is evaluated at $y$.
In these cases $|\zeta(y')-\zeta(y)|\le\min\{1,(|y'-y_0|_\infty+2)/\ell\}$, and the arguments in $\mathsf k_2$ are at least $\frac32\ell-2\ge\ell$; since $\kappa_k\le2\sqrt{13}$, only $C$ is altered.
Accordingly (B1) gives $\|\beta\|\le C\kappa_k^{1/2}(\kappa_k\ell)^{-1}$, and in the same manner the bound $C\kappa_k^{-1/2}(\kappa_k\ell)^{-1}$ for the second factor, which is formed with $\tilde S^{(k)}$.
Consequently
\begin{equation*}
\Xi_k=-q^2\bigl[\zeta^2\xi_k^\infty+\mathrm{err}_k\bigr],
\end{equation*}
where $\mathrm{err}_k$ has size $C(\kappa_k\ell)^{-1}$, and $\sum_k(\kappa_k\ell)^{-1}\le2$.

\emph{Step 24} (truncation; the leading part).
Fix $m\le k$, and let $\mathsf R:=\lceil\kappa_k^{-1}2^{(k-m)/2}\rceil$; then $\mathsf R\le\ell$.
A far part has as its cochain the $n$-periodisation of $\mathsf K(\cdot-y)1[|\cdot-y|_\infty>\mathsf R]$.
The bound (Ker) gives $|\mathsf K(z)|\le C\mathsf f_1\kappa_k^3(1+\kappa_k|z|_\infty)^{-8}$, where $\mathsf f_1=\kappa_k^{-1}$ for the first kernel and $\mathsf f_1=\kappa_k^{-2}$ for the second.
Applying Minkowski's inequality to the sum over the images $n\mathsf m$, $\mathsf m\in\mathbb Z^3$, one bounds the norm of this cochain in $C^2(\Gamma)$ by
\begin{multline*}
\|\mathsf K1_{>\mathsf R}\|_{\ell^2(\mathbb Z^3)}+\sum_{\mathsf m\ne0}Cn^{3/2}\mathsf f_1\kappa_k^3\bigl(\tfrac12\kappa_kn|\mathsf m|_\infty\bigr)^{-8}\\
\le C\mathsf f_1\kappa_k^{3/2}\bigl[(\kappa_k\mathsf R)^{-13/2}+(\kappa_kn)^{-13/2}\bigr]\le C\kappa_k^{\pm1/2}(\kappa_k\mathsf R)^{-6},
\end{multline*}
the last inequality resting on $n>10\mathsf R$ and $\kappa_k\mathsf R\ge1$.
The second kernel decays around $y+e$ rather than around $y$, and from $|\sigma-y|_\infty>\mathsf R$ one gets $|\sigma-y-e|_\infty>\mathsf R-1$.
Since $1+\kappa_k|z-e|_\infty\ge(1+\kappa_k)^{-1}(1+\kappa_k|z|_\infty)$ and $\kappa_k\le2\sqrt{13}$, this affects only $C$.
Likewise for the first kernel at $e=-e_\nu$, whose centre of decay is $y-e_\nu$, and for the second kernel at $e'=-e_\nu$, whose centre of decay is $y+e-e_\nu$: here $1+\kappa_k|z-v|_\infty\ge(1+2\kappa_k)^{-1}(1+\kappa_k|z|_\infty)$ for $|v|_\infty\le2$.
The far parts of the two kernels thus have $\|\beta\|\le C\kappa_k^{\pm1/2}(\kappa_k\mathsf R)^{-6}$, the bounds (G1) and (G2) being valid for every $\beta\in C^2(\Gamma)$, closed or not.
It follows that $\xi_k^\infty-\xi_k$ has size $C2^{-3(k-m)}$.
Let $\mathsf X_\nu:=\allowbreak \mathrm{Vx}_\nu(0)=\allowbreak -i(\mathrm{Ehp}_{e_\nu}-\mathrm{Ehp}_{-e_\nu})$ \cite[Proposition~37.1\,(d)]{QEDIII}.
Up to a lattice difference, $\sum_ea^{[k]}(y,y+e)\mathrm{Ehp}_e$ equals $i\sum_\nu a^{[k]}_\nu(y)\mathsf X_\nu$.
Hence $\xi_k^\infty=\mathsf L_k+\varrho_k$, with
\begin{equation*}
\mathsf L_k(y):=-\sum_{\mu,\nu,z}a^{[k]}_\mu(y)\,a^{[k]}_\nu(y+z)\,\mathsf X_\mu\tilde S^{(k)}(z)\mathsf X_\nu
\end{equation*}
and with $\varrho_k$ of size $C\kappa_k$.

\emph{Step 25} (the even sector).
Put $\operatorname{Ad}_\iota:=\allowbreak \operatorname{Ad}_1\operatorname{Ad}_2\operatorname{Ad}_3$.
Then $\mathsf P^{+++}=\allowbreak \mathsf P^{+++}\,\allowbreak \frac12(1+\operatorname{Ad}_\iota)$, and by Step~19 the map $\operatorname{Ad}_\iota$ carries $\mathsf X_\nu$ into $-\mathsf X_\nu$ and $\tilde S^{(k)}(z)$ into $\tilde S^{(k)}(-z)$.
Hence $\frac12(1+\operatorname{Ad}_\iota)\mathsf L_k$ is the expression $\mathsf L_k$ in which the kernel $\tilde S^{(k)}$ receives the symbol
\begin{equation*}
\varphi_k\,\tfrac12\bigl(\mathrm{Rsy}(p)+\mathrm{Rsy}(-p)\bigr)=\varphi_k\Bigl[J\,\frac{\mathbf m'}{\mathbf m'^2+|p|^2}\,J^*+\tfrac12\bigl(\mathrm{Er}(p)+\mathrm{Er}(-p)\bigr)\Bigr]
\end{equation*}
\cite[Proposition~37.1\,(f)]{QEDIII}, whose norm does not exceed $4\mathbf m'\kappa_k^{-2}+C_E$.
Being odd in $p$, the singular part $-iJ\frac{\gamma\cdot p}{\mathbf m'^2+|p|^2}J^*$ cancels.
It follows that $\mathsf P^{+++}\xi_k^\infty$ has size $C(\kappa_k+\mathbf m'/\kappa_k)$.
Since $\varphi_k=0$ for $k\ge\mathsf J+2$ and $\kappa_{\mathsf J+1}\le2\sqrt{13}$, one has
\begin{equation*}
\sum_{k=3}^{\mathsf J+1}\Bigl(\kappa_k+\frac{\mathbf m'}{\kappa_k}\Bigr)\le4\sqrt{13}+\frac14.
\end{equation*}

\emph{Step 26} (the seven odd sectors).
Let $\varepsilon_\lambda=-1$, and let $\mathsf o_y$ denote the real part or the imaginary part of an entry of $(\mathsf P^\varepsilon\xi_k)(y;\cdot)$.
By Steps~20 to~22, these functionals satisfy (L), (O) and (S) of Lemma~5.15 with $\mathsf Z=B^{++}$ and $\mathsf t=C$.
Here $m\le k$ are the integers of Step~24, with $3\le m\le\mathsf J$, as in (H), and $k\le\mathsf J+1$, as in the definition of $Y_m$ in Step~17.
Let $x\in\mathbb T'$ be the point that occurs in condition (H) of Step~18; the bounds of this step hold uniformly in $x$.
Denote by the same letter a representative of $x$ in $\mathbb Z^3$, chosen arbitrarily, so that the argument $x+\mathsf n\mathsf m-y$ of the kernel of $\Phi_m$ on $\mathbb Z^3$ is a point of $\mathbb Z^3$.
Another representative has the form $x+\mathsf n\mathsf m_0$ with $\mathsf m_0\in\mathbb Z^3$; the summand of index $\mathsf m$ formed with it equals the summand of index $\mathsf m+\mathsf m_0$ formed with $x$, so that the sums over $\mathsf m\in\mathbb Z^3$ in this step do not depend on the choice of the representative.
For $\mathsf m\in\mathbb Z^3$ set $w=w_{\mathsf m}:=\Phi_m(x+\mathsf n\mathsf m-\cdot)\,\zeta^2$, and let $x^*=x^*_{\mathsf m}$ be the point of the projection of $\mathsf Z$ that is nearest to $(x+\mathsf n\mathsf m)_\lambda$.
This point is unique, because $(x+\mathsf n\mathsf m)_\lambda$ is an integer, and it is an integer, as Lemma~5.15 requires, because $\mathsf Z\subset\mathbb Z^3$.
As the symbol of $\Phi_m$ is real and even, its kernel is real, and so is $\zeta$; moreover $\zeta$ vanishes off $B^{++}$, a box of $\mathbb Z^3$ on which the quotient map $\mathbb Z^3\to\mathbb T'$ is injective (Step~9), and the kernel of $\Phi_m$ on $\mathbb T'$ is the periodisation of its kernel on $\mathbb Z^3$ ((K8) and the paragraph on filters).
Hence $\sum_{\mathsf m}\sum_yw\,\mathsf o_y$ is the real part or the imaginary part of an entry of $\bigl[\Phi_m*(\zeta^2\mathsf P^\varepsilon\xi_k)\bigr](x)$, namely of the entry and the part that define $\mathsf o_y$.
Lemma~5.15 is applied to each $\mathsf m$ separately, and Minkowski's inequality over $\mathsf m$ combines the resulting bounds.
The cut-off $\zeta$ vanishes wherever $\operatorname{dist}_\infty(\cdot,B^+)\ge\ell$, that is, on the two outermost layers of $B^{++}$ and off $B^{++}$.
Extended by zero, $w$ therefore has no jump across the faces of $\mathsf Z$, and
\begin{equation*}
|\partial w|\le|\nabla_\lambda\Phi_m|+\frac2\ell\,|\Phi_m|,
\end{equation*}
the kernel of the second term being taken at the shifted point; as in Step~24, the shift affects only the constant $C$.
Apply (Ker) at $\kappa=\kappa_m$ to $\Phi_m$ and, with $\mathsf f_0=C\kappa_m$, to the symbol $\Phi_m(p)(1-\exp(\mp ip_\lambda))$ of $\nabla_\lambda\Phi_m$.
Together with $\ell^{-1}\le\mathbf m'\le\kappa_m$ this yields
\begin{equation*}
|\partial w(y)|\le C\kappa_m^4\bigl(1+\kappa_m|x+\mathsf n\mathsf m-y|_\infty\bigr)^{-8}.
\end{equation*}
Moreover $|y_\lambda-x^*|\le|y_\lambda-(x+\mathsf n\mathsf m)_\lambda|$, because $x^*$ is the point nearest to $(x+\mathsf n\mathsf m)_\lambda$ of an interval to which $y_\lambda$ belongs.
The sets $x+\mathsf n\mathsf m-\mathsf Z$, $\mathsf m\in\mathbb Z^3$, are mutually disjoint, as $7\ell<\mathsf n$.
Consequently
\begin{equation*}
\begin{split}
\sum_{\mathsf m}\sum_y\bigl(|y_\lambda-x^*|+1+a_0\bigr)^{1/2}|\partial w|&\le\sum_{z\in\mathbb Z^3}\bigl(|z|_\infty+1+a_0\bigr)^{1/2}C\kappa_m^4\bigl(1+\kappa_m|z|_\infty\bigr)^{-8}\\
&\le C\bigl(\kappa_m^{-1}+a_0\bigr)^{1/2}\kappa_m,
\end{split}
\end{equation*}
where, owing to $\kappa_k\le2\sqrt{13}$,
\begin{equation*}
a_0\kappa_m\le2\cdot2^{-(k-m)/2}+6\kappa_k2^{-(k-m)}\le45.3\cdot2^{-(k-m)/2}.
\end{equation*}
Lemma~5.15 then gives
\begin{equation*}
\Bigl\|\sum_{\mathsf m}\sum_yw\,\mathsf o_y\Bigr\|_{L^2}\le C\,2^{-(k-m)/4},
\end{equation*}
while (G2) bounds the norm of the same sum in $L^{2\mathsf p}$ by $C\mathsf p$.
Lyapunov's inequality (K9) shows that its norm in $L^{\mathsf p}$ is at most $C\mathsf p\,2^{-(k-m)/(4(\mathsf p-1))}$, and the sum over $k$ of these numbers does not exceed $C\mathsf p^2$.

\emph{Step 27} ((H) for one box and one level).
In $\xi_k$ take $\mathsf R=\mathsf R(k,m)$, the number of Step~24.
From Steps~17 and~23 and from $\sum_\varepsilon\mathsf P^\varepsilon=1$ one obtains
\begin{multline*}
Y_m=-q^2\sum_{k=m}^{\mathsf J+1}\Bigl[\Phi_m*\bigl(\zeta^2\mathsf P^{+++}\xi_k^\infty\bigr)+\sum_{\varepsilon\ne+++}\Phi_m*\bigl(\zeta^2\mathsf P^\varepsilon\xi_k\bigr)\\
+\sum_{\varepsilon\ne+++}\Phi_m*\bigl(\zeta^2\mathsf P^\varepsilon(\xi_k^\infty-\xi_k)\bigr)+\Phi_m*\mathrm{err}_k\Bigr];
\end{multline*}
here the even sector enters without truncation, and the odd sectors enter truncated.
For each real component of the four terms, Steps~25, 26, 24 and~23 give, in this order, the bounds
\begin{equation*}
C\mathsf p\Bigl(\kappa_k+\frac{\mathbf m'}{\kappa_k}\Bigr),\qquad C\mathsf p\,2^{-(k-m)/60},\qquad C\mathsf p\,2^{-3(k-m)},\qquad C\mathsf p\,(\kappa_k\ell)^{-1}
\end{equation*}
on the $L^{\mathsf p}$-norm.
Summed over $k$, these bounds are at most $C\mathsf p$ times $4\sqrt{13}+\frac14$ (because $m\ge3$), $(1-2^{-1/60})^{-1}<87.07$, $\frac87$ and $2$, respectively.
Minkowski's inequality over $k$, over the sectors and over the $2\cdot192^2$ components, combined with Steps~23 to~26 and with $\|\Phi_m\|_1\le C$, consequently yields (H), the constant being $\mathsf C_{\mathsf p}=C\mathsf p^2$.
Fix $\epsilon_3>0$ and $c_*\in(0,c_6]$, sufficiently small for Steps~14 and~18 to imply (C2), and let $r:=c_*\mathbf m'/q^2-1$.
Define
\begin{equation*}
\mathcal G=\mathcal G(B,\mathbf m'):=\begin{cases}\mathcal G^{\mathrm{lin}}(r)\cap\mathcal G^{\mathrm{qu}}&\text{if }r\ge1,\\\emptyset&\text{otherwise.}\end{cases}
\end{equation*}
If $r\ge1$, Steps~11 and~18 and the inequality
\begin{equation*}
\exp(-r)\le\exp(1)\,16!\,\Bigl(\frac{q^2}{c_*\mathbf m'}\Bigr)^{16}
\end{equation*}
lead to
\begin{equation*}
P_n(\mathcal G^c)\le C_{11}(\ell\mathbf m')^3\Bigl(\frac{q^2}{\mathbf m'}\Bigr)^{16},
\end{equation*}
with a constant $C_{11}\ge(2/c_*)^{16}$ depending on $\rho$ alone.
If $r<1$, then $q^2/\mathbf m'>c_*/2$; as $\ell\mathbf m'\ge1$, it follows that $C_{11}(\ell\mathbf m')^3(q^2/\mathbf m')^{16}\ge1\ge P_n(\mathcal G^c)$, so that the same bound is valid.
On $\mathcal G$, it follows from Step~3 above and from Step~2 of the proof of Proposition~5.11 that the number $N_B(\lambda')$ of eigenvalues not exceeding $\lambda'$ of the form $\|D_A\phi\|^2$ on the spinors over $B$ satisfies
\begin{equation*}
N_B(\lambda')\le C_3|B|\bigl(\sqrt{\lambda'}+\mathbf m'\bigr)^3.
\end{equation*}
The gauge transformation and unitary identification of Step~9 identify
this count with the count for the original field on $B$.

\emph{Step 28} (levels).
Let $\mathbf m$, $\mathbf m'$ and $n$ be as in the theorem, and set $\ell:=\lceil1/\mathbf m'\rceil$.
In the case $\ell<4$ one has $N_A\le192n^3\le192\cdot27\,\mathbf m'^3n^3$.
In the case $\ell\ge4$ one has $10\ell<n$, and Step~1 of the proof of Proposition~5.11 gives $N_A(\lambda)\le\sum_aN_a(\lambda')$, where $\lambda'=2\lambda+c_J\ell^{-2}$ and the number of boxes $B_a$ is at most $(n/\ell)^3$.
For each integer $k\ge0$ with $\mathbf m'_k:=2^k\mathbf m'\le1$ let $\mathcal G_{a,k}:=\mathcal G(B_a,\mathbf m'_k)$, formed with the same $\ell$.
As $\ell\mathbf m'\le2$ and $q^2\le\mathbf m'$, Step~27 yields, also when $\mathcal G_{a,k}=\emptyset$,
\begin{equation*}
P_n(\mathcal G_{a,k}^c)\le C_{11}(\ell\mathbf m'_k)^3\Bigl(\frac{q^2}{\mathbf m'_k}\Bigr)^{16}=C_{11}(\ell\mathbf m')^3\Bigl(\frac{q^2}{\mathbf m'}\Bigr)^{16}2^{-13k}\le8C_{11}2^{-13k}.
\end{equation*}
Here Steps~1 to~27 are applied to the pair $(\ell,\mathbf m'_k)$.
This pair satisfies Setting~(S), because $\mathbf m\le\mathbf m'_k\le1$ and $\ell\mathbf m'_k\ge1$.
The conditions of Setting~(S) bound $\ell\mathbf m'$ from below only, and for the largest admissible $k$ one has $\ell\mathbf m'_k>\ell/2$.
Every bound of Steps~1 to~27 in which $\ell$ and $\mathbf m'$ occur together (Setting~(S) enumerates these bounds, in the paragraph on the dependence on~$\ell$) holds as soon as $\ell\mathbf m'\ge1$, or becomes better as $\ell\mathbf m'$ increases, with the exception of the factor $(\ell\mathbf m')^3$ of Steps~8, 11 and~18.
This factor is retained; in the estimate above it appears as $(\ell\mathbf m'_k)^3=(\ell\mathbf m')^32^{3k}$.

\emph{Step 29} (summation over the levels).
Denote by $k_{\max}$ the largest integer $k$ for which $2^k\mathbf m'\le1$.
On the event that $\mathcal G_{a,k}$ is the first of $\mathcal G_{a,0},\dots,\mathcal G_{a,k_{\max}}$ to occur, $N_a(\lambda')$ is estimated by Step~27 at the level $k$.
On the event that none of them occurs, $N_a\le192(4\ell)^3(2^{k_{\max}+1}\mathbf m')^3$.
With $C_{12}:=\max\{C_3,192\}\max\{1,8C_{11}\}$ it follows that
\begin{equation*}
E_nN_a(\lambda')\le C_{12}(4\ell)^3\bigl(\sqrt{\lambda'}+\mathbf m'\bigr)^3\Bigl[1+\sum_{k\ge1}2^{3k}2^{-13(k-1)}\Bigr],
\end{equation*}
which gives the first bound in the display of Theorem~5.14.
The second bound is a consequence of the formula
\begin{equation*}
\mathsf T_n=2\int_{\mathbf m^2}^\infty\lambda^{-3}n^{-3}E_nN_A(\lambda)\,\mathrm d\lambda
\end{equation*}
and of the inequality $(\xi+\zeta)^3\le4(\xi^3+\zeta^3)$, as in Step~8 of the proof of Proposition~5.11.
The third bound results from $\mathsf c_n\le3\pi^2\mathsf h_3^4q^2\mathsf T_n$, which is part of Theorem~5.7.

\emph{Step 30} (continuum).
Choose $\mathbf m:=\boldsymbol\mu_N$ and $\mathbf m':=\max\{\boldsymbol\mu_N,e^2\eta\}\le1$.
Then $\mathbf m'/\mathbf m\le\mathsf g_*$ and $g\le e^2/(c_{\boldsymbol\mu}\bar m)$.
Assertion (a) follows from Lemma~5.10\,(a) together with $B_n(t_0)\le3\mathsf S_n$; the latter inequality holds since the terms of $B_n(t_0)$ are non-negative and since, by Lemma~5.9\,(a), the corresponding sum over all $y$ equals $3\mathsf S_n$.
For (b), Proposition~5.6\,(a) and Theorem~5.7 give
\begin{equation*}
K_n(u,u)=I_n(u)-\|du\|^2\le\frac{\pi^2}{4}\,\mathsf c_n\|du\|^2.
\end{equation*}
Assertion (c) is Proposition~5.13\,(e), applied with $C_M/3$ in place of $C_M$.
\end{proof}
\subsection{The tuned-mass sum rule}
\begin{named}{Theorem 5.16}
\emph{The tuned masses of \cite{QEDIII}: a bound from above, a sum rule, a bound from below.}
Let $0<\rho<1$, let the weights be standard, let $\bar m\in(0,1]$, and let $e>0$ be such that $e_K=e\,h^{1/2}\le e_1$;
the data of \cite[Theorem~44]{QEDIII} satisfy this condition.
Let $N\ge N_{\min}(R_{\mathrm{top}}(e),e)$, and denote by $(\varepsilon^{\mathrm v,t}_k,m_k,E_k)$ the trajectory of the tuned pair $P_N(e)$ on $R_{\mathrm{top}}(e)$ \cite[Definition before Proposition~28.3]{QEDIII};
in particular $m_{k_\star}=0$, where $k_\star:=k_{\mathrm{top}}$.
The numbers $L$, $M$, $K_0$, $c_{\mathrm{fl}}$, $C_0$, $\varepsilon$, $b=1$, $b_k$, $b_L$ and $c_b$ are those of \cite{QEDIII}.
The index $\star$ refers to the level $k_\star$:
the numbers $a_\star\in\{h,h/L\}$, $e_\star=e\,a_\star^{1/2}\le e_K$, $\mu_\star=\bar m\,a_\star\in(L^{-2},1]$ and $\ell_\star$ are independent of $N$, and $L^{-k_\star}=\eta/a_\star$ \cite[Section~G4.0.2, (Sc1), (Sc3), (Sc5); Section~I.6.1]{QEDIII}.
For $0\le j\le k_\star$ put $a_j:=L^{j-N}$;
the running charge of level $j$ is $e_j:=e\,a_j^{1/2}=q\,L^{j/2}$ \cite[Section~I.6.1]{QEDIII}.
Throughout the theorem and its proof the symbol $e_j$ denotes this number, not a unit vector, and in a formula written for a general index $j$ it bears this meaning at $j=0,1,2$ too, where it equals $q$, $L^{1/2}q$ and $Lq$ \cite[Section~I.2.5]{QEDIII};
the symbol $e_1$ in the condition $e_K\le e_1$, on the other hand, is the ceiling of the charge given by line~(C9) of \cite[Section~I.8.2]{QEDIII}.
Put
\begin{equation*}
d_\star:=\bigl(b_\star^{-1}+\mu_\star^{-1}\bigr)^{-1},\qquad Z_0:=b_\star^2\,(b_\star+\mu_\star)^{-2},\qquad c_\star:=d_\star-\frac{d_\star^2}{b_\star}=\mu_\star Z_0 .
\end{equation*}
Here $b_\star=b_{k_\star}=b_-\,(1-\eta/a_\star)^{-1}$, with $b_-:=b\,(1-L^{-1})$ \cite[Section~I.3.4]{QEDIII}.
As $a_\star$ is independent of $N$ and $\eta=L^{-N}$, the number $b_\star$ decreases to $b_-$ as $N\to\infty$;
thus $b_\star$, $d_\star$, $Z_0$ and $c_\star$ depend on $N$.
With the number $Z_0^\infty:=b_-^2\,(b_-+\mu_\star)^{-2}>0$, which does not depend on $N$, one has
\begin{equation*}
Z_0^\infty\le Z_0\le1,\qquad\mu_\star Z_0^\infty\le c_\star\le\mu_\star,
\end{equation*}
and $Z_0$ decreases to $Z_0^\infty$ as $N\to\infty$.
In this theorem and its proof, $n:=R_{\mathrm{top}}(e)/\eta=L^{k_\star}\ell_\star=\ell_0$, an odd number \cite{QEDIII};
this $n$ differs from the number $n=8hL^{T+N}$ of the notation fixed above and of Theorem~5.14, whereas $\eta$, $q$ and $\boldsymbol\mu_N$ have the same meaning as there.

Let $\Omega$ be the constant bond function on $T^0_{k_\star}$ with $e_\star\Omega_\mu\ell_\star=\pi$.
For $\zeta=(\zeta_j)_{j<k_\star}\in\prod_j\mathrm{Box}_j$ define $A_{k_\star}:=\Omega$ and, for $j<k_\star$, $A_j:=A^{\min}_j(A_{j+1})+\mathrm{Sq}_j\zeta_j$ \cite{QEDIII}.
The letter $D_A$ stands, in the theorem and in its proof, for the operator \cite[(I.2.3)]{QEDIII}, which contains no mass, and $\nabla^*_A\nabla_A$ for that of \cite[(I.2.4)]{QEDIII}.
The operators $D_{A_0}$, $D_0$ and $D_{\Omega_0}$, where $D_0$ is $D_A$ at $A=0$ \cite[Section~I.5.7]{QEDIII} and $\Omega_0$ is defined in Step~13 of the proof, are those of level $0$, with mesh $1$ and charge $q$, on $\mathfrak S^{\mathrm{ch}}(1,n)$;
the operator $D_\Omega$ is that of level $k_\star$, with mesh $L^{-k_\star}$ and charge $e_\star$.
At level $0$ and for real $\boldsymbol\mu$, the operator $D_A+\boldsymbol\mu$ is the operator $D_A$ of \cite[Definition~2.10]{QEDII} with the lattice mass $\mathbf m=\boldsymbol\mu$ \cite[(I.2.5); Section~I.9.3]{QEDIII}.
On $\mathcal H_0:=\mathfrak S^{\mathrm{ch}}(1,n)$ let $\mathcal D_0:=D_{A_0}+\boldsymbol\mu_N$.
With $\mathcal H_j:=\mathfrak S(1,\ell_j)$, the map $\mathcal R=\allowbreak \mathcal R(\zeta)\colon\mathcal H_0\to\mathcal H_{k_\star}$ is the product of the averaging maps, defined in Step~1 of the proof.
Set
\begin{equation*}
B:=L^{k_\star}\mathcal D_0+b_\star\,\mathcal R^*\mathcal R,\qquad\mathcal D:=b_\star-b_\star^2\,\mathcal R B^{-1}\mathcal R^*,
\end{equation*}
let $\mathsf P_N$ be the measure with
\begin{equation*}
\mathrm d\mathsf P_N\propto\det B\,\prod_j\mathbf 1_{\mathrm{Box}_j}\,\mathrm d\mu_I(\zeta_j),
\end{equation*}
and let $\mathsf E_N$ denote its mean.
Finally, with $u:=(1,0,0,0)$, let
\begin{gather*}
\Psi_u(y):=\ell_\star^{-3/2}\exp(-ie_\star\Omega\cdot y)\,u,\qquad\psi:=b_\star B^{-1}\mathcal R^*\Psi_u,\qquad\Phi:=\mathcal D\Psi_u,\\
Z_N:=\mathsf E_N\|\psi\|^2,\qquad\mathsf W_N:=\mathsf E_N\,\frac\rho2\,(\psi,\nabla^*_{A_0}\nabla_{A_0}\psi)^{\mathrm{ch}}\ge0;
\end{gather*}
the number $\mathsf W_N$ is the Wilson energy of the dressed zero mode.
\begin{itemize}
\item[(a)] \emph{The flow density as a mixture of Gaussians with a non-negative weight.}
Whatever the real value of $\boldsymbol\mu_N$, one has $\det B\ge0$, the measure $\mathsf P_N$ is a probability measure, and
\begin{equation*}
\frac{\mathscr F_{k_\star}(\Omega;\Psi)}{\mathscr F_{k_\star}(\Omega;0)}=\mathsf E_N\exp\bigl(-\langle\bar\Psi,\mathcal D\Psi\rangle\bigr),\qquad\mathscr F_{k_\star}:=\mathscr F_{k_\star}\bigl[\varepsilon^{\mathrm v,t}_{k_\star},0,E_{k_\star}\bigr].
\end{equation*}
If $\boldsymbol\mu_N>0$, then $\operatorname{Re}\mathcal D(\zeta)\ge\bigl(b_\star^{-1}+(L^{k_\star}\boldsymbol\mu_N)^{-1}\bigr)^{-1}$ for every $\zeta$.
\item[(b)] \emph{From above.}
For all $N\ge0$,
\begin{equation*}
m_{0,N}(e)\le10^{-4}\,e_\star^{3/4-10\varepsilon}\,\bar m\,\eta .
\end{equation*}
Consequently $\limsup_N\boldsymbol\mu_N/\eta\le\bar m\bigl(1+10^{-4}e_K^{3/4-10\varepsilon}\bigr)$;
the tuned assignment is therefore physical from above, that is, $\boldsymbol\mu_N\le C\eta$ for large $N$.
\item[(c)] \emph{A sum rule, for every real value of $\boldsymbol\mu_N$.}
At every $\zeta$ with $\det B\ne0$ the identity \textbf{(SR)}
\begin{equation}
\tag{SR}
\operatorname{Re}(\Psi_u,\mathcal D\Psi_u)=b_\star^{-1}\|\Phi\|^2+L^{k_\star}\Bigl[\boldsymbol\mu_N\|\psi\|^2+\frac\rho2\,(\psi,\nabla^*_{A_0}\nabla_{A_0}\psi)^{\mathrm{ch}}\Bigr]
\end{equation}
holds.
In the mean,
\begin{equation*}
\boldsymbol\mu_NZ_N+\mathsf W_N=\frac\eta{a_\star}\,c_\star\,(1+\vartheta_N),\qquad|\vartheta_N|\le e_\star^{3/4-10\varepsilon},\qquad Z_N\ge\tfrac12Z_0 .
\end{equation*}
\item[(d)] \emph{From below.}
For all large $N$,
\begin{equation*}
\boldsymbol\mu_N\ge-2\rho\,C_5\,q^{2-\varepsilon},\qquad C_5:=108\,c_b^{-1}+6 .
\end{equation*}
For comparison, the bound $|m_{0,N}(e)|\le\frac18q^{3/4-8\varepsilon}$ of \cite[Proposition~28.3\,(d)]{QEDIII} gives $\boldsymbol\mu_N\ge-\frac18q^{3/4-8\varepsilon}$.
\end{itemize}
\end{named}
\begin{proof}[Proof of Theorem~5.16]
\emph{Step 1} (averaging maps).
By \cite[Propositions~5.31, 5.52\,(i), (iii)]{QEDI} the functions $\mathscr B_j:=H_jA^{\min}_j(A_{j+1})$ and $\omega_j:=\omega_j(A_{j+1})$ are real.
The map $(\mathsf S_jv)(y'):=L^{-3/2}v(y'/L)$ is unitary from $\mathcal H_{j+1}$ onto $\mathfrak S(L,\ell_j)$, the latter space carrying the inner product $\langle\cdot,\cdot\rangle_L$, and $v_L=L^{1/2}\mathsf S_jv$ \cite{QEDIII}.
Let
\begin{equation*}
R_j:=\mathsf S_j^{-1}\exp(-ie_j\omega_j)\,Q(\mathscr B_j)\quad(0<j<k_\star),\qquad R_0:=\mathsf S_0^{-1}\exp(-iq\omega_0)\,Q(\mathscr B_0)J^*,
\end{equation*}
and $\mathcal R:=R_{k_\star-1}\cdots R_0$.
It follows from \cite[(I.J1), (I.J3), (I.Q1), (I.Q2)]{QEDIII} that
\begin{equation*}
R_j^*=Q^T(-\mathscr B_j)\exp(ie_j\omega_j)\,\mathsf S_j\quad(0<j<k_\star),\qquad R_0^*=JQ^T(-\mathscr B_0)\exp(iq\omega_0)\,\mathsf S_0,
\end{equation*}
that $R_jR_j^*=1$, and that every $R_j$ commutes with the spin matrices.

\emph{Step 2} (a single Gaussian integration).
Suppose that $\mathcal D_j$ acts on $\mathcal H_j$ and that $G_j:=\mathcal D_j+b_LR_j^*R_j$ is invertible.
In the exponent of $\mathbf T_j$ \cite{QEDIII} one has $\langle Q(-\mathscr B)\bar\Psi,\allowbreak v'\rangle_L=\allowbreak \langle\bar\Psi,\allowbreak Q^T(-\mathscr B)v'\rangle$.
Hence \cite[Proposition~22.1\,(G2)$^{\mathrm w}$]{QEDIII} yields
\begin{equation*}
\int\mathrm D\Psi_j\,\exp\bigl\{-b_L\langle\cdots\rangle_L-\langle\bar\Psi_j,\mathcal D_j\Psi_j\rangle\bigr\}=\det G_j\,\exp\bigl(-\langle\bar\Psi_{j+1},\mathcal D_{j+1}\Psi_{j+1}\rangle\bigr),
\end{equation*}
where $\mathcal D_{j+1}:=L\bigl[b_L-b_L^2R_jG_j^{-1}R_j^*\bigr]$.
When $\mathcal D_j$ is invertible, (K10) gives $\mathcal D_{j+1}^{-1}=b^{-1}+L^{-1}R_j\mathcal D_j^{-1}R_j^*$.

\emph{Step 3} (accretivity).
If $\operatorname{Re}\mathcal D_j\ge d>0$, then $\operatorname{Re}\mathcal D_{j+1}\ge\bigl(b^{-1}+(Ld)^{-1}\bigr)^{-1}$.
Indeed, let $\Phi'=\mathcal D_{j+1}\Psi'$ and $\phi=\mathcal D_j^{-1}R_j^*\Phi'$.
Then $\Psi'=b^{-1}\Phi'+L^{-1}R_j\phi$ and
\begin{equation*}
\operatorname{Re}(\Psi',\Phi')=b^{-1}\|\Phi'\|^2+L^{-1}\operatorname{Re}(\phi,\mathcal D_j\phi)\ge b^{-1}\|\Phi'\|^2+L^{-1}d\,\|\phi\|^2,
\end{equation*}
while the Schwarz inequality gives
\begin{equation*}
\|\Psi'\|^2\le\bigl(b^{-1}+(Ld)^{-1}\bigr)\bigl(b^{-1}\|\Phi'\|^2+L^{-1}d\,\|\phi\|^2\bigr).
\end{equation*}
One has $\operatorname{Re}\mathcal D_0\ge\boldsymbol\mu_N$ \cite[(I.O2)]{QEDIII}, while $b_1=b$ and $b_{j+1}^{-1}=b^{-1}+L^{-1}b_j^{-1}$ for $j\ge1$ \cite[Section~I.3.4]{QEDIII}.
For $\boldsymbol\mu_N>0$ the implication above, used at $j=0$ with $d=\boldsymbol\mu_N$ and then by induction on $j\ge1$, therefore gives for every $\zeta$ and $1\le k\le k_\star$
\begin{equation*}
\operatorname{Re}\mathcal D_k\ge\bigl(b_k^{-1}+(L^k\boldsymbol\mu_N)^{-1}\bigr)^{-1}.
\end{equation*}
Applying Step~2 with the same start, at $j=0$ with $b_1=b$, then by induction on $j\ge1$, and using $R_jR_j^*=1$, one finds $\mathcal D_{k_\star}^{-1}=\allowbreak b_\star^{-1}+L^{-k_\star}\mathcal R\mathcal D_0^{-1}\mathcal R^*$;
by (K10) this means $\mathcal D_{k_\star}=\mathcal D$.

\emph{Step 4} (the identity of a step, iterated).
At every level $j<k_\star$ the data satisfy (OS), $|m_j|\le w^{\mathrm m}_j$, $\|\mathscr E_j\|^x_j\le w^{\mathrm E}_j$, $\mathscr E_j|_0=E_j$ and $\mathfrak P_j(\cdot,0)=0$ \cite[Propositions~26.6\,(b), 28.1, 28.3\,(a), (b)]{QEDIII}.
Constant bond functions belong to $\mathrm{Dom}_k$, and $A_j\in\mathrm{Dom}_j$ \cite[Proposition~5.38\,(i), (ii)]{QEDI}, \cite{QEDIII}.
Starting from
\begin{equation*}
\mathscr F_0=\det(D_0+\bar m\eta)^{-1}\exp\bigl\{-\langle\bar\Psi_0,\mathcal D_0\Psi_0\rangle^{\mathrm{ch}}-\varepsilon^{\mathrm v}_{\mathrm{in}}n^3\bigr\}
\end{equation*}
\cite{QEDIII} and iterating \cite[Proposition~24.1\,(c)]{QEDIII} at $\mathscr C'=0$, one therefore obtains with \cite[(G5)$^{\mathrm w}$]{QEDIII}
\begin{equation*}
\mathscr F_{k_\star}(\Omega;\Psi)=\int\prod_j\mathbf 1_{\mathrm{Box}_j}\,\mathrm d\mu_I(\zeta_j)\;\mathcal G(\zeta;\Psi);
\end{equation*}
the integrand $\mathcal G$ arises by iterating the Berezin integration of Step~2, the factors $c_j$ included.

\emph{Step 5} (the weight).
Let $\boldsymbol\mu_N>0$.
By Steps~2 and~3, $\mathcal G=\mathsf w\exp(-\langle\bar\Psi,\mathcal D\Psi\rangle)$ with $\mathsf w\propto\prod_j\det G_j$.
Integration over $\Psi$, together with $k_\star$ applications of \cite[Proposition~22.2\,(a)]{QEDIII}, gives
\begin{equation*}
\mathsf w\det\mathcal D=\exp\bigl(-\varepsilon^{\mathrm v}_{\mathrm{in}}n^3\bigr)\,\frac{\det\mathcal D_0}{\det(D_0+\bar m\eta)},
\end{equation*}
and by (K10), $\det\mathcal D\,\allowbreak \det B=\allowbreak b_\star^{4n_{k_\star}}\det\bigl(L^{k_\star}\mathcal D_0\bigr)$.
Hence $\mathsf w=C_N\det B$ with a constant $C_N>0$.

\emph{Step 6} (all real masses; proof of (a)).
At a given $\zeta$, each coefficient of $\mathcal G$ is a polynomial function of the mass, whereas each coefficient of $C_N\det B\exp(-\langle\bar\Psi,\mathcal D\Psi\rangle)$ is a rational function of it;
they coincide for $\boldsymbol\mu_N>0$ by Step~5, and so, by (K11), for every real $\boldsymbol\mu_N$ at which $\det B\ne0$.
By \cite[(I.O2), (I.O3)]{QEDIII} and Step~1, each of $X=B$ and $X=\mathcal D$ satisfies
\begin{align}
&[X,\Sigma_3]=0,\tag{P1}\\
&X^*=\gamma_5X\gamma_5,\tag{P2}
\end{align}
and $\gamma_5$ interchanges the eigenspaces of $\Sigma_3$;
hence $\det B=|\det B^+|^2\ge0$, where $B^+$ is the block of $B$ on the eigenspace of $\Sigma_3$ for the eigenvalue $1$.
The function $\det B$ is real analytic in $\zeta$, and it does not vanish identically, because $\mathscr F_{k_\star}(\Omega;0)\ne0$;
by (K11), $\det B\ne0$ almost everywhere.
Since $\Sigma_3u=u$, the function
\begin{equation*}
\det B\,\|B^{-1}\mathcal R^*\Psi_u\|^2=\|\operatorname{adj}(B^+)\,\mathcal R^*\Psi_u\|^2
\end{equation*}
is bounded, and therefore $\|\psi\|^2$, $\|\Phi\|^2$ and $(\psi,\nabla^*\nabla\psi)$ are integrable.

\emph{Step 7} (the vector $\Psi_u$).
One has $H_{k_\star}\Omega=\Omega\in\frac14\tilde R_{k_\star}$ \cite[Proposition~5.31]{QEDI}, \cite{QEDIII}.
The function $f(c,g):=\ell_\star^{-3/2}\exp(-ie_\star\Omega\cdot c)\,u$ is antiperiodic and covariantly constant with respect to $\Omega$;
hence $D_\Omega f=0$ \cite[(I.O8)]{QEDIII}, and $Q_{k_\star}(\Omega)J^*f=\Psi_u$, $JQ^T_{k_\star}(-\Omega)\Psi_u=f$ \cite{QEDIII}.
The definitions of $S_{k_\star}(\Omega)$ and $D_{k_\star}(\Omega)$ in \cite{QEDIII} then give $S_{k_\star}(\Omega)f=(\mu_\star+b_\star)^{-1}f$ and $D_{k_\star}(\Omega)\Psi_u=d_\star\Psi_u$.

\emph{Step 8} (degrees two and four).
Since $m_{k_\star}=0$, one has by \cite{QEDIII}
\begin{equation*}
\mathscr F_{k_\star}(\Omega;\Psi)=\mathscr F_{k_\star}(\Omega;0)\exp\bigl\{-\langle\bar\Psi,D_{k_\star}(\Omega)\Psi\rangle+\mathcal A_2+\mathcal A_4+\dots\bigr\},
\end{equation*}
where $\mathcal A_{2i}$ denotes the part of degree $(i,i)$ of $\sum_XE_{k_\star}\bigl(X;\Omega,H^\#_{k_\star}(\Omega)\Psi\bigr)$.
Write $\mathcal A_2=:\langle\bar\Psi,K_E\Psi\rangle$ and $\delta\mathcal D:=\mathcal D-\mathsf E_N\mathcal D$.
A comparison of the logarithms of the two sides of the identity in (a) gives
\begin{equation*}
\mathsf E_N\mathcal D=D_{k_\star}(\Omega)-K_E .
\end{equation*}
Let $a,b$ be indices with $\Sigma_3=+1$ and $c,d$ indices with $\Sigma_3=-1$.
As $\delta\mathcal D_{ad}=\delta\mathcal D_{cb}=0$ by (P1), the coefficient of $\bar\Psi_a\Psi_b\bar\Psi_c\Psi_d$ in $\mathcal A_4$ is $\pm\mathsf E_N\bigl[\delta\mathcal D_{ab}\,\allowbreak \delta\mathcal D_{cd}\bigr]$.

\emph{Step 9} (sizes).
From \cite[Proposition~20.1\,(c)]{QEDIII}, from \cite[(k6)]{QEDIII}, from $c_\kappa\le K_0$ \cite{QEDIII} and from the bound $\|E_{k_\star}\|_{k_\star}\le4e_\star^{1/4-6\varepsilon}$ of \cite[Proposition~28.3\,(b)]{QEDIII} one obtains
\begin{equation*}
\Bigl|\sum_XE_{k_\star}\bigl(X;\Omega,H^\#_{k_\star}(\Omega)\Psi\bigr)\Bigr|_{\mathrm{aw}_2}\le4e_\star^{1/4-6\varepsilon}K_0\,\Bigl(\frac{\ell_\star}{M}\Bigr)^3,\qquad\mathrm{aw}_2=C_0^{-1}e_\star^{-1/4+\varepsilon}.
\end{equation*}
Therefore the moduli of the entries of $K_E$ have a sum not larger than $\ell_\star^3\delta_\star$, and $|\mathcal A_4|_1\le\ell_\star^3\delta_2$, with
\begin{equation*}
\delta_\star:=4C_0^2K_0M^{-3}e_\star^{3/4-8\varepsilon},\qquad\delta_2:=4C_0^4K_0M^{-3}e_\star^{5/4-10\varepsilon}.
\end{equation*}
By \cite[(T2)]{QEDIII}, valid by \cite[(Sc6)]{QEDIII}, and by $c_{\mathrm{fl}}=3294K_0$ \cite{QEDIII},
\begin{equation*}
\ell_\star^3\delta_\star\le10^{-5}e_\star^{3/4-10\varepsilon}d_\star,\qquad\delta_2\le10^{-6}e_\star^{5/4-12\varepsilon}d_\star^2 .
\end{equation*}
With Step~7 it follows that the number $\bar z:=(\Psi_u,\mathsf E_N\mathcal D\Psi_u)$ satisfies $|\bar z-d_\star|\le\delta_\star$, and that $\bigl|\,\allowbreak \|\mathsf E_N\mathcal D\Psi_u\|-d_\star\bigr|\le\ell_\star^3\delta_\star$.

\emph{Step 10} (proof of (b)).
If $\boldsymbol\mu_N\le0$, there is nothing to prove, and for $N<N_{\min}$ one has $m_{0,N}=0$ by definition.
Let $\boldsymbol\mu_N>0$.
Taking the mean in (a), one finds
\begin{equation*}
\bigl(b_\star^{-1}+(L^{k_\star}\boldsymbol\mu_N)^{-1}\bigr)^{-1}\le\operatorname{Re}\bar z\le d_\star(1+t),\qquad t:=\delta_\star/d_\star .
\end{equation*}
Hence
\begin{equation*}
(L^{k_\star}\boldsymbol\mu_N)^{-1}\ge d_\star^{-1}(1-t)-b_\star^{-1}=\mu_\star^{-1}\bigl(1-t(1+\mu_\star/b_\star)\bigr)\ge\mu_\star^{-1}(1-3t),
\end{equation*}
so that $L^{k_\star}\boldsymbol\mu_N\le\mu_\star(1+6t)$;
moreover $L^{-k_\star}\mu_\star=\bar m\eta$.

\emph{Step 11} (the identity (SR)).
Let $B$ be invertible.
Then $\Phi=b_\star(\Psi_u-\mathcal R\psi)$, and $\mathcal R^*\Phi=\allowbreak B\psi-b_\star\mathcal R^*\mathcal R\psi=\allowbreak L^{k_\star}\mathcal D_0\psi$.
Consequently
\begin{equation*}
(\Psi_u,\Phi)=b_\star^{-1}\|\Phi\|^2+(\mathcal R\psi,\Phi)=b_\star^{-1}\|\Phi\|^2+L^{k_\star}(\psi,\mathcal D_0\psi),
\end{equation*}
and (SR) results on taking real parts by means of \cite[(I.O2)]{QEDIII}.

\emph{Step 12} (proof of (c)).
Denote by $a\mapsto\hat a$ the permutation of the indices given by $\gamma_5$.
By (P2), $\overline{\delta\mathcal D_{ab'}}=\allowbreak \delta\mathcal D_{\hat b'\hat a}$, and Steps~8 and~9 give
\begin{equation*}
\mathsf E_N\|\delta\mathcal D\Psi_u\|^2=\sum_{a,b,b'}\mathsf E_N\bigl[\delta\mathcal D_{ab}\,\delta\mathcal D_{\hat b'\hat a}\bigr]\,\Psi_u(b)\,\overline{\Psi_u(b')}\le\ell_\star^{-3}|\mathcal A_4|_1\le\delta_2 .
\end{equation*}
Averaging (SR), one obtains
\begin{equation*}
L^{k_\star}\bigl(\boldsymbol\mu_NZ_N+\mathsf W_N\bigr)=\operatorname{Re}\bar z-b_\star^{-1}\bigl(\|\mathsf E_N\mathcal D\Psi_u\|^2+\mathsf E_N\|\delta\mathcal D\Psi_u\|^2\bigr).
\end{equation*}
By Step~9 and the inequalities $b_\star^{-1}\le2$, $d_\star\le1$ and $c_\star\ge\frac13d_\star$, the right side differs from $c_\star=d_\star-d_\star^2/b_\star$ by at most $e_\star^{3/4-10\varepsilon}c_\star$;
in addition $\mathsf E_N\|\Phi\|^2\le2$.
Finally, $\|\psi\|\ge|(\Psi_u,\allowbreak \mathcal R\psi)|=\allowbreak |1-b_\star^{-1}(\Psi_u,\allowbreak \mathcal D\Psi_u)|$, whence
\begin{equation*}
Z_N\ge|1-\bar z/b_\star|^2\ge\bigl(Z_0^{1/2}-2\delta_\star\bigr)^2 .
\end{equation*}

\emph{Step 13} (a gauge with a small potential).
Put $\mathscr A_j:=H_jA_j$.
By \cite[Proposition~5.52\,(iii)]{QEDI},
\begin{equation*}
\mathscr A_j=(\mathscr A_{j+1})_L-d\omega_j+H_j\mathrm{Sq}_j\zeta_j,
\end{equation*}
with $H_0=1$ and $\mathscr A_{k_\star}=\Omega$.
It follows that $A_0=\Omega_0+\tilde A-d\lambda$ with a real $\lambda$, where
\begin{equation*}
\Omega_0:=L^{-k_\star/2}\Omega,\qquad\tilde A:=\sum_{j<k_\star}\bigl(H_j\mathrm{Sq}_j\zeta_j\bigr)_{L^j}.
\end{equation*}
From $\mathrm{Sq}_j=CC_j^{1/2}$ \cite{QEDI}, from \cite[Proposition~5.38\,(iii)]{QEDI} and from $A_L(b)=L^{-1/2}A(b/L)$ one obtains $|\tilde A|_\infty\le\sum_jL^{-j/2}\,\frac12e_j^{-\varepsilon/2}$.
Since $e_j=qL^{j/2}\ge q$ and $L\ge5$, the latter being (L5) of \cite[Proposition~14.2\,(b)]{QEDIII},
\begin{equation*}
|\tilde A|_\infty\le\tfrac12\bigl(1-L^{-1/2}\bigr)^{-1}q^{-\varepsilon/2}\le q^{-\varepsilon/2}.
\end{equation*}
Therefore, for $x=(c,g)$,
\begin{equation*}
\varsigma(x):=\bigl|\exp\bigl(iq\tilde A(c,c_3)\bigr)-1\bigr|\le q^{1-\varepsilon/2}.
\end{equation*}

\emph{Step 14} (the Wilson energy for fixed $\zeta$).
By \cite[(I.O4)]{QEDIII} one may take $A_0=\Omega_0+\tilde A$.
Since $w_3=6$ and every site lies on a single bond of type~$3$, the bond formula of \cite[(I.O2)]{QEDIII} gives
\begin{equation*}
(\psi,\nabla^*_{A_0}\nabla_{A_0}\psi)\le2\,(\psi,\nabla^*_{\Omega_0}\nabla_{\Omega_0}\psi)+12\,\|\varsigma\psi\|^2 .
\end{equation*}
From \cite[Proposition~1.3\,(b)]{QEDIII}, applied to the free operator with mesh~$1$ and character $(\pi,\pi,\pi)$, from \cite[(I.O6)]{QEDIII} and from $\|\mathcal D_0\psi\|=L^{-k_\star}\|\Phi\|$ one gets
\begin{equation*}
(\psi,\nabla^*_{\Omega_0}\nabla_{\Omega_0}\psi)\le c_b^{-1}\|D_{\Omega_0}\psi\|^2\le3c_b^{-1}\bigl(L^{-2k_\star}\|\Phi\|^2+\boldsymbol\mu_N^2\|\psi\|^2+36\,\|\varsigma\psi\|^2\bigr).
\end{equation*}
Together these give the inequality \textbf{(WE)}:
\begin{equation}
\tag{WE}
\frac\rho2\,(\psi,\nabla^*_{A_0}\nabla_{A_0}\psi)\le3\rho c_b^{-1}\bigl(L^{-2k_\star}\|\Phi\|^2+\boldsymbol\mu_N^2\|\psi\|^2\bigr)+\rho\,C_5\,\|\varsigma\psi\|^2 .
\end{equation}

\emph{Step 15} (proof of (d)).
Let $\boldsymbol\mu_N<0$.
From (c), (WE) and Steps~12 and~13,
\begin{multline*}
|\boldsymbol\mu_N|Z_N=\mathsf W_N-\frac\eta{a_\star}\,c_\star(1+\vartheta_N)\\
\le\Bigl[6\rho c_b^{-1}\frac{\eta^2}{a_\star^2}-\frac\eta{2a_\star}\,c_\star\Bigr]+\bigl(3\rho c_b^{-1}\boldsymbol\mu_N^2+\rho\,C_5\,q^{2-\varepsilon}\bigr)Z_N .
\end{multline*}
Since $c_\star\ge\mu_\star Z_0^\infty$ and $\eta\to0$, once $N$ is sufficiently large the square bracket takes a negative value and, by \cite[Proposition~28.3\,(d)]{QEDIII}, one has $3\rho c_b^{-1}|\boldsymbol\mu_N|\le\frac12$.
Division by $Z_N>0$ then yields (d), for every $N\ge N_{\min}(R_{\mathrm{top}}(e),e)$ such that
\begin{equation*}
12\rho\,\eta<c_bc_\star a_\star\qquad\text{and}\qquad6\rho\bigl(\bar m\eta+\tfrac18q^{3/4-8\varepsilon}\bigr)\le c_b ;
\end{equation*}
the first of these conditions holds whenever $12\rho\,\eta<c_b\,\mu_\star Z_0^\infty a_\star$, an inequality whose right side does not depend on $N$.
In the estimate above,
\begin{equation*}
|\vartheta_N|\le e_\star^{3/4-10\varepsilon}\le e_\star^{3(1/4-9\varepsilon)}\le\tfrac1{512},
\end{equation*}
on account of $e_\star\le1$ and of the third inequality of \cite[(T2)]{QEDIII}.
\end{proof}


\subsection{A clustering estimate at a fixed spacing}

Here $P_\lambda$ denotes conditional expectation onto the
$\mathbb Z^3$-invariant events. The local algebra $\mathcal O(Q)$ and
its $L^2$ closure are those of \cite[Section~10.1]{QEDII}.
\begin{named}{Proposition 5.17}
(A clustering rate at a fixed spacing.)
Consider a law $\lambda$ on $\Omega$ which is invariant under the group $W$ and which is reflection positive at each mirror belonging to one of the nine families; by \cite[Theorem~1.2]{QEDII}, every weak limit point of $\mu_{C(n_i)}$, $n_i\to\infty$, is such a law when $0<\rho<1$, $w_i>0$, $M\in\mathbb R$ and $q>0$.
Let $\mathsf r\ge1$ be an integer, put
\begin{equation*}
Q_{\mathsf r}:=\{x:|x_j+\tfrac12|<\mathsf r,\ j=1,2,3\},
\end{equation*}
let $G$ and $G'$ belong to the closure of $\mathcal O(Q_{\mathsf r})$ in $L^2(\lambda)$, and let $a\in\mathbb Z^3$ be such that $\mathsf a:=|a|_\infty-2\mathsf r\ge1$.
Then, with the scalar product and the norm of $L^2(\lambda)$,
\begin{equation*}
\bigl|\langle G,\tau_aG'\rangle_2-\langle G,P_\lambda G'\rangle_2\bigr|\le C_\flat\,\|G\|_2\,\|G'\|_2\Bigl(\frac{\mathsf r}{\mathsf a}+\frac{\mathsf r^2}{\mathsf a^2}\Bigr),
\end{equation*}
where
\begin{equation*}
C_\flat:=\frac{128(1+\sqrt2)}{4\pi^2}\,C''<2\cdot10^4,\qquad C'':=200+64\Bigl(\frac{32}{\pi}+\frac{256}{\pi^2}\Bigr).
\end{equation*}
\end{named}

\begin{proof}
The argument is that of \cite[Lemmas~12.36--12.38, Theorem~12.39]{QEDII}, carried out with the lattice objects.
Both sides of the inequality depend continuously on $G$ and $G'$ in $L^2(\lambda)$, so that it suffices to treat $G,G'\in\mathcal O(Q_{\mathsf r})$.
\begin{enumerate}
\item[1.] \emph{Pictures.}
The reflection $\theta_j:=\theta_{x_j,-1/2}$ maps $Q_{\mathsf r}$ onto itself, and $\tau_{\mathsf re_j}Q_{\mathsf r}\subset\{x_j>-\frac12\}$.
Let $\{i,i'\}=\{1,2,3\}\setminus\{j\}$.
In the axis picture $j$ one has $0\le T_j\le1$, and the operators $T_j$, $U_{j,i}$, $U_{j,i'}$ commute, the latter two being unitary \cite[Lemma~10.9, Definition~10.11]{QEDII}.
By the joint spectral theorem they possess a joint spectral measure $\mathfrak E_j$, written in the variables
\begin{equation*}
t=\exp(-\mathsf e),\quad\mathsf e\in[0,\infty],\qquad u=(\exp ik_i,\exp ik_{i'}),\quad k\in(-\pi,\pi]^2.
\end{equation*}
The vector $\psi^j_G:=[\tau_{\mathsf re_j}G]_j$ satisfies $\|\psi^j_G\|\le\|G\|_2$.
As $\theta_j\tau_{\mathsf re_j}\theta_j=\tau_{-\mathsf re_j}$, one has for $y_j\ge2\mathsf r$
\begin{equation*}
\langle G,\tau_yG'\rangle_2=\bigl\langle\psi^j_{\theta_jG},T_j^{\,y_j-2\mathsf r}U^{y_\perp}\psi^j_{G'}\bigr\rangle,\qquad U^{y_\perp}:=U_{j,i}^{y_i}U_{j,i'}^{y_{i'}}.
\end{equation*}
\item[2.] \emph{The cone.}
By \cite[Theorem~10.12\,(b)]{QEDII}, the measure $\mathfrak E_j$ is carried by the set where $(1-t)^2\ge t\,|1-u_i|^2$, and likewise by the corresponding set for $i'$.
Thus $t=1$ forces $u=1$, and $\sinh\frac{\mathsf e}2\ge|\sin\frac{k_i}2|$ whenever $0<t<1$.
The function $x\mapsto\sin\frac x2-\sinh\frac x4$ is concave on $[0,\pi]$, zero at $0$ and positive at $\pi$; it follows that $\mathsf e\ge\frac12|k|_\infty$.
\item[3.] \emph{The kernel.}
Suppose that $\theta_jG=\pm G$, and let $\varrho:=\langle\psi^j_G,\mathfrak E_j(\cdot)\psi^j_G\rangle$.
With $z_m$, $b_m$, $B_z$ as in \cite[Lemma~12.35]{QEDII}, and with $\Upsilon$ the function written $\Psi$ in that lemma, define
\begin{equation*}
K(y):=\sum_mz_mb_m\,\langle\theta_jG,\tau_{2\mathsf rme_j+y}G\rangle_2,\qquad y\in\mathbb Z^2.
\end{equation*}
Step~1, the unitarity of $\theta_j$ and $\tau_v$, and items (1) to (5) of the proof of \cite[Lemma~12.37]{QEDII} give
\begin{equation*}
K(y)=\int d\varrho\,\Upsilon(t^{2\mathsf r})\,u^y,\qquad K(-y)=\overline{K(y)},\qquad|K|\le8\|G\|_2^2.
\end{equation*}
\item[4.] \emph{Decay.}
Let $y_i\ge2\mathsf r$.
In the picture $i$, as in item (6) of the same proof,
\begin{equation*}
K(y)=\pm\bigl\langle\psi^i_{\theta_iG},T_i^{\,y_i-2\mathsf r}\,\mathsf B\,U_{i,i'}^{y_{i'}}\psi^i_G\bigr\rangle,\qquad\mathsf B=\int d\mathfrak E_i\,B_z(2\mathsf rk_j).
\end{equation*}
Since $B_z(2\mathsf rk_j)\ne0$ forces $|k_j|>\frac{\pi}{8\mathsf r}$ \cite[Lemma~12.35]{QEDII}, Step~2 gives $\mathsf e>\alpha:=\frac{\pi}{16\mathsf r}$ wherever $B_z(2\mathsf rk_j)\ne0$, and so
\begin{equation*}
|K(y)|\le8\exp\bigl(-\alpha(|y|_\infty-2\mathsf r)\bigr)\|G\|_2^2\qquad\text{for }|y|_\infty\ge2\mathsf r.
\end{equation*}
\item[5.] \emph{Density.}
For $s\ge1$ there are $8s$ points $y\in\mathbb Z^2$ with $|y|_\infty=s$, and $\sum_{m\ge1}(2\mathsf r+m)\exp(-\alpha m)\le2\mathsf r\alpha^{-1}+\alpha^{-2}$.
Consequently
\begin{equation*}
\sum_y|K(y)|\le\bigl[8(4\mathsf r+1)^2+64\bigl(2\mathsf r\alpha^{-1}+\alpha^{-2}\bigr)\bigr]\|G\|_2^2\le C''\mathsf r^2\|G\|_2^2.
\end{equation*}
By Fourier inversion for an absolutely summable sequence, the image of the measure $\Upsilon(t^{2\mathsf r})\,d\varrho$ under $(t,u)\mapsto k$ is $h\,dk$, where
\begin{equation*}
0\le h\le\frac{C''\mathsf r^2\|G\|_2^2}{4\pi^2}.
\end{equation*}
Since $\Upsilon>0$ on $[0,1)$, this implies $\varrho(\{t<1,\ k=0\})=0$.
\item[6.] \emph{Cube vectors.}
The inequality $\Upsilon(\exp(-x))\ge\tanh x/(1+\sqrt2)$ \cite[Lemma~12.35]{QEDII}, Step~2, the monotonicity of $x\mapsto\exp(-\mathsf ax)/\tanh(2\mathsf rx)$ and the bound $\coth x\le1+\frac1x$ yield, since $t^{\mathsf a}=0$ at $t=0$,
\begin{multline*}
\int_{\{t<1\}}t^{\mathsf a}\,d\varrho\le(1+\sqrt2)\int dk\,h\,\frac{\exp(-\mathsf a|k|_\infty/2)}{\tanh(\mathsf r|k|_\infty)}\\
\le(1+\sqrt2)\,\frac{C''}{4\pi^2}\,\|G\|_2^2\Bigl(\frac{32\mathsf r^2}{\mathsf a^2}+\frac{16\mathsf r}{\mathsf a}\Bigr).
\end{multline*}
For a general $G=G_++G_-$, with $\theta_jG_\pm=\pm G_\pm$, the bound acquires a factor $4$ \cite[Lemma~12.38]{QEDII}.
\item[7.] \emph{Conclusion.}
Let $a_j=|a|_\infty$; if no coordinate of $a$ equals $|a|_\infty$, one takes complex conjugates.
Put $\hat P:=\mathfrak E_j(\{t=1\})$ and $S:=T_j^{\mathsf a}-\hat P$, so that $T_j^{\mathsf a}=\hat P+S$.
By Step~2, $\hat PU^{a_\perp}=\hat P$; moreover
\begin{equation*}
\bigl\langle\psi^j_{\theta_jG},\hat P\psi^j_{G'}\bigr\rangle=\bigl\langle\tau_{-\mathsf re_j}G,P_j\tau_{\mathsf re_j}G'\bigr\rangle_2=\langle G,P_\lambda G'\rangle_2,
\end{equation*}
where $P_j$ is the projection onto $\{H:\tau_{e_j}H=H\}$ \cite[proof of Theorem~10.12\,(c), Step~1; Theorem~10.12\,(c)]{QEDII}.
The term containing $S\ge0$ is estimated by the Schwarz inequality and Step~6.
\end{enumerate}
\end{proof}


\section{Bounds on the tuned lattice masses}

Given a power $h$ of $L$, one calls \emph{class} the interval $I_h:=(\frac1{Lh},\frac1h]$.
For $m>0$ put $h(m):=L^{\lfloor\log_L(1/m)\rfloor}$;
then $I_h$ consists of those $m>0$ for which $h(m)=h$, so that $m\in I_{h(m)}$, and the number $h$ of the notation fixed above is $h(\bar m)$.
Where $h$ is the index of a class, as here and in Proposition~6.3(c), (d), the starred numbers, $\epsilon_e$ and $R_{\mathrm{top}}(e)$ are formed with this power of $L$;
where $h$ enters together with the parameter $\bar m$, as in Theorem~6.1 and in the number $\epsilon_e$ of that theorem, it is the number $h=h(\bar m)$ of the notation fixed above, and $\bar m\in I_h$.
Let $m\in I_h$, $e\,h^{1/2}\le e_1$ and $N\ge N_{\min}(R_{\mathrm{top}}(e),e)$.
Then $m_{0,N}(e;m)$ denotes the tuned shift formed with $\bar m=m$ \cite[Definition before Proposition~28.3]{QEDIII}, and one puts $\varphi_N(m):=m+m_{0,N}(e;m)/\eta$, so that $\boldsymbol\mu_N=\varphi_N(\bar m)\eta$.
The index $\star$ refers to the level $k_\star=k_{\mathrm{top}}$ on $R_{\mathrm{top}}(e)=hM\,r_K(e)$, where $r_K$ is the smallest power $L^i$ with $L^i\ge\log^r(1/(e\,h^{1/2}))$.
More generally, for a power $h'\ge1$ of $L$ with $e\,h'^{1/2}\le e_1$ let $R_{\mathrm{top}}(e;h'):=h'M\,r'$, where $r'$ is the smallest power $L^i$ with $L^i\ge\log^r(1/(e\,h'^{1/2}))$.
Thus $r'$ is the number $r_k$ of \cite[Section~I.6.1]{QEDIII} with $k=N+\log_Lh'$, and $R_{\mathrm{top}}(e;h')$ is the side $R_{\mathrm{top}}(e)$ of \cite[Section~I.6.2]{QEDIII} formed at any parameter $\bar m\in I_{h'}$;
for the class $I_h$ at hand, $R_{\mathrm{top}}(e)=R_{\mathrm{top}}(e;h)$.
The numbers $a_\star\in\{h,h/L\}$, $e_\star=e\,a_\star^{1/2}\le e\,h^{1/2}$, $\ell_\star$ and $N_{\min}$ depend on $(e,h)$, but neither on $N$ nor on $m\in I_h$ \cite[Section~G4.0.2, (Sc0), (Sc1), (Sc5)]{QEDIII};
the inclusion $a_\star\in\{h,h/L\}$ results from $k_\star\in\{K,K-1\}$, which is (Sc5) there, and from $a_K=h$ at the level $K=N+\log_Lh$ \cite[Section~I.6.1]{QEDIII}.
On the other hand, $b_\star=b_{k_\star}=b_-/(1-\eta/a_\star)$ depends on $N$, and $b_-:=1-\frac1L\le b_\star\le1$.
Further, $\mu_k=ma_k$, and
\begin{equation*}
d_k(\mu):=\frac{b_k\mu}{b_k+\mu},\qquad Z_k(\mu):=\frac{b_k^2}{(b_k+\mu)^2}=d_k'(\mu),\qquad\delta^\flat_k:=4K_0M^{-3}e_k^{2\varepsilon}w_k.
\end{equation*}
Let the flow data (FL) of \cite[Sections~I.1.8 and~G4.0.1]{QEDIII} hold on $R$ at $\bar m=m$;
the set $\mathrm{Tor}(e)$ to which $R$ belongs, the side $R_{\mathrm{top}}(e)$, the number $N_{\min}(R,e)$ and the stop level $k_*(R)$ are those of \cite[Section~I.6.2]{QEDIII}.
The \emph{flow of $(m;t)$ on $R$} is then the trajectory \cite[Section~G4.0.5]{QEDIII} of a real pair with $m_{\mathrm{in}}=(t-m)\eta$, that is, of a real pair with the lattice mass $t\eta$.
Its entries $m_k=m_k(t;m)$ and $E_k$ do not depend on $\varepsilon^{\mathrm v}_{\mathrm{in}}$.
Let $h_J$ be the largest power of $L$ with $e\,h_J^{1/2}\le e_1$, and $m_\flat:=1/h_J$; then $m_\flat<Le^2/e_1^2$.
The operator $\mathcal K_j:=H_jCC_j^{1/2}$ has the kernel $\mathcal K_j(b,\Upsilon)$, where $b$ is a fine bond of $T_j$.
Put
\begin{equation*}
\kappa_1:=\sup\sum_\Upsilon|\mathcal K_j(b,\Upsilon)|,\qquad\kappa_\infty:=\sup\eta_j^3\sum_b|\mathcal K_j(b,\Upsilon)|,
\end{equation*}
the first supremum being taken over the levels $j\ge0$, the sides and $b$, the second over the levels $j\ge0$, the sides and $\Upsilon$,
where a side is an element of $\mathrm{Cl}=\{\ell:\ell/(LM)\text{ is an odd integer}\ge13\}$ \cite[Section~I.1.1]{QEDIII}, \cite[Section~5.2.3]{QEDI}.
Here $\eta_j:=L^{-j}$ is the mesh of the level $j$ \cite[Section~I.1.1]{QEDIII}, and at the level $k$ it is written $\eta_k=L^{-k}$.
With the place $i$ in a sequence $(N_i)$ or $(N_i,T_i)$ of cutoffs as index, $\eta_i:=L^{-N_i}$ is the spacing at the cutoff $N_i$, as in Proposition~5.6 and Lemma~6.4(a).
With the constant $C_5=108\,c_b^{-1}+6$ of Theorem~5.16, let
\begin{gather*}
C_{\mathsf W}:=1200\,\rho\,C_5\,\kappa_1^2,\qquad C_\natural:=C_{\mathsf W}+3Le_1^{-2},\\
\epsilon_e:=\min\Bigl\{\frac1{12},\,4\bigl(L^{1/2}e\,h^{1/2}\bigr)^{\varepsilon}\Bigr\},\qquad e_7:=\Bigl(\frac{\bar m}{4C_\natural}\Bigr)^{1/2}.
\end{gather*}

\subsection{Background estimates for the mass comparison}
\begin{itemize}
\item[(K22)] (\v{S}id\'ak 1967 \cite{Sidak1967}.)
Let $(Y_1,\dots,Y_d)$ be a centred jointly Gaussian vector, and let $c_1,\dots,c_d>0$.
Then
\begin{equation*}
P\bigl(|Y_i|\le c_i\ \text{for every }i\bigr)\ge\prod_iP\bigl(|Y_i|\le c_i\bigr).
\end{equation*}
\item[(K23)] (Schur.)
For a matrix $K$,
\begin{equation*}
\|K\|^2_{\ell^2\to\ell^2}\le\sup_i\sum_j|K_{ij}|\cdot\sup_j\sum_i|K_{ij}|.
\end{equation*}
If $X$ is a $d\times d$ matrix with $\|X\|<1$, then $|\det(1+X)|\ge(1-\|X\|)^d\ge1-d\|X\|$.
\item[(K24)] Every coclosed one-cochain $A$ of mean zero on $\mathbb Z^3/\ell\mathbb Z^3$ satisfies
\begin{equation*}
\|dA\|^2\ge4\sin^2\frac\pi\ell\,\|A\|^2\ge16\ell^{-2}\|A\|^2.
\end{equation*}
If $\omega$ is a function on a cube of $d_1^3$ sites and $\bar\omega$ is its mean, then $\sum|\omega-\bar\omega|^2\le\frac{d_1^2}4\sum_b|d\omega(b)|^2$.
(Kato.) For $a,b\in\mathbb C^4$ and $|u|=1$ one has $\bigl||a|-|b|\bigr|\le|ua-b|$.
\item[(K28)] (Stone--Weierstrass.) Let $X$ be a locally compact Hausdorff space.
A subalgebra of $C_0(X;\mathbb R)$ which separates the points of $X$ and has no common zero is dense.
\end{itemize}

\begin{named}{Theorem 6.1}
\emph{Two-sided bounds on the tuned lattice masses.}
Suppose that $\bar m\in(0,1]$ and $e>0$ satisfy $e\,h^{1/2}\le e_1$; the data of \cite[Theorem~44]{QEDIII} satisfy this condition.
Then there is a number $N_1(e)<\infty$, one and the same for all such $\bar m$, such that every $N\ge N_1(e)$ satisfies
\begin{equation*}
\bigl(\bar m(1-\epsilon_e)-C_\natural e^2\bigr)\,\eta\le\boldsymbol\mu_N=\bar m\eta+m_{0,N}(e)\le\bar m\,\eta\,\bigl(1+10^{-4}e_\star^{3/4-10\varepsilon}\bigr).
\end{equation*}
\begin{itemize}
\item[(a)] Let $0<e\le e_7$.
Then $e\,h^{1/2}\le e_1$, and
\begin{equation*}
\tfrac12\bar m\,\eta\le\boldsymbol\mu_N\le2\bar m\,\eta\qquad\text{for }N\ge N_1(e).
\end{equation*}
Hence $\liminf_N\boldsymbol\mu_N/\eta\ge\frac12\bar m>0$, and the
tuned assignment satisfies the mass hypothesis of Theorem~A.
Here $e_7=c\,\bar m^{1/2}$, and the number $c=(4C_\natural)^{-1/2}$ does not depend on $\bar m$.
\item[(b)] For $N\ge N_1(e)$ one has $|m_{0,N}(e)|\le(C_\natural e^2+\epsilon_e\bar m)\,\eta$.
Since $\bar m\le1$ and $e\le1$, the right side does not exceed $(C_\natural+4L^{\varepsilon/2})\,e^{\varepsilon}\eta$, uniformly in $\bar m$.
This bound does not reach the power $e^2$, on account of the term $\epsilon_e\bar m$;
in Steps~2 and~3 of the proof this term bounds from above the errors of the slope, which come from (M) of Proposition~6.3\,(c), together with the errors at the junctions, which come from (J) of Proposition~6.3\,(d) combined with (M).
\item[(c)] If $e\le e_7$ and $N\ge N_1(e)$, then $Z_N\le2.1\,Z_0$.
\end{itemize}
\end{named}

\begin{named}{Proposition 6.2}
\emph{A bound from below at every sign of the mass, and the second moment that yields it.}
Let $m\in(0,1]$ with $e\,h(m)^{1/2}\le e_1$, and form the objects of Theorem~5.16 at $\bar m=m$.
Put
\begin{equation*}
v_0:=\log\frac1q,\qquad j_1:=\lceil2p_0\log_Lv_0\rceil,\qquad\mathsf l:=L^{2j_1},\qquad\tilde a_j(b):=L^{-j/2}(\mathcal K_j\zeta_j)(L^{-j}b),
\end{equation*}
so that $\tilde A=\sum_{j<k_\star}\tilde a_j$ by Step~13 of the proof of Theorem~5.16.
Let $\mathcal G$ be the set of those $\zeta$ for which
\begin{equation*}
\sum_{b\cap\mathsf Q\ne\emptyset}\tilde a_j(b)^2\le16L^{-j}\kappa_1^2\mathsf l^3
\end{equation*}
holds for every $j<j_1$ and every block $\mathsf Q$ of side $\mathsf l$ of $\mathbb Z^3/n\mathbb Z^3$.
The set $\mathcal G$ and the blocks are formed only for the $N$ with $2j_1\le k_\star$, which are the large $N$;
for them the levels $j<j_1$ exist, and $\mathsf l$ divides $n=L^{k_\star}\ell_\star$.
\begin{itemize}
\item[(a)] There is a number $N_A(e,m)$, defined in Step~10 of the proof, which is a function of $(e,h(m),\mu_\star)$ and of constants, such that $\varphi_N(m)\ge-C_{\mathsf W}e^2$ for all $N\ge N_A(e,m)$.
By comparison, Theorem~5.16\,(d) gives the lower bound $-2\rho C_5e^{2-\varepsilon}\eta^{-\varepsilon/2}$ for $\varphi_N(m)$.
\item[(b)] Let $2j_1\le k_\star$.
On $\mathcal G\cap\prod_j\mathrm{Box}_j$, every block $\mathsf Q$ satisfies
\begin{equation*}
\sum_{b\cap\mathsf Q\ne\emptyset}\tilde A(b)^2\le200\,\kappa_1^2\mathsf l^3;
\end{equation*}
this bounds a second moment by a number of the order one, whereas the bound on the supremum is $q^{-\varepsilon}$.
Let $\bar\mu$ denote the zero-flux law of \cite[Proposition~31.1]{QEDIII}, taken on the torus with $n=R_{\mathrm{top}}/\eta$ sites per side and at the lattice mass $\varphi_N(m)\eta$.
Then
\begin{equation*}
\bar\mu(\mathcal G^c)\le j_1n^3\exp\Bigl(-\frac{2\kappa_1}{\kappa_\infty}v_0^{6p_0}\Bigr)\longrightarrow0\qquad(N\to\infty).
\end{equation*}
\end{itemize}
\end{named}

\begin{named}{Proposition 6.3}
\emph{Running mass versus bare mass; slope one; junctions.}
Assume $e\,h^{1/2}\le e_1$, where $h=h(m)$ in (a) and (b).
\begin{itemize}
\item[(a)] \emph{One parameter.}
Suppose that (FL) holds on $R$ at $\bar m=m$, that $k\le k_*(R)$, and that the flows of $(m;t)$ and of $(m;t')$ on $R$ satisfy $|m_j|\le w_j$ for $j<k$.
Then
\begin{equation*}
m_k(t;m)-m_k(t';m)=a_k(t-t')(1+\theta),\qquad|\theta|\le\tfrac1{15}e_k^{\varepsilon}.
\end{equation*}
Thus, on the set of those $t$ whose flows stay in these windows, the running mass is a strictly increasing function of the bare mass.
\item[(b)] \emph{Reading.}
Suppose that (FL) holds on $R$ at $\bar m=m$, that $1\le k\le k_*(R)$, and that $|m_j(t;m)|\le w_j$ for $j<k$.
Then there is a complex number $\mathsf z_k$, defined in Step~2 of the proof, which depends on $(e,N,R,k,t)$ alone and not on $m$:
it is one and the same number for every $m$ with the two properties just required, although for each such $m$ the flow of $(m;t)$ is formed with the primitives of (FL) at $(m,e,R,N)$;
these primitives depend on $m$ through $\mu_k=m\,a_k$ and through the scheme of $h(m)$, that is, through $R_{\mathrm{top}}(e;h(m))$ and, on a torus $R<R_{\mathrm{top}}(e;h(m))$, through the projections with $\mathrm{Jr}=R_{\mathrm{top}}/R$ \cite[Section~G4.0.1]{QEDIII};
in Step~7 of the proof, for instance, which proves (d) by the argument of Step~5, the flows of $(m_b;t)$ and of $(m_c;t)$ with $t=\varphi_N(m_c)$, where $m_b\in I_{Lh}$ and $m_c\in I_h$, are taken on one torus, $R=R_{\mathrm{top}}(e;h)$, which satisfies $R\le R_{\mathrm{top}}(e;Lh)$ and belongs to the tori of the scheme of $Lh$ as well;
they are formed with the schemes of $Lh$ and of $h$, respectively;
apart from $(e,N,R,k,t)$, the number $\mathsf z_k$ involves only constants and the spinor $u=(1,0,0,0)$ of the statement of Theorem~5.16.
This number satisfies
\begin{equation*}
\mathsf z_k=d_k(\mu_k)+Z_k(\mu_k)\,\bigl(m_k(t;m)-\varkappa_k\bigr),\qquad|\varkappa_k|\le\delta_k^\flat.
\end{equation*}
\item[(c)] \textbf{(M)}
Let $N\ge N_{\min}(R_{\mathrm{top}}(e),e)$, and let $m'<m''$ be points of $I_h$.
Then
\begin{equation*}
\bigl|\varphi_N(m'')-\varphi_N(m')-(m''-m')\bigr|\le e_\star^{\varepsilon}\,\bigl[(m''-m')+w_\star/a_\star\bigr].
\end{equation*}
In particular $\varphi_N(m'')>\varphi_N(m')$ whenever $m''-m'\ge2e_\star^\varepsilon w_\star/a_\star$.
\item[(d)] \textbf{(J)}
Suppose, in addition, that $e\,(Lh)^{1/2}\le e_1$.
Let $N\ge N_{\min}(R_{\mathrm{top}}(e;h),e)$, let $m_b:=\frac1{Lh}$ be the upper end point of $I_{Lh}$, and let $m_c\in I_h$ satisfy $m_c-m_b\le\frac{w_\star}{2a_\star}$.
Then
\begin{equation*}
\varphi_N(m_c)-\varphi_N(m_b)\ge-\frac{w_\star}{6a_\star}.
\end{equation*}
\end{itemize}
\end{named}

\begin{proof}[Proof of Proposition~6.3]
\emph{Step 1} (one density under two names).
By \cite[Section~I.5.10, (I.L0.22), (I.L0.7), (I.L0.2)]{QEDIII}, in the form written out for the flow in \cite[Chapter~29, P1, (Q4)]{QEDIII}, the density $\mathscr F_0[\varepsilon^{\mathrm v}_{\mathrm{in}},\allowbreak m_{\mathrm{in}},\allowbreak 0](A_0;\Psi)$ equals a positive number multiplied by
\begin{equation*}
\exp\bigl\{-\langle\bar\Psi,(D_{A_0}+t\eta)\Psi\rangle^{\mathrm{ch}}\bigr\}.
\end{equation*}
Each map $\mathbf T_j$ is linear, and the objects from which it is built, namely $c_j$, $\mathrm{Box}_j$, $\mu_I$, $b_L$, $\omega_j$, $Q(\mathscr B_j)$, $A^{\min}_j$ and $\mathrm{Sq}_j$, are formed with $(e_j,\ell_j)$ and constants \cite[Sections~I.3.2--I.3.4, I.5.7]{QEDIII}, \cite[Section~5.2.5]{QEDI};
thus neither a fermion mass nor a primitive enters $\mathbf T_j$.
Indeed, $H_j$ (which occurs in $\mathscr B_j=H_jA^{\min}_j(A')$), $C_j^{1/2}$, $\mathrm{Sq}_j$, $A^{\min}_j$ and $\omega_j$ are the exact photon objects of \cite[Section~5.2.5]{QEDI};
they are formed with $(j,\ell_j,L)$ and with the basis $\{e_i^0\}$, which is fixed once for each $L$.
Assume that $|m_j|\le w_j$ for $j<k$.
At the levels $j<k$, \cite[Proposition~28.1]{QEDIII} then furnishes the hypotheses of \cite[Proposition~24.1\,(c)]{QEDIII} at $\mathscr C'=0$, the primitives being those of the flow data (FL) at $(m,e,R,N)$;
these primitives are determined by the quadruple $(m,e,R,N)$ \cite[Sections~I.1.8, G4.0.1; Proposition~26.6\,(e)]{QEDIII}, and together with them the data of the levels $j<k$ are (OS) \cite[Proposition~26.6\,(b)]{QEDIII}.
Combined with \cite[Proposition~5.38\,(i), (ii)]{QEDI}, this gives
\begin{equation*}
\mathscr F_k\bigl[\varepsilon^{\mathrm v,t}_k,m_k,E_k\bigr]=\mathbf T_{k-1}\cdots\mathbf T_0\,\mathscr F_0\qquad\text{on }\mathrm{Dom}_k .
\end{equation*}
Consequently $\mathscr F_k(A;\cdot)$ is, up to a positive factor, a function of $(e,N,R,k,t)$ and of $A$ only.

\emph{Step 2} (the constant mode).
Let $\Omega_k$ be the constant bond function with $e_k\Omega_{k,\mu}\ell_k=\pi$, let
\begin{equation*}
\Psi_u(y):=\ell_k^{-3/2}\exp(-ie_k\Omega_k\cdot y)\,u,\qquad f(c,g):=\ell_k^{-3/2}\exp(-ie_k\Omega_k\cdot c)\,u,
\end{equation*}
and define the number $\mathsf z_k$ by the condition that the component of degree $(1,1)$ of
\begin{equation*}
\log\frac{\mathscr F_k(\Omega_k;\Psi)}{\mathscr F_k(\Omega_k;0)},
\end{equation*}
evaluated at $(\overline{\Psi_u},\Psi_u)$, equals $-\mathsf z_k$.
As in Step~7 of the proof of Theorem~5.16, whose definitions \cite[(I.O8); Sections~I.3.2, I.3.3, I.3.5]{QEDIII} are the same at every level $k\ge1$ (compare \cite[Proposition~30.3\,(f)]{QEDIII}), one finds the following.
One has $H_k\Omega_k=\Omega_k\in\frac14\tilde R_k$ \cite[Propositions~5.31, 5.38\,(i)]{QEDI}, \cite[Section~I.5.7, paragraph ``Settings'']{QEDIII}, so that $\Omega_k$ lies in the set over which the supremum in \cite[(k6)]{QEDIII} is taken.
Here $H_k\Omega_k$ is the photon minimiser applied to $\Omega_k$, whereas $H_k(A)$, whose argument is a connection, is the fermion operator \cite[Section~I.5.7, first paragraph]{QEDIII};
by \cite[Section~I.3.5]{QEDIII},
\begin{equation*}
H_k(A)=b_kS_k(A)JQ_k^T(-A),\qquad\bar H_k(A)=\mathsf C^{-1}H_k(-A)\mathsf C,
\end{equation*}
where $\mathsf C=i\gamma_2$ is the real matrix of \cite[Section~I.2.1, (I.c3)]{QEDIII}.
By \cite[Section~I.5.4, paragraph ``Fields to be inserted'']{QEDIII}, the fields inserted into an activity at the connection $\Omega_k$ are
\begin{equation*}
\psi=H_k(\Omega_k)\Psi,\qquad\chi=\delta_{\alpha,\Omega_k}\psi,\qquad\bar\psi=\bar H_k(\Omega_k)\bar\Psi,\qquad\bar\chi=\delta_{\alpha,-\Omega_k}\bar\psi .
\end{equation*}
The maps $\delta_{\alpha,\pm\Omega_k}$ are the H\"older quotients of \cite[(I.1.7)]{QEDIII}, and their exponent $\alpha$ is the H\"older order;
this is a constant, chosen in the line (C1) of \cite[Section~I.8.2]{QEDIII} as at stage~2 of \cite[Definition~5.97]{QEDI}, with $1-\bar\varepsilon<\alpha<1$, where $\bar\varepsilon=\frac1{200}$ is fixed in the line (C0) there.
Furthermore $D_{\Omega_k}f=0$, $D_k(\Omega_k)\Psi_u=d_k(\mu_k)\Psi_u$ and $H_k(\Omega_k)\Psi_u=Z_k^{1/2}f$.
The vector $\mathsf C\overline{\Psi_u}$ and the function $\mathsf C\bar f$ are $\Psi_u$ and $f$ formed at $-\Omega_k$ with the spinor $\mathsf C\bar u$ in the place of $u$;
the same computation at $-\Omega_k$ therefore yields $H_k(-\Omega_k)\mathsf C\overline{\Psi_u}=\allowbreak Z_k^{1/2}\mathsf C\bar f$, that is, $\bar H_k(\Omega_k)\overline{\Psi_u}=Z_k^{1/2}\bar f$.
Since $U_{\Omega_k}(x,x')f(x')=f(x)$ for all pairs $(x,x')$, it follows that $\delta_{\alpha,\Omega_k}f=0=\delta_{\alpha,-\Omega_k}\bar f$ \cite[(I.1.6), (I.1.7)]{QEDIII}.
At $(\bar\Psi,\Psi)=(\overline{\Psi_u},\Psi_u)$ the fields $\psi$, $\chi$, $\bar\psi$, $\bar\chi$ therefore take the values
\begin{equation*}
\psi=Z_k^{1/2}f,\qquad\chi=Z_k^{1/2}\delta_{\alpha,\Omega_k}f=0,\qquad\bar\psi=Z_k^{1/2}\bar f,\qquad\bar\chi=Z_k^{1/2}\delta_{\alpha,-\Omega_k}\bar f=0.
\end{equation*}
The definition of the density in \cite[Section~I.5.7, paragraph ``The density and the block transformation'']{QEDIII}, with these fields inserted, then gives
\begin{equation*}
\mathsf z_k=d_k+Z_km_k-Z_k\sum_X\int\!\!\int\bar f^{T}E_{k,20}(X;\Omega_k)\,f .
\end{equation*}

\emph{Step 3} (proof of (b)).
From \cite[Section~I.5.1, paragraphs ``The norms of the scheme'' and ``Matrix notation'']{QEDIII} and from the display following (k7) in \cite[Section~I.5.4]{QEDIII} one has
\begin{equation*}
e_k^{-1/2}\int\!\!\int\sum_{\beta\beta'}\bigl|E_{20}(X)_{\beta\beta'}\bigr|=\bigl\|E^{[2,0]}(X)\bigr\|_{h_k}\le\|E(X)\|_{h_k},
\end{equation*}
and $|f|^2=\ell_k^{-3}$.
Moreover, \cite[(k6)]{QEDIII} holds with a rate $\kappa\ge7\kappa_0+7\ge\kappa_0$ \cite[Section~I.8.2, line~(C2)]{QEDIII};
by \cite[Proposition~5.58\,(t2)]{QEDI}, \cite[Chapter~24, Section~0.6, (t2)]{QEDIII}, applied to the whole torus, which is the union of $(\ell_k/M)^3$ $M$-cubes,
\begin{equation*}
\sum_{X\in D_k}\exp\bigl(-\kappa_0\,d_M(X)\bigr)\le K_0\Bigl(\frac{\ell_k}M\Bigr)^3 ;
\end{equation*}
and, for the flow of $(m;t)$, $\|E_k\|_k\le4e_k^{1/4-6\varepsilon}$ by \cite[Proposition~28.1]{QEDIII}.
It follows that the sum in the last term of Step~2,
\begin{equation*}
\varkappa_k:=\sum_X\int\!\!\int\bar f^{T}E_{k,20}(X;\Omega_k)\,f,
\end{equation*}
satisfies
\begin{equation*}
|\varkappa_k|\le4K_0M^{-3}e_k^{3/4-6\varepsilon}=\delta^\flat_k .
\end{equation*}
By Step~1, $\mathsf z_k$ is determined by $(e,N,R,k,t)$, by the spinor $u$ and by constants, the basis $\{e_i^0\}$ among them, so that it does not depend on $m$;
the identity of Step~2, read with $m_k=m_k(t;m)$ and with $\varkappa_k$ formed from $E_k$ of the flow of $(m;t)$, and the bound on $\varkappa_k$ prove (b).
The number $\delta_\star$ of Step~9 of the proof of Theorem~5.16, formed with the bound on $H_{k_\star}$ in the supremum norm, equals $C_0^2e_\star^{-2\varepsilon}\delta^\flat_\star$;
here $H_k(\Omega_k)\Psi_u$ is known exactly, and this makes the accuracy $\delta^\flat_k$ with which $\mathsf z_k$ determines $m_k(t;m)$ finer than the window $w_k$.

\emph{Step 4} (proof of (a)).
Let $t\ne t'$, and let primed letters refer to the flow of $(m;t')$.
Put
\begin{equation*}
x_j:=m_j-m_j',\qquad Y_j:=\|E_j-E_j'\|_j ;
\end{equation*}
one has $Y_0=0$.
The recursion of the flow \cite[Section~I.5.7, paragraph ``Flows'']{QEDIII}, in whose difference the term $LE^{\det}_j$ cancels, and \cite[Propositions~48.2\,(a), (b), 48.5\,(a), (b), 24.1\,(e)]{QEDIII} give
\begin{gather*}
x_{j+1}=L\bigl(x_j+m(E_j-E_j')\bigr),\qquad\bigl|m(E_j-E_j')\bigr|\le c_{\mathrm{fl}}\,e_j^{1/2}\,Y_j,\\
Y_{j+1}\le\bigl(c_{\mathrm{fl}}L^{-1/4+2\varepsilon}+e_j^{\varepsilon}\bigr)Y_j+e_j^{-1/2+2\varepsilon}|x_j| .
\end{gather*}
From $c_{\mathrm{fl}}L^{-3/8+6\bar\varepsilon}\le\frac1{16}$ and $L\ge5$ \cite[Proposition~14.2\,(b), rows (L2), (L5)]{QEDIII}, and from $(c_{\mathrm{fl}}+2)e_j^{\varepsilon}\le\frac1{16}$ \cite[(fl)]{QEDIII}, induction on $j$ shows that $x_j\ne0$ and
\begin{equation*}
u_j:=\frac{e_j^{1/2}\,Y_j}{|x_j|}\le e_j^{2\varepsilon},
\end{equation*}
the inductive step being
\begin{equation*}
|x_{j+1}|\ge\tfrac{15}{16}L|x_j|,\qquad u_{j+1}\le\tfrac{16}{15}L^{-3/4}\bigl(c_{\mathrm{fl}}L^{-1/4+2\varepsilon}+e_j^{\varepsilon}+1\bigr)e_j^{2\varepsilon}.
\end{equation*}
Hence $x_{j+1}=Lx_j(1+\theta_j)$ with $|\theta_j|\le c_{\mathrm{fl}}e_j^{2\varepsilon}$, and by \cite[(fl)]{QEDIII}
\begin{equation*}
\mathsf s:=\sum_{j<k}|\theta_j|\le\tfrac1{16}e_k^{\varepsilon}.
\end{equation*}
Because $\mathsf s\le\frac1{16}$, one has
\begin{equation*}
\Bigl|\prod_{j<k}(1+\theta_j)-1\Bigr|\le\exp(\mathsf s)-1\le\tfrac{16}{15}\,\mathsf s .
\end{equation*}
As $x_0=(t-t')\eta$ and $L^k\eta=a_k$, this proves (a).

\emph{Step 5} (one link).
Let $m_1<m_2=m_1+s$ be points of $I_h$ with $s\le\frac{w_\star}{2a_\star}$;
put $t_2:=\varphi_N(m_2)$, $R=R_{\mathrm{top}}(e)$ and $\mu_k^{(i)}:=m_ia_k$.
By \cite[Proposition~28.2]{QEDIII}, the flow of $(m_2;t_2)$ satisfies $|m_k|\le\frac18w_k$ and $m_{k_\star}=0$.
The flow of $(m_1;t_2)$ is estimated by induction on $k$.
Let $|m_j(t_2;m_1)|\le w_j$ for $j<k$.
Both flows have the same $t=t_2$, so that (b) attaches to them one and the same number $\mathsf z_k$;
with the numbers $\varkappa_k^{(1)}$, $\varkappa_k^{(2)}$ that (b) gives for the two flows, it follows that
\begin{equation*}
Z_k(\mu_k^{(1)})\bigl(m_k(t_2;m_1)-\varkappa_k^{(1)}\bigr)=d_k(\mu_k^{(2)})-d_k(\mu_k^{(1)})+Z_k(\mu_k^{(2)})\bigl(m_k(t_2;m_2)-\varkappa_k^{(2)}\bigr).
\end{equation*}
Since $Z_k=d_k'$ is decreasing, $0\le d_k(\mu_k^{(2)})-d_k(\mu_k^{(1)})\le s\,a_kZ_k(\mu_k^{(1)})$;
since $a_k/w_k$ increases with $k$, one obtains
\begin{equation*}
|m_k(t_2;m_1)|\le\bigl(\tfrac12+\tfrac18\bigr)w_k+2\delta_k^\flat\le\tfrac34w_k .
\end{equation*}
For $k=0$ the same bound follows from the relation $m^{(1)}_{\mathrm{in}}=m^{(2)}_{\mathrm{in}}+s\eta$ between the initial masses of the two flows.
At $k=k_\star$, where $m_{k_\star}(t_2;m_2)=0$, the identity above gives
\begin{equation*}
m_{k_\star}(t_2;m_1)=s\,a_\star\vartheta'+r,\qquad\vartheta'=\frac{b_\star+\mu_\star^{(1)}}{b_\star+\mu_\star^{(2)}}\in[1-w_\star,1],\qquad|r|\le2\delta^\flat_\star .
\end{equation*}

\emph{Step 6} (proof of (c)).
Clause (a), applied to the flow of $(m_1;t_2)$ and to the tuned pair of $m_1$, gives
\begin{equation*}
m_{k_\star}(t_2;m_1)=a_\star\bigl(\varphi_N(m_2)-\varphi_N(m_1)\bigr)(1+\theta).
\end{equation*}
By the third inequality of \cite[(T2)]{QEDIII}, $w_\star\le e_\star^{\varepsilon}\bigl(e_\star^{1/4-9\varepsilon}\bigr)^3\le\frac1{512}e_\star^{\varepsilon}$.
With Step~5 it follows that
\begin{equation*}
\bigl|\varphi_N(m_2)-\varphi_N(m_1)-s\bigr|\le\tfrac18e_\star^{\varepsilon}s+2.2\,\frac{\delta^\flat_\star}{a_\star}.
\end{equation*}
Let $[m',m'']$ be divided into $\lceil2a_\star(m''-m')/w_\star\rceil$ intervals of equal length, each of which is a link as in Step~5.
By \cite[(fl)]{QEDIII},
\begin{equation*}
\frac{\delta^\flat_\star}{w_\star}=4K_0M^{-3}e_\star^{2\varepsilon}\le10^{-4}e_\star^{\varepsilon},
\end{equation*}
and the sum over the links of the inequality for one link yields (c).

\emph{Step 7} (proof of (d)).
Let $R:=R_{\mathrm{top}}(e;h)$;
then $R=Lh\cdot L^T$, where $L^T:=Mr_K/L\ge1$.
Applied to $Lh$, \cite[(Sc0)]{QEDIII} gives $v_K^r<Lv_{K+1}^r\le Lr_{K+1}$, and $Lr_{K+1}$ is a power of $L$;
hence $r_K\le Lr_{K+1}$ and $R\le R_{\mathrm{top}}(e;Lh)$.
Thus $R$ belongs to the tori of the scheme of $Lh$, while the numbers $N_{\min}(R,e)$, $k_*(R)=k_\star$, $\ell_k$ and $e_k$ do not depend on $h$ \cite[Sections~I.6.1, I.6.2]{QEDIII}.
By \cite[Section~G4.0.2, (Sc2)]{QEDIII}, $N\ge N_{\min}(R,e)\ge N_{\min}(R_{\mathrm{top}}(e;Lh),e)$, so that the shift entering $\varphi_N(m_b)$ is the tuned one.
By \cite[Proposition~28.3\,(a), (b)]{QEDIII}, the tuned pair of $m_b$ has $|m_k|\le\frac18w_k$ and is controlled on $R$.
The argument of Step~5, with $m_b$ and $m_c$ in the places of $m_1$ and $m_2$, shows that the flow of $(m_b;\varphi_N(m_c))$ on $R$ is controlled and has $m_{k_\star}\ge-2\delta^\flat_\star$.
Clause (a) on $R$ therefore gives
\begin{equation*}
a_\star\bigl(\varphi_N(m_c)-\varphi_N(m_b)\bigr)(1+\theta)\ge-2\delta^\flat_\star-\tfrac18w_\star,
\end{equation*}
and (d) follows.
\end{proof}

\begin{proof}[Proof of Proposition~6.2]
Throughout the proof put
\begin{gather*}
\mathrm{mup}:=\prod_{j<k_\star}\mathbf 1_{\mathrm{Box}_j}\,\mu_I,\qquad g_\star:=b_\star^{-1}+\mu_\star^{-1}=d_\star^{-1},\qquad g_-:=b_-^{-1}+\mu_\star^{-1},\\
c_-:=\mu_\star\,b_-^2\,(b_-+\mu_\star)^{-2},\qquad\varkappa_D:=\bigl(16\,\ell_\star^3\,g_-\bigr)^{-1},\qquad r_1:=\frac{\mu_\star^2}{24}\,\min\Bigl\{\varkappa_D,\frac{c_-}{28}\Bigr\},\qquad\boldsymbol\mu:=\varphi_N(m)\,\eta .
\end{gather*}
Since $b_-\le b_\star$ and the map $b\mapsto\mu b^2(b+\mu)^{-2}$ is increasing, one has $g_\star\le g_-$ and $c_-\le c_\star$.
The numbers $g_-$, $c_-$, $\varkappa_D$ and $r_1$ depend on $(e,h,\mu_\star)$ but not on $N$, whereas $g_\star$, $d_\star$, $c_\star$ and $Z_0$ depend on $N$ through $b_\star$.

\emph{Step 1} (Theorem~5.16 at a general block field).
Denote by $V_{\mathrm t}(T^0_{k_\star})$ the space of the coclosed bond functions with mean zero on the unit lattice $T^0_{k_\star}$ of the level $k_\star$ \cite[Section~G4.0.6]{QEDIII}, and put $p_\star:=p_{k_\star}=v_\star^{p}$, where $v_\star=-\log e_\star$ and $p$ is the integer of the line (C7) of \cite[Section~I.8.2]{QEDIII}, as in \cite[Section~I.6.1]{QEDIII}.
Let $\theta\in\mathbb R^3$, let $A\in V_{\mathrm t}(T^0_{k_\star})$ satisfy $|dA|_\infty\le p_\star$, and put $A':=A+c_\theta$, where $c_\theta$ is the constant bond function with the twist $\theta$.
Then $A'\in\mathrm{Dom}_{k_\star}$ by \cite[Proposition~32.1]{QEDIII}.
The objects $A_j$, $\mathcal R$, $B$, $\mathcal D$, $\mathcal D_0$, $\psi$ and $\Phi$ are formed as in Theorem~5.16, with $A_{k_\star}:=A'$ and with the same vector $\Psi_u$, and $z:=(\Psi_u,\Phi)$.
Let $\mathsf P^{A'}$ and $\mathsf P^{0,A'}$ be the measures proportional to $\det B\,\mathrm{mup}$ and to $\det\mathcal D_0\,\mathrm{mup}$, respectively, and let $\mathsf E^{A'}$ be the mean under $\mathsf P^{A'}$.
Put $\mathscr A':=H_{k_\star}A'$;
by \cite[Proposition~5.31]{QEDI}, $\mathscr A'=H_{k_\star}A+c_\theta$.
The following parts of the proof of Theorem~5.16 are carried out at $A_{k_\star}=A'$:
Steps~1 to~6;
Step~8, with $\mathscr A'$ in the place of $H_{k_\star}\Omega=\Omega$;
in Step~9, the bound on $\sum_XE_{k_\star}$ in the norm $|\cdot|_{\mathrm{aw}_2}$, the bounds on $K_E$ and on $\mathcal A_4$ in terms of $\delta_\star$ and $\delta_2$, and the two inequalities for $\ell_\star^3\delta_\star$ and $\delta_2$ derived from \cite[(T2)]{QEDIII};
Step~11;
in Step~12, the permutation $a\mapsto\hat a$ of the indices, the bound $\mathsf E_N\|\delta\mathcal D\Psi_u\|^2\le\delta_2$ and the identity obtained by averaging (SR);
Step~13, with $\mathscr A_{k_\star}=\mathscr A'$.
These arguments are valid with any real element of $\mathrm{Dom}_{k_\star}$ in the place of $\Omega$ (\cite[Chapter~29, Lemma~2\,(t4)]{QEDIII}; the bounds of Step~9 of that proof are suprema over $\tilde R_{k_\star}$ \cite[Proposition~20.1\,(c)]{QEDIII}), and with any antiperiodic unit vector $\Psi_u$ such that $\Sigma_3u=u$ and $|\Psi_u(y)|^2=\ell_\star^{-3}$.
Step~7 of that proof establishes the relation $D_{k_\star}(\Omega)\Psi_u=d_\star\Psi_u$;
on this relation rest the bounds of its Step~9 on $\bar z=(\Psi_u,\mathsf E_N\mathcal D\Psi_u)$ and on $\|\mathsf E_N\mathcal D\Psi_u\|$, its Step~10, and, in its Step~12, the estimate comparing the right side of the averaged identity (SR) with $c_\star$, the bound $\mathsf E_N\|\Phi\|^2\le2$ and the inequality $|1-\bar z/b_\star|^2\ge(Z_0^{1/2}-2\delta_\star)^2$.
Step~8 below takes the place of these parts and of Step~7 of that proof.
Consequently, whatever the sign of $\boldsymbol\mu$, the following hold.
\begin{itemize}
\item[(i)] One has $\det B\ge0$;
the objects $\mathcal D$, $\psi$, $\Phi$ and $z$ are defined off the null set $\{\det B=0\}$, and the relations below are understood off this set.
Both measures can be normalised, as $\mathscr F_{k_\star}(A';0)\ne0$ and $Y_*(A')>0$ \cite[Proposition~32.1]{QEDIII}.
Moreover $\det\mathcal D\ge0$, and with the constant $\mathsf C_{\det}>0$ of Step~5 of the proof of Theorem~5.16, which is denoted $C_N$ there,
\begin{equation*}
\mathscr F_{k_\star}(A';\Psi)=\mathsf C_{\det}\int\mathrm d\,\mathrm{mup}\;\det B\,\exp\bigl(-\langle\bar\Psi,\mathcal D\Psi\rangle\bigr),\qquad y_0\bigl(\mathrm{Af}_0(A',\zeta)\bigr)=\mathsf C_{\det}\,\det B\,\det\mathcal D .
\end{equation*}
In the second identity, $\mathrm{Af}_0(A',\zeta)$ is the last field $A_0$ of the recursion $A_j=A^{\min}_j(A_{j+1})+\mathrm{Sq}_j\zeta_j$ begun at $A_{k_\star}=A'$ \cite[Sections~I.5.7, G4.0.6]{QEDIII}, and $y_0$ is the function
\begin{equation*}
y_0(\breve A):=\exp\bigl(-\varepsilon^{\mathrm v}_{\mathrm{in}}n^3\bigr)\,\frac{\det(D_{\breve A}+\boldsymbol\mu)}{\det(D_0+m\eta)}
\end{equation*}
of a real one-cochain $\breve A$ on $\mathbb Z^3/n\mathbb Z^3$;
it is the function $\mathrm{wdn}$ of \cite[Chapter~29, (29.1)]{QEDIII}, and on $\mathrm{Dom}_0$ it agrees with $y_0[\varepsilon^{\mathrm v}_{\mathrm{in}},m_{\mathrm{in}},0]$ \cite[Section~I.5.7; Chapter~29, Lemma~1\,(iv)]{QEDIII}.
For $A_0=\mathrm{Af}_0(A',\zeta)$ one has $\det\mathcal D_0=\det(D_{A_0}+\boldsymbol\mu)$;
the second identity of the display is therefore the relation $\mathsf C_{\det}\det B\cdot\det\mathcal D=\allowbreak \exp(-\varepsilon^{\mathrm v}_{\mathrm{in}}n^3)\det\mathcal D_0/\det(D_0+m\eta)$ obtained in Steps~5 and~6 of the proof of Theorem~5.16.
\item[(ii)] With $\delta\mathcal D:=\mathcal D-\mathsf E^{A'}\mathcal D$ one has
\begin{equation*}
\mathsf E^{A'}\mathcal D=D_{k_\star}(\mathscr A')-K_E,
\end{equation*}
where the moduli of the entries of $K_E$ have a sum not exceeding $\ell_\star^3\delta_\star$;
furthermore $\mathsf E^{A'}\|\delta\mathcal D\Psi_u\|^2\le\delta_2$, and summation of the identities of Steps~8 and~12 of that proof over all indices yields $\mathsf E^{A'}\|\delta\mathcal D\|^2_{\mathrm{HS}}\le2\,\allowbreak \ell_\star^3\delta_2$.
\item[(iii)] The identity (SR) of Theorem~5.16\,(c) holds.
\item[(iv)] One has $A_0=H_\theta+\mathbf a+\tilde A-d\lambda$, where $H_\theta$ is the constant one-cochain of the twist $\theta$, as in \cite[Section~I.6.6]{QEDIII}, and $\mathbf a:=(H_{k_\star}A)_{L^{k_\star}}$;
moreover $q|\mathbf a|_\infty=e_\star L^{-k_\star}|H_{k_\star}A|_\infty$.
\end{itemize}

\emph{Step 2} (the complement of an event costs at most $\frac14$ and carries no weight $\|\psi\|^2$).
Put
\begin{equation*}
W:=\frac\rho2\,(\psi,\nabla^*_{A_0}\nabla_{A_0}\psi),\qquad c(A'):=\mathsf E^{A'}\bigl[\operatorname{Re}z-b_\star^{-1}\|\Phi\|^2\bigr].
\end{equation*}
As $\operatorname{Re}z\le\|\Phi\|$ and $x-x^2/b_\star\le b_\star/4\le\frac14$ for real $x$, the identity (SR) gives, for every Borel set $\mathcal B$,
\begin{equation*}
c(A')\le L^{k_\star}\,\mathsf E^{A'}\bigl[\boldsymbol\mu\|\psi\|^2+W;\mathcal B\bigr]+\tfrac14\,\mathsf P^{A'}(\mathcal B^c).
\end{equation*}

\emph{Step 3} (Kato's inequality and a Poincar\'e inequality).
Let $\mathsf Q$ be a cube of $\mathsf l^3$ photon sites, let $f\in\mathcal H_0$, and let $\tau\ge0$ be a function on the sites over $\mathsf Q$, whose mean is $\langle\tau\rangle_{\mathsf Q}$.
In this step $\|\cdot\|_{\mathsf Q}$ denotes the norm of $\mathcal H_0$ restricted to the $48\,\mathsf l^3$ sites over $\mathsf Q$;
the symbols $P_{\mathsf Q}$ and $\langle\cdot\rangle_{\mathsf Q}$ both stand for averaging over these sites;
and $(f,\nabla^*_{A_0}\nabla_{A_0}f)_{\mathsf Q}$ collects those terms of \cite[(I.O2)]{QEDIII} that belong to bonds with both end points over $\mathsf Q$.
The classical fact (K24) of Section~6, the bond formula \cite[(I.O2)]{QEDIII} and \cite[Proposition~1.4]{QEDIII}, applied to $(|f|,0,0,0)$, yield
\begin{equation*}
\bigl\||f|-P_{\mathsf Q}|f|\bigr\|^2_{\mathsf Q}\le C_{\mathrm{Po}}\,\mathsf l^2\,(f,\nabla^*_{A_0}\nabla_{A_0}f)_{\mathsf Q}.
\end{equation*}
Together with $|f|^2\le2(P_{\mathsf Q}|f|)^2+2(|f|-P_{\mathsf Q}|f|)^2$, it follows that
\begin{equation*}
\|\tau^{1/2}f\|^2_{\mathsf Q}\le2\,\langle\tau\rangle_{\mathsf Q}\,\|f\|^2_{\mathsf Q}+2\,|\tau|_\infty\,C_{\mathrm{Po}}\,\mathsf l^2\,(f,\nabla^*_{A_0}\nabla_{A_0}f)_{\mathsf Q}.
\end{equation*}

\emph{Step 4} (the good set $\mathcal G$; the first half of (b)).
In the suprema that define $\kappa_1$ and $\kappa_\infty$, the side ranges, as in \cite[Section~I.1.1]{QEDIII}, over the class $\mathrm{Cl}$ of those $\ell$ for which $\ell/(LM)$ is an odd integer not smaller than $13$ \cite[Section~5.2.3]{QEDI}.
The letter $\alpha$ stands for the H\"older order of \cite[Section~2.7]{QEDI}, the exponent of the H\"older quotients $\delta_{\alpha,\cdot}$ of \cite[(I.1.7)]{QEDIII};
it is the constant $\alpha\in(1-\bar\varepsilon,1)$, $\bar\varepsilon=\frac1{200}$, of the line (C1) of \cite[Section~I.8.2]{QEDIII}.
The numbers $\kappa_1$ and $\kappa_\infty$ are finite, and each of them is at most a number that depends on $L$ and $\alpha$ but not on $e$, $m$, $N$, $j$ or the side.
Indeed, for $j\ge0$ and for sides in $\mathrm{Cl}$, $|H_j(b,b')|$ and $|C_j^{1/2}(\Upsilon,\Upsilon')|$ are at most $C_P\exp(-\gamma_Pd)$, where $C_P$ and $\gamma_P$ depend on $L$ and $\alpha$, and $C$ has range at most $2L$ and row sums at most $c_C=c_C(L)$ \cite[Section~5.4, Propositions~5.32, 5.34, 5.35]{QEDI};
the sides $\ell_j(R_{\mathrm{top}})$, $j<k_\star$, at which these bounds are applied lie in $\mathrm{Cl}$ \cite[Section~G4.0.2, (Sc4)]{QEDIII};
by the locality stated in \cite[Proposition~5.35]{QEDI}, each column of $C$ has at most $C(L)$ non-zero entries, each of modulus at most $c_C$;
and the lattice sums of $\exp(-\gamma_Pd)$ over the unit bonds and over the fine bonds with the weight $\eta_j^3$ are bounded by numbers depending on $\gamma_P$, while the sum over $\Upsilon$ is bounded by a number depending on $\gamma_P$ and $L$.
For large $N$ one has $2j_1\le k_\star$, so that the blocks $\mathsf Q$ tile $\mathbb Z^3/n\mathbb Z^3$;
at most $4\,\mathsf l^3$ bonds meet a block $\mathsf Q$.
On $\prod_j\mathrm{Box}_j$ every coordinate satisfies $|\zeta_j(\Upsilon)|\le p_{0,j}=v_j^{p_0}$, where $v_j=-\log e_j$ \cite[Section~5.2.5]{QEDI}, \cite[Section~I.6.1]{QEDIII};
since $e_j=qL^{j/2}\ge q$, one has $v_j\le v_0$, whence $|\tilde a_j|_\infty\le\kappa_1L^{-j/2}v_0^{p_0}$, and therefore
\begin{equation*}
|\tilde A|_\infty\le2\,\kappa_1v_0^{p_0},\qquad\sum_{j\ge j_1}|\tilde a_j|_\infty\le1.81\,\kappa_1 .
\end{equation*}
On $\mathcal G$, the inequality $\bigl(\sum_{j<j_1}\tilde a_j\bigr)^2\le3.02\sum_jL^{j/4}\tilde a_j^2$ leads to
\begin{equation*}
\sum_{b\cap\mathsf Q\ne\emptyset}\tilde A(b)^2\le2\,\bigl(3.02\cdot16\cdot1.43+4\cdot3.3\bigr)\,\kappa_1^2\,\mathsf l^3\le200\,\kappa_1^2\,\mathsf l^3 .
\end{equation*}
Hence the function $\varsigma(x)=|\exp(iq\tilde A(c,c_3))-1|$ (eight labels for each direction) satisfies $\langle\varsigma^2\rangle_{\mathsf Q}\le67\,\kappa_1^2q^2$.

\emph{Step 5} (the Wilson energy on $\mathcal G$).
Step~14 of the proof of Theorem~5.16, with $H_\theta$ in the place of $\Omega_0$ (\cite[Proposition~1.3\,(b)]{QEDIII} holds for every character), and the function
\begin{equation*}
\varsigma'(x):=\bigl|\exp\bigl(iq(\tilde A+\mathbf a)(c,c_3)\bigr)-1\bigr|\le\varsigma(x)+q|\mathbf a|_\infty
\end{equation*}
give
\begin{equation*}
X:=(\psi,\nabla^*_{A_0}\nabla_{A_0}\psi)\le6\,c_b^{-1}\bigl(L^{-2k_\star}\|\Phi\|^2+\boldsymbol\mu^2\|\psi\|^2\bigr)+2\,C_5\,\|\varsigma'\psi\|^2 .
\end{equation*}
The bounds that follow are asserted for the pairs $(\theta,A)$ with $q|\mathbf a|_\infty\le L^{-k_\star}$, among which, by Step~8, are those of the set $\mathsf N$ of Step~7.
Step~3, with $\tau=\varsigma^2$ and summed over the blocks $\mathsf Q$, together with Step~4, gives
\begin{equation*}
\|\varsigma\psi\|^2\le134\,\kappa_1^2q^2\|\psi\|^2+2\,|\varsigma|^2_\infty\,C_{\mathrm{Po}}\,\mathsf l^2X,
\end{equation*}
and $|\varsigma|_\infty^2\,\allowbreak \mathsf l^2\le4\,\allowbreak \kappa_1^2L^4q^2v_0^{10p_0}$, a quantity which tends to $0$ as $N\to\infty$.
For large $N$ it follows that on $\mathcal G$
\begin{gather*}
\boldsymbol\mu\|\psi\|^2+W\le6\rho\,c_b^{-1}L^{-2k_\star}\|\Phi\|^2+\beta_N\|\psi\|^2,\\
\beta_N:=\boldsymbol\mu+6\rho\,c_b^{-1}\boldsymbol\mu^2+4\rho\,C_5\,L^{-2k_\star}+540\,\rho\,C_5\,\kappa_1^2q^2 .
\end{gather*}

\emph{Step 6} (Gaussian domination and large deviations; the second half of (b)).
Here $n=R_{\mathrm{top}}(e)/\eta=L^{k_\star}\ell_\star$ is the number of Theorem~5.16, which differs from the number $n$ of the notation fixed above, and $C^1(n)$, $C^2(n)$ are the spaces of the real one-cochains and two-cochains on $\mathbb Z^3/n\mathbb Z^3$.
The exact two-cochains form the subspace $F_N:=dC^1(n)$ of $C^2(n)$ \cite[Sections~I.7.1, I.7.3]{QEDIII}.
Its scalar product is the pairing $\langle w,\allowbreak w'\rangle=\allowbreak \sum_{\mathrm{pl}}w(\mathrm{pl})\,\allowbreak w'(\mathrm{pl})$, summed over the plaquettes $\mathrm{pl}$ of the torus, as in \cite[Definition~2.6]{QEDII} and \cite[Section~I.7.1]{QEDIII};
for $w\in F_N$ the norm $\langle w,w\rangle^{1/2}$, written $\|w\|$ in Section~5 and in \cite[Section~I.7.1]{QEDIII}, is denoted $|w|$.
By \cite[Chapter~29, P1, (Q7.4)]{QEDIII}, the Euclidean space so defined is the space $F_N$ of the first paragraph of \cite[Section~14.5.2]{QEDI} (compare \cite[Section~5.2.1]{QEDI}).
One has $\zeta_j(\Upsilon)=\langle\phi_{j,\Upsilon},F^{\mathrm u}\rangle$, the $\phi_{j,\Upsilon}$ forming an orthonormal family in $F_N$ \cite[Theorem~14.17\,(c)]{QEDI};
hence $\bar\mu(\exp\langle v,\zeta\rangle)\le\exp\frac12|v|^2$ at every real mass \cite[Proposition~31.1\,(d)]{QEDIII}.
From (K7) of Section~5 and Tonelli's theorem one obtains, for every real symmetric $\mathsf M\ge0$ and every $s\ge0$ with $4s\|\mathsf M\|\le1$,
\begin{equation*}
\bar\mu\exp\bigl(s\langle\zeta,\mathsf M\zeta\rangle\bigr)\le\det(1-2s\mathsf M)^{-1/2}\le\exp\bigl(2s\operatorname{tr}\mathsf M\bigr).
\end{equation*}
Let $\mathsf M$ be the matrix with $\sum_{b\cap\mathsf Q\ne\emptyset}\tilde a_j(b)^2=\allowbreak \langle\zeta_j,\allowbreak \mathsf M\zeta_j\rangle$;
then $\operatorname{tr}\mathsf M\le4\,\mathsf l^3L^{-j}\kappa_1^2$, and $\|\mathsf M\|\le L^{2j}\kappa_1\kappa_\infty$ by (K23) of this section.
Markov's inequality, with $s=(4L^{2j}\kappa_1\kappa_\infty)^{-1}$, bounds the $\bar\mu$-probability that this sum exceeds $16L^{-j}\kappa_1^2\mathsf l^3$ by
\begin{equation*}
\exp\Bigl(-\frac{2\kappa_1}{\kappa_\infty}\,\bigl(\mathsf lL^{-j}\bigr)^3\Bigr);
\end{equation*}
moreover $\mathsf lL^{-j}\ge v_0^{2p_0}$ and $\log n\le2v_0+C(e)$.

\emph{Step 7} (from $\bar\mu$ to $\mathsf P^{0,A'}$).
Steps~1 to~3 of the proof of \cite[Proposition~29\,(c)]{QEDIII} are carried out with $\mathbf 1_{\mathcal B}(\zeta)$ in the place of $\Phi(A)$.
By \cite[Chapter~29, Lemma~2\,(t6)--(t8)]{QEDIII}, by boundedness and by Fubini's theorem, Steps~1 and~3 of that proof hold with this function.
In its Step~2, part (2c) gives the normaliser
\begin{equation*}
\int\mathrm d\,\mathrm{mup}\;y_0\bigl(\mathrm{Af}_0(A+c_\theta,\cdot)\bigr)=Y_*(A+c_\theta),
\end{equation*}
where $y_0$, the function of Step~1\,(i) above, is denoted $\mathrm{wdn}$ in that proof;
at $\mathrm{Af}_0(A+c_\theta,\zeta)$ it is $\det\mathcal D_0$ times $\exp(-\varepsilon^{\mathrm v}_{\mathrm{in}}n^3)/\det(D_0+m\eta)$, a number which is positive by \cite[Proposition~31.1\,(a)]{QEDIII};
the reasoning of (2b), for its part, estimates the difference of the integrands uniformly in $\zeta$, so that the factor $\mathbf 1_{\mathcal B}$ leaves it intact, the integral with this factor being continuous, hence measurable, in $(\theta,A)$.
These steps give, for every Borel set $\mathcal B$,
\begin{equation*}
\bar\nu\bigl(\{\zeta\in\mathcal B\}\cap\mathfrak G\bigr)=\int\frac{d\theta}{(2\pi)^3}\int\mathrm d\mathsf m_*(A)\,\chi_*(A)\,Y_*(A+c_\theta)\,\mathsf P^{0,A+c_\theta}(\mathcal B).
\end{equation*}
In this step $\bar\nu$ is the average over the twist
\begin{equation*}
\bar\nu:=\int_{[-\pi,\pi]^3}\frac{d\theta}{(2\pi)^3}\,\nu^\theta,\qquad\nu^\theta:=\mathrm{dens}^{\theta,0}_N\,\gamma_N,
\end{equation*}
where $\gamma_N$ is the standard Gaussian measure of $F_N$ and the density $\mathrm{dens}^{\theta,0}_N$ of \cite[(I.7.2)]{QEDIII} is formed with the tuned pair of $m$ on $R_{\mathrm{top}}(e)$ \cite[Section~I.7.3]{QEDIII};
by \cite[Proposition~31.1\,(c)]{QEDIII}, $\bar\mu=\bar\nu/\bar\nu(1)$.
Further, $A_*$ is the field of the level $k_\star$ that \cite[Definition~5.11]{QEDI} attaches to $F^{\mathrm u}$, $\mathsf m_*$ is the law of $A_*$ under $\gamma_N$, a measure carried by $V_{\mathrm t}(T^0_{k_\star})$, and $\chi_*(A):=\mathbf 1\bigl[|dA|_\infty\le p_\star\bigr]$;
the set $\mathfrak G$ of this step is the small-field set
\begin{equation*}
\mathfrak G:=\bigl\{F^{\mathrm u}\in F_N:\ |\zeta_j(\Upsilon)|\le p_{0,j}\text{ for every }\Upsilon\text{ and every }j<k_\star,\ |dA_*|_\infty\le p_\star\bigr\}
\end{equation*}
of \cite[Section~G4.0.6]{QEDIII} and \cite[Definition~5.11]{QEDI}.
Here $Y_*(A+c_\theta)=\mathrm{Rtw}_\theta\,(1+\mathrm{vrem})$ with $|\mathrm{vrem}|\le\varepsilon_1$ (\cite[Proposition~32.1]{QEDIII}; $\varepsilon^{\mathrm v,t}_\star=0$ by \cite[Proposition~28.2]{QEDIII}), and
\begin{equation*}
|\log\mathrm{Rtw}_\theta|\le c_{\mathrm{tw}}:=2\varepsilon_5\bigl(\tfrac1L\bigr)+C_\Sigma
\end{equation*}
by \cite[Proposition~30.3\,(e)]{QEDIII} and \cite[Section~G4.0.2, (Sc3)]{QEDIII};
moreover $\bar\nu(1)\le\bar\nu(\mathfrak G)/(1-\varepsilon_3)$ \cite[Proposition~32.2\,(7.2)]{QEDIII}.
Let $\mathsf N$ be the set of the pairs $(\theta,A)$ with $\theta\in[\pi-r_1,\pi]^3$ and $|dA|_\infty\le r_1$;
on $\mathsf N$ one has $\chi_*=1$, since $r_1\le1\le p_\star$, the second inequality because $v_\star\ge8$ \cite[(fl)]{QEDIII}.
By \cite[Theorem~14.17\,(d)]{QEDI}, to each plaquette $\mathrm{pl}$ of $T^0_{k_\star}$ there corresponds an element $\phi_{\mathrm{pl}}\in F_N$ with $|\phi_{\mathrm{pl}}|\le1$ and $(dA_*)(\mathrm{pl})=\allowbreak \langle\phi_{\mathrm{pl}},\allowbreak F^{\mathrm u}\rangle$;
as $\gamma_N$ is the standard Gaussian measure of $F_N$ \cite[Section~I.7.3]{QEDIII}, the variables $(dA_*)(\mathrm{pl})$ are centred Gaussian under $\gamma_N$, with variance $|\phi_{\mathrm{pl}}|^2\le1$;
so (K22) of Section~6 gives
\begin{equation*}
\mathsf m_*\bigl(|dA|_\infty\le r_1\bigr)\ge\operatorname{erf}\bigl(r_1/\sqrt2\bigr)^{3\ell_\star^3}.
\end{equation*}
(A lower bound that does not depend on $N$ holds also without (K22).
One has $|dA|_\infty\le|dA|_2\le|dH_{k_\star}A|_2$, the first inequality because each plaquette of $T^0_{k_\star}$ carries the weight one, the second by \cite[Proposition~14.8\,(v)]{QEDI};
by \cite[Theorem~14.17\,(e)]{QEDI}, $\mathsf m_*$ has a density proportional to $\exp(-\frac12|dH_{k_\star}A|_2^2)$ on $V_{\mathrm t}(T^0_{k_\star})$, whose dimension is $2\ell_\star^3-2$ \cite[Section~14.1.1]{QEDI}, so that under $\mathsf m_*$ the variable $|dH_{k_\star}A|_2^2$ is $\chi^2$-distributed with $2\ell_\star^3-2$ degrees of freedom;
hence $\mathsf m_*(|dA|_\infty\le r_1)\ge(\frac12r_1^2)^{\ell_\star^3-1}\exp(-\frac12r_1^2)/\Gamma(\ell_\star^3)$.)
Taking the mean over $\mathsf N$, one finds a pair $(\theta,A)\in\mathsf N$ with $\mathsf P^{0,A+c_\theta}(\mathcal G^c)\le\varpi\,\allowbreak \bar\mu(\mathcal G^c)$, where
\begin{equation*}
\varpi:=\frac{(1+\varepsilon_1)\exp(2c_{\mathrm{tw}})}{(1-\varepsilon_1)(1-\varepsilon_3)}\Bigl(\frac{2\pi}{r_1}\Bigr)^3\operatorname{erf}\Bigl(\frac{r_1}{\sqrt2}\Bigr)^{-3\ell_\star^3}.
\end{equation*}
\emph{The number $\varpi$ depends on $(e,h,\mu_\star)$ and not on $N$}:
$\varepsilon_1$ and $\varepsilon_3$ depend on $(e_\star,\ell_\star,p_{0\star},p_\star)$ \cite[Section~G4.0.3]{QEDIII}, where $p_{0\star}:=p_{0,k_\star}=v_\star^{p_0}$ and $p_\star=v_\star^{p}$ are functions of $e_\star$ and of constants, and neither $e_\star$ nor $\ell_\star$ depends on $N$ \cite[Section~G4.0.2, (Sc1)]{QEDIII};
the estimate of \cite[Proposition~30.3\,(e)]{QEDIII} holds uniformly with respect to the mesh $L^{-k_\star}$ and to $b_\star\in[b_-,1]$, so that $c_{\mathrm{tw}}$ is a constant;
finally, $r_1$ is defined through $b_-$.

\emph{Step 8} (near $\Omega$).
For real $\mathscr A_1$ and $\mathscr A_2$,
\begin{equation*}
\|D_k(\mathscr A_1)-D_k(\mathscr A_2)\|\le12\,\mu_k^{-2}e_k\,|\mathscr A_1-\mathscr A_2|_\infty .
\end{equation*}
This follows from the definitions of $S_k(\mathscr A)$ and $D_k(\mathscr A)$ in \cite[Section~I.3.5]{QEDIII} and from the classical resolvent identity
\begin{equation*}
S_k(\mathscr A_1)-S_k(\mathscr A_2)=S_k(\mathscr A_1)\Bigl(D_{\mathscr A_2}-D_{\mathscr A_1}+b_k\,J\bigl(P_k(\mathscr A_2)-P_k(\mathscr A_1)\bigr)J^*\Bigr)S_k(\mathscr A_2),
\end{equation*}
with $b_k\le1$, $\mu_k\le1$ and $\|S_k(\mathscr A)\|\le\mu_k^{-1}$ for real $\mathscr A$ \cite[(I.O2), (I.Q2)]{QEDIII}, together with the two bounds
\begin{equation*}
\bigl\|\bigl(Q_k(\mathscr A_1)-Q_k(\mathscr A_2)\bigr)J^*\bigr\|\le\tfrac32\,e_k|\mathscr A_1-\mathscr A_2|_\infty,\qquad\|D_{\mathscr A_1}-D_{\mathscr A_2}\|\le6\,e_k|\mathscr A_1-\mathscr A_2|_\infty;
\end{equation*}
the first of them results from \cite[(I.3.3)]{QEDIII}, from the bound $\eta_k|(\tau_k\mathscr A)(y,c)|\le\frac32|\mathscr A|_\infty$ \cite[Section~I.9.2, Lemma (Bound for the axial functional)]{QEDIII} and from the Schwarz inequality on $\Delta_y$, and the second is \cite[(I.O6)]{QEDIII}.
The terms add up to $b_k^2\bigl[3\mu_k^{-1}+(6+3b_k)\mu_k^{-2}\bigr]e_k|\mathscr A_1-\mathscr A_2|_\infty$.
Moreover $\|D_k(\mathscr A)^{-1}\|\le g_k:=b_k^{-1}+\mu_k^{-1}$ by (K10) of Section~5.
On $\mathsf N$, the fact (K24), the first clause of \cite[(T2)]{QEDIII} and the bound $|H_{k_\star}A|_\infty\le C_0|A|_\infty$, which is the kernel bound quoted in \cite[(32.8)]{QEDIII} and, being linear, holds for every real $A$, give
\begin{equation*}
e_\star|H_{k_\star}A|_\infty\le e_\star C_0\,\ell_\star^{5/2}r_1\le r_1 .
\end{equation*}
Consequently $q|\mathbf a|_\infty\le L^{-k_\star}$ and $e_\star|\mathscr A'-\Omega|_\infty\le2r_1$.
It follows that $D_{k_\star}(\mathscr A')$ differs in norm from $D_{k_\star}(\Omega)$ and from $D_{k_\star}(c_\theta)$ by at most
\begin{equation*}
12\,\mu_\star^{-2}\cdot2r_1=\min\Bigl\{\varkappa_D,\frac{c_-}{28}\Bigr\}\le\min\Bigl\{\varkappa_D,\frac{c_\star}{28}\Bigr\}.
\end{equation*}
From $D_{k_\star}(\Omega)\Psi_u=d_\star\Psi_u$, from Step~1\,(ii) and from the sizes obtained in Steps~9 and~12 of the proof of Theorem~5.16 one concludes
\begin{equation*}
c(A')\ge c_\star-7\Bigl(\frac{c_\star}{28}+\ell_\star^3\delta_\star\Bigr)-2\delta_2\ge\tfrac12c_\star,\qquad\mathsf E^{A'}\|\Phi\|^2\le2 .
\end{equation*}

\emph{Step 9} (from $\mathsf P^{0,A'}$ to $\mathsf P^{A'}$).
By Step~1\,(i), $\mathsf E^{A'}[\det\mathcal D;\mathcal B]=\allowbreak \mathsf E^{A'}[\det\mathcal D]\,\allowbreak \mathsf P^{0,A'}(\mathcal B)$.
Let $E_{\star,00}(X):=E_{k_\star}(X;\mathscr A',0,0)$ be the component of degree $(0,0)$ of $E_{k_\star}\bigl(X;\mathscr A',\allowbreak H^\#_{k_\star}(\mathscr A')\Psi\bigr)$.
Since $m_{k_\star}=0=\varepsilon^{\mathrm v,t}_{k_\star}$, the definition of $\mathscr F_{k_\star}$ in \cite[Section~I.5.7]{QEDIII} yields $\mathscr F_{k_\star}(A';0)=\allowbreak \det D_{k_\star}(0)^{-1}\exp\sum_XE_{\star,00}(X)$, and therefore
\begin{equation*}
\mathsf E^{A'}\det\mathcal D=\frac{Y_*(A')}{\mathscr F_{k_\star}(A';0)}=\det D_{k_\star}(c_\theta)\,(1+\mathrm{vrem})\exp\Bigl(-\sum_XE_{\star,00}\Bigr)\le1.21\,\det D_{k_\star}(c_\theta),
\end{equation*}
by \cite[Proposition~30.3\,(f)]{QEDIII} and because $|\sum_XE_{\star,00}|\le e_\star^{1/4-7\varepsilon}\le\frac18$.
Here the bound of Step~9 of the proof of Theorem~5.16, taken at $\mathscr A'$, gives $|\sum_XE_{\star,00}|\le4e_\star^{1/4-6\varepsilon}K_0(\ell_\star/M)^3$, and the second and third clauses of \cite[(T2)]{QEDIII}, with $c_{\mathrm{fl}}=3294K_0$, give the two inequalities.
On the event $\{\|\mathcal D-D_{k_\star}(\mathscr A')\|\le\varkappa_D\}$, two applications of (K23), together with $g_\star\varkappa_D\le g_-\varkappa_D=(16\ell_\star^3)^{-1}$, give $\det\mathcal D\ge\frac9{16}\det D_{k_\star}(c_\theta)$.
By the first clause of \cite[(T2)]{QEDIII}, which reads $8C_0^2\ell_\star^6g_-\le e_\star^{-\varepsilon/2}$, and by \cite[(fl)]{QEDIII}, one has $\|K_E\|\le\frac12\varkappa_D$;
by Chebyshev's inequality,
\begin{equation*}
\mathsf P^{A'}\bigl(\|\delta\mathcal D\|_{\mathrm{HS}}>\tfrac12\varkappa_D\bigr)\le8\,\varkappa_D^{-2}\,\ell_\star^3\,\delta_2\le e_\star^{5/4-12\varepsilon}.
\end{equation*}
Hence $\mathsf P^{A'}(\mathcal B)\le e_\star^{5/4-12\varepsilon}+3\,\allowbreak \mathsf P^{0,A'}(\mathcal B)$.

\emph{Step 10} (conclusion).
Let $(\theta,A)$ be the pair of Step~7.
Steps~2 (with $\mathcal B=\mathcal G$), 5, 8 and~9 give
\begin{equation*}
\tfrac12c_\star\le12\rho\,c_b^{-1}L^{-k_\star}+L^{k_\star}\beta_N\,\mathsf E^{A'}\bigl[\|\psi\|^2;\mathcal G\bigr]+\tfrac14\bigl(e_\star^{5/4-12\varepsilon}+3\varpi\,\bar\mu(\mathcal G^c)\bigr).
\end{equation*}
By \cite[(T2)]{QEDIII}, $c_\star\ge\frac13d_\star\ge\frac83e_\star^{\varepsilon/2}$ and $e_\star^{5/4-12.5\varepsilon}\le e_\star^{1/4-9\varepsilon}\le\frac18$, so that $\frac14e_\star^{5/4-12\varepsilon}\le\frac1{16}c_\star$.
For large $N$ the first term on the right is at most $\frac18c_\star$, and $\frac34\varpi\,\bar\mu(\mathcal G^c)\le\frac1{16}c_\star$, since $c_\star\ge c_-$ and neither $c_-$ nor $\varpi$ depends on $N$ (Steps~6 and~7);
thus the first and the last term are at most $\frac18c_\star$ each.
As $\mathsf E^{A'}\|\psi\|^2<\infty$ by Step~6 of the proof of Theorem~5.16, it follows that $\beta_N>0$.
If $\boldsymbol\mu<0$, then for large $N$ one has $6\rho\,c_b^{-1}|\boldsymbol\mu|\le\frac12$ \cite[Proposition~28.3\,(d)]{QEDIII} and $4L^{-2k_\star}\le60\,\kappa_1^2q^2$, whence $\frac12\boldsymbol\mu+600\,\rho\,C_5\,\kappa_1^2q^2>0$.
Accordingly, $N_A(e,m)$ is the least $N_0\ge N_{\min}(R_{\mathrm{top}}(e),e)$ such that every $N\ge N_0$ satisfies
\begin{gather*}
2j_1\le k_\star,\qquad64\,C_5\,C_{\mathrm{Po}}\,\kappa_1^2L^4q^2v_0^{10p_0}\le1\quad\text{(Step~5)},\qquad96\rho\,\eta\le c_bc_-a_\star,\\
12\,\varpi\,j_1n^3\exp\Bigl(-\frac{2\kappa_1}{\kappa_\infty}\,v_0^{6p_0}\Bigr)\le c_-,\qquad3\rho\,q^{3/4-8\varepsilon}\le2c_b,\qquad\eta\le15\,\kappa_1^2e^2a_\star^2 .
\end{gather*}
This number is a function of $(e,h,\mu_\star)$ and of constants.
\end{proof}

\begin{proof}[Proof of Theorem~6.1]
Put $h_j:=L^jh$ for $0\le j\le J'$, where $h_{J'}=h_J$;
a subscript $j$ marks the numbers with the index $\star$ that belong to the class $I_{h_j}$.
Let $N_1(e)$ be the largest of the numbers $N_A(e,m_\flat)$ and $N_{\min}(R_{\mathrm{top}}(e;h'),e)$, where $h'$ ranges over the finite set of the powers of $L$ in $[1,h_J]$, that is, of the numbers $h(m)$ with $m\in(0,1]$ and $e\,h(m)^{1/2}\le e_1$.

\emph{Step 1}.
The upper bound is Theorem~5.16\,(b).
By Proposition~6.2\,(a), $\varphi_N(m_\flat)\ge-C_{\mathsf W}e^2$.
Let $h=h_J$.
Then (M) of Proposition~6.3\,(c), applied on $[\bar m,m_\flat]$ together with $w_\star/a_\star\le\frac18m_\flat$ (the bound on $w_j/a_j$ of Step~3), gives $\varphi_N(\bar m)\ge-C_{\mathsf W}e^2-2m_\flat$;
this is at least $\bar m-C_\natural e^2$, because $\bar m\le m_\flat<Le^2/e_1^2$.

\emph{Step 2} (a chain of classes, $h<h_J$).
Let $h<h_J$, and for $j=J'-1,\dots,0$ put
\begin{equation*}
m_b^{(j)}:=\frac1{Lh_j},\qquad\mathrm{top}_j:=\frac1{h_j}\ \ (j\ge1),\qquad\mathrm{top}_0:=\bar m,\qquad m_c^{(j)}:=\min\Bigl\{m_b^{(j)}+\frac{w_j}{2a_j},\,\mathrm{top}_j\Bigr\}.
\end{equation*}
Adding the inequalities (J) of Proposition~6.3\,(d), from $m_b^{(j)}=\mathrm{top}_{j+1}$ to $m_c^{(j)}$, and (M), from $m_c^{(j)}$ to $\mathrm{top}_j$, one obtains
\begin{equation*}
\varphi_N(\bar m)-\varphi_N(m_\flat)\ge(\bar m-m_\flat)-\sum_je_j^\varepsilon\bigl(\mathrm{top}_j-m_b^{(j)}\bigr)-\sum_j\frac{w_j}{a_j}\Bigl(\frac23+e_j^\varepsilon\Bigr).
\end{equation*}

\emph{Step 3} (the error terms).
One has $\mathrm{top}_0-m_b^{(0)}\le\bar m$ and $\mathrm{top}_j<L^{1-j}\bar m$;
moreover $e_j\le L^{j/2}e\,h^{1/2}$, and $e_j^\varepsilon\le\frac1{48}$ by \cite[(fl)]{QEDIII}.
By \cite[(T2)]{QEDIII}, $\mu_\star^{-1}\le\frac18e_\star^{-\varepsilon/2}$ for every parameter of the class;
letting the parameter tend to $\frac1{Lh_j}$, one obtains
\begin{equation*}
\frac{w_j}{a_j}\le\frac1{8Lh_j}\,e_j^{3/4-9\varepsilon}<\frac{\bar m}8\,L^{-j}e_j^{3/4-9\varepsilon},
\end{equation*}
and $e_j^{1/4-9\varepsilon}\le\frac18$.
Therefore the two sums of Step~2 are together at most $\epsilon_e\bar m$.

\emph{Step 4} (proof of (a) to (c)).
If $e\le e_7$, then $C_\natural e^2\le\frac14\bar m$ and $e^2Lh\le e_1^2/12$, so that $\varphi_N(\bar m)\ge(\frac{11}{12}-\frac14)\bar m$.
For (b), one has $h^{\varepsilon/2}\bar m\le1$.
For (c): by Theorem~5.16\,(c), since $\mathsf W_N\ge0$ and $c_\star=\bar m\,a_\star Z_0$,
\begin{equation*}
Z_N\le\frac\eta{a_\star}\,c_\star\,\frac{1+\vartheta_N}{\boldsymbol\mu_N}\le2.1\,Z_0 .
\end{equation*}
The numbers $Z_0$ and $c_\star$ depend on $N$ through $b_\star$.
\end{proof}

\subsection{One law}

\begin{named}{Lemma 6.4}
\emph{What the atom needs; the measure $\mathrm{kl}$ and the components.}
Let $e>0$, let the mass assignment be arbitrary, and let $\mu$ be a volume-first limit point with the ergodic decomposition $\mu=\int\mu_\xi\,d\nu(\xi)$.
\begin{itemize}
\item[(a)] The bounds of Theorem~5.14 are statements at a single cutoff, and the conclusion of Theorem~5.14\,(c) is obtained along one defining sequence, in the following form.
Let $c>0$, let $\mu=\lim_i\mu^{(N_i)}$ with $\mu^{(N_i)}\in\mathcal V_{N_i}$, and suppose that $c\,\eta_i\le\boldsymbol\mu_{N_i}\le1$ for all large $i$, where $\eta_i=L^{-N_i}$.
Then, along a defining sequence of $\mu$, every probe $\beta$ satisfies
\begin{equation*}
\bar I(\beta)\le\Bigl(1+\frac{C_F'e^2}{c}\Bigr)\|f\|^2,\qquad C_F':=15\,\pi^4\mathsf h_3^4\,C_\sharp\,\max\{1,e^2/c\}^3,
\end{equation*}
and
\begin{equation*}
\mathrm{kl}_\mu(\{0\})\ge\Bigl(1+\frac{C_F'e^2}{c}\Bigr)^{-1}.
\end{equation*}
Consequently, if $\limsup_N\boldsymbol\mu_N/\eta>0$, the atom is obtained at the limit points taken along those $N_i$ for which $c\,\eta_i\le\boldsymbol\mu_{N_i}\le1$, and nothing is asserted at the limit points taken along other sequences.
\item[(b)] The map $\xi\mapsto\mathrm{kl}_\xi$ is measurable, and $\mathrm{kl}_\mu=\int\mathrm{kl}_\xi\,d\nu(\xi)$; in particular $\mathrm{kl}_\mu(\{0\})=\int\mathrm{kl}_\xi(\{0\})\,d\nu(\xi)$.
\end{itemize}
\end{named}

\begin{proof}
\emph{Step 1} (proof of (a)).
The first display of Theorem~5.14 is a statement about one quadruple $(q,\mathbf m,\mathbf m',n)$, and Step~30 of the proof of Theorem~5.14 takes $\mathbf m=\boldsymbol\mu_N$ and $\mathbf m'=\max\{\boldsymbol\mu_N,e^2\eta\}$.
Theorem~5.14\,(b) therefore holds, with $c_{\boldsymbol\mu}\bar m=c$, at every cutoff at which $c\,\eta\le\boldsymbol\mu_N\le1$, $e^2\eta\le1$ and $n\boldsymbol\mu_N\ge20$.
By Step~5 of the proof of Proposition~5.13, $\mu$ has defining sequences $(N_i,T_i)$ whose $T_i$ are as large as one wishes.
The assertions of (a) then follow from \cite[Proposition~47.1\,(b1), (b5)]{QEDIII} and Proposition~5.6\,(d).

\emph{Step 2} (proof of (b)).
By \cite[Theorem~1.4, Part~I\,(d)--(f)]{QEDII}, almost every $\mu_\xi$ has the form (C7) of \cite[Theorem~44]{QEDIII}.
Let $\mathsf D_\xi$ and $\mathsf D_\mu$ be the functions $\mathsf D$ of Step~9 of the proof of Proposition~5.6, formed with $\mathrm{kl}_\xi$ and with $\mathrm{kl}_\mu$ respectively.
The function $\mathsf D_\xi(r)$ is continuous in $r$, and it arises as the limit of $S_2^\xi(g_n,g_n)$ along forms $g_n$ whose weights accumulate on the sphere $|p|^2=r$; it is therefore measurable in $\xi$ \cite[Proposition~45.6\,(d0)]{QEDIII}.
The identity $S_2^\mu=\int S_2^\xi\,d\nu$ \cite[Proposition~45.6\,(a$'$), (a$''$)]{QEDIII} and the bounds $0\le\mathsf D_\xi\le1$ give $\mathsf D_\mu(r)=\int\mathsf D_\xi(r)\,d\nu(\xi)$.
By partial fractions, the closed linear span of the functions $s\mapsto(r+s)^{-1}$, $r>0$, is an algebra, and by the theorem of Stone and Weierstrass (K28) it is the whole of $C_0([0,\infty))$.
Hence $\int\mathrm{kl}_\xi\,d\nu$ is a measure of mass at most one that represents $S_2^\mu$.
Such a measure is unique, at every $e>0$ and for every mass assignment, by \cite[Theorem~1.4, Part~I\,(d)]{QEDII}, and also by the density just established.
For the atoms one lets $r\downarrow0$ and applies dominated convergence.

\end{proof}


\section{Scores and ergodic components}

The notation $(\mathrm{FI}_\lambda)$ with constant $C$ means
$\lambda(\partial_f\Phi)^2\le C\|f\|^2\lambda(\Phi^2)$ for every
probe and every smooth bounded cylinder function with all derivatives bounded.
Equivalently $C(\lambda)\le C$.


\emph{Conventions.}
Unless stated otherwise, $\mu$ denotes a law that is invariant under $\mathbb R^3$ and has the property (G) of the notation fixed above;
the $\sigma$-algebra $I$, the kernel $\omega\mapsto\mu_\omega$ and the measure $\nu$ are those of \cite[Proposition~45.6]{QEDIII}, and ``almost every $\xi$'' is understood in the sense of \cite[(n7)]{QEDIII}.
Put $L^2_{\mathrm c}:=\Pi_{\mathrm c}L^2(\mathbb R^3;\Lambda^2)$.
For a law $\mathrm{lw}$, the symbol $\mathcal A(\mathrm{lw})$ stands for the class of the functions $\Phi=\phi(X(g_1),\dots,X(g_k))$ of Proposition~5.6\,(c), regarded as elements of $L^2(\mathrm{lw})$, with $\partial_f\Phi$ defined as there;
its subclass $\mathcal A_{\mathrm c}$ is formed by the $\Phi$ with $g_j\in C_c^\infty$, which are functions on $\mathcal S'$ and constitute an algebra, and on which $U(a)\Phi=\phi(X(\tau_a\cdot g_1),\dots,X(\tau_a\cdot g_k))$.
For a law $\mathrm{lw}$ with (G) and for $f\in L^2_{\mathrm c}$ set
\begin{equation*}
I_{\mathrm{lw}}(f):=\sup_{\Phi\in\mathcal A(\mathrm{lw})}\bigl[2\,\mathrm{lw}(\partial_f\Phi)-\mathrm{lw}(\Phi^2)\bigr]\in[0,\infty],\qquad C(\mathrm{lw}):=\sup_{\beta:\,dd^*\beta\ne0}\frac{I_{\mathrm{lw}}(dd^*\beta)}{\|dd^*\beta\|^2}.
\end{equation*}
Replacing $\Phi$ by $t\Phi$, with real $t$, one finds that $I_{\mathrm{lw}}(f)\le\mathsf c$ is equivalent to $\mathrm{lw}(\partial_f\Phi)^2\le\mathsf c\,\mathrm{lw}(\Phi^2)$ for all $\Phi\in\mathcal A(\mathrm{lw})$;
thus $(\mathrm{FI}_{\mathrm{lw}})$ with constant $C$ is equivalent to $C(\mathrm{lw})\le C$.
One writes $I_\xi:=I_{\mu_\xi}$, $C_\xi:=C(\mu_\xi)$ and $\mathsf z_\xi:=\mathrm{kl}_\xi(\{0\})$.
For the probes considered, $f=dd^*\beta\ne0$;
with $\beta^{(l)}:=l^{1/2}\beta(\cdot/l)$ and $f_l:=dd^*\beta^{(l)}$ one has $\|f_l\|=\|f\|$.
A \emph{law of the class} means a member of $\mathcal K$.

Put
$\mathcal K(C)=\{\lambda\in\mathcal K:C(\lambda)\le C\}$.
The space $\mathcal D_0\subset C_c^\infty$ is the countable
$\mathbb Q$-vector space of \cite[Lemma~8.11(i)]{QEDI}.
The Gaussian laws $\gamma^\varrho$ and the operation $\ast_\epsilon$
are defined in Section~4. Scores are denoted by $\mathsf s_f$;
$T_f\omega=\omega+f$. The constants $C_\natural$, $e_7$ and $N_1(e)$
are those of Section~6, and $C_\flat$ is the lattice clustering
constant in Proposition~5.17.
\emph{Classical facts.}
\begin{itemize}
\item[(K12)] (Riesz.)
Let a linear functional be defined on a dense subspace of a Hilbert space and have norm $\kappa$.
Then it equals $\langle s,\cdot\rangle$ for exactly one vector $s$, and $\|s\|=\kappa$.
\item[(K13)] The theorems of monotone and of dominated convergence, Fatou's lemma and Tonelli's theorem; Jensen's inequality for the function $x\mapsto1/x$.
(Monotone class theorem.)
On some set, consider a vector space $\mathrm{Mv}$ of bounded real functions which contains the constant functions and is stable under pointwise limits of increasing sequences that are uniformly bounded, and a subset $\mathrm{Mu}\subset\mathrm{Mv}$ that is closed under products.
Then every bounded $\sigma(\mathrm{Mu})$-measurable function belongs to $\mathrm{Mv}$.
\item[(K14)] If $(Y_0,\dots,Y_k)$ is a centred Gaussian vector and $\phi\in C^1_b$, then
\begin{equation*}
E\bigl[Y_0\,\phi(Y_1,\dots,Y_k)\bigr]=\sum_j\operatorname{Cov}(Y_0,Y_j)\,E[\partial_j\phi].
\end{equation*}
\item[(K15)] (von Neumann.)
For a unitary representation of $\mathbb Z^3$, the averages over the sets $\{|x|_\infty\le\ell\}$ converge strongly to the projection onto the invariant vectors.
\item[(K16)] (Stone, Naimark, Ambrose, Godement.)
Every strongly continuous unitary representation of $\mathbb R^3$ has the form $\int\exp(i\langle a,\mathsf p\rangle)\,dE(\mathsf p)$ with a unique projection valued measure $E$.
A closed subspace that is invariant under a bounded self-adjoint operator reduces that operator.
The identity theorem.
(Transitive actions.)
Suppose that a locally compact second countable group $G$, with left Haar measure $\lambda_G$, acts continuously and transitively on a second countable, locally compact Hausdorff space $Y$.
If $m$ is a $\sigma$-finite Borel measure on $Y$ and $S\subset Y$ is a Borel set with $m(gS\,\triangle\,S)=0$ for all $g\in G$, then either $m(S)=0$ or $m(Y\setminus S)=0$.
\begin{proof}
The set $N:=\{(g,y):1_S(gy)\ne1_S(y)\}$ is Borel, and its section at $g$ is $g^{-1}S\,\triangle\,S$, which is $m$-null.
Tonelli's theorem shows that for $m$-almost every $y$ the set $\{g:1_S(gy)\ne1_S(y)\}$ is $\lambda_G$-null.
Assume that $m(S)>0$ and $m(Y\setminus S)>0$.
Then there are points $y\in S$ and $y'\notin S$ with this property, and some $g_0\in G$ with $y'=g_0y$.
It follows that $gy\in S$ for $\lambda_G$-almost every $g$, and that $gg_0y\notin S$ for $\lambda_G$-almost every $g$.
The second assertion means that $g'y\notin S$ for every $g'$ in a right translate of a conull set;
since a right translate of a $\lambda_G$-null set is null, by the modular function, this is incompatible with $\lambda_G\ne0$.
\end{proof}
\end{itemize}

\emph{Four facts on one law with (G).}
In (S1)--(S4), $\mathrm{lw}$ is a law with (G).
\begin{itemize}
\item[(S1)] The functions $\cos X(g)$ and $\sin X(g)$, $g\in\mathcal D_0$, span over $\mathbb R$ a subspace of $\mathcal A_{\mathrm c}$ that is dense in the real space $L^2(\mathrm{lw})$ \cite[Lemma~8.12\,(i)]{QEDI}.
\item[(S2)] If $I_{\mathrm{lw}}(f)<\infty$, then $\mathrm{lw}(\partial_f\Phi)$ is a linear functional of the class of $\Phi$, of norm $I_{\mathrm{lw}}(f)^{1/2}$.
Hence, by (S1) and (K12), there is a unique $\mathsf s_f\in L^2(\mathrm{lw})$, called the \emph{score}, such that $\mathrm{lw}(\partial_f\Phi)=\mathrm{lw}(\Phi\,\mathsf s_f)$ for every $\Phi\in\mathcal A(\mathrm{lw})$;
moreover $\|\mathsf s_f\|_2^2=I_{\mathrm{lw}}(f)$.
For the $\mathsf s_\beta$ of Proposition~5.6\,(c) this gives $\mathsf s_\beta=\mathsf s_f$ and $I_\mu(f)\le\bar I(\beta)$.
\item[(S3)] Let $\mathcal F_k\subset C_c^\infty(\mathbb R^k)$ be a countable set that is dense in $C_c^\infty(\mathbb R^k)$ for the $C_b^2$ norm, and let $\mathcal A_0\subset\mathcal A_{\mathrm c}$ be the countable set of those $\Phi$ for which $g_j\in\mathcal D_0$ and $\phi\in\mathcal F_k$.
Each $\Phi\in\mathcal A(\mathrm{lw})$ is the limit of a uniformly bounded sequence of $\Phi_n\in\mathcal A_0$ such that $\Phi_n\to\Phi$ and $\partial_f\Phi_n\to\partial_f\Phi$ in $L^2(\mathrm{lw})$ for every $f$.
Such a sequence is obtained by cutting $\phi$ off at the scale $n$, approximating in $C^2_b$, choosing $g^{(n)}_j\to g_j$ in $L^2$, and using the bound $\|X(g)\|_2\le\|g\|$ \cite[Lemma~14.3\,(a)]{QEDII}.
It follows that the supremum defining $I_{\mathrm{lw}}$ may be restricted to $\mathcal A_0$, that $\xi\mapsto I_\xi(f)$ is measurable \cite[Proposition~45.6\,(d0)]{QEDIII}, and that an identity $\mathrm{lw}(\partial_f\Phi)=\mathrm{lw}(\Phi\,\mathsf s)$ with $\mathsf s\in L^2(\mathrm{lw})$ extends from $\mathcal A_0$ to $\mathcal A(\mathrm{lw})$.
\item[(S4)] On $L^2_{\mathrm c}$, the function $f\mapsto I_{\mathrm{lw}}(f)$ is a supremum of continuous affine functions and is therefore lower semicontinuous;
moreover the forms $dd^*\beta$, $\beta\in\mathcal D_0$, are dense in $L^2_{\mathrm c}$, by Step~7 of the proof of Proposition~5.6 and the density of $\mathcal D_0$ in $\mathcal S$.
Hence $C(\mathrm{lw})$ equals the supremum over $\beta\in\mathcal D_0$, the map $\xi\mapsto C_\xi$ is measurable, and $I_{\mathrm{lw}}\le C(\mathrm{lw})\,\|\cdot\|^2$ on $L^2_{\mathrm c}$.
\end{itemize}
\subsection{Disintegration of the score}
\begin{named}{Theorem 7.1}
\emph{The atom in almost every component.}
Assume that $0<\rho<1$, that the weights are standard, that $\bar m\in(0,1]$ and that $e>0$;
the mass assignment is arbitrary, every physical one and the tuned ones being admitted.
\begin{itemize}
\item[(a)] \emph{One probe; limits in any order.}
Fix a limit point $\mu$ together with a defining sequence, a probe $\beta$ with $\bar I(\beta)<\infty$, and a $\Sigma$-measurable representative $\mathsf s$ of $\mathsf s_\beta$.
Then, for almost every $\xi$, the function $\mathsf s$ belongs to $L^2(\mu_\xi)$ and serves as the score of $\mu_\xi$ along $f$;
furthermore
\begin{equation*}
\int I_\xi(f)\,d\nu=I_\mu(f)\le\bar I(\beta),\qquad\|f\|^4\le S_2^\xi(f,f)\,I_\xi(f),\qquad\int\frac{\|f\|^4}{S_2^\xi(f,f)}\,d\nu(\xi)\le\bar I(\beta).
\end{equation*}
Consequently $S_2^\xi(f,f)>0$, and the components equal to the zero field form a $\nu$-null set.
\item[(b)] \emph{The atom.}
Let $\mu$ be volume-first, and let $(N_i,T_i)$ be a defining sequence of $\mu$ as a limit point;
such a sequence exists, since by \cite[Theorem~1.3\,(a)]{QEDII} each iterated limit is a limit point;
compare Step~9 of the proof of Proposition~5.6.
Let $\bar I$ be formed along this sequence, let $\beta$ be \emph{one} probe, and let $l_k\to\infty$ be \emph{one} sequence.
Then
\begin{equation*}
\int\mathsf z_\xi^{-1}\,d\nu\le\liminf_{k\to\infty}\frac{\bar I(\beta^{(l_k)})}{\|f\|^2}.
\end{equation*}
This inequality holds for every defining sequence, and the hypothesis of Proposition~5.6\,(d) is understood along one of them, as in Lemma~6.4(a) and Theorem~7.1(b).
If this hypothesis is satisfied, then
\begin{equation*}
\int\mathrm{kl}_\xi(\{0\})^{-1}\,d\nu(\xi)\le1+c.
\end{equation*}
Hence \textbf{$\mathrm{kl}_\xi(\{0\})>0$ for $\nu$-almost every $\xi$}, and $\nu\{\mathsf z_\xi\le t\}\le ct/(1-t)$ for $0<t<1$.
Since $\mathrm{kl}_\mu(\{0\})=\int\mathsf z_\xi\,d\nu$, an identity contained in the equality $\mathrm{kl}_\mu=\int\mathrm{kl}_\xi\,d\nu$ of measures in Lemma~6.4\,(b), Jensen's inequality gives back the conclusion of Proposition~5.6\,(d).
\item[(c)] \emph{Laws.}
Part~(a) extends to every law $\mu$ that has (G) and is invariant under $\mathbb R^3$, with $I_\mu(f)$ replacing $\bar I(\beta)$, and to every $f\in L^2_{\mathrm c}\setminus\{0\}$ of compact support such that $I_\mu(f)<\infty$;
without the last condition, the identity $\int I_\xi(f)\,d\nu=I_\mu(f)$ holds in $[0,\infty]$.
Part~(b) extends to all laws of the class, with $I_\mu(f_{l_k})$ replacing $\bar I(\beta^{(l_k)})$.
Accordingly $\int\mathsf z_\xi^{-1}\,d\nu\le C$ for $\mu\in\mathcal K(C)$.
\end{itemize}
\end{named}

\begin{proof}
\emph{Step 1} (reduction to (c)).
Every limit point is invariant under $\mathbb R^3$ and has (G), and a volume-first limit point belongs to the class \cite[Theorem~1.3\,(a), (b)(i)]{QEDII}.
Moreover (S2) gives $I_\mu(f)=\|\mathsf s_\beta\|_2^2\le\bar I(\beta)$.

\emph{Step 2} (an invariant function is a constant for $\partial_f$).
Consider $f\in L^2_{\mathrm c}$ of compact support with $I_\mu(f)<\infty$, a function $\Phi\in\mathcal A_{\mathrm c}$, and $\mathsf v\in L^2(\mu)$ such that $U(a)\mathsf v=\mathsf v$ for all $a$.
It is claimed that
\begin{equation*}
\mu(\mathsf v\,\partial_f\Phi)=\mu(\mathsf v\,\Phi\,\mathsf s_f).
\end{equation*}
Given $\epsilon>0$, (S1) furnishes $\Psi=\psi(X(g_1'),\dots,X(g_{k'}'))\in\mathcal A_{\mathrm c}$ with $\|\Psi-\mathsf v\|_2<\epsilon$;
choose $a$ such that none of the sets $a+\operatorname{supp}g_i'$ meets $\operatorname{supp}f$.
The function $\Psi_a:=U(a)\Psi$ lies in $\mathcal A_{\mathrm c}$ and satisfies $\partial_f\Psi_a=0$;
moreover $\|\Psi_a-\mathsf v\|_2=\|U(a)(\Psi-\mathsf v)\|_2<\epsilon$, because $U(a)$ is unitary \cite[Proposition~45.1\,(5),~(11)]{QEDIII}.
As $\Phi\Psi_a\in\mathcal A_{\mathrm c}$, the Leibniz rule and (S2) yield
\begin{equation*}
\mu(\Psi_a\,\partial_f\Phi)=\mu\bigl(\partial_f(\Phi\Psi_a)\bigr)=\mu(\Psi_a\,\Phi\,\mathsf s_f).
\end{equation*}
Replacing $\Psi_a$ by $\mathsf v$ changes the first member of this chain by at most $\epsilon\sup|\partial_f\Phi|$ and the last one by at most $\epsilon\sup|\Phi|\,\|\mathsf s_f\|_2$.

\emph{Step 3} (components).
For $\mathrm{Ie}\in I$, Step~2 with $\mathsf v=1_{\mathrm{Ie}}$ expresses the equality $E_\mu[\partial_f\Phi\,|\,I]=E_\mu[\Phi\,\mathsf s\,|\,I]$.
Consequently $\mu_\omega(\partial_f\Phi)=\mu_\omega(\Phi\,\mathsf s)$ outside a $\mu$-null set of $\omega$;
moreover $\int\mu_\omega(\mathsf s^2)\,d\mu=\mu(\mathsf s^2)$ by \cite[Proposition~45.6\,(a$'$)]{QEDIII}.
By \cite[Proposition~45.6\,(d0)]{QEDIII} the set
\begin{equation*}
\mathcal B_\Phi:=\{\xi:\ \xi(\mathsf s^2)<\infty,\ \xi(\partial_f\Phi)=\xi(\Phi\,\mathsf s)\}
\end{equation*}
is measurable, and \cite[Proposition~45.6\,(a$''$)]{QEDIII} gives $\nu(\mathcal B_\Phi)=1$.
By \cite[Proposition~45.6\,(c)]{QEDIII}, read with \cite[(n7)]{QEDIII}, some set of the $\sigma$-algebra of \cite[(n7)]{QEDIII} has $\nu$-measure $1$ and consists of laws with (G);
let $\Xi_f$ be the intersection of such a set with the sets $\mathcal B_\Phi$, $\Phi\in\mathcal A_0$.
If $\xi\in\Xi_f$, then by (S3) and (S2) the function $\mathsf s$ is the score of $\mu_\xi$, and $I_\xi(f)=\xi(\mathsf s^2)$;
the integral of the latter with respect to $\nu$ is $\mu(\mathsf s^2)=I_\mu(f)$.

\emph{Step 4} (the identity in $[0,\infty]$).
Assume that $\int I_\xi(f)\,d\nu<\infty$.
For $\Phi\in\mathcal A_{\mathrm c}$ the Schwarz inequality gives
\begin{equation*}
\mu(\partial_f\Phi)^2\le\Bigl(\int I_\xi(f)^{1/2}\,\xi(\Phi^2)^{1/2}\,d\nu\Bigr)^2\le\int I_\xi(f)\,d\nu\cdot\mu(\Phi^2).
\end{equation*}
By (S3) it follows that $I_\mu(f)<\infty$, and Step~3 applies.

\emph{Step 5} (the Cram\'er--Rao inequality within a component).
For $\xi\in\Xi_f$ put $\Phi_r:=r\tanh(X(f)/r)$;
then $\|f\|^2\,\allowbreak \xi\bigl(\cosh^{-2}(X(f)/r)\bigr)=\allowbreak \xi(\Phi_r\,\allowbreak \mathsf s)$.
Since $|\Phi_r|\le|X(f)|$, letting $r\to\infty$ and using (K13) one obtains
\begin{equation*}
\|f\|^2=\xi\bigl(X(f)\,\mathsf s\bigr)\le S_2^\xi(f,f)^{1/2}\,I_\xi(f)^{1/2}.
\end{equation*}
Division by $S_2^\xi(f,f)$ and integration over $\xi$ by means of Step~3 give (a).

\emph{Step 6} (scales).
Let $\mu$ be a law of the class, and let $\Xi_1$ be a measurable set with $\nu(\Xi_1)=1$ on which $\mu_\xi$ is a law of the class, is ergodic, clusters on $L^2(\mu_\xi)$, and has a two-point function $S_2^\xi$ of K\"all\'en--Lehmann form with the measure $\mathrm{kl}_\xi$ \cite[Theorem~1.3\,(b)(v); Theorem~1.4, Part~I\,(d)--(f)]{QEDII}.
For $\xi\in\Xi_1$ and $r>0$ put $\mathsf D_\xi(r):=\int d\mathrm{kl}_\xi(t)\,r/(r+t^2)$, the function $\mathsf D$ of Step~9 of the proof of Proposition~5.6 formed with $\mathrm{kl}_\xi$, as in Step~2 of the proof of Lemma~6.4;
let $\varpi$ be the probability measure on the momenta $k$ with $d\varpi\propto|\hat f|^2\,dk$, which is the Fourier probability measure used in Proposition~4.1.
The number
\begin{equation*}
x_l(\xi):=\frac{S_2^\xi(f_l,f_l)}{\|f\|^2}=\int\mathsf D_\xi(|k|^2/l^2)\,d\varpi(k)
\end{equation*}
is a non-increasing function of $l>0$ with limit $\mathsf z_\xi$, since the computation in Proposition~4.1 is valid for every law whose two-point function has the K\"all\'en--Lehmann form.
Hence $\xi\mapsto\mathsf z_\xi$ is measurable and $x_l^{-1}\uparrow\mathsf z_\xi^{-1}$.
Choose an increasing subsequence of $(l_k)$ along which $I_\mu(f_{l_k})/\|f\|^2$ tends to its lower limit, which may be taken finite, since otherwise there is nothing to prove.
Along it, (K13) and Step~5 at $f_{l_k}$ give
\begin{equation*}
\int\mathsf z_\xi^{-1}\,d\nu=\lim_k\int x_{l_k}^{-1}\,d\nu\le\liminf_k\frac{I_\mu(f_{l_k})}{\|f\|^2}.
\end{equation*}

\emph{Step 7} (consequences).
As $\mathsf z_\xi\le1$, the inequality $\int(\mathsf z_\xi^{-1}-1)\,d\nu\le c$ has a non-negative integrand, and the tail bound is Markov's inequality.
By \cite[Proposition~45.6\,(a$'$)]{QEDIII}, $S_2^\mu(f_l,f_l)=\int S_2^\xi(f_l,f_l)\,d\nu$, and the limit $l\to\infty$ gives
\begin{equation*}
\mathrm{kl}_\mu(\{0\})=\int\mathsf z_\xi\,d\nu\ge\Bigl(\int\mathsf z_\xi^{-1}\,d\nu\Bigr)^{-1}.
\end{equation*}
For volume-first limit points this equality is the case of the atom in Lemma~6.4\,(b).

\emph{Step 8} (the order of the choices).
First $\mu$, the kernel and $\nu$ are fixed;
then the set $\Xi_1$;
then a single probe;
then a countable family of scales $l=l_k$;
for every such $l$, a representative of $\mathsf s_{f_l}$;
then the countable class $\mathcal A_0$, which for (b) may be replaced by the functions $r\tanh(X(f_l)/r)$, $r\in\mathbb N$;
and last of all one exceptional null set, namely the complement of $\Xi_1\cap\bigcap_l\Xi_{f_l}$, which belongs to the $\sigma$-algebra of \cite[(n7)]{QEDIII}.
The function $\Psi$ and the translation $a$ of Step~2 serve an approximation in $L^2(\mu)$ and create no further null set.
The limits $r\to\infty$ and then $l\to\infty$ are taken inside a fixed component, and the integration over $\xi$ is performed last.
The null set depends on $\beta$, whereas the conclusion does not.
\end{proof}

\subsection{Sharpness for Gaussian mixtures}
\begin{named}{Proposition 7.2}
\emph{Optimality; what is false.}
The equality in (b) concerns the full class $\mathcal K(C)$:
Theorem~7.1(c) gives the upper-bound inclusion, and the Gaussian mixtures
in (a) give the reverse inclusion. The examples below are not asserted
to be QED3 cutoff limit points or their components.
\begin{itemize}
\item[(a)] Let $\pi$ be a Borel probability measure on $(0,1]$, and put
\begin{equation*}
\mu_\pi:=\int\gamma_a\,d\pi(a).
\end{equation*}
The law $\mu_\pi$ belongs to the class, and the map $\omega\mapsto-\omega$ leaves it invariant.
The components of $\mu_\pi$ are the laws $\gamma_a$;
its measure $\nu$ is the image of $\pi$ under $a\mapsto\gamma_a$;
the component $\xi=\gamma_a$ has $\mathsf z_\xi=a$ and $C_\xi=1/a$.
For every probe,
\begin{equation*}
I_{\mu_\pi}(f)=\|f\|^2\int a^{-1}\,d\pi(a).
\end{equation*}
\item[(b)] Let $C\ge1$.
As $\mu$ runs through $\mathcal K(C)$, the distribution of $\xi\mapsto\mathsf z_\xi$ under $\nu$ runs \emph{exactly} through the probability measures $\pi$ on $(0,1]$ such that $\int a^{-1}\,d\pi(a)\le C$.
Consequently, inside $\mathcal K(1+c)$, the law of $\mathsf z_\xi$ is subject to no restriction other than $\int\mathsf z_\xi^{-1}\,d\nu\le1+c$;
each $\mu_\pi$ satisfies $\int\mathsf z_\xi^{-1}\,d\nu=C(\mu_\pi)$;
and the tail bound $\nu\{\mathsf z_\xi\le t\}\le ct/(1-t)$ of Theorem~7.1\,(b) is attained for each $0<t<1$ such that $ct/(1-t)\le1$.
\item[(c)] Let $c>0$.
Each of the three statements below is false for some law in $\mathcal K(1+c)$.
The first is the inequality $\mathsf z_\xi\ge(1+c)^{-1}$ for almost every $\xi$:
if $0<a<a'\le1$, the law $\mu=\frac12(\gamma_a+\gamma_{a'})$ has
\begin{equation*}
C(\mu)=\int\mathsf z_\xi^{-1}\,d\nu=\frac{a+a'}{2aa'},\qquad\mathrm{kl}_\mu(\{0\})=\frac{a+a'}2,
\end{equation*}
and $\mathsf z_\xi=a<1/C(\mu)$ on a set of $\nu$-measure $\frac12$.
Take $a$ and $a'$ subject to $\frac{a+a'}{2aa'}=1+c$, for example $a'=1$ and $a=(1+2c)^{-1}$.
Then $C(\mu)=1+c$, whence $\mu\in\mathcal K(1+c)$ by (a), and $1/C(\mu)=(1+c)^{-1}$, so that $\mathsf z_\xi=a<(1+c)^{-1}$ on that set.
The second is a bound $\mathsf z_\xi\ge\delta(c)>0$ for almost every $\xi$, whatever the number $\delta(c)$.
The third, for a given $p>1$, is the integrability $\mathsf z_\xi^{-1}\in L^p(\nu)$.
\item[(d)] The conclusion of Proposition~5.6\,(d), a positive lower bound for $\mathrm{kl}_\mu(\{0\})$, does not imply that $\mathsf z_\xi>0$ for almost every $\xi$.
Take $M>0$ and $\mu=\frac12(\gamma+\gamma^{\delta_M})$;
then $\mathrm{kl}_\mu(\{0\})=\frac12$,
while the components of $\mu$ are $\gamma$, for which $\mathsf z=1$, and $\gamma^{\delta_M}$, for which $\mathsf z=0$.
At every probe the information of $\mu$ is finite,
\begin{equation*}
I_\mu(f_l)=\|f\|^2+\tfrac12M^2l^2\|d^*\beta\|^2,
\end{equation*}
and yet $C(\mu)=\int\mathsf z_\xi^{-1}\,d\nu=\infty$.
\item[(e)] A law $\mathrm{lw}$ may satisfy $(\mathrm{FI})$, in the form $(\mathrm{FI}_{\mathrm{lw}})$, without being Gaussian.
Examples are the laws $\mathrm{lw}=\mu_\xi^\epsilon=\mu_\xi\ast_\epsilon\gamma$ of Proposition~4.4, for the tuned masses, $e\le e_2$, $\xi\in\Xi_\mu$ and $0<\epsilon\le\epsilon_0$:
they are ergodic laws of the class, they are not Gaussian, and they satisfy $C(\mathrm{lw})\le1/\epsilon$;
their measure $\mathrm{kl}_{\mathrm{lw}}$ has at zero the atom $\epsilon+(1-\epsilon)\mathsf z_\xi\ge\epsilon$.
\end{itemize}
\end{named}

\begin{proof}
\emph{Step 1} (Gaussian laws).
Let $\gamma^\varrho$ be one of the laws of Proposition~4.4\,(d), let $\Phi=\phi(X(g_1),\dots,X(g_k))$ be an element of $\mathcal A(\gamma^\varrho)$, and let $g\in L^2_{\mathrm c}$.
Then (K14) yields
\begin{equation*}
\gamma^\varrho\bigl(X(g)\,\Phi\bigr)=\sum_jC_\varrho(g,g_j)\,\gamma^\varrho(\partial_j\phi).
\end{equation*}
Write $\mathsf D_\varrho(r):=\int d\varrho(t)\,r/(r+t^2)$;
this is the function $\mathsf D$ of Step~9 of the proof of Proposition~5.6 with $\varrho$ in the place of $\mathrm{kl}_\mu$,
and by the formula for $C_\varrho$ in Proposition~4.4\,(d) the multiplier of $C_\varrho$ is $\Pi(p)\,\mathsf D_\varrho(|p|^2)$.
Suppose that the form $g$ given by $\hat g:=\hat f/\mathsf D_\varrho(|p|^2)$ belongs to $L^2$.
Then $C_\varrho(g,g_j)=\langle f,g_j\rangle$ for every $j$,
so that $X(g)$ is the score of $\gamma^\varrho$ along $f$, and
\begin{equation*}
I_{\gamma^\varrho}(f)=(2\pi)^{-3}\int|\hat f|^2\,\mathsf D_\varrho(|p|^2)^{-1}\,dp.
\end{equation*}
For $\gamma_a$ this gives $I_{\gamma_a}(f)=\|f\|^2/a$;
for $\gamma^{\delta_M}$, use $\operatorname{div}(d^*\beta)=0$ and $\ast f=\operatorname{curl}(d^*\beta)$.
The Fourier formula in Step~1 then gives $I_{\gamma^{\delta_M}}(f)=\|f\|^2+M^2\|d^*\beta\|^2$.
Finally, $\|d^*\beta^{(l)}\|=l\,\|d^*\beta\|$.

\emph{Step 2} (the law $\mu_\pi$).
Tonelli's theorem shows that, for each $\mathrm{Ev}\in\Sigma$, the map $a\mapsto\gamma_a(\mathrm{Ev})=\allowbreak \gamma\{\omega:a^{1/2}\omega\in\mathrm{Ev}\}$ is measurable;
moreover
\begin{equation*}
\int\exp\bigl(\tfrac a2\|\Pi_{\mathrm c}g\|^2\bigr)\,d\pi(a)\le\exp\bigl(\tfrac12\|g\|^2\bigr).
\end{equation*}
Being linear in the law, condition (dF) of the notation fixed above, the invariances and the inequalities of reflection positivity of \cite[(n3)]{QEDIII} carry over to $\mu_\pi$ from the laws $\gamma_a$, which satisfy them \cite[Lemma~12.29]{QEDII}.
The laws $\gamma_a$ are ergodic,
and distinct values of $a$ give distinct laws;
the uniqueness statement of \cite[Proposition~45.6\,(a$''$)]{QEDIII} therefore identifies $\nu$ with the image of $\pi$.
By Proposition~4.4\,(d), $\mathrm{kl}_{\gamma_a}=a\delta_0$.
The formula for $I_{\mu_\pi}$ is a consequence of Step~1 and Theorem~7.1\,(c).

\emph{Step 3} (proof of (b) and (c)).
Theorem~7.1\,(c) gives one of the two inclusions in (b), and (a) gives the other.
Let $0<t<1$ be such that $p:=\frac{ct}{1-t}\le1$,
and let $\pi=(1-p)\delta_1+p\,\delta_t$;
then $\int a^{-1}\,d\pi(a)=1+c$.
For this measure the tail bound is attained,
and $\operatorname{ess\,inf}\mathsf z_\xi=t$, which can be taken as small as one wishes.
Let $p>1$ and $0<\theta\le\min\{1,c(p-1)\}$;
the measure $\pi=(1-\theta)\delta_1+\theta\,p\,a^{p-1}\,da$ on $(0,1]$ has
\begin{equation*}
\int a^{-1}\,d\pi(a)=1+\frac{\theta}{p-1}\le1+c,\qquad\int a^{-p}\,d\pi(a)=\infty,
\end{equation*}
and (a) places $\mu_\pi$ in $\mathcal K(1+c)$.

\emph{Step 4} (proof of (d) and (e)).
Proposition~4.4(d) shows that $\gamma^{\delta_M}$ is an ergodic law of the class whose measure $\mathrm{kl}$ is $\delta_M$;
with Steps~1 and~2, this proves (d).
For (e), take $F$ and $G$ independent, of the laws $\mu_\xi$ and $\gamma$, as in the definition of $\ast_\epsilon$ in Section~4,
and let $Y:=(1-\epsilon)^{1/2}F+\epsilon^{1/2}G$.
For $\Phi\in\mathcal A(\mathrm{lw})$ and $f\in L^2_{\mathrm c}$, (K14), applied to $G$ while $F$ is held fixed, gives
\begin{equation*}
\epsilon^{1/2}\,E\bigl[(\partial_f\Phi)(Y)\bigr]=E\bigl[\Phi(Y)\,G(f)\bigr];
\end{equation*}
consequently $\mathrm{lw}(\partial_f\Phi)^2\le\epsilon^{-1}\|f\|^2\,\allowbreak \mathrm{lw}(\Phi^2)$.
The other assertions of (e) are contained in Proposition~4.4\,(a)--(c).
\end{proof}

\subsection{Weighted Fisher criteria}
\begin{named}{Proposition 7.3}
\emph{The bound $(\mathrm{FI})$ in the components, deduced from $(\mathrm{FI}_\mu)$ alone.}
Let $\mu\in\mathcal K(C)$.
For example, at any charge $e>0$ and for any mass assignment, every volume-first limit point $\mu$ that satisfies the hypothesis of Proposition~5.6\,(d) belongs to $\mathcal K(1+c)$.
\begin{itemize}
\item[(a)] \emph{Each probe.}
Each $f\in L^2_{\mathrm c}$ satisfies
\begin{equation*}
\int I_\xi(f)\,d\nu=I_\mu(f)\le C\|f\|^2.
\end{equation*}
In particular, for almost every $\xi$, the inequality $\mu_\xi(\partial_f\Phi)^2\le I_\xi(f)\,\mu_\xi(\Phi^2)$ holds for all $\Phi\in\mathcal A(\mu_\xi)$;
and $\nu\{I_\xi(f)>t\|f\|^2\}\le C/t$ for $t>0$.
For a law $\mathrm{lw}$ with (G), a probe $\beta$, $a=d^*\beta$ and $f=da$ with $I_{\mathrm{lw}}(f)<\infty$, the current $J(a)$ is defined by
\begin{equation*}
e\,J(a):=X(f)-\mathsf s_f,
\end{equation*}
where $\mathsf s_f$ is the score of (S2) and $e>0$ is a fixed number;
the product $e\,J(a)$ does not depend on $e$.
If $\mathrm{lw}$ is a limit point at the charge $e$ and $\bar I(\beta)<\infty$, then (S2) yields $\mathsf s_\beta=\mathsf s_f$, and this current is the function $J(a)$ of Proposition~5.6\,(c).
For $f=da$ as above, let $J(a)$ be formed for $\mu$ with a $\Sigma$-measurable representative $\mathsf s$ of $\mathsf s_f$, as in Theorem~7.1\,(a).
Then, for almost every $\xi$, the field equation of Proposition~5.6\,(c) is valid in $\mu_\xi$ with this $J(a)$, and
\begin{equation*}
\int\|e\,J(a)\|^2_{L^2(\mu_\xi)}\,d\nu(\xi)\le(C-1)\,\|f\|^2.
\end{equation*}
For an arbitrary $f\in L^2_{\mathrm c}$, the field equation and this bound hold with $X(f)-\mathsf s$ in the place of $e\,J(a)$, where $\mathsf s$ is a $\Sigma$-measurable representative of $\mathsf s_f$.
Proposition~7.2 shows that integrability in $L^1(\nu)$ is the best possible here.
\item[(b)] \emph{Large scales.}
Fix a probe $\beta$ and define
\begin{equation*}
c_\xi(\beta):=\liminf_{\mathbb N\ni l\to\infty}\frac{I_\xi(f_l)}{\|f\|^2}.
\end{equation*}
Then $\int c_\xi(\beta)\,d\nu\le C$ and $\mathsf z_\xi\ge1/c_\xi(\beta)$.
Hence $c_\xi(\beta)<\infty$ for almost every $\xi$,
and almost every component satisfies the hypothesis of Proposition~5.6\,(d) in the form in which it enters Theorem~7.1\,(b).
\item[(c)] \emph{One null set for all probes; a stronger norm.}
Let $s>\frac32$, and put $\mathrm{ZH}_{2s}:=\sum_\kappa\lambda_\kappa^{-2s}$, with the eigenvalues $\lambda_\kappa$ of \cite[(n8)]{QEDIII}, as at the beginning of \cite[Section~45]{QEDIII};
in \cite[Section~7.1.5]{QEDI} this number is written $\mathsf Z_{2s}$.
One can choose a measurable function $\xi\mapsto T_\xi\ge0$, with $\int T_\xi\,d\nu\le C\,\mathrm{ZH}_{2s}$, such that
\begin{equation*}
I_\xi(f)\le T_\xi\,\inf_a\|\tau_af\|_s^2
\end{equation*}
holds for almost every $\xi$ and \emph{all} $f\in L^2_{\mathrm c}\cap H_s$ simultaneously;
these $f$ include $dd^*\beta$ for every probe $\beta$.
As a consequence, for almost every $\xi$ the law $\mu_\xi$ possesses, at every probe, a score depending linearly on $f$, the current $J(a)\in L^2(\mu_\xi)$ of (a), formed with $\mathrm{lw}=\mu_\xi$, and the bound on shifts of Proposition~5.6\,(c) with $I_\xi(f)$ instead of $\bar I(\beta)$.
Since $\|\cdot\|_s$ is not invariant under the scaling $\beta\mapsto\beta^{(l)}$, the bound in (c) gives no lower bound for $\mathsf z_\xi$.
\item[(d)] \emph{One constant: a criterion.}
One has $C_\xi\ge\mathsf z_\xi^{-1}$.
Moreover $\int C_\xi\,d\nu$ is the smallest constant $C'$ for which
\begin{equation*}
(\mathrm{FI}^I_\mu)\qquad\mu\Bigl(\sum_{k\le m}1_{B_k}\,\partial_{f_k}\Phi\Bigr)^2\le C'\,\mu(\Phi^2)\max_k\|f_k\|^2
\end{equation*}
holds for every finite partition $(B_k)\subset I$, all probes $\beta_k$ and every $\Phi\in\mathcal A_{\mathrm c}$;
for $m=1$ this inequality is $(\mathrm{FI}_\mu)$.
If $\nu(\{\xi_0\})>0$, then $C_{\xi_0}\le C/\nu(\{\xi_0\})$.
If the components of $\mu$ are laws of the form $\gamma^\varrho$, then $C_\xi=\mathsf z_\xi^{-1}$ and $\int C_\xi\,d\nu=C(\mu)$.
\item[(e)] It is not decided whether every $\mu\in\mathcal K(C)$ has $C_\xi<\infty$ for almost every $\xi$.
Invariance under all Euclidean maps, (G) of the notation fixed above, (dF), mixing of the components and $(\mathrm{FI}_\mu)$ do not imply it:
some mixture $\mu$ of Gaussian laws has all these properties, is \emph{not of the class, as it is reflection positive at no plane}, and satisfies $C(\mu)\le2$ and $C_\xi=\infty$ for every $\xi$.
\end{itemize}
\end{named}

\begin{proof}
\emph{Step 1} (proof of (a)).
By (S4), the score $\mathsf s_f$ depends linearly and boundedly on $f\in L^2_{\mathrm c}$, as an element of $L^2(\mu)$.
In Step~2 of the proof of Theorem~7.1 both members of the identity depend continuously on $f\in L^2_{\mathrm c}$;
consequently that step holds for every $f\in L^2_{\mathrm c}$, and so do Steps~3 to~5 of the same proof.
Read for $\mu$, Step~5 gives $\mu(X(f)\,\mathsf s_f)=\|f\|^2$;
expanding the square with the help of (S2), one obtains the equality in
\begin{equation*}
\mu\bigl((X(f)-\mathsf s_f)^2\bigr)=S_2(f,f)-2\|f\|^2+I_\mu(f)\le(C-1)\,\|f\|^2.
\end{equation*}
The inequality follows from $I_\mu(f)\le C\|f\|^2$, by (S4), and from $S_2(f,f)=\mu(X(f)^2)\le\|f\|^2$, a consequence of (G) \cite[Lemma~14.3\,(a)]{QEDII}.

\emph{Step 2} (proof of (b)).
Fatou's lemma, combined with (a), bounds $\int c_\xi(\beta)\,d\nu$.
Steps~5 and~6 of the proof of Theorem~7.1 give $x_l^{-1}\le I_\xi(f_l)/\|f\|^2$ and $x_l^{-1}\uparrow\mathsf z_\xi^{-1}$, whence $\mathsf z_\xi^{-1}\le c_\xi(\beta)$.

\emph{Step 3} (proof of (c)).
With the Hermite functions $h_\kappa$ and the eigenvalues $\lambda_\kappa$ of \cite[(n8)]{QEDIII}, put $e_\kappa:=\Pi_{\mathrm c}h_\kappa$, choose a representative $\varsigma_\kappa$ of $\mathsf s_{e_\kappa}$, and define
\begin{equation*}
T_\xi:=\sum_\kappa\lambda_\kappa^{-2s}\,\xi(\varsigma_\kappa^2).
\end{equation*}
From (a) one obtains $\int T_\xi\,\allowbreak d\nu\le C\sum_\kappa\lambda_\kappa^{-2s}=\allowbreak C\,\allowbreak \mathrm{ZH}_{2s}<\infty$, the number $\mathrm{ZH}_{2s}$ being finite because $2s>3$ \cite[Section~7.1.5]{QEDI};
compare \cite[Proposition~45.1\,(3)]{QEDIII}, where the same number occurs for the same range of $s$.
Take $\xi\in\bigcap_\kappa\Xi_{e_\kappa}$ with $T_\xi<\infty$, where $\Xi_f$ denotes the set of Step~3 of the proof of Theorem~7.1.
Every $f\in L^2_{\mathrm c}\cap H_s$ has the expansion $f=\sum_\kappa\langle h_\kappa,f\rangle\,e_\kappa$ in $L^2$,
and the Schwarz inequality together with \cite[(n8)]{QEDIII} yields
\begin{equation*}
\sum_\kappa|\langle h_\kappa,f\rangle|\,\|\varsigma_\kappa\|_{L^2(\xi)}\le\|f\|_s\,T_\xi^{1/2}.
\end{equation*}
The map $f\mapsto\xi(\partial_f\Phi)$ is linear and continuous for the norm of $L^2$;
hence $\sum_\kappa\langle h_\kappa,f\rangle\,\varsigma_\kappa$ is the score of $\xi$ along $f$.
Translation invariance of $\xi$ gives $I_\xi(\tau_af)=I_\xi(f)$.
For the bound on shifts, Step~8 of the proof of Proposition~5.6 is read in $\mu_\xi$.

\emph{Step 4} (proof of (d), from $(\mathrm{FI}^I_\mu)$).
Step~2 of the proof of Theorem~7.1 turns the left side of $(\mathrm{FI}^I_\mu)$ into $\mu\bigl(\Phi\sum_k1_{B_k}\mathsf s_{f_k}\bigr)^2$.
Taking the supremum over the $\Phi$ with $\mu(\Phi^2)=1$ and using (S1), \cite[Proposition~45.6\,(a$'$)]{QEDIII} and Step~3 of the proof of Theorem~7.1, one finds the value
\begin{equation*}
\sum_k\mu\bigl(1_{B_k}\mathsf s_{f_k}^2\bigr)=\sum_k\int_{B_k}I_{\mu_\omega}(f_k)\,d\mu(\omega).
\end{equation*}
Let $f_1,f_2,\dots$ be an enumeration of the normalised forms $dd^*\beta/\|dd^*\beta\|$ with $\beta\in\mathcal D_0$.
Let $B_k\in I$ be the set of the $\omega$ for which $k$ is the smallest index $j\le m$ at which $I_{\mu_\omega}(f_j)$ is maximal.
With this partition, $\int\max_{j\le m}I_\xi(f_j)\,d\nu\le C'$.
As $m\to\infty$, (K13) and (S4) yield $\int C_\xi\,d\nu\le C'$.

\emph{Step 5} (proof of the other assertions of (d)).
Conversely, $\int_{B_k}I_{\mu_\omega}(f_k)\,\allowbreak d\mu\le\|f_k\|^2\int_{B_k}C_{\mu_\omega}\,\allowbreak d\mu$ for each $k$.
The inequality $C_\xi\ge\mathsf z_\xi^{-1}$ follows from (b).
Since an atom of $\nu$ is contained in no null set, (a) gives $\nu(\{\xi_0\})\,I_{\xi_0}(f)\le C\|f\|^2$ for every $f$.
Let $\varrho(\{0\})>0$.
As $\mathsf D_\varrho\ge\varrho(\{0\})$, Step~1 of the proof of Proposition~7.2 shows $C(\gamma^\varrho)\le\varrho(\{0\})^{-1}$;
Theorem~7.1\,(c) and Proposition~7.3\,(a) then give
\begin{equation*}
\int C_\xi\,d\nu=\int\mathsf z_\xi^{-1}\,d\nu\le C(\mu)\le\int C_\xi\,d\nu.
\end{equation*}

\emph{Step 6} (proof of (e): the law $\mu$).
For $\xi\in[0,1)$ define
\begin{equation*}
m_\xi(p):=1+\sum_{n\ge0}2^n\,1\bigl[2^j\le|p|<2^{j+1},\ j=2^n+\lfloor2^n\xi\rfloor\bigr],
\end{equation*}
denote by $\gamma_{[\xi]}$ the centred Gaussian law whose covariance has the multiplier $\Pi(p)/m_\xi(p)$, and set
\begin{equation*}
\mu:=\int_0^1\gamma_{[\xi]}\,d\xi.
\end{equation*}
By \cite[Lemma~9.11\,(a)--(c)]{QEDI} and \cite[Remark~14.4]{QEDII}, every $\gamma_{[\xi]}$ exists, satisfies (G) and (dF), is invariant under all Euclidean maps, and is mixing;
distinct values of $\xi$ give distinct laws.
On $[0,1)\times\mathbb R^3$ the function $(\xi,p)\mapsto m_\xi(p)$ is Borel.
Indeed, $m_\xi(p)-1$ is the countable sum, taken over the pairs $(n,i)$ with $n\ge0$ and $0\le i<2^n$, of $2^n$ multiplied by the indicator function of the set
\begin{equation*}
[i2^{-n},(i+1)2^{-n})\times\bigl\{p:2^{2^n+i}\le|p|<2^{2^n+i+1}\bigr\}.
\end{equation*}
For each $g\in\mathcal S$, Tonelli's theorem therefore makes the function
\begin{equation*}
\xi\mapsto\operatorname{Cov}_\xi(g,g):=(2\pi)^{-3}\int|\Pi(p)\,\widehat{\ast g}(p)|^2\,m_\xi(p)^{-1}\,dp
\end{equation*}
Borel, and with it $\xi\mapsto\gamma_{[\xi]}(\exp iX(g))=\allowbreak \exp\bigl(-\frac12\operatorname{Cov}_\xi(g,\allowbreak g)\bigr)$.
Denote by $\mathrm{Mv}$ the class of the bounded real $\Sigma$-measurable functions $\Phi$ such that $\xi\mapsto\gamma_{[\xi]}(\Phi)$ is Borel.
This class is a linear space containing the constants, and it contains the limit of every bounded increasing sequence of its elements.
A finite product of functions $\cos X(g)$ and $\sin X(g)$, $g\in\mathcal S$, is a linear combination of functions of the same kind;
hence $\mathrm{Mv}$ contains the set $\mathrm{Mu}$ of all such products.
Moreover $\sigma(\mathrm{Mu})=\Sigma$, because $X(g)=\lim_nn\sin X(g/n)$.
The monotone class theorem (K13) shows that every bounded $\Sigma$-measurable function lies in $\mathrm{Mv}$.
Consequently $\mu$ is a mixture as defined in \cite[(n6)]{QEDIII},
and the map $\xi\mapsto\gamma_{[\xi]}$ is measurable with respect to the $\sigma$-algebra of \cite[(n7)]{QEDIII}.
Under this map, Lebesgue measure on $[0,1)$ has an image that is concentrated on ergodic laws and represents $\mu$;
by \cite[Proposition~45.6\,(a$''$)]{QEDIII} this image is $\nu$,
and the laws $\gamma_{[\xi]}$ are the components of $\mu$.

The law $\mu$ does not belong to the class.
The multiplier of $S_2^\mu$ is $\Pi(p)M(p)$, with
\begin{equation*}
M(p):=\int_0^1\frac{d\xi}{m_\xi(p)}=
\begin{cases}
1&\text{if }|p|<2,\\
\dfrac{2^n}{1+2^n}&\text{if }2^j\le|p|<2^{j+1}\text{ and }2^n\le j<2^{n+1}.
\end{cases}
\end{equation*}
If $\mu$ were a law of the class, \cite[Theorem~1.4, Part~I\,(d), (f)]{QEDII} would yield $M(p)=\mathsf D(|p|^2)$ for Lebesgue-almost all $p$, where $\mathsf D(r)=\int d\mathrm{kl}_\mu(t)\,r/(r+t^2)$ is a non-decreasing function of $r$;
whereas $M=1$ on $\{|p|<2\}$ and $M=\frac12$ on $\{2\le|p|<4\}$.
The law $\mu$ is invariant under all Euclidean maps and satisfies (G) and (dF);
among the conditions of (H) and of $\mathcal C_9$ \cite[(n3)]{QEDIII}, reflection positivity is therefore the only one that $\mu$ can fail.
Let $r_{\hat n,\mathrm{po}}$ be the reflection in the plane $\{\hat n\cdot y=\mathrm{po}\}$,
and let $\mathrm{Gi}$ be a Euclidean map that carries this plane onto the plane $\{\hat n'\cdot y=\mathrm{po}'\}$.
Because $r_{\hat n',\mathrm{po}'}=\allowbreak \mathrm{Gi}\,\allowbreak r_{\hat n,\mathrm{po}}\,\allowbreak \mathrm{Gi}^{-1}$ and $\mu$ is invariant, reflection positivity of $\mu$ at a single plane would entail reflection positivity at every plane.
Therefore $\mu$ is reflection positive at no plane.

\emph{Step 7} (proof of (e): the constants).
The computation of Step~1 of the proof of Proposition~7.2 gives $I_{\gamma_{[\xi]}}(f)=(2\pi)^{-3}\int m_\xi\,|\hat f|^2\,dp$ whenever $m_\xi\hat f\in L^2$.
For every $p$ one has $\int_0^1m_\xi(p)\,d\xi\le2$,
and Theorem~7.1\,(c) then gives $C(\mu)\le2$.
Choosing forms $f$ whose $\hat f$ is carried by the shell $\{2^j\le|p|<2^{j+1}\}$ with $j=2^n+\lfloor2^n\xi\rfloor$, and applying (S4), one obtains $C(\gamma_{[\xi]})\ge1+2^n$ for every $n$.
\end{proof}

\subsection{One constant for all probes}
\begin{named}{Theorem 7.4}
\emph{A single constant valid at every probe, for the volume-first limit points of the model.}
Let $0<\rho<1$, let the weights be standard, let $\bar m\in(0,1]$ and $e>0$ be arbitrary, and let the mass assignment be such that $\boldsymbol\mu_N>0$ for all large $N$.
Let $\mu$ be a volume-first limit point,
\begin{equation*}
\mu=\lim_i\mu^{(N_i)},\qquad\mu^{(N_i)}=\lim_jP_{N_i,T_{ij}},
\end{equation*}
and define
\begin{equation*}
c_W:=\liminf_{i\to\infty}\,3\pi^2\mathsf h_3^2\,q_i^2\,\liminf_{j\to\infty}\mathsf S_{n_{ij}},
\end{equation*}
where $\eta_i=L^{-N_i}$, $q_i=e\,\eta_i^{1/2}$ and $n_{ij}$ are the spacing, the lattice charge and the number of sites per side at the cutoff $(N_i,T_{ij})$ (the notation fixed above),
and $\mathsf S_{n_{ij}}$ is the number $\mathsf S_n$ of Section~5 formed with $q=q_i$ and $\mathbf m=\boldsymbol\mu_{N_i}$.
Then
\begin{equation*}
\int\mathsf z_\xi^{-1}\,d\nu\le\int C_\xi\,d\nu(\xi)\le1+c_W.
\end{equation*}
In other words, if $c_W<\infty$, then for $\nu$-almost every $\xi$ the component $\mu_\xi$ fulfils the condition $(\mathrm{FI}_{\mu_\xi})$ at every probe and for every $\Phi\in\mathcal A(\mu_\xi)$ with one and the same constant $C_\xi<\infty$, and $C_\xi\in L^1(\nu)$.
Among all laws of the class, the condition $(\mathrm{FI}^I_\mu)$ of Proposition~7.3\,(d) implies no integrability of $C_\xi$ beyond this one:
for each $p>1$ there is a law of the class which satisfies $(\mathrm{FI}^I_\mu)$ and for which $C_\xi\notin L^p(\nu)$ (Proposition~7.2\,(c), Proposition~7.3\,(d)).
For the volume-first limit points formed with a physical mass assignment, or with the tuned masses of \cite{QEDIII}, it is not decided whether $C_\xi$ belongs to $L^p(\nu)$ for some $p>1$.
If the hypothesis of Proposition~5.13\,(e) is satisfied, then $c_W\le\frac12C\,e^2/\bar m$, with the constant $C$ of that proposition.
By Theorem~5.14\,(a), for every charge $e>0$ and every mass assignment such that $c_{\boldsymbol\mu}\bar m\eta\le\boldsymbol\mu_N\le1$ for all large $N$ and some $c_{\boldsymbol\mu}>0$ (every physical mass assignment is of this kind),
\begin{equation*}
c_W\le\tfrac12C_F\,\frac{e^2}{\bar m};
\end{equation*}
here $C_F$ is the number of Theorem~5.14\,(b), whose dependence on $e^2/(c_{\boldsymbol\mu}\bar m)$ is described in Theorem~A.
\end{named}

\begin{proof}
Put $\lambda_i:=(J_{\eta_i})_*\mu^{(N_i)}$.
By \cite[Proposition~12.15\,(a)--(c)]{QEDII}, the laws $\mu_{C(n_{ij})}$ converge weakly on $\Omega$ to $\lambda_i$ as $j\to\infty$,
the law $\mu^{(N_i)}=(\sharp_{\eta_i})_*\lambda_i$ has the property (G) of the notation fixed above,
and $\lambda_i$ is invariant under $W=\mathbb Z^3\rtimes O_{\mathrm{cub}}$.
Throughout Steps~1 to~6 the number $N=N_i$ is held fixed, with $\mathbf m=\boldsymbol\mu_N>0$;
one writes $\lambda=\lambda_i$, and $n=n_{ij}$ runs through a subsequence along which $\mathsf S_{n_{ij}}$ converges to $\liminf_j\mathsf S_{n_{ij}}$.
For a real field $A$ on $C(n)$, one-cochains $u,u'$, links $b,b'$ and sites $x,y$ set
\begin{equation*}
k_A(u,u'):=-\partial_u\partial_{u'}\ell_n(A),\qquad k_A(b,b'):=k_A(\delta_b,\delta_{b'}),\qquad\mathsf n_A(x,y):=\|D_A^{-1}(x,y)\|^2_{\mathrm{HS}}.
\end{equation*}
A function is called \emph{local} if it depends smoothly on finitely many of the $F_\sigma$ and is bounded together with all its derivatives.
On $C(n)$ with $n$ large, local functions and the one-cochains $u$, $u_\alpha$ of finite support on $\mathbb Z^3$ are read in the manner of Step~4 of the proof of Proposition~5.3.
For $g\in W$ and a function $G$ one writes $(gG)(F):=G(g^*F)$.

\emph{Step 1} (field by field).
In each real field $A$ on $C(n)$ the following three statements hold.
\begin{itemize}
\item[(i)] There is a bound on $k_A(b,b')$ and on $\mathsf n_A(x,y)$ that does not depend on $n$ (Lemma~5.1; Step~1 of the proof of Lemma~5.8).
\item[(ii)] One has $\mathsf n_A(x,y)=\mathsf n_A(y,x)$, and whenever $b\ne b'$,
\begin{equation*}
|k_A(b,b')|\le\mathsf h_3^2q^2\sum_{x'\in V(b),\,y'\in V(b')}\mathsf n_A(x',y'),
\end{equation*}
which is the inequality of Steps~1 and~2 of the proof of Lemma~5.9 before the expectation is formed.
\item[(iii)] Whenever $\vartheta$ is a zero-cochain of finite support, $k_A(u,d\vartheta)=0$;
moreover, for each $g\in W$,
\begin{equation*}
k_{g^*A}(g^*u,g^*u')=k_A(u,u'),\qquad\mathsf n_{(g^{-1})^*A}(gx,gy)=\mathsf n_A(x,y)
\end{equation*}
(Step~2 of the proof of Lemma~5.8; Step~1 of the proof of Lemma~5.9).
\end{itemize}

\emph{Step 2} (a weighted identity).
For every local $h$,
\begin{equation*}
E_n[h\,s_u^2]=E_n[h]\,\|du\|^2+E_n[h\,k_A(u,u)]+E_n[\partial_u^2h].
\end{equation*}
To see this, take the relation $E_n[\partial_u\Phi]=-E_n[\Phi\,s_u]$ of Step~2 of the proof of Proposition~5.6 once with $\Phi=h\,s_u$, using $\partial_us_u=-\|du\|^2-k_A(u,u)$ (Step~2a\,(i) of that proof),
and once with $\Phi=\partial_uh$.

\emph{Step 3} (weights).
Consider local functions $\Psi_1,\dots,\Psi_m\ge0$ with $\sum_\alpha\Psi_\alpha\le1$, let $\Lambda_\ell:=\{x\in\mathbb Z^3:|x|_\infty\le\ell\}$, and define
\begin{equation*}
h_\alpha=h_{\alpha,\ell}:=\frac{1}{48\,|\Lambda_\ell|}\sum_{g_0\in O_{\mathrm{cub}}}\ \sum_{x\in\Lambda_\ell}(\tau_xg_0)\Psi_\alpha.
\end{equation*}
These functions satisfy $h_\alpha\ge0$ and $\sum_\alpha h_\alpha\le1$;
for each $g\in W$ the number $\sup|gh_\alpha-h_\alpha|$ tends to $0$ as $\ell\to\infty$;
and, the number of $x$ that contribute being bounded,
\begin{equation*}
\sup|\partial_uh_\alpha|+\sup|\partial_u^2h_\alpha|\le\frac{\mathsf c}{|\Lambda_\ell|}.
\end{equation*}
Let $G$ be local.
Applying Step~2 of the proof of Proposition~5.6 with $\Phi=h_\alpha G$, and then the Schwarz inequality together with $(\sum_\alpha h_\alpha s_{u_\alpha})^2\le\sum_\alpha h_\alpha s_{u_\alpha}^2$, one finds
\begin{equation*}
\Bigl|E_n\sum_\alpha h_\alpha\,\partial_{u_\alpha}G\Bigr|\le(E_nG^2)^{1/2}\Bigl(\sum_\alpha E_n[h_\alpha s_{u_\alpha}^2]\Bigr)^{1/2}+\frac{\mathsf c'}{|\Lambda_\ell|}.
\end{equation*}

\emph{Step 4} (limit kernels).
On account of (i) and of a diagonal argument, letting first $j\to\infty$ and then $\ell=\ell_r\to\infty$, one obtains for all links $b,b'$ of $\mathbb Z^3$ and all sites $x,y$ of $\mathbb X$ limits
\begin{equation*}
E_n[h_\alpha k_A(b,b')]\to k^{(\alpha)}(b,b'),\qquad E_n[h_\alpha\mathsf n_A(x,y)]\to\mathsf n^{(\alpha)}(x,y)\ge0.
\end{equation*}
The expectation $E_n$ is invariant under $W$ (Step~2 of the proof of Lemma~5.8), so that (iii) yields $E_n[h\,k_A(g^*u,g^*u')]=E_n[(gh)\,k_A(u,u')]$;
in view of Step~3, the kernels $k^{(\alpha)}$ and $\mathsf n^{(\alpha)}$ are therefore invariant under $W$.
From (ii) and (iii) it follows that $k^{(\alpha)}$ is real and symmetric, that it annihilates each $d\delta_y$, and that
\begin{equation*}
|k^{(\alpha)}(b,b')|\le\mathsf h_3^2q^2\sum_{x'\in V(b)}\ \sum_{y'\in V(b')}\mathsf n^{(\alpha)}(x',y')\qquad(b\ne b').
\end{equation*}
The group $W$ acts transitively on the sites, and hence
\begin{equation*}
\sigma^{(\alpha)}:=\sum_y|v(x,y)|^2\,\mathsf n^{(\alpha)}(x,y)
\end{equation*}
is the same number for every site $x$.
By Fatou's lemma, by $\sum_\alpha h_\alpha\le1$ and by Lemma~5.9\,(a),
\begin{equation*}
\sum_\alpha\sigma^{(\alpha)}\le\frac18\liminf_{j\to\infty}\mathsf S_{n_{ij}}.
\end{equation*}

\emph{Step 5} (Lemma~5.9\,(c) and Lemma~5.8 on $\mathbb Z^3$).
With the properties of $k^{(\alpha)}$ and $\mathsf n^{(\alpha)}$ obtained in Step~4 and the geometry of the links, the argument of Steps~3 and~4 of the proof of Lemma~5.9 applies to $k^{(\alpha)}$ and gives
\begin{equation*}
\mathsf c^{(\alpha)}:=\max_b\sum_{b'}|s(b,b')|^2\,|k^{(\alpha)}(b,b')|\le192\,\mathsf h_3^2q^2\sigma^{(\alpha)}.
\end{equation*}
Form $\hat k_{\mu\nu}=\hat k^{(\alpha)}_{\mu\nu}$ and $m_{\mu\nu}$ from $k^{(\alpha)}$ as in Steps~4 and~5 of the proof of Lemma~5.8.
Plancherel's formula on $\mathbb Z^3$, contained in (K8) of Section~5, takes in polarised form the shape
\begin{equation*}
k^{(\alpha)}(u,u)=(2\pi)^{-3}\int_{[-\pi,\pi]^3}\hat u^*\,\hat k^{(\alpha)}\,\hat u\,dp
\end{equation*}
for every $u$ of finite support.
With this formula, Steps~3 to~6 of the proof of Lemma~5.8 give four facts:
the matrix $\hat k^{(\alpha)}(p)$ is real and symmetric;
$\hat k^{(\alpha)}\mathsf s=0$;
$|\hat k_{\mu\nu}(p)-\hat k_{\mu\nu}(0)|\le\frac12|p|^2m_{\mu\nu}$;
and $\hat k_{\mu\nu}(te_\nu)=0$ whenever $0<t\le\pi$.
Continuity gives $\hat k^{(\alpha)}(0)=0$, whence $|\hat k_{\mu\nu}(p)|\le\frac12|p|^2m_{\mu\nu}$ at every $p$;
on a torus, where the mode $p=0$ contributes, Step~5 of the proof of Lemma~5.8 yields at $p\ne0$ the bound $|p|^2m_{\mu\nu}$, twice as large.
As $|p|\le\frac\pi2|\mathsf s(p)|$, it follows that
\begin{equation*}
|k^{(\alpha)}(u,u)|\le\frac{\pi^2}{8}\,\mathsf c^{(\alpha)}\|du\|^2\le24\pi^2\mathsf h_3^2q^2\sigma^{(\alpha)}\|du\|^2.
\end{equation*}

\emph{Step 6} (fixed $N$).
In the inequality of Step~3 let $j\to\infty$, local functions being bounded and continuous on $\Omega$, and afterwards let $\ell=\ell_r\to\infty$.
With
\begin{equation*}
\bar\Psi_\alpha:=\frac1{48}\sum_{g_0\in O_{\mathrm{cub}}}g_0\Psi_\alpha,
\end{equation*}
(K15) shows that $h_{\alpha,\ell}\to P_\lambda\bar\Psi_\alpha$ in $L^2(\lambda)$, where $P_\lambda$ stands, as in the Conventions of Section~5 and in Proposition~5.17, for the projection of $L^2(\lambda)$ onto the vectors invariant under $\mathbb Z^3$;
as $\tau_x$ corresponds to $U(\eta x)$ under $\Psi\mapsto\Psi\circ\sharp_\eta$ (Step~7), (K15), used in $L^2(\lambda)$ and in $L^2(\mu^{(N_i)})$, gives $P_\lambda(\Psi\circ\sharp_\eta)=(P_i\Psi)\circ\sharp_\eta$ for $\Psi\in\mathcal A_{\mathrm c}$, with the projection $P_i$ of Step~8.
Steps~2, 4 and~5 thus lead to
\begin{equation}
\begin{split}
\lambda\Bigl(\sum_\alpha(P_\lambda\bar\Psi_\alpha)\,\partial_{u_\alpha}G\Bigr)^2
&\le\lambda(G^2)\sum_\alpha\bigl[\lambda(\Psi_\alpha)+24\pi^2\mathsf h_3^2q^2\sigma^{(\alpha)}\bigr]\,\|du_\alpha\|^2\\
&\le\lambda(G^2)\Bigl(1+3\pi^2\mathsf h_3^2q^2\liminf_{j\to\infty}\mathsf S_{n_{ij}}\Bigr)\max_\alpha\|du_\alpha\|^2.
\end{split}
\tag{7.1}\label{eq:weighted-lattice}
\end{equation}

\emph{Step 7} (data for $N\to\infty$).
Let $\Psi_\alpha=\psi_\alpha(X(g_1'),\dots)$ and $\Phi=\phi(X(g_1),\dots,X(g_k))$ be elements of $\mathcal A_{\mathrm c}$, so that $\psi_\alpha$, like $\phi$, is smooth and bounded together with all its derivatives, and suppose that $\psi_\alpha\ge0$ and $\sum_\alpha\psi_\alpha\le1$,
all the forms $g_c'$ and $g_j$ lying in $C_c^\infty((-\mathrm{vp},\mathrm{vp})^3;\Lambda^2)$;
let further $\beta_\alpha$ be probes, and $f_\alpha:=dd^*\beta_\alpha$.
Since $X(g)\circ\sharp_\eta=\langle g_{[\eta]},\cdot\rangle$, the functions $\Psi_\alpha\circ\sharp_\eta$ and $G:=\Phi\circ\sharp_\eta$ are local.
The operators $\hat w$ by which the elements $w\in O_{\mathrm{cub}}$ act on functions on $\mathcal S'$ correspond to $w$, and the operator $U(\eta x)$ corresponds to $\tau_x$ \cite[Lemma~12.6\,(a)]{QEDII}.
The one-cochain $u_\alpha:=\eta^{1/2}a^\circ_\eta[\beta_\alpha]$ is finitely supported, and $du_\alpha=(f_\alpha^{(\eta)})_{[\eta]}$ \cite[Proposition~47.1\,(b1), (b2$^\circ$)]{QEDIII}, \cite[Lemma~12.5\,(b)]{QEDII}.
Consequently, as $N\to\infty$, one has $\|du_\alpha\|\le\|f_\alpha^{(\eta)}\|\to\|f_\alpha\|$ and
\begin{equation*}
\langle g_{[\eta]},du_\alpha\rangle=\langle\mathsf P_\eta g,f_\alpha^{(\eta)}\rangle\to\langle g,f_\alpha\rangle,
\end{equation*}
with $\mathsf P_\eta$ the projection onto the forms which are constant on each cube $Q_\eta(\sigma)$ \cite[Proposition~47.1\,(b3)]{QEDIII}.
Hence, uniformly,
\begin{equation*}
\partial_{u_\alpha}G=(\partial_{f_\alpha}\Phi)\circ\sharp_\eta+o(1).
\end{equation*}

\emph{Step 8} (an ergodic theorem uniform in the spacing).
Denote by $P_i$ the projection of $L^2(\mu^{(N_i)})$ onto the vectors invariant under $\eta_i\mathbb Z^3$,
and, for $\mathsf L\in\mathbb N$, by $A^\eta_{\mathsf L}$ the average of the operators $U(a)$ over the points $a\in\eta\mathbb Z^3$ with $|a|_\infty\le\mathsf L$.
Then $\|A^\eta_{\mathsf L}\bar\Psi_\alpha-P_i\bar\Psi_\alpha\|_2^2$ equals the average, over the pairs $(a,a')$ of such points, of the numbers
\begin{equation*}
\langle\bar\Psi_\alpha,U(a'-a)\bar\Psi_\alpha\rangle-\langle\bar\Psi_\alpha,P_i\bar\Psi_\alpha\rangle.
\end{equation*}
Proposition~5.17, taken with $\mathsf r=\lceil\mathrm{vp}/\eta\rceil+2$, bounds each of these numbers by
\begin{equation*}
\min\Bigl\{2,\ \frac{3C_\flat(\mathrm{vp}+3)}{|a-a'|_\infty}\Bigr\}.
\end{equation*}
For $\mathsf L\ge\mathrm{vp}+3$ the average is therefore at most $\mathsf c\,(\mathrm{vp}+3)/\mathsf L$, where the constant $\mathsf c$ is absolute; the bound holds \emph{uniformly in $i$}.

\emph{Step 9} (the limit).
Let $A_{\mathsf L}$ be the average of $U(a)$ over $a\in[-\mathsf L,\mathsf L]^3$.
Condition (G) for $\mu^{(N_i)}$ gives
\begin{equation*}
\|(A^\eta_{\mathsf L}-A_{\mathsf L})\bar\Psi_\alpha\|_2\le\operatorname{Lip}(\psi_\alpha)\sum_c\ \sup_{|y|_\infty\le\eta}\|\tau_yg_c'-g_c'\|+\frac{\mathsf c\,\eta}{\mathsf L},
\end{equation*}
a quantity which tends to $0$ as $i\to\infty$;
moreover the function $(A_{\mathsf L}\bar\Psi_\alpha)\,\partial_{f_\alpha}\Phi$ is bounded and continuous on $H_{-4}$, since $|\omega(g)|\le\|\omega\|_{-4}\,\|g\|_4$ whenever $\omega\in H_{-4}$ and $g\in\mathcal S$ \cite[Section~12.5]{QEDII}.
Write $P:=P^{123}_\mu$.
By \cite[Lemma~8.13\,(iii), (iv)]{QEDI} one has $P=E_\mu[\,\cdot\,|I]$ and $A_{\mathsf L}\bar\Psi_\alpha\to P\bar\Psi_\alpha$ in $L^2(\mu)$;
moreover $P\bar\Psi_\alpha=P\Psi_\alpha$, because $\nu$-almost every component is invariant under $O_{\mathrm{cub}}$ \cite[Theorem~1.3\,(b)(v)]{QEDII}.
As to the convergence, let $A^{123}_{2\mathsf L}$ be the average of $U(a)$ over $a\in[0,2\mathsf L]^3$ which enters that lemma;
then
\begin{equation*}
A_{\mathsf L}=U\bigl(-\mathsf L(e_1+e_2+e_3)\bigr)\,A^{123}_{2\mathsf L},\qquad U(y)P=P,
\end{equation*}
and consequently $\|A_{\mathsf L}\Psi-P\Psi\|_2=\allowbreak \|A^{123}_{2\mathsf L}\Psi-P\Psi\|_2\to0$ for each $\Psi\in L^2(\mu)$.
Together with Steps~7 and~8, an argument with three epsilons shows that the left side of \eqref{eq:weighted-lattice} at $N_i$ converges to $\mu\bigl(\sum_\alpha(P\Psi_\alpha)\,\allowbreak \partial_{f_\alpha}\Phi\bigr)^2$, and $\lambda_i(G^2)\to\mu(\Phi^2)$.
Letting $i\to\infty$ along a subsequence that realises $c_W$, one arrives at
\begin{equation}
\mu\Bigl(\sum_\alpha(P\Psi_\alpha)\,\partial_{f_\alpha}\Phi\Bigr)^2\le(1+c_W)\,\mu(\Phi^2)\max_\alpha\|f_\alpha\|^2.
\tag{7.2}\label{eq:weighted-continuum}
\end{equation}

\emph{Step 10} (partitions).
As a function of $(\Psi_\alpha)\in L^2(\mu)^m$, the left side of \eqref{eq:weighted-continuum} is continuous.
Let $(B_\alpha)\subset I$ be a finite partition; by (S1) one may choose $\Psi_\alpha'\in\mathcal A_{\mathrm c}$ converging to $1_{B_\alpha}$ in $L^2(\mu)$.
Fix smooth functions $\sigma_\delta\colon\mathbb R\to[0,1]$ with $|\sigma_\delta(t)-t|\le\delta$ on $[0,1]$ and $\zeta_\delta$ with $\max(1,t)\le\zeta_\delta(t)\le\max(1,t)+\delta$, and set
\begin{equation*}
\Psi_\alpha:=\frac{\sigma_\delta(\Psi_\alpha')}{\zeta_\delta\bigl(\sum_{\alpha'}\sigma_\delta(\Psi_{\alpha'}')\bigr)}.
\end{equation*}
These functions are admissible in Step~7, and, the map $(t_{\alpha'})_{\alpha'}\mapsto\sigma_\delta(t_\alpha)/\zeta_\delta(\sum_{\alpha'}\sigma_\delta(t_{\alpha'}))$ being bounded and continuous, $\limsup\|\Psi_\alpha-1_{B_\alpha}\|_2\le\mathsf c\,m\,\delta$.
Since $P1_{B_\alpha}=1_{B_\alpha}$, the inequality \eqref{eq:weighted-continuum} yields in the limit $(\mathrm{FI}^I_\mu)$ with $C'=1+c_W$.
Suppose first that $c_W<\infty$.
For $m=1$ and $B_1=\mathcal S'$, the condition $(\mathrm{FI}^I_\mu)$ is the inequality
\begin{equation*}
\mu(\partial_f\Phi)^2\le(1+c_W)\,\|f\|^2\,\mu(\Phi^2),
\end{equation*}
valid for every probe and every $\Phi\in\mathcal A_{\mathrm c}$.
According to (S3), in the supremum which defines $I_\mu$ the function $\Phi$ may be restricted to $\mathcal A_0\subset\mathcal A_{\mathrm c}$; thus $C(\mu)\le1+c_W$.
By Step~1 of the proof of Theorem~7.1, $\mu$ is a law of the class.
So $\mu\in\mathcal K(1+c_W)$, and Proposition~7.3\,(d) is applicable.
Suppose, finally, that $c_W=\infty$.
The right-hand inequality of the theorem then holds trivially.
The left-hand one is proved by means of Theorem~7.1\,(c), which is applicable since $\mu$ is a law of the class also when $c_W=\infty$:
for a probe $\beta$ and a sequence $l_k\to\infty$ it gives $\int\mathsf z_\xi^{-1}\,d\nu\le\liminf_kI_\mu(f_{l_k})/\|f\|^2$,
and in $[0,\infty]$
\begin{equation*}
I_\mu(f_{l_k})=\int I_\xi(f_{l_k})\,d\nu\le\|f\|^2\int C_\xi\,d\nu,
\end{equation*}
because $\|f_{l_k}\|=\|f\|$ and $I_\xi(f_{l_k})\le C_\xi\|f_{l_k}\|^2$ by the definition of $C_\xi$;
the two bounds together give $\int\mathsf z_\xi^{-1}\,d\nu\le\int C_\xi\,d\nu$.

\emph{Step 11} (the constants).
Under the lower bound $\boldsymbol\mu_N\ge c_{\boldsymbol\mu}\bar m\eta$, which is part of the hypotheses of Proposition~5.13\,(e) and of Theorem~5.14,
\begin{equation*}
q^2\mathsf S_n\le\frac{e^2}{c_{\boldsymbol\mu}\bar m}\,\boldsymbol\mu_N\mathsf S_n.
\end{equation*}
Under the hypothesis of Proposition~5.13\,(e), let $C_M$ be the number of that hypothesis, so that $\limsup_T\boldsymbol\mu_N\mathsf S_n\le C_M$ for all large $N$.
Since $T_{ij}\to\infty$ as $j\to\infty$, one obtains for all large $i$
\begin{equation*}
3\pi^2\mathsf h_3^2\,q_i^2\liminf_{j\to\infty}\mathsf S_{n_{ij}}\le\frac{3\pi^2\mathsf h_3^2C_M\,e^2}{c_{\boldsymbol\mu}\bar m}=\frac12\,C\,\frac{e^2}{\bar m},\qquad C=\frac{6\pi^2\mathsf h_3^2C_M}{c_{\boldsymbol\mu}},
\end{equation*}
and $C$ is the constant of Proposition~5.13\,(e); hence $c_W\le\frac12Ce^2/\bar m$.
For the mass assignments of Theorem~5.14, part~(a) of that theorem gives $\boldsymbol\mu_N\mathsf S_n\le\frac13C_M$ for all large $N$ and all large $T$, where $C_M$ denotes the number defined in that part and not the datum of Proposition~5.13\,(e);
finally, $2\pi^2\mathsf h_3^2C_M/c_{\boldsymbol\mu}=C_F$.
\end{proof}

\section{Proof of Theorem A and Corollary B}
\begin{proof}[Proof of Theorem A]
\emph{Step 1} (hypotheses).
The limit point $\mu$ has defining sequences whose $T_i$ are as large as one wishes (Step~5 of the proof of Proposition~5.13);
the lattice probes satisfy $\sum_xu(x)=0$ \cite[Proposition~47.1\,(b1)]{QEDIII};
and $\bar I(\beta)=\allowbreak \|f\|^2+\liminf_iK_{n_i}(\hat{\mathfrak a}_i,\allowbreak \hat{\mathfrak a}_i)$ (Proposition~5.6(a) and Step~3 of its proof).
Theorem~5.14\,(b) therefore yields the hypothesis of Proposition~5.6\,(d) with $c=c_F$, and Theorem~7.1\,(b) then gives the atom.
The bound $\bar I(\beta)\le(1+c_F)\|f\|^2$ at all probes, which is that hypothesis, is contained in Lemma~6.4\,(a), with $c_{\boldsymbol\mu}\bar m$ in the place of the constant $c$ of that statement.
Theorem~5.14\,(a) gives $c_W\le\frac12c_F$ in Theorem~7.4.
Put $\Xi^\gamma_\mu:=\Xi_1\cap\{C_\xi<\infty\}$, with the set $\Xi_1$ of Step~6 of the proof of Theorem~7.1. By Theorem~7.4 this set has full $\nu$-measure.
Fix $\xi\in\Xi^\gamma_\mu$ and put $\mathrm{lw}=\mu_\xi$ for the remaining steps.
\emph{Step 2} ((i) and (ii)).
Assertion (i) is \cite[Theorem~1.4, Part~I\,(e), (f), Part~II]{QEDII}.
Since $C_\xi<\infty$, Proposition~7.3(a) and Step~1 of its proof give
$\|f\|^4\le S_2^\xi(f,f)I_\xi(f)$ for every real $f\in L^2_{\mathrm c}$.
Applying this inequality to one probe gives $x_l(\xi)^{-1}\le I_\xi(f_l)/\|f\|^2\le C_\xi$ for all $l$, while $x_l^{-1}\uparrow\mathsf z_\xi^{-1}$ (Step~6 of the proof of Theorem~7.1);
hence $\mathsf z_\xi^{-1}\le C_\xi$.
The tail bound is Markov's inequality, applied to $\mathsf z_\xi^{-1}-1$.
The description of $\mathsf z_\xi$ as a best constant is Proposition~4.1, which holds for every law whose second moment has the K\"all\'en--Lehmann form.

\emph{Step 3} (infrared limit).
For each $\xi$ in this full-measure set, the law $\mu_\xi$ is ergodic
and belongs to $\mathcal K$. Apply Theorem~4.3(b), with
$c_0=\mathrm{kl}_\xi(\{0\})=z_\xi$, to obtain exactly the
characteristic-functional convergence in Theorem~A(iv).

Put
\[
\mathcal H_{\mathrm E}^{(1)}=
\overline{\{[X(g)]:g\in C_c^\infty(\{y_1>0\};\Lambda^2\mathbb C^3)\}},
\qquad
\mathcal H_{\mathrm E,\gamma}=E(\partial V)\mathcal H_{\mathrm E}^{(1)}.
\]
Step~11 identifies these Euclidean spaces with the spaces
$\mathcal H^{(1)}$ and $\mathcal H_\gamma$ in the statement.

\emph{Step 4} (shell and group).
The law $\mathrm{lw}$ has the properties listed in Proposition~3.1, and Theorems~3.2 and~3.3 hold for such laws (Remark after Proposition~3.1).
From $\mathcal U(0,\allowbreak \Lambda)\mathcal U(a,\allowbreak 1)\mathcal U(0,\allowbreak \Lambda)^{-1}=\allowbreak \mathcal U(\Lambda a,\allowbreak 1)$, from Theorem~3.2\,(a) and from (K16) one obtains
\begin{equation*}
\mathcal U(0,\Lambda)\,E(\Delta)\,\mathcal U(0,\Lambda)^{-1}=E\bigl((\Lambda^T)^{-1}\Delta\bigr),
\end{equation*}
and $\Lambda^T\in SO_0(1,2)$ preserves $\partial V$ and $\{0\}$.
By the property (Sp) of Proposition~3.1, $E$ is carried by $\{\omega\ge|\vec p|\}$;
hence $E(\{0\})=1_{\{0\}}(H)=|\Omega\rangle\langle\Omega|$ by the property (Vac) of Proposition~3.1.

\emph{Step 5} (the vectors $\psi_h$).
Proposition~4.2\,(a)--(c) is stated for every probability measure with (G) and (dF) of the notation fixed above, the conditions $(T_{123})$, $(R_1)$, $(RP_1)$ of Section~2, and the form (C7) of \cite[Theorem~44]{QEDIII};
the identities obtained in item~4 of its proof hold for every such measure whose spectral measure is carried by the cone.
Consequently $\langle\psi_h,E(\cdot)\psi_h\rangle$ is the image of the measure
\begin{equation*}
\mathrm{kl}(ds)\,\frac{|\hat h(k)|^2\,|k|^2}{2\,\omega_s(k)}\,\frac{d^2k}{(2\pi)^2}
\end{equation*}
under $(s,k)\mapsto(\omega_s(k),k)$, and $\omega_s(k)=|k|$ holds for $s=0$ alone.
This gives the norm of $E(\partial V)\psi_h$;
moreover $\|\psi_h\|^2=\allowbreak \int c_h\,\allowbreak d\mathrm{kl}\le c_h(0)=\allowbreak \frac12\|h\|_{1/2}^2$, and the measure has no atom at $k=0$.

\emph{Step 6} ($\mathcal H_{\mathrm E}^{(1)}$ is invariant).
The operators $\exp(-tH)$ with $t\ge0$ and $V(y)=\exp i(y_2P_2+y_3P_3)$, and the rotations about $e_1$, map $[X(g)]$ to $[X(\mathrm{Gi}\cdot g)]$, with the corresponding Euclidean map $\mathrm{Gi}$ (the property (Rec) of Proposition~3.1, Theorem~3.2\,(b));
by (K16), $\mathcal H_{\mathrm E}^{(1)}$ reduces $H$.
Let $g$ have compact support in $\{y_1>0\}$, and let $t_0>0$ be small.
Then $\zeta\mapsto\exp(\zeta L_j)[X(g)]$ is holomorphic on $|\operatorname{Re}\zeta|<t_0$, and at real $\zeta$ its value is $[X(R^{(j)}_\zeta\cdot g)]\in\mathcal H_{\mathrm E}^{(1)}$ (Theorem~3.2\,(c)).
Pairing with a vector $\phi\perp\mathcal H_{\mathrm E}^{(1)}$ and using the identity theorem, one finds $\exp(isL_j)[X(g)]\in\mathcal H_{\mathrm E}^{(1)}$.
Together with Step~4, it follows that $\mathcal H_{\mathrm E,\gamma}$ is closed and invariant under $\mathcal U$.

\emph{Step 7} (spectral measure of the vectors $[X(g)]$).
Take $g$ as in Step~6, but allow complex values;
the first entry of $C_\varrho$ is the conjugate linear one.
The properties (Rec) and (C7) give
\begin{equation*}
\bigl\langle[X(g)],\exp(-tH)\,V(y)[X(g)]\bigr\rangle=\int d\mathrm{kl}(s)\,C_{\delta_s}(\theta\cdot g,\tau_{te_1+y}\cdot g).
\end{equation*}
Repeating the computation in Step~11 of the proof of Proposition~4.4, with $\tau_{te_1+y}\cdot g$ in the second entry and with $E_s=(|k|^2+s^2)^{1/2}$, one obtains
\begin{equation*}
C_{\delta_s}(\theta\cdot g,\tau_{te_1+y}\cdot g)=\int\frac{d^2k}{(2\pi)^2}\,\frac{|v_2|^2+|v_3|^2-|v_1|^2}{2E_s}\,\exp(-tE_s-ik\cdot y),
\end{equation*}
whose density is non-negative.
Here $v$ is the vector in Step~11 of the proof of Proposition~4.4, with its mass parameter replaced by $s$.
For $s=0$ the density is $|a_g|^2/(2|k|)$, where
\[
a_g(k)=\int_0^\infty dt\int_{\mathbb R^2}d^2y\,
e^{-|k|t-ik\cdot y}\,\operatorname{div}(\ast g)(t,y).
\]
This is the linear amplitude of \cite[Lemma~12.26(a), Remark~12.27]{QEDII}, extended complex linearly.

\emph{Step 8} (a single particle).
According to \cite[Lemma~12.25]{QEDII}, the measures of Step~7, carried over by the map $(s,k)\mapsto(E_s(k),-k)$, give the spectral measure of $[X(g)]$;
on $\partial V$ only $s=0$ contributes, and the apex carries no mass.
Therefore $\iota\colon E(\partial V)[X(g)]\mapsto\mathsf z_\xi^{1/2}a_g$ is isometric, and $\iota\,\mathsf F(H,\vec P)=\mathsf F(|k|,-k)\,\iota$ for every bounded Borel function $\mathsf F$.

\emph{Step 9} ($\iota$ is onto).
As $\varepsilon\to0$, and then $t\downarrow0$, one has $X(u_{\varepsilon,t})\to B_t(h)\to B_0(h)$ in $L^2(\mathrm{lw})$ (Proposition~4.2\,(a) and item~2 of its proof).
Hence $\psi_h\in\mathcal H_{\mathrm E}^{(1)}$, and the integration by parts of \cite[Lemma~12.26\,(c)]{QEDII} gives $\iota E(\partial V)\psi_h=\mathsf z_\xi^{1/2}|k|\hat h$.
The span of these functions is dense, and $\iota$ is therefore onto.

\emph{Step 10} (irreducibility; spectra).
On $\mathcal H_{\mathrm E,\gamma}$ the operators $\mathsf F(H,\vec P)$ act as the maximal abelian algebra $L^\infty(\mathbb R^2)$, so that every closed $\mathcal U$-invariant subspace equals $\iota^{-1}L^2(S)$ for some Borel set $S\subset\mathbb R^2\setminus\{0\}$.
Denote by $m$ the measure on $\partial V$ obtained from $d^2k$ by the map $k\mapsto(|k|,-k)$.
Steps~8 and~9 show that $E(\Delta)$ vanishes on $\mathcal H_{\mathrm E,\gamma}$ if and only if $m(\Delta)=0$.
Suppose that $E(\tilde S)\mathcal H_{\mathrm E,\gamma}=\iota^{-1}L^2(S)$, where $\tilde S:=\{(|k|,-k):k\in S\}$, is invariant under $\mathcal U$.
Step~4 and $\mathcal U(0,\Lambda)\mathcal H_{\mathrm E,\gamma}=\mathcal H_{\mathrm E,\gamma}$ give
\begin{equation*}
E\bigl((\Lambda^T)^{-1}\tilde S\bigr)\mathcal H_{\mathrm E,\gamma}=\mathcal U(0,\Lambda)\,E(\tilde S)\mathcal H_{\mathrm E,\gamma}=E(\tilde S)\mathcal H_{\mathrm E,\gamma},
\end{equation*}
so that $m\bigl((\Lambda^T)^{-1}\tilde S\,\triangle\,\tilde S\bigr)=0$ for all $\Lambda$.
The group $SO_0(1,2)$ acts on $\partial V\setminus\{0\}$ by $\mathsf p\mapsto(\Lambda^T)^{-1}\mathsf p$;
this action is continuous, and the rotations about the axis together with the boosts make it transitive.
By (K16), either $\tilde S$ or its complement is an $m$-null set.
The assertions on the spectra follow from $H\cong|k|$ on $\mathcal H_{\mathrm E,\gamma}$ and from Step~4.

\emph{Step 11} (the field $F$ and the subspace $\mathcal H_{\mathrm E}^{(1)}$).
For $z\in\mathcal T$ and $b\in\Lambda^2\mathbb C^3$, denote by $\Psi_1(z;b)$ the vectors introduced in Step~6 of the proof of Theorems~3.2 and~3.3, and by $T(\zeta)$ the semigroup of Step~1 of that proof.
As functions of $z$ on the convex set $\mathcal T$, the $\Psi_1(z;b)$ are holomorphic in norm, and they depend linearly on $b$;
their closed span is denoted $\mathcal H^{\Psi_1}$.
(a) Step~7 of that proof, at $n=1$, gives $[X(g)]=\int dy\,\Psi_1(Jy;Jg(y))\in\mathcal H^{\Psi_1}$.
If $\phi\perp\mathcal H_{\mathrm E}^{(1)}$, then $\int dy\,\langle\phi,\Psi_1(Jy;Jg(y))\rangle=0$ for every $g\in C_c^\infty(\{y_1>0\};\Lambda^2)$.
For each $b$, the continuous function $y\mapsto\langle\phi,\Psi_1(Jy;Jb)\rangle$ therefore vanishes on $\{y_1>0\}$;
Fact~2 of that proof (with $L=J$) then shows $\phi\perp\mathcal H^{\Psi_1}$, and therefore $\mathcal H_{\mathrm E}^{(1)}=\mathcal H^{\Psi_1}$.
(b) Write $\mathcal H^F$ for the closure of $\{F(h')\Omega:h'\in\mathrm{Sw}_1\}$.
The rule (e4) of Step~18 of that proof, at $n=1$, gives $\exp(-tH)F(h')\Omega=\allowbreak \int dx\,\allowbreak \Psi_1(x+ite_0;h'(x))\in\mathcal H^{\Psi_1}$ for $t>0$, and the limit $t\downarrow0$ gives $\mathcal H^F\subset\mathcal H^{\Psi_1}$.
The rule (e3) of the same step and $T(-a)\Omega=\Omega$ give $T(a)F(h')\Omega=F(\tau_ah')\Omega$ for real $a$.
Thus $\mathcal H^F$ and its orthogonal complement are invariant under the unitary group $T(a)=\exp i(a^0H+\vec a\cdot\vec P)$, $a\in\mathbb R^3$.
Conjugation by the unitary operator $2Q-1$, where the projection $Q$ onto $\mathcal H^F$ commutes with this group, turns $E$ into a projection valued measure for the same group;
by the uniqueness in (K16) it is $E$ itself.
Consequently $\mathcal H^F$ reduces $H$, and $\exp(-tH)$ leaves it invariant.
For $\phi\perp\mathcal H^F$ one obtains $\int dx\,\langle\phi,\Psi_1(x+ite_0;h'(x))\rangle=0$ for every $h'\in\mathrm{Sw}_1$, whence $\langle\phi,\Psi_1(x+ite_0;b)\rangle$ vanishes for every real $x$ and every $b$.
The function $z\mapsto\langle\phi,\Psi_1(z+ite_0;b)\rangle$ is holomorphic on $\mathcal T-ite_0$, which contains $\mathbb R^3$, and Fact~2 with $L=1$ yields $\phi\perp\mathcal H^{\Psi_1}$.
Thus the three spaces $\mathcal H^F$, $\mathcal H^{\Psi_1}$ and $\mathcal H_{\mathrm E}^{(1)}$ coincide, and $\{E(\partial V)F(h')\Omega\}$ is dense in $\mathcal H_{\mathrm E,\gamma}$.
By the definition of $\mathcal H^{(1)}$ in Theorem~A,
$\mathcal H_{\mathrm E}^{(1)}=\mathcal H^{(1)}$ and
$\mathcal H_{\mathrm E,\gamma}=\mathcal H_\gamma$.
That the boosts leave $\mathcal H_{\mathrm E}^{(1)}$ invariant, shown in Step~6, can also be read off from $\mathcal U(0,\Lambda)\Psi_1(z;b)=\Psi_1(\Lambda z;\Lambda b)$ (Step~12 of the proof of Theorems~3.2 and~3.3).


\end{proof}


\begin{proof}[Proof of Corollary B]
The physical mass bounds are Theorem~6.1(a). The fourth-cumulant
bounds in the parent law and the components are
\cite[Theorem~44(B3), (D6)]{QEDIII}. Theorem~A gives $z_\xi>0$.
If $\mathrm{kl}_\xi((0,\infty))=0$, the rigidity theorem in
Section~4 gives $\mu_\xi=\gamma_{z_\xi}$, contradicting the positive
fourth cumulant. Since the total spectral mass is at most one,
$z_\xi<1$. The Maxwell limit follows from Theorem~A(iv).
\end{proof}

The construction gives limit points. Convergence to a unique law and
ergodicity of the parent limit point remain open. The mass hypothesis
in Theorem~A is essential to the Fisher estimate supplied here; at
nonpositive lattice mass or at mass $o(\eta)$ this argument supplies
no positive photon weight. At arbitrary charge the theorem does not
decide Gaussianity. It also gives neither a scattering theory nor a
description of possible additional massless sectors.


\appendix
\section{The finite-torus moment representation}
\label{app:finite-torus-proof}
We prove Proposition~5.12 using reflection positivity and a truncated Stieltjes moment representation.
The lattice quantities and the facts (K1)--(K11) are those of Section~5.

\begin{proof}[Proof of Proposition~5.12]
\emph{Step 1} (positivity).
The torus $C(n)$ is adapted to $x_1$, $x_2$ and $x_3$, with height period $n$ \cite[Section~1.1]{QEDII}, and $\mathrm{Adm}(x_1)=\frac12\mathbb Z$ \cite[Definition~2.2]{QEDII}.
Fix $c\in\frac12\mathbb Z$, and consider the sites $x$ of the open slab $\{c<x_1<c+\frac n2\}$.
For these sites $\mathfrak M(x)\in\mathcal B^+$ and $\Theta_c\mathfrak M(x)=\mathfrak M^\star(\theta_cx)$, while $\mathbb E[\mathfrak M^\star(y)\mathfrak M(x)]=\mathsf M_n(y,x)$ (Proposition~5.4\,(a), (c)).
Consequently \cite[Theorem~1.1\,(c)]{QEDII} shows that, on these sites, the kernel $(x,x')\mapsto\mathsf M_n(\theta_cx,x')$ is Hermitian and positive semidefinite.
This property is denoted \textbf{(RP)}.

\emph{Step 2} (proof of (a)).
Every site of the layer $L_\tau$ has $\tau-\frac12<x_1<\tau+\frac12$ \cite[Definition~2.4]{QEDII}, and $\theta_cL_\tau=L_{2c-\tau}$ \cite[Proposition~3.6\,(c)]{QEDII}.
Evaluated on the indicator functions of the layers that lie in the slab, (RP) together with translation invariance yields the positive semidefinite matrix $\bigl(n^2G_n(\tau+\tau'-2c)\bigr)_{\tau,\tau'}$.
When $c=\frac12$, the slab contains the layers $L_\tau$ with $1\le\tau\le k_A$; when $c=0$, it contains those with $1\le\tau\le\lfloor\frac{n-1}2\rfloor$, and $\lfloor\frac{n-1}2\rfloor\ge k_A-1$.

\emph{Step 3} (proof of (b), a truncated Stieltjes problem).
Put $\mathsf g_j:=G_n(j+1)$ and $k:=k_A-1$, write $\mathcal P_d$ for the space of real polynomials of degree at most $d$, and let $\mathsf L$ be the linear form on $\mathcal P_{2k}$ with $\mathsf L(\lambda^j):=\mathsf g_j$.
Part (a) shows that $(f,h)\mapsto\mathsf L(fh)$ is positive semidefinite on $\mathcal P_k$, and that $\mathsf b(f,h):=\mathsf L(\lambda fh)$ defines a positive semidefinite form on $\mathcal P_{k-1}$.
Whenever $h\in\mathcal P_{k-1}$ and $\mathsf L(h^2)=0$,
\begin{equation*}
\mathsf b(f,h)^2\le\mathsf L(\lambda^2f^2)\,\mathsf L(h^2)=0.
\end{equation*}
Hence $\mathsf b$ descends to the quotient $\mathcal P_{k-1}/\{h:\mathsf L(h^2)=0\}$, equipped with the inner product induced by $\mathsf L$.
There it takes the form $\mathsf b(f,h)=\langle[f],\mathsf J[h]\rangle$ with an operator $\mathsf J=\mathsf J^*\ge0$ satisfying $\mathsf J^i[1]=[\lambda^i]$ for $i\le k-1$.
Let $\sigma_n$ be the spectral measure of $\mathsf J$ at the vector $[1]$; then
\begin{equation*}
\int\lambda^{i+i'}\,\mathrm d\sigma_n=\mathsf g_{i+i'}\quad(i,i'\le k-1),\qquad\int\lambda^{2k-1}\,\mathrm d\sigma_n=\mathsf b(\lambda^{k-1},\lambda^{k-1})=\mathsf g_{2k-1}.
\end{equation*}

\emph{Step 4} (proof of (c)).
Decompose $\sigma_n$ at $\lambda_0:=\exp(-E_0)$: its part on $\{\lambda\le\lambda_0\}$, expressed in the variable $E$, is $\varrho_n$, and its part on $\{\lambda>\lambda_0\}$ produces $R_n(t)$.
If $\lambda>\lambda_0$ and $t\le k_A$, then $\lambda^{t-1}\le\lambda_0^{t-k_A}\lambda^{k_A-1}$, and therefore $R_n(t)\le\exp(E_0(k_A-t))\,G_n(k_A)$.
From (E1) one obtains
\begin{equation*}
G_n(k_A)\le\mathsf F(k_A)\le\exp\bigl(-2\varepsilon_*\lfloor k_A/2\rfloor\bigr)\,\mathsf F(0)\le96\,\mathbf m^{-2}\exp\bigl(-\varepsilon_*(k_A-1)\bigr).
\end{equation*}
Next, for standard weights and $\mathbf m\le(1+\rho)\,\mathrm{sw}$, put $\xi:=\frac{\mathbf m}{(1+\rho)\mathrm{sw}}\le1$; then
\begin{equation*}
\exp(4\vartheta)=(1+\xi)^4\le1+15\xi<1+4\,\frac{\mathrm{sw}}{w_3}\,\xi=\exp(2\varepsilon_*).
\end{equation*}
Finally, let $n\ge6$ be even.
By Lemma~5.9\,(a), $G_n(2)=G_n(n-2)$, and both arguments lie in the range of (b).
For a measure $\sigma\ge0$ on $[0,1)$, on the other hand, $\int\lambda^{n-3}\,\mathrm d\sigma<\int\lambda\,\mathrm d\sigma$ unless $\sigma((0,1))=0$.

\emph{Step 5} (proof of (d)).
The total mass of $\sigma_n$ is $G_n(1)\le96\,\mathbf m^{-2}$, and $R_n(t)\to0$ as $n\to\infty$ for every $E_0<\varepsilon_*$.
By (K5) there is a subsequence along which the restrictions of $\sigma_{n_i}$ to $[0,\exp(-E_0)]$ converge.
Their limit represents $G_\infty$, is independent of $E_0$, and is carried by $\bigcap_{E_0<\varepsilon_*}[0,\exp(-E_0)]$.
\end{proof}

\begin{thebibliography}{99}

\bibitem{OS1973}
K.~Osterwalder and R.~Schrader,
\emph{Axioms for Euclidean Green's functions},
Commun.\ Math.\ Phys.\ 31 (1973), 83--112.
\url{https://doi.org/10.1007/BF01645738}.

\bibitem{OS1975}
K.~Osterwalder and R.~Schrader,
\emph{Axioms for Euclidean Green's functions II},
Commun.\ Math.\ Phys.\ 42 (1975), 281--305.
\url{https://doi.org/10.1007/BF01608978}.

\bibitem{ReehSchlieder1961}
H.~Reeh and S.~Schlieder,
\emph{Bemerkungen zur Unit\"ar\"aquivalenz von Lorentzinvarianten Feldern},
Nuovo Cimento 22 (1961), 1051--1068.
\url{https://doi.org/10.1007/BF02787889}.

\bibitem{QEDI}
\emph{Noncompact lattice QED in three dimensions: limit points of the curvature field
and their fourth cumulant} (companion paper).

\bibitem{QEDII}
\emph{Noncompact lattice QED in three dimensions: chamber fermions,
reflection positivity at all mirrors and rotation invariance} (companion paper).

\bibitem{QEDIII}
\emph{Noncompact lattice QED in three dimensions: the chamber fermion,
the renormalisation group analysis and the fourth cumulant} (companion paper).

\bibitem{Sidak1967}
Z.~\v{S}id\'ak, \emph{Rectangular confidence regions for the means of multivariate normal distributions}, J.\ Amer.\ Statist.\ Assoc.\ 62 (1967) 626--633.

\bibitem{SteinWeiss1971}
E.~M.~Stein and G.~Weiss, \emph{Introduction to Fourier analysis on Euclidean spaces},
Princeton University Press, Princeton, NJ, 1971.

\bibitem{AxlerBourdonRamey}
S.~Axler, P.~Bourdon and W.~Ramey, \emph{Harmonic function theory},
second edition, Graduate Texts in Mathematics 162, Springer, 2001.
\url{https://www.axler.net/HFT.pdf}.


\bibitem{FageotAminiUnser2014}
J.~Fageot, A.~Amini and M.~Unser,
\emph{On the continuity of characteristic functionals and sparse stochastic modeling},
arXiv:1401.6850v2, Theorem~1.
\url{https://arxiv.org/abs/1401.6850}.

\end{thebibliography}

\end{document}
